% This file was converted to LaTeX by Writer2LaTeX ver. 0.5
% see http://www.hj-gym.dk/~hj/writer2latex for more info
\documentclass{article}
\usepackage[ascii]{inputenc}
\usepackage[T1]{fontenc}
\usepackage[ngerman]{babel}
\usepackage{amsmath,amssymb,amsfonts,textcomp}
\usepackage{makeidx}
\usepackage{multicol}
\usepackage{array}
\usepackage{supertabular}
\usepackage{hhline}
\usepackage{graphicx}
\usepackage{ooomath}
\makeindex
\setlength\tabcolsep{1mm}
\renewcommand\arraystretch{1.3}
% Non-floating captions
\makeatletter
\newcommand\captionof[1]{\def\@captype{#1}\caption}
\makeatother
\newcounter{Drawing}
\renewcommand\theDrawing{\arabic{Drawing}}
\title{Arbeitsstandards f\"ur den Sanit\"atsdienst, Version f\"ur \"Osterreich, Stand November 2008 (1)}
\begin{document}
\begin{center}
\tablehead{}\begin{supertabular}{m{3.0279999cm}m{10.0130005cm}m{3.372cm}}
\multicolumn{3}{m{16.813cm}}{\centering Der Sanit\"atsdienst\par

\begin{center}
\begin{minipage}{16.97cm}
\centering Arbeits- und Ausbildungsstandards\par

\centering und eine Einf\"uhrung in die Welt der Medizin und
Sanit\"atshilfe\par

\end{minipage}
\end{center}
\begin{center}
 [Warning: Image ignored] % Unhandled or unsupported graphics:
%\includegraphics[width=15.432cm,height=11.903cm]{AASdev20090228-img1.png}

\end{center}
}\\\hline
Wien, {}--.-{}-.2008

 &
 &
\raggedleft Ausgabe 0\par

\\\hline
   [Warning: Image ignored] % Unhandled or unsupported graphics:
%\includegraphics[width=1.36cm,height=3.962cm]{AASdev20090228-img2.png}
 

 &
Mission Statement {\textperiodcentered} Das Ausbildungszentrum (ABZ) ist
eine Dienst\-leistungs\-ein\-richtung des ASB Floridsdorf-Donau\-stadt
und orientiert sein Angebot an den Be\-d\"urf\-nissen unserer Partner
und an den strategischen Zielen des ASB. Wir sind sowohl f\"ur die 
medizinische und sanit\"atsdienstliche Ausbildung unserer Mitarbeiter
als auch f\"ur die Durchf\"uhrung von Schulungen und Lehrg\"angen zur
Schulung der Bev\"olkerung verantwortlich. In Zusammenarbeit mit der
\"Arztlichen Leitung und den \"ubergeordneten Gremien  entwickeln wir
Arbeitsstandards und Lehrmeinungen. 

Um die Qualit\"at der Arbeit des ASB auf hohem Niveau zu halten verfolgt
das ABZ eine aktive Informations- und Ausbildungspolitik. 

 &
\centering   [Warning: Image ignored]
% Unhandled or unsupported graphics:
%\includegraphics[width=2.182cm,height=2.208cm]{AASdev20090228-img3.png}
 \par

\centering Ausbildungszentrum\par

\centering Floridsdorf-Donaustadt\par

\\\hline
 &
\centering DRAFT -- ENTWURF -- DRAFT\par

 &
\\
\end{supertabular}
\end{center}
\setcounter{tocdepth}{1}
\renewcommand\contentsname{Inhalt - Schnell\"ubersicht}
\tableofcontents
\setcounter{tocdepth}{10}
\renewcommand\contentsname{Inhalt - Detail\"ubersicht}
\tableofcontents
Autoren und Mitwirkende

(in alphabetischer Reihenfolge)

\begin{multicols}{2}
\begin{center}
\begin{tabular}{m{6.813cm}}
Josef Emhofer

Mag. phil.

Ausbildungszentrum des ASB Floridsdorf-Donaustadt

Wallenberggasse 2, A-1220 Wien

\\
Gerhard Gabriel

Dr. med. univ. 

Facharzt f\"ur Innere Medizin

A-2103, Langenzersdorf

\\
Sebastian Gabriel

Medizinische Universit\"at Wien -- Besondere Einrichtung zur
Medizinischen Aus- und Weiterbildung

Spitalgasse 23, A-1090 Wien

\\
Lena Hirtler

Medizinische Universit\"at Wien -- Zentrum f\"ur Anatomie und
Zellbiologie, Abt. f\"ur systematische Anatomie

W\"ahringerstra{\ss}e 13, A-1090 Wien 

\\
Thomas Hochreiter

DGKP

Ausbildungszentrum des ASB Floridsdorf-Donaustadt

Wallenberggasse 2, A-1220 Wien

\\
Roman Koch

Ing.

Leiter Ausbildungszentrum des ASB Floridsdorf-Donaustadt

Wallenberggasse 2, A-1220 Wien

\\
Christof Koller

Ausbildungszentrum des ASB Floridsdorf-Donaustadt

Wallenberggasse 2, A-1220 Wien

\\
Michael Motal

Ausbildungszentrum des ASB Floridsdorf-Donaustadt

Wallenberggasse 2, A-1220 Wien

\\
Johannes Steuer

Dr. med. univ., M.Sc.

\"Arztlicher Leiter und Medizinisch-wissenschaftlicher Leiter der
Curricula

ASB Floridsdorf-Donaustadt

Wallenberggasse 2, A-1220 Wien

Sozialmedizinisches Zentrum Ost der Stadt Wien -- Donauspital

Abteilung f\"ur An\"asthesiologie und Intensivmedizin

Langobardenstra{\ss}e 122, A-1220 Wien

\\
Michael Witthofner

Dr. med, univ. 

Wien 

\\
\end{tabular}
\end{center}
\end{multicols}
Danksagungen

New Somerset Hospital, Kapstadt, S\"udafrika

Western Cape METRO Emergency Medical Service, Kapstadt, S\"udafrika

Sponsoren

Software

\begin{center}
\tablehead{}\begin{supertabular}{m{5.0740004cm}m{5.0740004cm}m{5.0740004cm}}
JabRef 2.4.2 

 &
\centering Literaturverzeichnis\par

 &
\\
OpenOffice 3.0

 &
\centering Textverarbeitung\par

 &
\raggedleft http://www.openoffice.com\par

\\
Ubuntu Linux 8.04, 8.10

 &
\centering Betriebssystem\par

 &
\raggedleft http://www.ubuntu.com\par

\\
 &
 &
\\
 &
 &
\\
\end{supertabular}
\end{center}
\section{T\"atigkeitsbild und Orientierung}
Autor fehlt

Derzeit in Vorbereitung.

Der Sanit\"ater in \"Osterreich -- zwischen Sch\"utzengraben und
Flusanit\"ater

Der Rettungsdienst in der gro{\ss}en weiten Welt 

Ein funktionierendes Sozialversicherungssystem ist Grundlage f\"ur
optimale Leistungen

\section{Betriebsinterne Regelungen}
Angepasst an: Arbeiter-Samariter-Bund Gruppe Floridsdorf-Donaustadt und 
Arbeiter-Samariter-Bund Floridsdorf-Donaustadt Rettungs-
Krankentransport und soziale Dieste gem. GesmbH.

Derzeit in Vorbereitung.

Siehe externe Beilage. 

\section{[RS III] Rechtliche Grundlagen}
Christof Koller [*]

Changelog:

2009-02-20: Fertig f\"ur Kurs 2009/02

2009-02-25: Edit Koller

\setcounter{tocdepth}{10}
\renewcommand\contentsname{\"Ubersicht}
\tableofcontents
\subsection{Allgemeine Rechtsgrundlagen}
[Warning: Draw object ignored]Struktur der \"osterreichischen
Rechtsordnung

\begin{center}
\tablehead{}\begin{supertabular}{m{2.299cm}|m{1.8669999cm}m{2.082cm}m{2.0839999cm}|m{2.082cm}m{2.0839999cm}|}
\multicolumn{1}{m{2.299cm}}{\centering [Warning: Draw object
ignored]EU-Recht\par

} &
\multicolumn{5}{m{10.999cm}}{\centering [Warning: Draw object
ignored]Struktur der \"osterreichischen Rechtsordnung\par

}\\\hhline{~-----}
 &
\multicolumn{3}{m{6.433cm}|}{\centering \"Offentliches Recht\par

} &
\multicolumn{2}{m{4.366cm}|}{\centering Privatrecht\par

}\\\hhline{~-----}
\centering Rechtsgebiete\par

 &
\multicolumn{1}{m{1.8669999cm}|}{\centering Strafrecht\par

} &
\multicolumn{1}{m{2.082cm}|}{\centering VwR\par

} &
\centering VfR\par

 &
\multicolumn{1}{m{2.082cm}|}{\centering Zivilrecht\par

} &
\centering Sonderprivat-rechte\par

\\\hhline{~-----}
\centering wichtigste Gesetzesbereiche\par

 &
\multicolumn{1}{m{1.8669999cm}|}{\centering StGB,StPO\par

} &
\multicolumn{1}{m{2.082cm}|}{\centering Berufsgesetze, KAKuG, UbG\par

} &
\centering BVG; StGG\par

 &
\multicolumn{1}{m{2.082cm}|}{\centering ABGB\par

} &
\centering Arbeitsrechts-gesetze, Unternehmens-gesetzbuch\par

\\\hhline{~-----}
\end{supertabular}
\end{center}
Die \"osterreichische Rechtsordnung ist ein zusammenh\"angendes Geflecht
bestehend aus zwei Hauptbereichen, dem \"offentlichen Recht und dem
Privatrecht. W\"ahrend das \"offentliche Recht die Beziehung zwischen
Staat und B\"urger regelt, bestimmt das Privatrecht die Verh\"altnisse
unter den B\"urgern (= Privaten) selber. 

\"Offentliches Recht: Staat ${\leftrightarrow}$ B\"urger

Privatrecht: B\"urger ${\leftrightarrow}$ B\"urger

Durch den Beitritt zur EU wurde ein weiterer Gesetzesgeber geschaffen,
der in beide Rechtsgebiete einwirken kann. Unabh\"angig davon l\"asst
sich jedes von beiden Rechts-gebieten in folgende Bereiche unterteilen:

\"Offentliches Recht

\begin{itemize}
\item Strafrecht: Erfasst jene Delikte, die der Gesetzgeber als derart
schwerwiegend erachtet hat, dass sie nicht nur (idR.) vom Staatsanwalt
vor Gericht zu verfolgen sind,  sondern auch (in den meisten F\"allen)
mit Haftstrafen sanktionierbar sind. Der Staatsanwalt vertritt hierbei
den Staat als Ankl\"ager. Die ein gehobene Geldstrafe erh\"alt der
Staat, nicht der Verletzte.

\item Verwaltungsrecht: Umfasst Regelungen die sich idR. nur an
bestimmte Personenkreise richten und von Verwaltungsbeh\"orden (nicht
von Gerichten) bei Versto{\ss} insofern sanktioniert werden, dass
(vorrangig) Geldstrafen verh\"angt werden.

\item Verfassungsrecht: Jener Teil der Rechtsordnung, der die
grundlegenden Bestimmungen \"uber Staatsgewalt, ihre Organe und ihre
Verh\"altnisse zu den Staatsb\"urgern sowie Grunds\"atze der
Rechtsordnung enth\"alt.

\end{itemize}
Privatrecht

\begin{itemize}
\item Zivilrecht: Enth\"alt jene allgemeinen Normen, die die
Rechtsverh\"altnisse zwischen gleichgestellten B\"urgern regeln. Bei
Verst\"o{\ss}en muss der Verletzte selber vor Gericht sein Recht
einfordern (kein Staatsanwalt), wodurch er vom Sch\"adiger
Ersatzleistungen f\"ur den Schaden fordern kann (keine Strafe).

\item Sonderprivatrechte: eigenst\"andige Rechtsgebiete, die mittels
selbstst\"andigen Normen besondere Rechtsverh\"altnisse zwischen den
B\"urgern regeln. Zu den wichtigsten z\"ahlen das Arbeits- und das
Unternehmensrecht.

\end{itemize}
Vertragsverh\"altnis im Rettungswesen

\begin{flushleft}
\tablehead{}\begin{supertabular}{m{3.3cm}m{5.6210003cm}m{4.1730003cm}}
 &
\centering Heilbehandlungsvertrag \par

\centering Gesch\"aftsf\"uhrung ohne Auftrag im Notfall\par

 &
\\
\raggedleft Patient\par

 &
\centering [Warning: Draw object ignored]\par

 &
Rettungsorganisation

(bzw. Krankenanstalt)

\\
\raggedleft Erf\"ullungsgehilfe\par

 &
[Warning: Draw object ignored]  [Warning: Draw object ignored]

 &
Arbeitsvertrag

Vereinsmitgliedschaft

Zivildienst

\\
 &
\centering (Not)Arzt/Sanit\"ater\par

 &
\\
\end{supertabular}
\end{flushleft}
Das Rechtsverh\"altnis im Rettungswesen besteht grunds\"atzlich aus drei
Parteien: dem Patienten, der Rettungsorganisation und dem (Not)Arzt
oder Sanit\"ater. (In einer Krankenanstalt ist das Verh\"altnis
\"ahnlich, mit dem Unterschied, dass an die Stelle der
Rettungsorganisation die Krankenanstalt und an jene des Sanit\"aters
der Krankenpfleger tritt.)

Als wichtigstes Verh\"altnis in dieser Dreiecksbeziehung ist das
Verh\"altnis zwischen Patient und Rettungsorganisation zu bezeichnen,
da direkter (Vertrags)Partner zum Pat. niemals der Sanit\"ater, sondern
stets die von ihm vertretene Organisation ist. Man nimmt n\"amlich an,
dass ein Pat., wenn er eine {\quotedblbase}Rettung{\textquotedblleft}
verst\"andigt, nicht mit dem einzelnen Einsatzpersonal in rechtlichen
Kontakt treten m\"ochte, sondern mit der Gesamtheit einer Organisation.
(Er bestellt ja nicht dieses eine Fahrzeug mit jenem speziellen
Personal.) Das Rechtsverh\"altnis kann nun auf zwei verschiedene Arten
zustande kommen:

\begin{itemize}
\item Heilbehandlungsvertrag: der Pat. selber w\"unscht behandelt zu
werden

\item Gesch\"aftsf\"uhrung ohne Auftrag im Notfall: der Pat. befindet
sich in einem derartigen Notfall, dass er sich nicht mehr \"au{\ss}ern
kann und dringend eine Behandlung ben\"otigt (Bsp.: Bewusstlosigkeit)

\end{itemize}
Aus beiden Verh\"altnissen schuldet die Rettungsorganisation die
h\"ochstm\"ogliche Sorgfalt, aber niemals einen bestimmten
Behandlungserfolg. Weiters hat sie vorrangig nach dem Willen des Pat.
vorzugehen und wenn dieser fehlt, nach seinem hypothetischen Willen.

Im Verh\"altnis zwischen Patient und (Not)Arzt/Sanit\"ater fungiert der
(Not)Arzt/Sanit\"ater blo{\ss} als Gehilfe zur Erf\"ullung der
Pflichten seiner Organisation gegen\"uber dem Pat., wodurch man ihn
auch als Erf\"ullungsgehilfen bezeichnet. Er ist nicht Vertragspartner
des Pat., muss aber mit der gleichen Sorgfalt, wie es von der
Organisation verlangt wird, arbeiten. Kommt es nun zu Sch\"adigungen
des Pat. durch Handlungen des (Not)Arztes/ Sanit\"aters, kann der
Gesch\"adigte deliktisch gegen den Erf\"ullungsgehilfen und vertraglich
gegen die Organisation vorgehen. (IdR. wird der Pat. die Organisation
verklagen, da ihm hierbei aufgrund des Vertrages mehrere gesetzliche
Erleichterungen zustehen.)

Das Verh\"altnis zwischen Rettungsorganisation und (Not)Arzt/Sanit\"ater
l\"asst sich auf drei verschiedene Arten begr\"unden:

\begin{itemize}
\item Arbeitsvertrag: Der (Not)Arzt/Sanit\"ater ist bei der Organisation
als Hauptamtlicher berufsm\"a{\ss}ig besch\"aftigt. 

\item Vereinsmitgliedschaft: Der (Not)Arzt/Sanit\"ater ist bei der
Organisation als Vereinsmitglied und somit als
Ehrenamtlicher/Freiwilliger t\"atig. 

\item Zivildienst: Der Sanit\"ater leistet seinen Zivildienst bei der
Organisation.

\end{itemize}
Aus diesem Verh\"altnis heraus m\"ussen manche der Pflichten, welche die
Organisation gegen\"uber dem Pat. treffen, auch durch den
(Not)Arzt/Sanit\"ater erf\"ullt werden. Gleichzeitig treffen auch die
Organisation und den Sanit\"ater Rechte und Pflicht im Umgang
miteinander.

Strafbarkeitspr\"ufung

\begin{itemize}
\item[] 1. Tatbestandserf\"ullung bzw. Schaden

2. Kausalit\"at

3. Rechtswidrigkeit

\end{itemize}
\begin{itemize}
\item[] 3.1. Rechtfertigungsgr\"unde:

\end{itemize}
\begin{itemize}
\item Heilbehandlung

\item Einwilligung

\item Notwehr

\item rechtfertigender Notstand

\item Anhalterecht Privater

\end{itemize}
\begin{itemize}
\item[] 4. Verschulden

\end{itemize}
\begin{itemize}
\item[] 4.1. Entschuldigungsgr\"unde:

\end{itemize}
\begin{itemize}
\item Notwehr\"uberschreitung im Affekt

\item entschuldigender Notstand

\item [Unzurechenbarkeit]

\end{itemize}
Damit es zur Haftung f\"ur eine gesetzte Handlung kommt, m\"ussen die
oben \"uberblicksartig dargestellten und folgend n\"aher beschriebenen
Voraussetzungen vollst\"andig erf\"ullt sein. Mangelt es auch nur an
einer von ihnen, kommt es zu keiner Haftung. Am h\"aufigsten kommen die
oben aufgez\"ahlten Rechtfertigungs- oder Entschuldigungsgr\"unde ins
Spiel.

Tatbestandserf\"ullung 

Erf\"ullung der rechtlichen Beschreibung einer Lebenssituation auf die
eine Rechtsfolge/Strafe angeordnet wird

Kausalit\"at

Die gesetzte Handlung war der Grund f\"ur den Eintritt der
Tatbestandserf\"ullung

Rechtswidrigkeit

Die gesetzte Handlung war objektiv falsch ${\rightarrow}$ Urteil \"uber
Tat ${\rightarrow}$ Versto{\ss} gegen gesetzliche Verbote bzw. Gebote,
vertraglichen Bestimmungen oder die guten Sitten 

 Rechtfertigungsgr\"unde

 Der begangenen Tat kann kein Urteil gemacht werden, da sich ein
ma{\ss}gerechter Mensch in  derselben Situation gleich verhalten
h\"atte

Verschulden

Die gesetzte Handlung war subjektiv falsch ${\rightarrow}$ Urteil \"uber
T\"ater

 Entschuldigungsgr\"unde

 \"Uber den T\"ater kann kein Urteil gemacht werden, da er sich nach
seinen pers\"onlichen   F\"ahigkeiten und Kenntnissen in dieser
Situation nicht anders verhalten h\"atte k\"onnen

Die einzelnen Rechtfertigungs- und Entschuldigungsgr\"unde werden in der
Ausarbeitung {\quotedblbase}Ein langer Dienst{\textquotedblleft}
separat behandelt.

\subsection{Ein langer Dienst}
Eine Geschichte aus 17 Rechtsf\"allen

Zweck

Behandlung rechtlicher Fragen anhand von Praxisbeispielen, die in einer
zusammen-h\"angenden Geschichte erz\"ahlt werden. Dabei wird
insbesondere versucht auf die Besonderheiten des Rettungswesens Bezug
zu nehmen und die sich daraus ergebenden rechtlichen Probleme einfach
darzustellen. Die Darstellung beinhaltet aber nur einen \"Uberblick
\"uber die behandelten Rechtsgebiete sowie eine Konkretisierung auf die
am h\"aufigsten gestellten Rechtsfragen. 

Aufbau

Jeder einzelne Rechtsfall wird in der Geschichte nach folgendem Aufbau
gestaltet:

\begin{itemize}
\item[] 1) Sachverhalt (= SV): Schilderung der Geschichte

2) Rechtsfrage (= RFr): Erkennung des sich stellenden Problems

3) Rechtslage (= RL): rechtliche Beurteilung des Problems

4) Rechtsfolgen (= RF): drohende Konsequenzen und
Ausweichm\"oglichkeiten

\end{itemize}
Vorkommende Personen

\begin{center}
\tablehead{}\begin{supertabular}{m{6.649cm}m{6.651cm}}
Alex, Fahrer

Sepp, Transportf\"uhrer

Bernd, Zivildiener

NAtascha, Not\"arztin

 &
Hr. Hauzu, 1. Patient

Fr. Uninformiert, 2. Patientin

Fr. Finito, 3. Patientin

Hr. Schmerz, 4. Patient

Fr. Aufgeregt, Angeh\"orige

Dr. Wasbringst, Aufnahmearzt

Fr. Hilflos, 5. Patientin

\\
\end{supertabular}
\end{center}
Der Tag beginnt ...

Unsere Geschichte beginnt an einem warmen Fr\"uhjahrsmorgen auf den die
drei Sanit\"ater Alex, Sepp und Bernd schon lange gehofft hatten, da
sie endlich nach dem langen Winter einen angenehm warmen Sonnentag
genie{\ss}en wollen. Doch wie immer an einem solchen Freudentag, ruft
die Arbeit mit einem ganz normalen Sanit\"atsdienst und seinen
typischen Problemen.

Teil A: Gefahrenbek\"ampfung

\paragraph{Fall 1. Randalierender Patient}
Sachverhalt

Die erste Ausfahrt beginnt nat\"urlich direkt nach Dienstbeginn, ohne
dass wenigstens ein Kaffee getrunken werden h\"atte k\"onnen. Der auf
unsere RTW-Mannschaft wartende Patient, Hr. Hauzu, hingegen hat die
ganze Nacht durchgemacht und ist schon weitl\"aufig bekannt f\"ur seine
tobend aggressiven Wutausbr\"uche im Zuge seiner psychiatrischen
Erkrankung. Als Alex, Sepp und Bernd eintreffen, schmei{\ss}t er mit
mehreren Gegenst\"anden um sich und droht sich oder einen anderen zu
verletzen, falls man sich ihm auch nur n\"ahern w\"urde.

Rechtsfrage

Darf man Hr. Hauzu festhalten? Muss Hr. Hauzu eingewiesen werden? 

Rechtslage

Eine Festhaltung ist erlaubt nach folgenden Paragraphen:

\begin{itemize}
\item {\S} 3 StGB -- Notwehr: Abwehr eines Angriffes auf ein eigenes
oder fremdes Rechtsgut, wenn:

\end{itemize}
\begin{itemize}
\item gegenw\"artig oder unmittelbar drohender Angriff,

\item rechtswidrig,

\item von Menschen gesetzt,

\item auf ein notwehrf\"ahiges Gut (= Leben, Gesundheit, Freiheit oder
Verm\"ogen) und

\item das Gelindeste der zur verl\"asslichen Abwehr verf\"ugbaren Mittel
verwendet wird

\end{itemize}
\begin{itemize}
\item {\S}~3 Abs. 2 StGB -- Notwehr\"uberschreitung im Affekt:

\end{itemize}
\begin{itemize}
\item Sollten Alex, Sepp und Bernd f\"alschlich einen Agriff angenommen
haben, sind sie entschuldigt, wenn sie aus Furcht gehandelt haben und
eine ma{\ss}gerechte Person sich gleich verhalten h\"atte. 

\end{itemize}
\begin{itemize}
\item {\S} 86 Abs. 2 StPO -- Anhalterecht Privater: Festhaltung bei
Verdacht auf eine gerichtlich strafbare Handlung, welche im
vorliegenden Fall die Gefahr einer K\"orperverletzung (KV) iSd.
{\S}{\S} 83 ff StGB und die Sachbesch\"adigung iSd. {\S} 125 StGB
darstellt. Dabei darf nur auf angemessene Weise (N\"otigung und leichte
KV) angehalten werden, wenn unverz\"uglich f\"ur eine Meldung an die
Polizei gesorgt wird.

\end{itemize}
Eine Unterbringung muss erfolgen, nach folgenden Gesichtspunkten:

\begin{itemize}
\item {\S} 3 UbG:

\end{itemize}
\begin{itemize}
\item[] a) psychiatrische Krankheit (= KH) 

b) Selbst- und/oder Fremdgef\"ahrdung

c) Verh\"altnism\"a{\ss}igkeitsgrundsatz

\end{itemize}
\begin{itemize}
\item[] zu der Feststellung von a) und b) ben\"otigt es das Gutachten
eines Amtsarztes, welcher der Bezirksverwaltungsstelle (in Wien die MA
15) unterstellt ist und \"uber die Polizei zu verst\"andigen ist. Der
Amtsarzt darf aber nur die {\quotedblbase}Zuweisung{\textquotedblleft}
bis zur Krankenanstalt (= KA) bestimmen, da erst in der KA von 2
Fach\"arzten der Psychiatrie die eigentliche
{\quotedblbase}Einweisung{\textquotedblleft} verf\"ugt wird. W\"ahrend
der Untersuchung und dem Transport darf auch Zwang angewendet werden,
aber nur durch die gelindesten zur Erfolgsherbeif\"uhrung m\"oglichen
Mittel bei Beachtung der Mittel-Ziel-Relation (= 
Verh\"altnism\"a{\ss}igkeitsgrundsatz).

\end{itemize}
Rechtsfolge

\begin{itemize}
\item Alex, Sepp und Bernd haben somit das Recht Gewaltausbr\"uche des
Hr. Hauzu gegen sich und andere abzuwehren. Dabei k\"onnen sie das
verh\"altnism\"a{\ss}ig gelindeste Zwangsmittel w\"ahlen, sind aber
verpflichtet unverz\"uglich die Polizei und \"uber diese den Amtsarzt
zu verst\"andigen. Trotzdem sollte Gewalt durch Sanit\"ater immer als
letzte Option verwendet werden, auch wegen dem Risiko sich durch das
Anhalten selbst zu verletzen. Daher der Grundsatz: Selbstschutz geht
vor Fremdschutz (bzw. \"ubertriebenen Heldenmut).

\end{itemize}
\paragraph{Fall 2. Gefahr durch Hund}
Sachverhalt

Kaum hat Hr. Hauzu geh\"ort, dass die Polizei und der Amtsarzt
verst\"andigt wurden, st\"urmt er zu einem kleinen Nebenzimmer in
seiner Gemeindebauwohnung, \"offnet die T\"ur und l\"asst einen
{\quotedblbase}kleinen{\textquotedblleft} (65~kg schweren)
Dobermann-Kampfhund heraus. Mit fletschenden Z\"ahnen n\"ahert er sich
den Sanit\"atern und droht jeden Moment anzugreifen.

Rechtsfrage

Was darf gegen den Hund getan werden?

Rechtslage

\begin{center}
\tablehead{}\begin{supertabular}{|m{6.649cm}|m{6.651cm}|}
\hline
Rechtfertigender Notstand: 

\begin{itemize}
\item in keinem Gesetz aufgezeichnet; wird aber einheitlich angenommen

\item anzuwenden wenn die Gefahr auch tats\"achlich (objektiv) vorliegt

\end{itemize}
 &
Entschuldigender Notstand - {\S} 10 StGB; {\S} 1306a ABGB:

\begin{itemize}
\item anzuwenden wenn man die Gefahr blo{\ss} subjektiv annimmt 

\end{itemize}
\\
\begin{itemize}
\item Notstandssituation: Gefahr ist:

\end{itemize}
\begin{itemize}
\item kein menschlicher Angriff

\item gegenw\"artig oder unmittelbar drohend

\item auf beliebiges Rechtsgut gerichtet

\item droht bedeutender Nachteil

\end{itemize}
\begin{itemize}
\item Notstandshandlung:

\end{itemize}
\begin{itemize}
\item Abwehr als schonendstes Mittel

\item Zurechnungsprinzip (=~Abwehr geht in Sph\"are des
Gefahrenverursachers)

\item geringes Risiko bei hoher Rettungschance 

\end{itemize}
 &
\begin{itemize}
\item Notstandssituation: Gefahr ist:

\end{itemize}
\begin{itemize}
\item kein menschlicher Angriff

\item gegenw\"artig oder unmittelbar drohend

\item auf beliebiges Rechtsgut gerichtet

\item droht bedeutender Nachteil

\end{itemize}
\begin{itemize}
\item Notstandshandlung: 

\end{itemize}
\begin{itemize}
\item Abwehr als schonendstes Mittel

\item Zurechnungsprinzip (= Abwehr geht in Sph\"are des
Gefahrenverursachers)

\item Drohender Nachteil darf geringf\"ugig kleiner als das Ergebis der
Abwehrhandlung sein.

\end{itemize}
\\\hline
\end{supertabular}
\end{center}
Rechtsfolge

\begin{itemize}
\item Alex, Sepp und Bernd d\"urfen sich mit den schonendsten Mitteln
verteidigen, um ihr eigenes Leben oder ihre Gesundheit zu sch\"utzen.

\end{itemize}
Praxistipp: Schon vor Betreten der Wohnung sollte ein Hund von seinem
Besitzer in einen  separaten Raum weggesperrt werden.

Teil B: Heilbehandlung

\paragraph{Fall 3. Blutzuckermessung = Heilbehandlung}
Sachverhalt

Durch die rechtzeitig eintreffende Polizei, kann die Gefahr durch den
Hund beseitigt und Hr. Hauzu festgehaltenen werden. Trotz der Zuweisung
durch den etwas sp\"ater eingetroffenen Amtsarzt, m\"ochte Bernd,
gewissenhaft wie er ist, einen vollen Neuro-Check mit einem
Blutzuckertest durchf\"uhren.

Rechtsfrage

Wann liegt eine Heilbehandlung (= HB) vor? Darf der Blutzuckertest
durchgef\"uhrt werden?

Rechtslage

Heilbehandlung:

Wird gewonnen aus allgemeinen Rechtsgrunds\"atzen und der systematischen
Betrachtung der gesamten Rechtsordnung.

Voraussetzung:

\begin{itemize}
\item medizinisch indiziert (= HB ist zur Diagnose oder Linderung von
Erkrankungen und/oder Verletzungen notwendig)

\item lege artis durchgef\"uhrt (= HB entspricht dem geltenden Stand der
Wissenschaft)

\item durchzuf\"uhren nach Aufkl\"arung des Pat. (${\rightarrow}$ siehe
ad 4.)

\end{itemize}
${\rightarrow}$ das Vorliegen einer HB wird als Rechtfertigung f\"ur
alle Einwirkungen am Pat. gewertet

${\rightarrow}$ wer die Voraussetzungen der HB nicht einh\"alt macht
sich strafbar nach:

\begin{itemize}
\item Strafrecht:

\end{itemize}
\begin{itemize}
\item[] a) K\"orperverletzung - {\S}{\S} 83 ff StGB: Sonderfall der
HB-Voraussetzungen: es ben\"otigt nur med. Indikation und lege artis
Durchf\"uhrung; die Aufkl\"arung ist hierf\"ur irrelevant

b) Eigenm\"achtige Heilbehandlung - {\S} 110 StGB: siehe ad 4.

\end{itemize}
\begin{itemize}
\item Zivilrecht:

\end{itemize}
\begin{itemize}
\item[] a) Schadenersatzanspr\"uche - {\S}{\S} 1293 ff ABGB: die HB muss
alle drei Voraussetzungen erf\"ullen (auch Aufkl\"arung); Umfang des
Schadenersatzes siehe ad 11.

\end{itemize}
Blutzuckermessung:

Die Messung des Blutzucker mit einer Lanzette war fr\"uher umstritten,
da manche Rettungsorganisationen davon ausgingen, dass es sich um einen
nicht dem Sanit\"ater erlaubten Eingriff am menschlichen K\"orper
handelt. Dieser Streit wurde durch eine Novellierung des SanG gel\"ost,
in dem man die Blutzuckermessung f\"ur den Sanit\"ater ausdr\"ucklich
zu lies.

Rechtsfolge

\begin{itemize}
\item Da ein kompletter Neuro-Check bei einem neurologisch auffallenden
Pat. medizinisch indiziert und gesetzlich erlaubt ist, darf jeder der
drei Sanit\"ater eine solche Ma{\ss}nahme lege artis durchf\"uhren.
Eine Aufkl\"arung ist aufgrund der mangelnden Einsichts- und
Urteilsf\"ahigkeit nicht m\"oglich.

\end{itemize}
Teil C: Selbstbestimmungsrecht

\paragraph{Fall 4. Aufkl\"arung und Einwilligung}
Sachverhalt

Nachdem nun Hr. Hauzu geb\"andigt, zur G\"anze untersucht und
eingewiesen wurde, freut sich die RTW-Mannschaft schon auf den
wohlverdienten Kaffee am St\"utzpunkt. Doch wie es kommen musste,
k\"onnen Alex, Sepp und Bernd nicht einr\"ucken, da sie einen
Folgetransport zugewiesen bekommen. Die vorgefundene Pat. Fr.
Uninformiert ist eine in der 30. Wo. schwangere Mutter, die erst letzte
Woche ihren 18. Geburtstag feierte. Grund f\"ur ihren Anruf sind starke
Kopfschmerzen und leichte Seh- und H\"oreinschr\"ankungen. Bei der
weiteren Untersuchung durch Sepp werden geschwollene Kn\"ochel, ein
Blutdruck von 140/90 mm/Hg und nach beil\"aufiger Aussage der Pat.
starker Harndrang mit s\"u{\ss}lichem Geruch diagnostiziert.

Schnell ist Sepp klar, dass es sich um eine EPH-Gestose handelt, die zur
Eklampsie f\"uhren kann und einen raschen, aber m\"oglichst reizarmen
Transport auf die Gyn\"akologie ben\"otigt. Da er aber wenig von jungen
M\"uttern h\"alt, erachtet er es auch nicht f\"ur notwendig Fr.
Uninformiert \"uber die n\"achsten Geschehnisse zu informieren, sondern
ordnet nur harsch Befehle an, die von Fr. Uninformiert
{\quotedblbase}gef\"alligst{\textquotedblleft} durchzuf\"uhren sind. Er
redet \"uberhaupt bis zur Aufnahme nur das n\"otigste mit ihr.

Rechstfrage

Besteht eine Aufkl\"arungspflicht und was umfasst sie? Muss eine
Einwilligung erfolgen und wenn ja, in welcher Form?

Rechtslage

\begin{multicols}{2}
Aufkl\"arungspflicht:

\begin{itemize}
\item Jeder Pat. ist aufzukl\"aren,

\item vor Setzung der Ma{\ss}nahme,

\item wenn er einsichts- und urteilsf\"ahig ist.

\end{itemize}
Einsichts- und Urteilsf\"ahigkeit:

\begin{itemize}
\item grunds\"atzlich: bei jedem vollj\"ahrigen Pat. gegeben 

\item fehlend bei: 

\end{itemize}
\begin{enumerate}
\item grunds\"atzlich unm\"undigen Minderj\"ahrigen (0-14 LJ.), wobei es
vorrangig auf die Einsichts- und Urteilsf\"ahigkeit ankommt

\item bei m\"undig Minderj\"ahrigen (14.-18 LJ.) vorhanden, solange
entsprechende Einsichts- und Urteilsf\"ahigkeit nicht fehlt

\item psychisch erkrankter Pat. mit fehlender Einsichts- und
Urteilsf\"ahigkeit

\item Pat. im Schockzustand

\item stark berauschter Pat. (ab ca. 2,5{\textperthousand})

\item Bewusstlosen

\end{enumerate}
\begin{itemize}
\item[] ad 1) \& 2): {\S} 146c ABGB:

\end{itemize}
\begin{itemize}
\item grunds\"atzlich entscheidend ist die vom Einsatzpersonal
getroffene Beurteilung der Einsichts- und Urteilsf\"ahigkeit des
Minderj\"ahrigen

\item nur im Zweifel ist erst der m\"undige Minderj\"ahrige
selbstst\"andig entscheidungsbefugt (es ben\"otigt keine Zustimmung der
Eltern) 

\item nur bei schwerwiegenden Angelegenheiten muss ein Elternteil
zustimmen ${\rightarrow}$ stimmt ein Elternteil zu, der andere
verweigert aber, so ist  entweder jener Teil entscheidungsbefugt der
die gr\"o{\ss}ere Rolle in der Erziehung des Pat. spielt (im Zweifel
die Mutter) oder das Pflegschaftsgericht ist anzurufen

\item lehnen beide ab, obwohl es sich um eine notwendige HB handelt,
muss das Pflegschaftsgericht angerufen werden

\end{itemize}
Umfang der Aufkl\"arung:

\begin{itemize}
\item Grundsatz: Die Aufkl\"arung hat in jenem Rahmen zu erfolgen, der
einer ma{\ss}gerechten Figur unter Ber\"ucksichtigung der individuellen
Besonderheiten des Pat. und der Situation gerecht wird.

\item Detail:

\end{itemize}
\begin{itemize}
\item Umfang und Tragweite der HB

\item Folgen, Gefahren und Nebenwirkungen der HB

\item \"uber Umst\"ande die f\"ur verst\"andigen Pat. ins Gewicht fallen
w\"urden

\item umgekehrte Proportion zur Dringlichkeit der HB

\end{itemize}
Unterlassen der Aufkl\"arung:

Unterbleibt die Aufkl\"arung wird angenommen, dass dem Pat. die
Tragweite der Ma{\ss}nahmen nicht bewusst ist und er daher nicht in die
HB einwilligt. 

Besonderheit Revers:

\begin{itemize}
\item Was ist der Revers?

\end{itemize}
\begin{itemize}
\item Der Revers ist eine schriftliche Dokumentierung \"uber die Aussage
des Patient, eine vom Einsatz\-personal empfohlenen Ma{\ss}nahme
abzulehnen. Die Verweigerung kann auch vom zust\"andigen Vertreter des
Patienten abgegeben werden.

\end{itemize}
\begin{itemize}
\item Welche Voraussetzungen hat der Revers?

\end{itemize}
\begin{itemize}
\item  auszuf\"ullen nach umfassender Aufkl\"arung des Patienten

\item  Schriftlichkeit (Transportschein, Reversschein oder beigelegtes
Dokument)

\item  Unterschrift des Patienten (oder seines zust\"andigen Vertreters)

\end{itemize}
\begin{itemize}
\item Wann und warum ist der Revers notwendig?

\end{itemize}
\begin{itemize}
\item  wenn das Einsatzpersonal eine Ma{\ss}nahme f\"ur notwendig
erachtet, der Patienten  (oder sein zust\"andiger Vertreter) diese
Ma{\ss}nahme aber verweigert

\item  zur Einhaltung der gesetzlichen Dokumentationspflicht und zur
Abwendung von  klagsweise geltendmachbaren Anspr\"uchen des Patienten

\end{itemize}
\end{multicols}
Diese Punkte f\"uhrt zu folgenden Rechtsfolgen:

Rechtsfolgen

Aufgrund der unterlassenen Aufkl\"arung machte sich Sepp strafbar nach
folgenden Bestimmungen:

\begin{itemize}
\item[] Strafrecht:

\item Eigenm\"achtige Heilbehandlung ({\S} 110 StGB): 

\end{itemize}
\begin{itemize}
\item nur bei {\quotedblbase}Gefahr in Verzug{\textquotedblleft}
ben\"otigt es keine Aufkl\"arung bzw. Einwilligung und es d\"urfen alle
nach Ansicht der Einsatzkr\"afte direkt notwendigen Ma{\ss}nahmen
gesetzt werden

\item Gefahr in Verzug: unmittelbar drohende oder gegenw\"artige Gefahr,
die ein Abwarten nicht zul\"asst bzw. unmittelbares Handeln verlangt

\end{itemize}
\begin{itemize}
\item K\"orperverletzung ({\S}{\S} 83 ff StGB):

\end{itemize}
\begin{itemize}
\item keine Strafbarkeit, solange die HB medizinische indiziert und lege
artis durchgef\"uhrt wurde

\end{itemize}
\begin{itemize}
\item[] Zivilrecht:

\item Schadenersatz wegen K\"orperverletzung ({\S}{\S} 1293 ff ABGB):

\end{itemize}
\begin{itemize}
\item f\"ur das Zivilrecht ist die Einwilligung aufgrund einer
Aufkl\"arung notwendige Voraussetzung f\"ur die rechtfertigende HB
${\rightarrow}$ mangelt es an ihr, kann Schadenersatz verlangt werden

\end{itemize}
\paragraph[Fall 5. Patientenverf\"ugung]{Fall 5. Patientenverf\"ugung}
Sachverhalt

Nach Ablieferung der Fr. Uninformiert dr\"angt Alex seine Kollegen zu
einer schnellen Abfahrt, um endlich seinen Kaffee zu bekommen. Doch
macht ihm der Journaler schon wieder einen Strich durch die Rechnung,
als er das Team zu einem weiteren BO beruft. 

Dieser liegt in einem kleinen Pflegeheim, indem eine aufgeregte
Heimpflegerin auf Alex, Sepp und Bernd zust\"urmt und sie zu der 70
j\"ahrigen, regungslosen Pat., Fr. Finito, bringt. Gleichzeitig h\"alt
die Heimpflegerin den Sanit\"atern eine Patientenverf\"ugung vor, die
Fr. Finito vor \"uber 5 Jahren auf einer Serviette unterschrieben hat
und nach der {\quotedblbase}bei ausbleibenden Kreislauf keine
Wiederbelebungsma{\ss}nahmen gesetzt werden{\textquotedblleft}
d\"urfen. Deshalb hatte die Pflegerin es bis jetzt auch unterlassen mit
einer Reanimation zu beginnen. Da sie aber ebenso nicht wei{\ss}, ob
die Verf\"ugung bindend ist, rief sie die Rettung und setzte sich
wartend neben die Pat..

Rechtsfrage

Was ist eine Patientenverf\"ugung? Ab wann und wie weit ist sie bindend?

Rechtslage

\begin{multicols}{2}
Allgemeines zur Patientenverf\"ugung:

\begin{itemize}
\item Willenserkl\"arung einer einsicht-, urteils- und
\"au{\ss}erungsf\"ahigen Person

\item erst von Bedeutung, wenn der Pat. seinen Willen nicht mehr
ausdr\"ucken kann

\item nur Ablehnung einer bestimmten HB kann Inhalt der Verf\"ugung sein

\item unbeachtlicher Inhalt:

\end{itemize}
\begin{itemize}
\item Ausschluss einer Grundversorgung

\item med. nicht indizierte Behandlung

\end{itemize}
zwei Arten der Patientenverf\"ugung:

\begin{itemize}
\item[] 1) verbindliche Patientenverf\"ugung:

\end{itemize}
\begin{itemize}
\item schriftlich vor Notar, Rechtsanwalt oder rechtskundigen
Patientenvertreter

\item eigenh\"andig unterschrieben und datiert

\item aufgrund einer umfassenden Aufkl\"arungen durch einen Arzt

\item Inhalt muss eine konkrete HB ablehnen ${\rightarrow}$ allgemeine
oder widerspr\"uchliche Formulierungen f\"uhren zur Unwirksamkeit

\item Wirkung nur f\"ur 5 Jahre

\end{itemize}
\begin{itemize}
\item[] eine Verf\"ugung die alle Elemente enth\"alt ist absolut bindend
und darf nicht umgangen werden

2) beachtliche Patientenverf\"ugung:

\end{itemize}
\begin{itemize}
\item grunds\"atzlich: alle Verf\"ugungen, die nicht die obigen
Voraussetzungen erf\"ullen

\item Umfang der Beachtlichkeit: blo{\ss} relative Beachtlichkeit:

\end{itemize}
\begin{itemize}
\item Verf\"ugung soll in Behandlungsentscheidungen ber\"ucksichtigt
werden, aber

\item das Wohl des Pat. bleibt oberstes Gebot 

\item daher kann man sich bei besonderer Begr\"undung \"uber die
Verf\"ugung hinwegsetzen

\end{itemize}
\begin{itemize}
\item absolute Unwirksamkeit/Unbeachtlichkeit bei:

\end{itemize}
\begin{itemize}
\item zu allgemeine oder widerspr\"uchliche Formulierungen

\item rechtswidriger Inhalt

\item erkennbare Willensm\"angel des Pat. bei Abgabe der Verf\"ugung
(nicht ernst gemeint, Irrtum, List, Drohung oder T\"auschung)

\end{itemize}
Suche nach Patientenverf\"ugung - {\S} 12 PatVG:

Bei Notfallversorgungen muss nicht nach einer Verf\"ugung gesucht
werden, wenn die Suche das Leben oder die Gesundheit des Pat ernstlich
gef\"ahrden w\"urde. Daher wird f\"ur das Rettungswesen idR. keine
Verpflichtung zur Suche, sondern zur Behandlung bestehen.

\end{multicols}
Rechtsfolge

\begin{itemize}
\item Die Patientenverf\"ugung gilt mit \"uber 5 Jahren, der fehlenden
\"arztlichen Aufkl\"arung und dem mangelnden Beiziehen eines
Rechtsverst\"andigen nicht mehr als verbindlich und w\"are somit als
blo{\ss} beachtlich einzustufen. 

\item Doch aufgrund des zu allgemeinen Inhalts, der Beil\"aufigkeit der
Verf\"ugung (Serviette als Unterlage) und wegen des hohen Alters der
Verf\"ugung kann kein konkreter und aktueller Wille der Pat. erkannt
werden, was dazu f\"uhrt, dass das Wohl der Pat. das oberste Gebot
bleibt. Mit einer Reanimation ist zu beginnen (wenn \"uberhaupt noch
erfolgreich).

\end{itemize}
Teil D: Rechte und Pflichten des Sanit\"aters

\paragraph{Fall 6. Unterlassung der Notarztnachforderung}
Sachverhalt

Endlich am St\"utzpunkt angekommen, ist Alex auch schon der Erste der
sich zum Kaffeeautomaten begibt. Entkr\"aftet durch die einsetzende
{\quotedblbase}Hypokoffeinomie{\textquotedblleft} kramt er nach
Kleingeld und kaum als er es einwurfbereit in H\"anden h\"alt, t\"ont
die n\"achste Ausfahrt durch die Lautsprecher, mit dem Hinweis einer
dringenden Berufung. Alex geht aber in seiner typisch ruhigen Manie,
gl\"ucklich durch die enteral erlangte Sedierung, zum Fahrzeug, mit dem
Spruch auf den Lippen: {\quotedblbase}immerhin konn des
Wagl{\textquoteright} eh net ohne mie foahrn!{\textquotedblleft}

Nach nur kurzer Anfahrtszeit begegnet Alex, Sepp und Bernd im ersten
Stock eines kleinen Reihenhaus der 55 j\"ahrige, adip\"ose Hr. Schmerz
mit Schmerzen hinter der Brust, die in den linken Arm ausstrahlen und
ihm selbst in Ruhelage mit Atemnot und kalten Schwei{\ss}ausbr\"uchen
plagen. F\"ur Sepp ist sofort klar, dass es sich um einen MCI handelt
und m\"ochte schleunigst einen NEF nachfordern. Doch Alex meint
hingegen, dass er keine Lust h\"atte schon wieder auf einen
{\quotedblbase}Boder{\textquotedblleft} zu warten und der Pat. auch
selbst runter zum Fahrzeug gehen sollte, da er es im Kreuz h\"atte und
sowieso noch eine ca. 20 Min. lange Anfahrtszeit zur Krankenanstalt
durch den schon einsetzende Morgenstau bew\"altigen m\"usste.

Rechtsfrage

Besteht die Pflicht f\"ur Sanit\"ater zur Nachforderung eines NA?

Rechtslage

Unterlassene Hilfeleistung - {\S} 95 StGB

\begin{itemize}
\item jeder Durchschnittsmensch ist zur Durchf\"uhrung von ihm
zumutbaren Hilfeleistungen verpflichtet. Die Zumutbarkeit bestimmt sich
dabei nach dem individuellen Stand an Ausbildung, F\"ahigkeiten und
Kenntnissen. Aber schon zu den Grundlagen einer Ersten Hilfe, die von
jedem Menschen mit Besuch eines F\"uhrerscheinkurses erwartet werden
kann, geh\"ort das Absetzen eines Notrufes zur  Anforderung von
geeigneten Rettungskr\"aften.

\end{itemize}
{\S} 4 SanG

\begin{itemize}
\item ein jeder Sanit\"ater, unabh\"angig von seinem Ausbildungsstand,
ist zur Durchf\"uhrung von lebensrettenden Sofortma{\ss}nahmen
verpflichtet, was gleichzeitig beinhaltet, dass er bei Notwendigkeit
eine unverz\"ugliche Anforderung des Notarztes oder sonstigen Arztes zu
veranlassen hat.

\end{itemize}
Rechtsfolge

\begin{itemize}
\item W\"urde Alex die zweckm\"a{\ss}ige Nachforderung eines NA
unterlassen, macht er sich gleichzeitig nach strafrechtlichen ({\S} 95
StGB) und verwaltungsrechtlichen Bestimmungen ({\S} 9 SanG) strafbar.

\end{itemize}
\paragraph{Fall 7. Sanit\"aterkompetenzen}
\begin{multicols}{2}
Sachverhalt

Nach langem hin und her, k\"onnen Sepp und Bernd gemeinsam den Alex
\"uberreden doch einen NA nachzufordern und Sepp, der gerade frisch
gebackener NFS geworden ist, sieht endlich seine M\"oglichkeit
umfangreich am Hr. Schmerz zu werken. Dabei m\"ochte er neben der
Verabreichung von Nitrolingual-Spray auch gleich einen ven\"osen Zugang
legen, da er das sowieso als ehemaliger Fleischhauer besser als jeder
Arzt kann.

Rechtsfrage

Was sind die Kompetenzen der einzelnen Sanit\"aterstufen
(RS/NFS/NKA/NKV/NKI)? Darf man die Kompetenzgrenzen (bei Gefahr in
Verzug) \"uberschreiten?

Rechtslage

Kompetenzstufen gem. {\S}{\S} 9-12 SanG:

\begin{itemize}
\item RS - {\S} 9 SanG:

\end{itemize}
\begin{itemize}
\item Z 1: selbstst\"andige und eigenverantwortliche Versorgung und
Betreuung des Pat.

\item Z 2: \"Ubernahme und \"Ubergabe des Pat. im Rahmen eines
Transportes

\item Z 3: Hilfestellung bei Akutsituationen; Verabreichung von O2

\item Z 4: Durchf\"uhrung von lebensrettenden Sofortma{\ss}nahmen

\item Z 5: Durchf\"uhrung von Sondertransporten

\item Z 6: Blutzuckermessung

\end{itemize}
\begin{itemize}
\item NFS - {\S} 10 SanG:

\end{itemize}
\begin{itemize}
\item Z 1: T\"atigkeiten gem. {\S} 9

\item Z 2: Unterst\"utzung des Arztes; Betreuung von Transporten

\item Z 3: Verabreichung der Arzneimittelliste 1

\item Z 4: eigenverantwortliche Betreuung von berufsspezifischen
Medizinprodukten und Arzneimittel

\item Z 5: Mitarbeit in Forschung

\end{itemize}
\begin{itemize}
\item[] Arzneimittelliste 1: Wird festgelegt von dem leitenden Arzt der
Rettungsorganisation. Die Verabreichung kann durchgef\"uhrt werden ohne
dass ein NA informiert werden muss oder es ein solcher anordnet.

\item NKA - {\S} 11 Abs. 1 Z 1 SanG: zus\"atzlich zu den oben genannten
Kompetenzen darf der NKA jene Medikamenten der Arzneimittelliste 2
verabreichen, die ohne ven\"osen Zugang verabreicht werden k\"onnen.

\item NKV - {\S} 11 Abs. 1 Z 2 SanG: zus\"atzlich zu den oben genannten
Kompetenzen darf der NKV jene Medikamenten der Arzneimittelliste 2
verabreichen, die \"uber ven\"osen Zugang verabreicht werden k\"onnen.

Arzneimittelliste 2: Wird festgelegt vom leitenden Arzt der
Rettungsorganisation, und darf erst verabreicht werden:

\end{itemize}
\begin{itemize}
\item auf Anweisung durch anwesenden Arzt (jeder mit jus practicandi --
nicht unbedingt NA) oder

\item nach Verst\"andigung eines Notarztes.

\end{itemize}
\begin{itemize}
\item NKI - {\S} 12 SanG: Durchf\"uhrung einer endotrachealen Intubation
samt Beatmung, aber ohne Pr\"amedikation

\end{itemize}
\begin{itemize}
\item NKI-Kompetenz muss alle 2 Jahre rezertifiziert werden

\item Durchf\"uhrung der Intubation erst:

\end{itemize}
\begin{itemize}
\item auf Anweisung durch anwesenden Arzt (jeder mit jus practicandi%
%
 -- nicht unbedingt NA) oder

\item nach Verst\"andigung eines Notarztes 

\end{itemize}
\"Uberschreiten der Kompetenzstufen:

Die Durchf\"uhrung von Ma{\ss}nahmen au{\ss}erhalb des eigenen
Kompetenzstandes, f\"uhrt zu einer Verwaltungsstrafe nach SanG bzw.
einer gerichtlichen Strafe wegen K\"orperverletzung und
zivilrechtlichen SchEA. Selbst bei Wissen \"uber Gefahren, Risiken und
Durchf\"uhrungsart einer kompetenzm\"a{\ss}ig h\"oheren Ma{\ss}nahme,
gibt es keine Rechtfertigung durch {\S} 95 StGB (Unterlassene
Hilfeleistung).

Dies bedeutet weiters, dass einfache Ma{\ss}nahmen, die zwar das
Berufsgesetz nicht aufz\"ahlt, aber von jedem Laien (mitunter nach
kurzer Einweisung) durchgef\"uhrt werden d\"urfen, auch von
Sanit\"atern gesetzt werden d\"urfen bzw. iSd. {\S} 95 StGB gesetzt
werden m\"ussen.

Rechtsfolge

\begin{itemize}
\item Sepp darf keinen ven\"osen Zugang legen oder sogar Medikamente
dar\"uber verabreichen, selbst  wenn Gefahr in Verzug vorliegt. 

\end{itemize}
\end{multicols}
\paragraph{Fall 8. rechtswidrige Anordnung durch Notarzt}
Sachverhalt

Sepp konnte sich nun dazu \"uberreden lassen keinen Venflon zu legen und
m\"ochte wenigstens Nitro verabreichen. Gerade noch rechtzeitig bemerkt
er eine offene Packung von Viagra am Tisch liegen. Hr. Schmerz meint,
als er Sepp die Packung betrachten sieht, dass das
{\quotedblbase}Zeug{\textquotedblleft} nichts nutzt, da er vor zwei
Stunden eine Tablette eingenommen h\"atte und noch immer
{\quotedblbase}nichts steht{\textquotedblleft}. Kurze Zeit sp\"ater
trifft der NEF mit der Not\"arztin Natascha ein, die den Pat. flapsig
untersucht und eine sofortige Nitro-Verabreichung zur Bek\"ampfung des
hohen Blutdrucks (200/110 mm Hg) anordnet, obwohl sie von Sepp \"uber
die gefundene Viagra-Packung informiert wurde.

Rechtsfrage

Ist der Sanit\"ater verpflichtet eine rechtswidrige bzw. gef\"ahrliche
Anordnung des NA (oder eines anderen Vorgesetzten) zu folgen?

Rechtslage

Ablehnung von rechtswidrigen Anordnungen:

\begin{itemize}
\item Zivildiener: {\S} 22 Abs. 2 Zivildienstgesetz:

\end{itemize}
\begin{itemize}
\item[] {\quotedblbase}Er darf die Befolgung einer Weisung nur dann
ablehnen, wenn [{\dots}] die Befolgung gegen strafgesetzliche
Vorschriften versto{\ss}en w\"urde.{\textquotedblleft}

\end{itemize}
\begin{itemize}
\item Haupt- und Ehrenamtliche: 

\end{itemize}
\begin{itemize}
\item {\S} 4 Abs. 1 S. 2 SanG

\item {\S} 7 iVm {\S} 879 Abs. 1 ABGB

\item {\S} 22 Abs. 2 ZDG analog:

\end{itemize}
\begin{itemize}
\item[] Aufgrund der im SanG fehlenden konkreten Bestimmung eines
Ablehnungsrechtes von Sanit\"atern, kann ein solches nur kompliziert
gewonnen werden aus:

\end{itemize}
\begin{itemize}
\item allgemeinen Verpflichtung zur Wahrung des Patientenwohls nach {\S}
4 SanG

\item nat\"urlichen Rechtsgrunds\"atzen iSd. {\S} 7 ABGB und dem
Nichtigkeitsgebot bei Sittenwidrigkeit des {\S} 879 ABGB

\item analoger Anwendung des ZDG 

\end{itemize}
Rechtsfolge

\begin{itemize}
\item Die RTW-Mannschaft darf daher die f\"ur den Pat. sch\"adigende
Anordnung ablehnen. Natascha hingegen macht sich strafbar wegen
K\"orperverletzung und riskiert uU. zivilrechtliche SchEA
(${\rightarrow}$ Nitro-Gabe nicht lege artis).

\end{itemize}
\paragraph{Fall 9. Dokumentationspflicht}
Sachverhalt

Genervt von der Auseinandersetzung der Not\"arztin mit den Sanit\"atern,
beginnt Hr. Schmerz wild vor Schmerzen die Beteiligten anzubr\"ullen,
dass sie doch endlich ihre {\quotedblbase}gottverdammte
Arbeit{\textquotedblleft} tun sollen. Der ebenso gereizte Bernd belehrt
daraufhin Hr. Schmerz, dass er sich doch gef\"alligst wie ein
{\quotedblbase}richtiger Pat.{\textquotedblleft} verhalten soll und
sich erst r\"uhren darf, wenn es ihm gesagt wird. Diese Aussage wird
Bernd aber zum Verh\"angnis, da sich Hr. Schmerz Wochen sp\"ater beim
Arbeitgeber von Alex, Sepp und Bernd \"uber das
{\quotedblbase}beleidigende Verhalten{\textquotedblleft} des Bernd
beschwert. Sepp seinerseits ist \"uberzeugt, dass er sich korrekt
verhalten hat, doch kann er sich so recht an den alten Fall nicht mehr
erinnern.

Rechtsfrage

Besteht eine Dokumentationspflicht und wenn ja, in welchem Umfang? 

Rechtslage

{\S}~5 SanG: {\quotedblbase}Sanit\"ater haben [{\dots}] die von ihnen
gesetzten Ma{\ss}nahmen zu dokumentieren.{\textquotedblleft}

Vertragsrechtliche Relevanz:

Im \"osterreichischen Recht hat grunds\"atzlich der Kl\"ager den
Schaden, die Rechtswidrigkeit und das Verschulden des Beklagten zu
beweisen. Besteht aber zwischen Kl\"ager und Beklagten ein Vertrag, so
wird gem. {\S} 1298 AGBG das Verschulden des Beklagten angenommen bzw.
der Beklagte muss beweisen, dass ihn kein Verschulden trifft (=
Beweislastumkehr). 

Im Rettungswesen besteht zu jedem Pat. ein Behandlungsvertrag, der eine
solche Beweislastumkehr bewirkt. Um nun f\"ur den Beklagten die
M\"oglichkeit zu schaffen, sein Unverschulden zu beweisen, wurde die
Dokumentationspflicht eingef\"uhrt, mit der bei ausreichender
Dokumentation wieder der Kl\"ager das Verschulden des Beklagten zu
beweisen hat.

Inhalt der Dokumentation:

\begin{itemize}
\item gesetzte med. Ma{\ss}nahmen

\item situationsrelevante Umst\"ande (beim Pat.; am BO; w\"ahrend des
Transportes; {\dots})

\item Verweigerungen oder besondere Anmerkungen des Pat. (kein Gurt;
Marcoumarpflichtig; {\dots})

\item uU. relevante Inhalte der Aufkl\"arung

\item sonstige f\"ur den Sanit\"ater entscheidend wirkende Umst\"ande 

\end{itemize}
Grundsatz: Alles darf dokumentiert werden, es darf aber nicht Nichts
dokumentiert werden. 

Durchf\"uhrung der Dokumentation:

Die Dokumentation soll vorrangig am Transportschein selbst oder wenn
nicht mehr ausreichend auf einem beigelegten Blatt Papier erfolgen.
Verweigert der Patient medizinische oder sonst notwendige
Ma{\ss}nahmen, ist von ihm ein Revers zu unterschreiben, der dem
Transportschein beizulegen ist. Dabei hat der Sanit\"ater zus\"atzlich
immer (!) durch ein Gespr\"ach auf m\"ogliche Risiken und Gefahren der
Verweigerung hinzuweisen.

Des Weiteren ist eine zus\"atzliche eigene Dokumentation f\"ur jeden
Sanit\"ater empfehlenswert.

Rechtsfolge

\begin{itemize}
\item Bernd hat keine ausreichende Dokumentation durchgef\"uhrt und kann
sich somit nicht gegen die Beschuldigungen zur Wehr setzen.

\end{itemize}
\paragraph{Fall 10. Verschwiegenheits- und Auskunftspflicht}
\begin{multicols}{2}
Sachverhalt

Durch die Aufregung verschlechtert sich allm\"ahlich der Zustand des Hr.
Schmerz und das RTW-Team samt NEF-Mannschaft sind sich einig, dass ein
schleunigster Abtransport durchzuf\"uhren ist. Noch bevor sie aber den
Pat. umlagern k\"onnen, erscheint Fr. Aufgeregt, die Angeh\"orige des
Hr. Schmerz, und verlangt \"uber alles informiert zu werden.

Rechtsfrage

Was umfasst die Verschwiegenheitspflicht der Sanit\"ater und \"Arzte?
Besteht auch eine Auskunftspflicht derselben?

Rechtslage

Verschwiegenheitspflicht:

\begin{itemize}
\item strafrechtliche Norm: {\S} 121 StGB

\item verwaltungsrechtliche Norm: {\S} 6 SanG

\item beiden gleich: Geheimnisse, die man aufgrund seiner T\"atigkeit
als Arzt/Sanit\"ater anvertraut bekommen hat, d\"urfen nicht an
unbefugte Dritte weitergegeben werden

\end{itemize}
\begin{itemize}
\item Geheimnisse: alle Informationen, die die Person des Pat., seinen
Gesundheitszustand oder die Umst\"ande am Einsatzort betreffen

\item daher darf man mit Kollegen oder sonstigen Personen \"uber die
erlebten Geschehnisse reden, solange aus dem Gesagten nicht erkennbar
ist, um welche Person es sich bei dem Pat. handelt 

\item Bsp: Name des Pat., Versicherungsnummer, Wohnadresse, {\dots}

\end{itemize}
Wann darf man Geheimnisse weitergeben?

\begin{itemize}
\item Identit\"at des Pat. bleibt gewahrt

\item bei Entbindung von der Schweigepflicht durch den Pat.

\end{itemize}
Wann muss man Geheimnisse weitergeben? = Auskunftspflicht - {\S} 7 SanG

\begin{itemize}
\item wenn es zur Behandlung des Pat. dient

\item wenn es gesetzlichen vorgeschrieben ist (Bsp: im Gerichtsverfahren
als Zeuge)

\item zum Schutz von h\"oherwertigen Interessen der \"offentlichen
Gesundheitspflege oder Rechtspflege (${\rightarrow}$ zur Meldepflicht
siehe ad G)

\end{itemize}
Wem gegen\"uber muss man Auskunft erteilen?

\begin{itemize}
\item Patienten

\item gesetzlichen Vertreter des Pat.

\item anderen Personen der Gesundheitsberufe, die den Pat. betreuen,
behandeln oder pflegen bzw. an denen der Pat. im Zuge des Transportes
\"ubergeben wird

\item Sozialversicherung und Krankenanstalt

\item vom Pat. als auskunftsberechtigt benannte Personen

\end{itemize}
Rechtsfolge

\begin{itemize}
\item Fr. Aufgeregt z\"ahlt als blo{\ss}e Angeh\"orige nicht zum Kreis
der Auskunftsberechtigten, solange der Pat. die Weitergabe von
Informationen nicht gestattet. 

\end{itemize}
\end{multicols}
\paragraph[Fall 11. Fortbildungs- und Rezertifizierungspflicht]{Fall 11.
Fortbildungs- und Rezertifizierungspflicht}
\begin{multicols}{2}
Sachverhalt

Nachdem nun auch Fr. Aufgeregt beruhigt wurde, wird Hr. Schmerz auf den
Sessel umgelagert, um \"uber die enge Treppe den ersten Stock zu
verlassen. Doch handelt es sich beim Tragsessel um einen neuen
Striker-Sessel, der dem Bernd noch vollkommen neu ist, da es den Sessel
zum Zeitpunkt seiner RS-Ausbildung noch gar nicht gegeben hat. Aber
genau wegen dieser Neuheit hatte sein Arbeitgeber eine interne
Fortbildung \"uber die {\quotedblbase}Neuen
Medizinprodukte{\textquotedblleft} abgehalten, zu der Bernd aber nicht
gekommen ist, da er meinte sowieso schon alles zu wissen und zu keinen
Fortbildungen verpflichtet ist.

Rechtsfrage

Besteht eine Fortbildungs- oder sogar Re\-zertifizierungspflicht?

Rechtslage

Fortbildungspflicht nach {\S}~50 SanG:

Jeder Sanit\"ater ist dazu verpflichtet sich innerhalb von 2 Jahren ab
seinem Stichtag im Umfang von 16 Stunden fortzubilden. Als Stichtag
wird der Monatserste nach bestandener Rettungssanit\"aterausbildung
gewertet, wobei sich der Stichtag auch bei Erreichen von h\"oheren
Kompetenzniveaus nicht \"andert. Zu welchem Zeitpunkt die 16
Fortbildungsstunden innerhalb des Rahmens von 2 Jahren besucht werden,
ist dem Sanit\"ater \"uberlassen. 

Man kann also entweder regelm\"a{\ss}ig Fortbildungen besuchen oder auf
einmal die Stundenzahl ansammeln. Weiters muss ein Sanit\"ater keine
bestimmten Themen als Fortbildungen besuchen, sondern kann nach seinem
Interesse vorgehen.

Fortbildungspflicht nach {\S}~83 MPG:

Die im Rettungswesen verwendeten Materialien gelten, abgesehen von
Arzneimitteln, als Medizinprodukte und d\"urfen daher iSd. MPG nur von
eingeschulten Personal angewendet werden. Diese Einschulungen m\"ussen
nicht nur alle Medizinprodukte der Rettungsorganisation beinhalten,
sondern auch von allen Sanit\"atern besucht werden. Daher kann ein
Sanit\"ater in Bezug auf Stundenanzahl oder Fortbildungsthemen nicht
frei w\"ahlen. Die Fortbildungsstunden f\"ur das MPG k\"onnen aber als
Fortbildungsstunden iSd. SanG gewertet werden. Zus\"atzlich ist zu
sagen, dass eine Einschulung der meisten Medizinprodukte schon durch
die Standardausbildung zum Rettungs\-sani\-t\"ater abgedeckt ist. F\"ur
die anderen, insbesondere h\"ohere technische Produkte, muss jede
Org\-anisation separate Fortbildungen anbieten.

Rezertifizierungspflicht gem.{\S}~51 SanG:

Innerhalb des 2 Jahre Rahmens der Fortbildungspflicht muss sich ein
Sanit\"ater, unabh\"angig von seinem Ausbildungsstand, pr\"ufen lassen
in der Herz-Lungen-Wiederbelebung mit halbautomatischen Defibrillator.
Der Rezertifizierungszeitpunkt innerhalb des 2 Jahre Rahmens kann vom
Sanit\"ater frei gew\"ahlt werden.

Sonderfall NKI: Ein Notfallsanit\"ater mit Kompetenzen der Intubation
muss sich zus\"atzlich innerhalb seines 2 Jahre Rahmens einmal \"uber
die Intubation pr\"ufen lassen.

Rechtsfolge

\begin{itemize}
\item Die f\"ur den vorliegenden Fall entscheidende Fortbildung gem. MPG
wurde von Bernd nicht besucht, was ihm nicht nur eine
Verwaltungsstrafe, sondern auch eine gerichtliche Strafe und
zivilrechtlichen SchEA einbringen kann. Denn sollte es im Umgang mit
dem Medizinprodukt zu K\"orperverletzungen kommen, wird aufgrund der
mangelnden Materialkenntnis die Behandlung als nicht-lege-artis
gewertet. Um einer solchen Strafe zu entgehen gen\"ugt es, dass Bernd
noch vor Ort, aber vor Anwendung, durch eine {\quotedblbase}geeignete
Person{\textquotedblleft} eingeschult wird. 

\item Als {\quotedblbase}geeignete Person{\textquotedblleft} gelten
vorrangig Medizinproduktebeauftragter,  Lehrsanit\"ater oder
Praxisanleiter sowie jene, die von der Rettungsorganisation als
geeignet genannt werden.

\item Da Sepp Praxisanleiter ist, darf er vorab Bernd einweisen.

\end{itemize}
\end{multicols}
Teil E: Schadenersatz

\paragraph{Fall 12. Verletzung des Pat. durch Sanit\"ater:}
Sachverhalt

Nachdem nun Sepp unseren Bernd kurzum erkl\"art hat, wie er mit den
Striker-Sessel umzugehen hat unterl\"asst er trotzdem das Anschnallen
des Pat., da er meint, dass {\quotedblbase}eh kana d\"as
mocht{\textquotedblleft}. Leider ist Hr. Schmerz kein Fliegengewicht
und der Treppenaufgang ist derart verwinkelt, dass es alle M\"uhe und
Not bedarf den Tragsessel ruhig zu halten. 

Fast am Ende des Abgangs angekommen, flitzt auf einmal die
aufgeschreckte Katze des Hr. Schmerz an den beiden
{\quotedblbase}Tr\"agern{\textquotedblleft} vorbei, wodurch der
vorangehende Bernd seine Konzentration verliert, ausrutscht und samt
dem Pat. seitlich nach vorne auf einen Kasten und die darunter sich
versteckende Katze f\"allt. Bernd kann sich beim Sturz wenigsten selber
abst\"utzen und bleibt (auch durch die abfedernde Wirkung der Katze)
unverletzt. Doch Hr. Schmerz prallt mit Sch\"adel und Brustkorb
ungesch\"utzt gegen den Kasten und wird unter dem hinter ihm
nachkommenden Tragsessel begraben. Er zieht sich neben einer RQW auf
der Stirn auch noch mehrere gebrochene Rippen zu. 

Rechtsfrage - Schadenersatz wegen K\"orperverletzung

Wann wird man schadenersatzpflichtig? Was umfasst der Schadenersatz? Wer
wird zur Haftung gezogen?

Rechtslage

Rechtliche Grundlage eines Schadenersatzanspruches (= SchEA):

\begin{itemize}
\item allgemein: {\S}{\S} 1293 ff ABGB

\item konkret f\"ur K\"orperverletzungen: {\S}{\S} 1325-1327 ABGB

\end{itemize}
Voraussetzungen der Schadenersatzpflicht:

\begin{itemize}
\item[] 1. Schaden = Nachteil am K\"orper, an der Gesundheit, im
Verm\"ogen

2. Kausalit\"at = bei Wegdenken der Sch\"adigerhandlung, bleibt auch der
Erfolg aus

3. Rechtswidrigkeit = objektiver Versto{\ss} gegen Gesetze oder
Vertragspflichten

4. Schuld = subjektiver Versto{\ss} gegen jene Sorgfaltspflichten, die
man einhalten h\"atte k\"onnen

\end{itemize}
Umfang des Schadenersatzes bei K\"orperverletzungen:

\begin{itemize}
\item Heilungskosten

\item Verdienstentgang

\item Schmerzengeld

\end{itemize}
Wer haftet wie:

Sch\"adigt der Sanit\"ater im Zuge seines Dienstes einen Pat. oder
dessen Angeh\"origen, kann der Gesch\"adigte nach {\S} 1313a ABGB
direkt gegen die Rettungsorganisation seine SchEA geltend machen. Denn
ein Sanit\"ater gilt in dieser Beziehung als Erf\"ullungsgehilfe der
Rettungsorganisation, wodurch er als verl\"angerter Arm der
Organisation deren vertragliche Pflichten (sorgf\"altige Behandlung)
erf\"ullt. Daher kann der Gesch\"adigte direkt gegen die Organisation
alle Anspr\"uche wie gegen den Sanit\"ater selber durchsetzen.

Rechtsfolge

\begin{itemize}
\item Im vorliegenden Fall ist der Schaden durch RQW und Frakturen neben
der Kausalit\"at gegeben. Weiters verst\"o{\ss}t es gegen die
gesetzlichen und vertraglichen Pflichten den Pat. ohne anzuschnallen zu
transportieren, weshalb eine Rechtswidrigkeit bejaht werden kann. 

\item Auch h\"atten sich unsere Sanit\"ater anders, n\"amlich den Pat.
anzuschnallen, verhalten k\"onnen, wodurch ebenso die Schuld gegeben
ist. Daher kann Hr. Schmerz SchEA gegen die Sanit\"ater oder,
wirtschaftlich vern\"unftiger, gegen die Rettungsorganisation geltend
machen.

\end{itemize}
\paragraph{Fall 13. Sachbesch\"adigung}
Sachverhalt

Leider sind die Verletzungen am Pat. nicht die einzigen Sch\"aden die
entstanden sind, denn der umgeschmissene Kasten war eine
Biedermeier-Rarit\"at mit einem Wert von stolzen 10.000 Euro. Auch die
Katze verwirkte als Auffangkissen f\"ur Bernd ihr Leben und hatte neben
dem Sachwert von 350 Euro auch einen besonderen emotionalen Wert f\"ur
Hr. Schmerz.

Rechtsfrage: Schadenersatz wegen Sachbesch\"adigung

Wann wird man schadenersatzpflichtig? Was umfasst der Schadenersatz? Wer
wird zur Haftung gezogen?

Rechtslage

Rechtliche Grundlage eines Schadenersatzanspruches:

\begin{itemize}
\item allgemein: {\S}{\S} 1293 ff ABGB

\item konkret f\"ur Sachsch\"aden: {\S} 1331 ABGB

\end{itemize}
Voraussetzungen der Schadenersatzpflicht:

\begin{itemize}
\item[] gleiche Voraussetzungen wie bei K\"orperverletzungen

\end{itemize}
Umfang des Schadenersatzes bei Sachbesch\"adigung:

\begin{itemize}
\item realer Schaden (= tats\"achlich eingetretener Schaden)

\item entgangener Gewinn (= in Zukunft drohender Schaden)

\end{itemize}
13.4. Rechtsfolge

\begin{itemize}
\item gleiche L\"osung wie unter ad 12.

\end{itemize}
Teil F: Behandlungspflicht

\paragraph{Fall 14. Behandlungspflicht von Krankenanstalten, \"Arzten,
Sanit\"atern und Durchschnittsb\"urgern:}
Sachverhalt

Vor Schmerzen und Trauer, um seine geliebte Katze
{\quotedblbase}Muschi{\textquotedblleft}, beschimpft und verflucht Hr.
Schmerz die Anwesenden mit wilden Gestikulierungen. Als er dann
wenigstens von Alex und Sepp vorbildlich versorgt wird, beruhigt sich
Hr. Schmerz und l\"asst sich sogar m\"uhelos ins Auto transportieren. 

Nach ausreichender Versorgung des Hr. Schmerz begibt sich der RTW auf
die Reise zur KA. Doch ist leider das einzig abbuchbare Bett, welches
neben einer Internen- auch eine Unfall-Ambulanz aufzuweisen hat, im
Wilhelminenspital, welches durch den gro{\ss}en Morgenstau gute 15 Min.
entfernt ist. Schon nach kurzer Fahrzeit klagt Hr. Schmerz \"uber
st\"arker werdende Atemprobleme und ein rasselndes, brodelndes
Atemger\"ausch t\"ont durch den gesamten Patientenraum. Es hat sich ein
kardiales Lungen\"odem gebildet und 

Natascha f\"uhlt sich, trotz Verabreichung von Diuretika, nicht mehr in
der Lage den Pat. w\"ahrend der Fahrt ausreichend zu behandeln. Alex
soll so schnell wie m\"oglich das n\"achste Spital anfahren. Durch eine
kurze Abweichung der Route, gelangt das Team zum AKH und prescht mit
dem Pat. zur nicht vorinformierten Aufnahme. 

Der dortige Aufnahmearzt Dr. Wasbringst hat schon einen 22 Stunden
Dienst hinter sich und ist alles andere als erfreut \"uber das
{\quotedblbase}Mitbringsel{\textquotedblleft} des Teams. Wutentbrannt
br\"ullt er Natascha an, was sie sich eigentliche erlaube trotz
gesperrten Bettenkontingent seine Station anzufahren und verweigert ihr
die Aufnahme des Hr. Schmerz.

Rechtsfrage

Besteht eine Behandlungspflicht f\"ur Krankenanstalten? Besteht eine
Behandlungspflicht f\"ur medizinische Berufe? Besteht eine
ersthelferische Versorgungspflicht f\"ur den Durchschnittsb\"urger?

Rechtslage

\begin{itemize}
\item Behandlungspflicht f\"ur Krankenanstalten (= KA) - {\S}{\S} 22 und
40 KAKuG:

\end{itemize}
\begin{itemize}
\item Jede KA, ob \"offentlich oder privat, ist verpflichtet einen
behandlungspflichtigen Pat.  aufzunehmen und zumindest \"arztliche
Grundbehandlung zukommen zu lassen. 

\item Das in Wien geltende Bettenkontingentsystem, mit seinen
Zuweisungen zu bestimmten KA{\textquoteright}s, ist trotzdem
zul\"assig, da sich der Pat. w\"ahrend des Transportes schon in einem
pr\"aklinischen Umfeld (= Einsatzwagen mit Sanit\"atern und
Medizinprodukten) befindet. Sollte sich aber der Zustand des Pat.
w\"ahrend der Fahrt verschlechtern, so darf das n\"achstgelegene Spital
angefahren werden.  Erst wenn der Pat. wieder transportf\"ahig ist,
darf er mit geeigneten Krankentransportmitteln in eine andere KA
weitergeschickt werden.

\end{itemize}
\begin{itemize}
\item \"arztliche Behandlungspflicht - {\S} 48 \"ArzteG:

\end{itemize}
\begin{itemize}
\item \"Arzte, unabh\"angig ihrer Fachausbildung, sind dazu verpflichtet
hilfsbed\"urftigen Pat. \"arztliche Grundversorgung zukommen zu lassen
und wenn n\"otig weitere Einsatzkr\"afte nachzufordern oder den Pat. in
weitere medizinische Behandlung zu schicken.

\end{itemize}
\begin{itemize}
\item Behandlungspflicht f\"ur Sanit\"ater - {\S} 4 iVm {\S} 8 iVm {\S}
9 Abs. 1 Z 1 SanG:

\end{itemize}
\begin{itemize}
\item Jeder t\"atigkeits- oder berufsberechtigter Sanit\"ater,
unabh\"angig von seinem Kompetenzstand, ist zur selbstst\"andigen und
eigenverantwortlichen Versorgung einer hilsbed\"urftigen Person
verpflichtet. Dabei muss er, auch au{\ss}erhalb des Dienstes, zumindest
qualifizierte Erste Hilfe leisten und wenn n\"otig einen Arzt (jeder
mit ius practicandi) nachfordern.

\end{itemize}
\begin{itemize}
\item ersthelferische Versorgungspflicht f\"ur Durchschnittsb\"urger -
{\S} 95 StGB:

\end{itemize}
\begin{itemize}
\item Jeder B\"urger (und nicht mehr t\"atigkeitsberechtigter
Sanit\"ater) ist zur Setzung von Hilfsleistungen verpflicht, die seinem
Wissens- und Kenntnisstand entsprechen. Daher ist jeder mit einem
F\"uhrerschein-Erste-Hilfe-Kurs zur Setzung von zumutbaren Erste-Hilfe
Ma{\ss}nahmen verpflichtet.

\end{itemize}
\begin{itemize}
\item Unzumutbarkeit der Behandlung - {\S} 95 Abs. 2 StGB:

\end{itemize}
\begin{itemize}
\item F\"ur alle Berufe oder Durchschnittsb\"urger ist die Behandlung
dann nicht zumutbar, wenn:

\end{itemize}
\begin{itemize}
\item man das eigene Leben oder die Gesundheit gef\"ahrden w\"urde oder

\item h\"oherwertige Interessen verletzen m\"usste.

\end{itemize}
Rechtsfolge

\begin{itemize}
\item Auch wenn das Bettenkontingent der KA aufgebraucht ist, muss Dr.
Wasbringst den Pat. aufnehmen und zumindest soweit stabilisieren, dass
er ohne Gefahren in das eigentliche Zielspital weitertransportiert
werden kann.

\end{itemize}
Teil G: Meldepflichtenrecht

\paragraph{Fall 15. Ansteckende Krankheiten}
Sachverhalt

Als nun Hr. Schmerz nach langem hin und her auf der Internen aufgenommen
wurde, erholen sich Alex, Sepp und Bernd bei einem kurzen Plausch mit
der NEF-Mannschaft. Die Ruhe h\"alt aber nicht lang an, als auch schon
der n\"achste Einsatz durchgegeben wird, der beide Teams, RTW und NEF,
zum selben BO schickt. 

Dort angekommen finden sie in einer Lacke von Blut die 25 j\"ahrige Fr.
Hilflos liegen, die sich selber (querverlaufend) die Pulsadern
aufgeschnitten hat. Als Sepp sie zu Versorgung beginnt, erz\"ahlt Fr.
Hilflos, dass ihr Leben {\quotedblbase}sowieso keinen Sinn mehr
hat{\textquotedblleft}, da sie ihr Freund mit HIV angesteckt hat.

Rechtsfrage

Bei welchen Erkrankungen besteht eine Meldepflicht? In welchem
Ausma{\ss} besteht eine solche Meldepflicht?

Rechtslage

Meldepflichtige Krankheiten:

\begin{itemize}
\item grunds\"atzlich:

\end{itemize}
\begin{itemize}
\item leicht \"ubertragbare Krankheiten (= KH)

\item mit Gefahr von schweren Gesundheits- oder Todesfolgen

\item Meldepflicht nur wenn es Gesetz oder Verordnung anordnen

\end{itemize}
\begin{itemize}
\item nach Epidemiegesetz: 

\end{itemize}
\begin{itemize}
\item bei blo{\ss}en Verdacht oder sicherer Diagnose einer im Gesetz
genannten Erkrankung.

\end{itemize}
\begin{itemize}
\item nach Tuberkulosegesetz:

\end{itemize}
\begin{itemize}
\item bei blo{\ss}en Verdacht auf eine Erkrankung die mit Sicherheit
oder Wahrscheinlichkeit durch Tuberkulosebakterien verursacht wurde

\end{itemize}
\begin{itemize}
\item nach Geschlechtskrankheitengesetz:

\end{itemize}
\begin{itemize}
\item nur bei Bef\"urchtung, dass sich die KH weiterverbreiten w\"urde
oder sich der Pat. einer HB entziehen w\"urde.

\end{itemize}
\begin{itemize}
\item nach AIDS-Gesetz:

\end{itemize}
\begin{itemize}
\item Anonymisierte Meldepflicht erst ab: manifester Erkrankung mit AIDS
(= bei sichtbaren Krankheitszeichen  nicht schon bei HIV) oder
Todesfall mit erkennbaren AIDS typischen KHs

\item es gibt keine Untersuchungs- oder Behandlungspflicht daher darf
kein beh\"ordlicher Zwang ausge\"ubt werden

\end{itemize}
\begin{itemize}
\item nach Verordnungen des Gesundheitsministerium:

\end{itemize}
\begin{itemize}
\item da sich st\"andig neue \"ubertragbare KH entwickeln, gibt das
Gesundheitsministerium in unregelm\"a{\ss}igen Abst\"anden Verordnungen
heraus, die eine Meldepflicht beinhalten.

\end{itemize}
Wer ist zur Meldung verpflichtet?

Nicht jeder Durchschnittsb\"urger ist zur Meldung verpflichtet, sondern
nur besondere Personengruppen worunter aber alle Gesundheitsberufe
fallen.

An wen ist zu melden?

\begin{itemize}
\item bei AIDS: Gesundheitsministerium

\item ansonsten: Bezirksverwaltungsbeh\"orde bzw. in Wien an MA 15

\end{itemize}
15.4. Rechtsfolge

\begin{itemize}
\item Nun ist zwar unser RTW-Team nicht zur Meldung von AIDS an eine
Beh\"orde verpflichtet, doch sehr wohl zur Meldung an die
weiterbehandelnden Gesundheitsberufe aufgrund der schon erw\"ahnten
Auskunftspflicht gem. {\S} 7 SanG.

\end{itemize}
\paragraph{Fall 16. Strafbare Handlungen}
Sachverhalt

Die Schnittwunden am Handgelenk sind nicht die einzigen Verletzungen der
Fr. Hilflos Mehrere H\"amatome \"uberziehen den K\"orper und ein
gro{\ss}es Feilchen ziert das rechte Auge der jungen Pat.. Als Natascha
auf diese weiteren Verletzungen eingeht, beginnt Fr. Hilflos weinend zu
berichten, dass ihr Freund sie \"ofters geschlagen hat und er ihr mit
dem Tod gedroht hat, falls sie die Polizei verst\"andigen w\"urde.

Rechtsfrage

Besteht eine Pflicht zur Anzeige von strafbaren Handlungen? Wenn ja,
f\"ur wen besteht diese Anzeigepflicht?

Rechtslage

Anzeigepflicht

Nur bestimmte Personengruppen sind dazu verpflichtet bei Kenntnis \"uber
strafbare Handlungen dies der Polizei oder Staatsanwaltschaft per
Anzeige zu melden.

\begin{itemize}
\item \"arztliche Anzeigepflicht - {\S} 54 Abs. 4 und 5 \"ArzteG:

nur bei Vorliegen folgender Delikte (Verdacht reicht aus) und
unabh\"angig vom Willen des Betroffenen:

\end{itemize}
\begin{itemize}
\item T\"otung bzw. schwere K\"orperverletzung

\item  Misshandelung, Qu\"alung, Vernachl\"assigung oder sexueller
Missbrauch von:

\end{itemize}
\begin{itemize}
\item Minderj\"ahrigen oder

\item Vollj\"ahrigen, die nicht alleine in der Lage sind ihre Interessen
wahrzunehmen

\end{itemize}
\begin{itemize}
\item polizeiliche Anzeigepflicht - {\S} 24 StPO

f\"ur alle Delikte

\end{itemize}
Anzeigerecht - {\S} 86 Abs. 1 StPO:

Jeder Durchschnittsb\"urger und auch sonstige Gesundheitsberufe sind
nicht verpflichtet zur Anzeige von Strafdelikten. Sie haben aber das
Recht eine strafbare Handlung anzuzeigen. 

Trotz der fehlenden Verpflichtung f\"ur den Sanit\"ater sollte bei jedem
Einsatz, bei dem der Verdacht einer strafbaren Handlung vorliegt,
zumindest die Polizei nachgefordert werden. Die Nachalarmierung kann
man in weiter Analogie mit der aus {\S} 4 SanG entstammenden Pflicht
f\"ur den Sanit\"ater {\quotedblbase}das Wohl der Patienten [{\dots}]
zu wahren{\textquotedblleft} begr\"unden.

Rechtsfolge

\begin{itemize}
\item Natascha muss aufgrund ihrer gesetzlichen Anzeigepflicht die
Polizei nachfordern, die zur Einbringung einer Anzeige den Sachverhalt
aufnehmen wird. Dies muss sogar dann geschehen, wenn die Pat. selber
keine Anzeige einbringen m\"ochte.

\end{itemize}
Teil H: Stra{\ss}enordnung

\paragraph{Fall 17. Blaulichtverwendung}
Sachverhalt

Nach Durchf\"uhrung der Behandlung und Anzeige gegen den aggressiven
Freund, begibt sich das RTW-Team ersch\"opft zum St\"utzpunkt zur\"uck.
Hungrig und m\"ude qu\"alt sich die RTW-Mannschaft durch das
allt\"agliche Wiener Verkehrschaos, als Sepp auf einmal meint:
{\quotedblbase} Geh, schalt doch das H\"orndl und die Blitzer ei, damit
dee uns endlie Platz machen und wir einr\"ucken
k\"annan!{\textquotedblleft}

Rechtsfrage

Wann darf das Blaulicht samt Folgetonhorn verwendet werden? Wie darf
bzw. muss man sich bei einer Blaulichtfahrt verhalten?

Rechtslage

Rechtsquellen:

{\S} 26 StVO und {\S} 107 KFG

Verwendung von Blaulicht und Folgetonhorn:

\begin{itemize}
\item grunds\"atzlich bei: Gefahr in Verzug

\item im n\"aheren bei:

\end{itemize}
\begin{itemize}
\item Fahrt zum BO

\item Haltezeit am BO

\item Fahrt zum AO

\end{itemize}
Verhalten bei Blaulichtfahrt:

\begin{itemize}
\item Verkehrsgebote und --beschr\"ankungen gelten nicht

\item Geschwindigkeitsh\"ochstgrenzen gelten nicht

\item Ampeln gelten als Stopptafeln

\item absoluter Vorrang vor Nicht-Einsatzfahrzeugen

\item es d\"urfen weder Personen noch Sachen gef\"ahrdet oder
besch\"adigt werden

\item es darf kein Nicht-Einsatzfahrzeug dem Rettungswagen unmittelbar
nachfahren

\item Vertrauensgrundsatz ist ung\"ultig (${\rightarrow}$ man darf bei
Blaulichtfahrten nicht mehr darauf vertrauen, dass sich die anderen
Verkehrsteilnehmer richtig verhalten daher ist defensives und
voraussehendes Fahren verpflichtend)

\end{itemize}
Vorrang der Einsatzfahrzeuge untereinander:

\begin{itemize}
\item[] 1. Rettungsfahrzeuge

2. Feuerwehrfahrzeuge

3. Polizeifahrzeuge

4. sonstige Einsatzfahrzeuge

\end{itemize}
Missbrauch von Blaulicht und Folgetonhorn:

Wer Blaulicht oder Folgetonhorn in Situationen verwendet, bei denen
nicht Gefahr in Verzug vorliegt, macht sich gerichtlich strafbar nach
{\S} 1 des {\quotedblbase}Bundesgesetzes gegen den Missbrauch von
Notzeichen{\textquotedblleft}.

Gebot bei Blaulichtfahrten:

Nur wer noch gr\"o{\ss}ere Sorgfalt und Vorsicht wahrt als bei normaler
Fahrt, verh\"alt sich richtig! Beinhaltet mitunter eine geringere
Fahrgeschwindigkeit als die h\"ochst zul\"assige. 

Exkurs: Anschnallpflicht - {\S} 106 KFG:

Der Fahrer eines jeden Fahrzeuges ist dazu verpflichtet zu achten, dass
er und seine Passagiere angeschnallt sind. Es steht aber jedem
Passagier frei zu sagen, dass er unangeschnallt mitfahren m\"ochte,
wodurch eine Schuld beim Fahrer entf\"allt. F\"ur Kranken- und
Rettungstransport besteht zus\"atzlich gem. {\S} 106 Abs. 3 Z 3 KFG die
Ausnahme, dass ein Passagier (Personal wie Pat.) unangeschnallt
mitfahren darf, wenn es ansonsten {\quotedblbase}mit dem Zweck der
Fahrt unvereinbar ist{\textquotedblleft}. Trotzdem sollte der Pat.
aufgefordert werden sich anzuschnallen bzw. sollte vom Sanit\"ater
angeschnallt werden. Verweigert der Pat., muss diese Verweigerung
dokumentiert werden (am besten mittels Revers, den der Pat.
unterschreibt).

Wird man nun im Zuge eines (sogar unverschuldeten) Verkehrsunfalls
verletzt und war man dabei nicht angeschnallt, verliert man einen Teil
des SchEA, weil das Fahren ohne angelegten Gurt als Mitverschulden
gewertet wird.

Rechtsfolge

\begin{itemize}
\item W\"urde unser RTW-Team tats\"achlich Blaulicht und Folgetonhorn
einschalten nur um schneller einr\"ucken zu k\"onnen, w\"urden sie sich
f\"ur die falsche Blaulichtfahrt gerichtlich strafbar machen. Bei einem
auch  unverschuldeten Unfall m\"ussten unsere drei Sanit\"ater die
volle Haftung tragen.

\end{itemize}
Der weitere Tag verl\"auft ruhig mit den \"ublichen Patienten.

\subsection{Abk\"urzungsverzeichnis}
\begin{multicols}{3}
ABGB:  Allgemeines B\"urgerliches Gesetzbuch

AKH:  Allgemeines Krankenhaus

AIDS:  acquired immune deficiency syndrome

AO:  Abgabeort

AR:  Arbeitsrecht

BO:  Berufungsort

DNHG:  Dienstnehmer\-haft\-pflicht\-gesetz

EPH-Gestose: 
\"Odem-Protein\-urie-Hypertonie-Schwanger\-schafts\-erkrankung

HB:  Heilbehandlung

HIV:  human immundeficiency virus

KA:  Krankenanstalt

KAKuG:  Krankenanstalt und Kur\-anstaltengesetz

KFG:  Kraftfahrgesetz

KH:  Krankheit

KTW:  Krankentransportwagen

MA:  Magistratsabteilung

MC:  Myocardinfarkt = Herz\-infarkt

MONA:  Morphin + O2 + Nitro\-lingual-Spray + Aspirin/Aspisol
medikament\"ose Herzinfarktbehandlung

MPG:  Medizinproduktegesetz

NA:  Notarzt

NaCl:  Natrium-Chlorid

NEF:  Notarzteinsatzfahrzeug

NFS:  Notfallsanit\"ater

NKA:  Notfallsanit\"ater mit Kompetenzen der Arzneimittel

NKI:  Notfallsanit\"ater mit Kompetenzen der Intubation

NKV:  Notfallsanit\"ater mit Kompetenzen des ven\"osen Zugangs

OGH:  Oberster Gerichtshof

OSR:  Oberster Sanit\"atsrat

Pat.:  Patient

PatVG:  Patienten\-ver\-f\"ugungs\-gesetz

RF:  Rechtsfolge

RFr:  Rechtsfrage

Rl:  Rechtslage

RQW:  Rissquetschwunde

RS:  Rettungssanit\"ater

RTW:  Rettungstransportwagen

SanG:  Sanit\"atergesetz

SchE:  Schadenersatz

SchEA:  Schadenersatzanspruch

StGB:  Strafgesetzbuch

StGG:  Staatsgrundgesetz

StPO:  Strafprozessordnung

StrR:  Strafrecht

StVO:  Stra{\ss}enverkehrsordnung

SV:  Sachverhalt

TB:  Tatbestand

UbG:  Unterbringungsgesetz

VfGH:  Verfassungsgerichtshof

VfR:  Verfassungsrecht

VwGH:  Verwaltungsgerichtshof

VwR:  Verwaltungsrecht

ZDG:  Zivildienstgesetz

ZR:  Zivilrecht

\end{multicols}
\section{[RS I] Erste Hilfe und erweiterte Erste Hilfe}
Autor fehlt

\begin{itemize}
\item \"Uberpr\"ufung der Atemt\"atigkeit (bei  \"uberstreckter HWS)
durch

\end{itemize}
\begin{itemize}
\item SEHEN 

\item H\"OREN

\item F\"UHLEN

\item F\"ur mindestens 10  Sekunden!! 

\end{itemize}
\paragraph{Freimachen der Atemwege}
\begin{itemize}
\item Freimachen der Atemwege

\end{itemize}
\begin{itemize}
\item \"Offnen des Mundes -{\textgreater} ESMARCH-Handgriff

\item Inspektion

\item Ausr\"aumen -{\textgreater}Fremdk\"orper/Erbrochenes durch
Seitw\"artsdrehen des Kopfes

\item (falls n\"otig) Absaugen

\item einzuf\"uhrende Katheterl\"ange: Distanz Ohrl\"appchen/Mundwinkel

\item \"Uberstrecken der HWS

\item Kinn/Stirn-Griff

\item Zungengrund wird angehoben -{\textgreater} Atemwege  sind frei!

\end{itemize}
\subsection{AED}
\section{[RS IX] Ger\"atelehre und Sanit\"atstechnik}
Thomas Hochreiter [R?]; BEMAW, Sebastian Gabriel

Derzeit in Vorbereitung. 

Kapitelliste nicht vollst\"andig

Blutdruckmessung

Grundlagen

Fallstricke

Die Blutdruckmanschette soll nach M\"oglichkeit nicht auf der Seite

\begin{itemize}
\item eines Shuntarmes bei Dialysepatienten

\end{itemize}
\begin{itemize}
\item Ein Dialyse-Shunt ist ein k\"unstlicher Kurzschluss zwischen einer
Armarterie und einer Armvene. Dadurch wird die Vene erweitert und sie
pulsiert, mann kann viel mehr Durchfluss f\"ur die Dialyse erreichen.
Die Nahtstelle zwischen Vene und Arterie ist jedoch sehr sensibel, wenn
gestaut wird kann sie aufplatzen ${\rightarrow}$ 

\item ${\rightarrow}$ Gefahr einer massiven Shuntblutung

\end{itemize}
\begin{itemize}
\item eines Arm\"odems

\item einer Brust-OP mit Lymphknoten-Entfernung (\"Odemgefahr)

\item einer Halbseitenschw\"ache oder -l\"ahmung 

\end{itemize}
angelegt werden. 

Shunt-Arme sind tabu f\"ur die Blutdruckmessung!

Bitten Sie im Praktikum am KTW Dialysepatienten h\"oflich Ihnen ihren
Shunt-Arm zu zeigen!

Ma{\ss}nahmen

\begin{itemize}
\item Funktion des Blutdruckmessger\"ates und des Stethoskops pr\"ufen

\end{itemize}
\begin{itemize}
\item Kontrolle der Dichte des Ventils, der Verbindungsschl\"auche und
der Manschette (Ventil schlie{\ss}en, Manschette etwas aufpumpen und
Zusammenpressen, Manschette ganz entl\"uften)

\item Kontrolle der G\"angigkeit des Ventils (es soll sich leicht auf
und zuschrauben lassen)

\item Kontrolle des Manometers auf korrekte Eichung: Die Nadel muss sich
im Ruhezustand genau bei 0mmHg befinden

\item Ohrst\"opsel des Stethoskops zeigen schr\"ag nach vorne

\item Trichterfunktion durch Ber\"uhren der Membran pr\"ufen

\end{itemize}
\begin{itemize}
\item Kontaktaufnahme mit dem Patienten

\end{itemize}
\begin{itemize}
\item Standardisierte Blutdruckmessung in Ruhe = konstante Bedingungen +
vor der Messung: kein Nikotin ( 30 min) -- Koffein ( 2 h) -- Alkohol (8
h)  {}-- k\"orperliche Anstrengung (1h)

\end{itemize}
\begin{itemize}
\item Patientenlagerung

\end{itemize}
\begin{itemize}
\item Arm freimachen und entspannt auf einen Tisch legen in leichter
Flexion und Supination

\item Oberarm auf Herzh\"ohe

\end{itemize}
\begin{itemize}
\item Auswahl der geeigneten Manschettengr\"o{\ss}e

\end{itemize}
\begin{itemize}
\item Manschette laut Herstellerangaben w\"ahlen z.B Standardmanschette:
Oberarmumfang max. 31 cm

\item Messwertverf\"alschung vermeiden! Der aufblasbare Teil der
Manschette soll 80\% des Oberarmumfangs umfassen, die Breite der
Manschette mindestens 1/3 des Oberarmumfangs betragen

\item Zu kleine Manschetten f\"uhren zum Messen zu hoher Werte, zu weite
Manschetten k\"onnen nicht am Arm fixiert werden

\end{itemize}
\begin{itemize}
\item Handhabung des Stethoskops

\item Bestimmen der Anlegestelle f\"ur Manschette und der
Auskultationsstelle durch Palpation

\end{itemize}
\begin{itemize}
\item Arteria brachialis medial in der Ellenbeuge aufsuchen

\end{itemize}
\begin{itemize}
\item Anlegen der Manschette

\end{itemize}
\begin{itemize}
\item Manschette muss auf unbekleideter Haut aufliegen

\item m\"a{\ss}ig straff

\item Abstand zur Ellenbeuge 2,5 cm

\item aufblasbaren Teil \"uber A. brachialis zentrieren

\item Schl\"auche nicht im Weg,

\item CAVE: Halbseitenschw\"ache/-l\"ahmung, Shuntarm bei
Dialysepatienten, Arm\"odem

\end{itemize}
\begin{itemize}
\item Puls f\"uhlen beim Aufpumpen

\end{itemize}
\begin{itemize}
\item Schnelles Aufpumpen

\item Tasten des Radialispulses mit Finger 2-4

\item Aufpumpen bis 30 mmHg \"uber jenen Druck, ab dem kein Radialispuls
mehr getastet wird

\end{itemize}
\begin{itemize}
\item Messvorgang

\end{itemize}
\begin{itemize}
\item Kopf des Stethoskops auf a. brachialis legen

\item Druck langsam ablassen 2 mmHg/sec

\item Systolischer Wert: Beginn der turbulenten Korotkoffschen
Str\"omungsger\"ausche

\item diastolischer Wert: Verschwinden der Korotkoffschen Ger\"ausche

\item Wenn Ger\"ausche \"uberhaupt nicht verschwinden (z.B. bei Kindern,
Schwangeren etc.), so gilt als diastolischer Wert der Messpunkt, an dem
die Ger\"ausche deutlich leiser werden

\item Wenn Ger\"ausche verschwunden sind, noch weitere 10mmHg langsam
ablassen -- wenn dann keine Ger\"ausche mehr h\"orbar sind, kann die
Manschette z\"ugig entleert werden

\item CAVE: auskultatorische L\"ucke

\end{itemize}
\begin{itemize}
\item Protokollierung der Blutdruckwerte

\end{itemize}
\begin{itemize}
\item Immer sofort nach Messung protokollieren

\end{itemize}
Das EKG

Das Elektrokardiogramm (EKG) zeigt mittels eines entsprechenden
Ger\"ates die elektrische Herzaktivit\"at des Reizleitungssystems an.
Dazu werden auf den K\"orper nach einem bestimmten Muster Elektroden
aufgeklebt und mittels eines Kabels mit dem EKG-Ger\"at verbunden. 

EKG -- Ablauf einer Erregung

\begin{itemize}
\item P-Welle \~{} Vorhoferregung

\item PQ-Zeit \~{} AV-\"Uberleitung

\item QRS-Komplex \~{} Kammererregung

\item ST-Strecke + T-Welle \~{} Erregungsr\"uckbildung

\end{itemize}
Das EKG gibt nur die elektrische Leitung, nicht aber die tats\"achliche
Muskelarbeit wieder! Es erlaubt keine Aussage \"uber die Herzleistung! 

Grundsatz: Selbst ein toter Patient kann ein
{\quotedblbase}sch\"ones{\textquotedblleft} EKG haben. (allerdings
nicht lange)

\subparagraph{EKG -- Ableitungen}
Grunds\"atzlich unterscheidet man zwischen den Extremit\"atenableitungen
und den Brustwandableitungen. 

\begin{itemize}
\item Extremit\"atenableitungen: I, II, III, aVL, aVR, aVF

\item Brustwandableitungen: V1 -- V6

\end{itemize}
Extremit\"atenableitungen

Die Extremit\"atenableitungen werden mittels 4 Elektroden abgeleitet. Je
eine Elektrobe wird an den Enden der Extremit\"aten angebracht. Zur
Vereinfachung k\"onnen die Elektroden auch am \"Ubergang vom Rumpf zu
den Extremit\"aten angebracht werden (z.B. statt dem rechten Arm an der
rechten Schulter).

Aus den Informationen der 4 Elektroden erzeugt das EKG-Ger\"at 6
Ableitungen:  I, II, III, aVL, aVR, aVF 

  [Warning: Image ignored] % Unhandled or unsupported graphics:
%\includegraphics[width=12.304cm,height=8.975cm]{AASdev20090228-img4.png}
 

\begin{center}
\begin{tabular}{|m{4.367cm}|m{4.367cm}|m{4.367cm}|}
\hline
\centering Elektrode\par

 &
\centering Position\par

 &
\centering Alternative Postion\par

\\\hline
Rot

 &
Rechte Hand

 &
Rechte Schulter

\\\hline
Gelb

 &
Linke Hand

 &
Linke Schulter

\\\hline
Gr\"un

 &
Linker Fu{\ss}

 &
Linke Leiste

\\\hline
Schwarz

 &
Rechter Fu{\ss}

 &
Rechte Leiste

\\\hline
\end{tabular}
\end{center}
Brustwandableitungen

Die Brustwandableitungen werden mittels 6 Elektroden abgeleitet und
erlauben eine genauere Zuordnung von St\"orungen zu der jeweiligen
Region des Herzens. Aus den Informationen der 6 Elektroden erzeugt das
EKG-Ger\"at 6 Ableitungen: V1 -- V6. 

\begin{center}
\begin{tabular}{|m{6.649cm}|m{6.651cm}}
\hline
\centering Elektrode\par

 &
\centering Position\par

\\\hline
V1

 &
4. Zwischenrippenraum rechts vom Sternum

\\\hline
V2

 &
4. Zwischenrippenraum links vom Sternum

\\\hline
V3

 &
Mitte zwischen V2 und V4

\\\hline
V4

 &
5. Zwischenrippenraum links in der Verl\"angerung der
Schl\"usselbein-Mittellinie

\\\hline
V5

 &
Mitte zwischen V4 und V6

\\\hline
V6

 &
6. Zwischenrippenraum in der mittleren Achsellinie

\\\hline
\end{tabular}
\end{center}
\subsection{Traumafertigkeiten}
Michael Motal

Helmabnahme

immer.

Material

\begin{itemize}
\item 2 Helfer

\item 1 stifneck f\"ur sp\"ater!

\end{itemize}
Vorgehen

\begin{itemize}
\item 1 Helfer fixiert den Helm von hinten und l\"asst nicht mehr los,
bis der Helm unten ist. Zwischen Helm und Schritt eine Helmh\"ohe
Abstand lassen, damit der Helm in einer Bewegung abgezogen werden kann

\item 1 Helfer spricht den Patienten von vorne an, \"offnet Visier und
Kinnriemen

\item der seitliche Helfer stabilisiert den Kopf: eine Hand fasst mit
Daumen und Zeigefinger das Unterkiefer, die andere Hand st\"utzt den
Kopf von unten. W\"ahrend der Helm abgezogen wird rutscht diese Hand
etwas nach

\item der hintere Helfer zieht den Helm in einer S-f\"ormigen Bewegung
ab: den Helm zuerst kippen bis die Nasenspitze sichtbar ist, dann die
Unterkante in der gegengleichen Bewegung zum Helfer bewegen

\item der hinere Helfer zieht den Helm sanft fertig ab und legt ihn
sicher ab

\item der hintere Helfer \"ubernimmt den Kopf im doppelten C-Griff:
Daumen und Zeigefinger um ein Ohr, Kopf leicht auf Zug halten.

\item der seitliche Helfer legt ein stifneck an (siehe unten)

\end{itemize}
Stifneck

zur Stabilisierung der Halswirbels\"aule bei jedem Verdacht auf eine
Verletzung; auf jeden Fall bei jeder Helmabnahme!

Material

\begin{itemize}
\item 2 Helfer

\item 1 stifneck

\end{itemize}
Vorgehen

\begin{itemize}
\item Pat. auffordern, den Kopf nicht zu bewegen

\item nach Helmabnahme: Kopf weiterhin leicht auf Zug halten

\item der hintere Helfer h\"alt Kopf/HWS in Neutralstellung

\item der seitliche Helfer misst die Entfernug von Schulter bis
Unterkiefer in Fingerbreiten aus

\item stifneck entsprechend einstellen: gemessene Entfernung auftragen:
Hartplastikrand, der an der Schulter aufliegt unten, Markierung am
verschiebbaren Teil bzw H\"ohe der Kinnst\"utze oben)

\item Verschiebemechanismus des stifneck blockieren, ab nun darf kein
Verstellen mehr m\"oglich  sein 

\item stifneck zuerst am Kinn anlegen/anpassen und festhalten

\item Klettverschlu{\ss}band so einfalten, dass es nicht am Boden
scheifen kann (h\"alt sonst nicht mehr)

\item Nackenteil unter Pat. durchf\"uhren (Vorsicht:
Klettverschlussband)

\item fest schlie{\ss}en

\end{itemize}
\begin{itemize}
\item[] Kinnspitze sollte ein kurzes St\"uck \"uber Kante hinausschauen
(kein Durchrutschen) aber nicht allzu weit (falsch eingestellt)

\end{itemize}
Schaufeltrage

zusammengesetzt aus 2 L\"angsh\"alften; Fu{\ss}teil jeder H\"alfte
l\"angenverstellbar, Kopfteil fixe L\"ange.

Zur beinahe bewegungsfreien Rettung bei Verdacht auf Wirbels\"aulen-,
Becken- oder Oberschenkelverletzungen. Letztendlich soll der Patient
auf der Vakuummatratze gelagert werden. Sollte der Patient gr\"o{\ss}er
als die Trage sein m\"ussen Kopf und Oberk\"orper auf jeden Fall
trotzdem auf der Trage zu liegen kommen. Ausnahme: Beinverletzung, bei
der eine Wirbels\"aulenverletzung ausgeschlossen werden kann.

Material

\begin{itemize}
\item 2 Helfer

\item Schaufeltrage

\item Gurte f\"ur Schaufeltrage

\end{itemize}
Vorgehen

\begin{itemize}
\item Schaufeltrage neben Patienten in zusammengebautem Zustand auf die
richtige L\"ange einstellen. Sollte ein Patient zu gro{\ss} sein
m\"ussen Kopf und Oberk\"orper auf jeden Fall trotzdem auf der Trage
liegen! Vorsicht: Markierung der max. L\"ange bei Schraubmodell

\item \"Offnen der Trage: Schnappverschl\"usse; je nach Methode nur den
unteren oder beide \"offnen.

\item Aufschaufeln, 1. Methode: Fu{\ss}seite offen, Pat. wird von oben
kommend wie auf eine Schere aufgeladen, die Trage unter ihm
geschlossen. Kontrolle des sicheren Schlu{\ss}es!

\item Aufschaufeln: 2. Methode: beide Teile getrennt, Pat. leicht an
Schulter und H\"ufte zur Seite gedreht und jew. Tragenteil
daruntergeschoben. H\"alften zuerst am Kopf-, dann am Fu{\ss}ende
schliessen. Kontrolle des sicheren Schlu{\ss}es!

\item Angurten: 2 Gurte; jeden Gurt zuerst um den Handgriff der einen,
dann der anderen H\"alfte f\"uhren, dann \"uber dem Patienten
schlie{\ss}en und festziehen.

\item Umlagern auf Vakuummatratze: Trage mit vorbeireiteter Matratze
(s.u.) neben Patienten stellen um diesen einen m\"oglichst geringen Weg
tragen/bewegen zu m\"ussen

\item Gurte \"offnen und entfernen.

\item H\"alften trennen: beide Schnappverschl\"usse \"offnen

\item H\"alften parallel zum Patienten flach wegziehen. Der Patient soll
hierbei nicht bewegt werden!

\item Sollten die Schnappschl\"osser klemmen kann man die Trage wie bei
der Scherentechnik an nur einer Seite \"offnen und in V-Form (je ein
Helfer auf einer Seite) entfernen.

\end{itemize}
Vakuummatratze:

zur Ruhigstellung des gesamten K\"orpers, bei Verdacht auf m\"ogliche
Wirbels\"aulenverletzungen, Becken- oder Oberschenkelfrakturen (lockere
Indikationsstellung!)

Am Ende soll die Matratze an die K\"orperform der Patienten eng
angepasst sein und somit ein Bewegen verhindern.

Material

\begin{itemize}
\item 2 Helfer

\item Fahrtrage

\item Vakuummatratze

\item Accuvac

\item Leintuch

\item ev. Adapter, je nach Matratzentyp

\end{itemize}
Vorgehen

\begin{itemize}
\item Fahrtrage auf tiefste Stufe stellen (ganz unten)

\item Vorbereiten: Vakuummatratze auf Fahrtrage legen. Die weichere
Seite ist die Patientenseite, Ventil beim Kopf. Kugeln so verteilen,
dass der gesamte K\"orper gut sowie die betroffene Stelle besonders gut
fixiert werden kann.

\item Vorabsaugen mit Accuvac, bis gut formbar, aber nicht allzu fest.

\item Vorformen, Ventil verschlie{\ss}en. Wenn vorhanden Leintuch auf
die Matratze legen.

\item Angegurteten Patienten mit Schaufeltrage auf die Vakuummatratze
legen, Schaufeltrage entfernen (siehe
{\quotedblbase}Schaufeltrage{\textquotedblleft})

\item Absaugen mit Accuvac, dabei z\"ugig Matratze an den Patienten
anpassen. Besonderes Augenmerk gilt der Fixierung der verletzten
Regionen sowie dem Kopf ({\quotedblbase}Ohren{\textquotedblleft}
formen: Ecken neben dem Kopf hochformen).

\item Ist kein weiteres Absaugen mehr m\"oglich und die Matratze gut
angepasst (Matratze muss jetzt bretthart sein) Ventil schlie{\ss}en,
Accuvac abdrehen.

\item Pat. an Fahrtrage angurten.

\end{itemize}
Vakuum-Beinschiene

Schienung einer Fu{\ss}-, Knie- oder Unterschenkelverletzung. Achtung:
Oberschenkel: Vakuummatratze

1 Helfer seitlich, 1 Helfer beim Fu{\ss}ende

Material

\begin{itemize}
\item 2 Helfer

\item Vakuumbeinschiene

\item Accuvac

\item ev. Adapter

\end{itemize}
Vorgehen

\begin{itemize}
\item L\"ange anpassen, gegebenenfalls oben umfalten.
Klettverschlu{\ss}b\"ander herrichten.

\item Kugeln gut verteilen. Auch am Fu{\ss}ende m\"ussen genug Kugeln
sein, um dem Fu{\ss} sowohl unter der Sohle als auch am Fu{\ss}r\"ucken
gut Halt geben zu k\"onnen

\item Schiene leicht vorabsaugen, damit Kugeln nicht mehr verrutschen
k\"onnen

\item Schuh am gesunden Bein ausziehen: Finden der schonensten Methode.
Ein Helfer h\"alt das Bein auf Zug, der andere \"offnet den Schuh
(B\"ander ausf\"adeln, zur Not aufschneiden)

\item vorsichtiges Wiederholen zu zweit am verletztem Bein, Socken
ausziehen (DMS)

\item Stiefelgriff: gleich nach Ausziehen von Schuh/Socken nimmt der
Helfer am Fu{\ss}ende durch das Loch in der Schiene die Ferse des
verletzten Fu{\ss}es, die andere Hand h\"alt den Fu{\ss}r\"ucken. Das
Bein wird so auf dem Boden liegend auf Zug gehalten. Der Helfer
verbleibt so, bis die Schiene fest sitzt!

\item Der seitliche Helfer richtet die Schiene vorsichtig unter dem
Patienten aus und fixiert die Schiene mit den Klettverschlussb\"andern.
Um den Fu{\ss} werden Ohren gebogen und kreuzweise festgeklettet: der
Fu{\ss} soll von allen Seiten so gut wie m\"oglich umschlossen sein!

\item Der seitliche Helfer saugt ab und zieht dabei die Klettb\"ander
immer wieder fest nach. Der andere Helfer h\"alt nach wie vor den
Fu{\ss} im Stiefelgriff.

\item Fertig abgesaugt: Ventil schliessen, Accuvac abschalten. Der
Stiefelgriff kann jetzt losgelassen werden.

\end{itemize}
Die Schiene soll bretthart sein, das Bein leicht auf Zug fixiert. Der
Fu{\ss} soll weder seitlich noch zum K\"orper hin beweglich sein.

Sam-Splint

Unterarmschienung

Material

\begin{itemize}
\item 2 Helfer

\item Sam-Splint

\item Peha-Haft, Mullbinde + Leukoplast oder 2-3 Dreiecktuchkravatten

\item ein weiteres Dreiecktuch

\end{itemize}
Vorgehen

\begin{itemize}
\item Sam-Splint in der H\"alfte falten, an gesundem Arm der L\"ange
nach anpassen, Knick beim Ellbogen

\item Anlegen an verletztem Arm, sanft anpassen. Fixierung mit
Peha-Haft, Mullbinden oder Dreiecktuchkravatten

\item Fingerspitzen m\"ussen sichtbar sein (DMS), gegebenenfalls Enden
umfalten.

\item Lagerung des geschienten Armes im Armtragetuch

\end{itemize}
Druckverband

zur Stillung einer starken, d.h. potentiell lebensbedrohlichen Blutung

Material

\begin{itemize}
\item 1-2 Helfer

\item sterile Wundauflage je nach Wundgr\"o{\ss}e

\item saugf\"ahiger Druckk\"orper: zB Mullbinde

\item Dreiecktuchkravatte

\end{itemize}
Vorgehen

\begin{itemize}
\item sofort bei Erkennen einer starken Blutung durch direkten Druck auf
die Wunde und Hochhalten die Blutung verlangsamen (Pat. kann dies je
nach Zustand selbst tun!)

\item sterile Wundauflage auf die Wunde legen

\item Druckk\"orper dem Wundverlauf folgend \"uber die sterile
Wundauflage legen, weiterhin dr\"ucken.

\item Druckk\"orper mit einer Dreiecktuchkravatte fest und m\"oglichst
vollst\"andig umschlie{\ss}en

\item Dreiecktuchkravatte festziehen und mit einem festen Knoten bei
ausreichendem Druck direkt \"uber der Wunde fixieren

\item 2. Knoten zur Sicherung des ersten

\end{itemize}
sollte die Blutung zu einem sp\"ateren Zeitpunkt wieder st\"arker werden
oder der Druckverband zu locker angelegt sein einen weiteren
Druckk\"orper mit einer weiteren Dreiecktuchkravatte noch fester
fixieren. Einen schon angelegten Druckverband nie \"offnen!

\subsection{Assistenzt\"atigkeiten}
\begin{center}
\tablehead{}\begin{supertabular}{|m{6.649cm}|m{6.651cm}|}
\hline
Teile des Textes und der Handlungsanleitungen enstammen dem Skriptum
{\quotedblbase}\"Arztliche Grundfertigkeiten{\textquotedblleft} der 
Besonderen Einrichtung zur medizinischen Aus- und Weiterbildung --
Medizinische Universit\"at Wien, Abt. Methodik und Entwicklung, welches
in Zusammenarbeit  mit dem Klinischem Institut f\"ur Hygiene und
Medizinische Mikrobiologie entstanden ist. 

Wir danken Hrn. Univ.Prof. Dr. Michael Schmidts und seinem Team f\"ur
die Unterst\"utzung!

 &
Powered by

\centering   [Warning: Image ignored]
% Unhandled or unsupported graphics:
%\includegraphics[width=4.551cm,height=2.17cm]{AASdev20090228-img5.gif}
 \par

\\\hline
\end{supertabular}
\end{center}
Aufziehen eines Medikamentes aus einer Ampulle

Ampulle mit rotem Punkt: Sollbruchstelle, Daumen muss beim aufbrechen
auf dem roten Punkt sein. 

Ampulle mit {\quotedblbase}Halsring{\textquotedblleft}: Ampulle kann
nach alle Richtungen aufgebrochen werden. 

\begin{center}
\begin{minipage}{6.556cm}
Warum mach{\textquotesingle} cih mir wegen der Verletzungsgefahr Sorgen?
Ist ja nur ein kleiner Schnitt, ist ja meine Sache {\dots}

Nein! Unter Normalbedingungen, wenn ich in aller Ruhe mein Medi
aufziehen kann ist es tats\"achlich mein eigenes Bier. Was aber wenn
ich mich bei einer Reanimation an der Ampulle schneide? Die anderen
Kollegen sind besch\"aftigt (Dr\"ucken, blasen, Defi {\dots}) und dann
steh ich da mit meinem kleinen Schnitt und blute alles an. Und der
Notarzt braucht dringend das Adrenalin {\dots} Was jetzt?

Inakzeptabel!

\end{minipage}
\end{center}
Material 

\begin{itemize}
\item Spritze

\item Kan\"ule(n)

\item Ampulle mit Medikament

\item Trockene Tupfer

\end{itemize}
Aufziehen des Medikamentes

\begin{itemize}
\item Kontrolle Ablaufdatum

\item Kontrolle Inhalt: Verf\"arbung, Ausflockung?

\item \"Offnen der Ampulle

\end{itemize}
\begin{itemize}
\item  Glasampulle

\end{itemize}
\begin{itemize}
\item Ampullenspitze beklopfen, darauf achten dass sich keine
Fl\"ussigkeit im Apullenkopf mehr befindet

\item \"Offnen der Medikamenten-Ampulle mit Tupfer, roter Punkt weist
zum K\"orper

\end{itemize}
[Warning: Draw object ignored]

\begin{center}
 [Warning: Image ignored] % Unhandled or unsupported graphics:
%\includegraphics[width=4.764cm,height=3.145cm]{AASdev20090228-img6.png}

\end{center}
\begin{itemize}
\item Plastikampulle,
{\quotedblbase}Dreh-und-Drink{\textquotedblleft}-Verschlu{\ss}

\end{itemize}
\begin{itemize}
\item Den Verschlussknopf ohne Kontamination des Ampullenhalses
(Sollbruchstelle) abdrehen

\end{itemize}
\begin{itemize}
\item Ampulle abstellen

\end{itemize}
\begin{itemize}
\item Verpackung der Spritze aufsch\"alen, Spritze in Verpackung lassen

\item Verpackung der Nadel aufsch\"alen, Spritze auf Konus aufstecken

\item Entfernen der Kan\"uleabdeckung (nicht abdrehen, sondern
abziehen!)

\item Einf\"uhren der Nadel in die Ampulle

\end{itemize}
\begin{itemize}
\item Beim Einf\"uhren der Kan\"ule in die Ampulle den Ampullenrand
nicht ber\"uhren

\end{itemize}
\begin{itemize}
\item Aufziehen des Medikamentes

\item Vorsichtig entl\"uften

\item Variante {\quotedblbase}Verabreichung \"uber
Venflon{\textquotedblleft} 

\end{itemize}
\begin{itemize}
\item Nadel h\"andisch abziehen und in den Kontamed verwerfen. 

\item Arzt das Medikament \"ubergeben, Ampulle zeigen.

\item Wird das Medikamen nicht sofort verwendet spritze mit roten
Stoppel verschliessen und Ampulle ankleben.

\end{itemize}
\begin{itemize}
\item Variante {\quotedblbase}Zuspritzen in eine
Infusionsflasche{\textquotedblleft}

\end{itemize}
\begin{itemize}
\item Gummimembran desinfizieren

\item Kan\"ule nicht bis zum Schaft einstechen

\item Medikament zuspritzen

\item Beschriften der Flasche

\end{itemize}
\begin{itemize}
\item Variante {\quotedblbase}direktes Spritzen{\textquotedblleft}

\end{itemize}
\begin{itemize}
\item Arzt nach Nadelgr\"o{\ss}e fragen

\item Aufziehnadel h\"andisch entfernen

\item Spritznadel aufstecken

\item Arzt Spritze mit Verschlusskappe reichen, Ampulle zeigen.

\end{itemize}
Vorbereitung einer Infusion

Material

\begin{itemize}
\item \item Infusionsflasche

\item Infusionsbesteck

\item Aufh\"angevorrichtung

\item alkoholhaltiges Desinfektionsmittel

\item einige saubere Tupfer

\item Permanentmarker

\end{itemize}
Handlungsschritte

\begin{itemize}
\item H\"andedesinfektion         

\item Aufh\"angevorrichtung auf Flasche aufsetzen

\item Er\"offnen der Flaschenabdeckung (Membranschutz)

\item Wischdesinfektion der Gummimembran

\end{itemize}
\begin{itemize}
\item Alkoholgetr\"ankter Tupfer kann w\"ahrend weiterer Vorbereitung
auf Membran belassen werden

\end{itemize}
\begin{itemize}
\item Montage des Infusionsbestecks

\end{itemize}
\begin{itemize}
\item Radklemme schlie{\ss}en 

\item Einstich des Tropfkammerdorns in Infusionsflasche mit
geschlossenem L\"uftungsventil und geschlossener Radklemme

\item {\quotedblbase}No touch technique{\textquotedblleft} -- Dorn der
Tropfkammer und Gummimembran der Infusionsflasche d\"urfen nicht
ber\"uhrt werden

\end{itemize}
\begin{itemize}
\item Infusionsschlauch f\"ullen

\end{itemize}
\begin{itemize}
\item Flasche wenden und hochhalten

\item Tropfkammerluftventil \"offnen

\item Tropfkammer durch Zusammendr\"ucken bis zur H\"alfte f\"ullen

\item Konusabdeckung des Infusionsschlauchs am anderen Ende muss bei der
Bel\"uftung nicht entfernt werden

\item Infusionsflasche und Ende des Infusionsschlauchs mit nicht
dominanter Hand halten  (Wenn sich das Infusionsschlauchende auf
gleicher H\"ohe wie der Spiegel in der Tropfkammer befindet,
entl\"uftet sich der Schlauch problemlos --  Prinzip der
kommunizierenden Gef\"a{\ss}e, an kommunizierenden Gef\"a{\ss}en stellt
sich immer dieselbe Fl\"ussigkeitsspiegelh\"ohe ein.)

\item Radklemme \"offnen

\item Infusionsschlauch {\quotedblbase}ohne
vergie{\ss}en{\textquotedblleft} entl\"uften %
%
(Fl\"ussigkeit bis ca. 5cm vor Konnektor rinnen lassen, endg\"ultig
entl\"uften bei Patienten vor Verabreichung)

\item Radklemme schlie{\ss}en

\end{itemize}
Venenverweilkan\"ule ({\quotedblbase}Venflon{\textquotedblleft}):
Vorbereitung und Assistenz 

  [Warning: Image ignored] % Unhandled or unsupported graphics:
%\includegraphics[width=13.467cm,height=4.658cm]{AASdev20090228-img7.png}
 

Bild: BEMAW

Farbcodierung der peripheren Venenverweilkan\"ulen

\begin{flushleft}
\begin{tabular}{|m{1.5619999cm}|m{1.7229999cm}|m{1.7229999cm}|m{2.0509999cm}|m{1.7839999cm}|m{2.102cm}|}
\hline
\multicolumn{3}{|m{5.408cm}|}{\centering Nur zur Illustration!\par

} &
\centering Durchfluss [ml/min]\par

 &
\centering Farbcode\par

 &
\\\hline
\centering {\O} mm\par

 &
\centering G\par

 &
\centering Lmm\par

 &
 &
 &
\\\hline
\centering 0,5 x 0,8\par

 &
\centering 22\par

 &
\centering 25\par

 &
\centering 25\par

 &
\centering blau\par

 &
\centering Sehr d\"unn\par

\\\hline
\centering 0,7 x 1,0\par

 &
\centering 20\par

 &
\centering 32\par

 &
\centering 55\par

 &
\centering rosa\par

 &
\centering d\"unn\par

\\\hline
\centering 0,9 x 1,2\par

 &
\centering 18\par

 &
\centering 40\par

 &
\centering 90\par

 &
\centering gr\"un\par

 &
\centering normal\par

\\\hline
\centering 1,1 x 1,4\par

 &
\centering 17\par

 &
\centering 42\par

 &
\centering 135\par

 &
\centering wei{\ss}\par

 &
\centering dick\par

\\\hline
\centering 1,3 x 1,7\par

 &
\centering 16\par

 &
\centering 42\par

 &
\centering 170\par

 &
\centering grau\par

 &
\centering Sehr dick\par

\\\hline
\centering 1,6 x 2,1\par

 &
\centering 14\par

 &
\centering 42\par

 &
\centering 265\par

 &
\centering orange\par

 &
\centering Pferd\par

\\\hline
\end{tabular}
\end{flushleft}
Tabelle: BEMAW (modifiziert)

Material

\begin{itemize}
\item Nierentasse 

\item Staubinde

\item Venflon: 

\end{itemize}
\begin{itemize}
\item Die Gr\"o{\ss}e (= Farbe) ist vom Arzt vorher zu erfragen!)

\item Die Fl\"ugel werden hinunter geklappt

\item Die farbige Lasche des Zuspritzventils in rechten Winkel zur
Venflonachse stellen.

\end{itemize}
\begin{itemize}
\item Einmalverschlusskappe ({\quotedblbase}roter
Stoppel{\textquotedblleft})

\item Fixierungspflaster

\item einige saubere trockene Tupfer

\item Alko-Tupfer

\item (Einwegspritze mit 5ml Sp\"ulfl\"ussigkeit)

\item (Faserschreiber)

\item durchstichsicherer Entsorgungsbeh\"alter
({\quotedblbase}sharp{\textquotedblleft})

\end{itemize}
Assistenz

\begin{itemize}
\item Stauschlauch zureichen

\item ge\"offnete Verpackung des Alko-Tupfer hinreichen sodass dieser
entnommen werden kann

\item Kontamed in Reichweite des  Arztes / NKV bereitstellen sodass er
sp\"ater die Nadel abwerfen kann.

\item Venflon zureichen

\end{itemize}
\begin{itemize}
\item Venflon an der Schutzkappe halten sodass der Arzt / NKV den
Venflon abziehen kann.

\end{itemize}
\begin{itemize}
\item Trockene Tupfer zureichen.

\item Je nach Situation Infusionsschlauch, roten Stoppel oder gar nichts
zureichen.  

\item Fixierpflaster zureichen

\end{itemize}
Intubation: Vorbereitung und Assistenz 

\subsection[O2 \ - Sauerstoff]{O2  {}- Sauerstoff}
Sauerstoff (chem. Symbol O2) ist ein lebenswichtiges Gas welches in
unserer Atemluft vorkommt und au{\ss}erdem ein essentielles
Notfall{\textquotedblright}medikament{\textquotedblleft}, die
fr\"uhzeitige Verabreichung ist wichtig f\"ur die Prognose des
Patienten!

Sauerstoff darf und muss bei Bedarf vom RS verabreicht werden.

\begin{center}
 [Warning: Image ignored] % Unhandled or unsupported graphics:
%\includegraphics[width=4.522cm,height=5.628cm]{AASdev20090228-img8.png}
\captionof{figure}[: NEIN!]{: NEIN!}

\end{center}
\paragraph{Sauerstoff -- technische Grundlagen}
Sauerstoff wird in Druckflaschen  gelagert. Laut EU-Norm sind diese
Flaschen wei{\ss} und mit einem {\quotedblbase}N{\textquotedblleft}
gekennzeichnet. Das {\quotedblbase}N{\textquotedblleft} steht f\"ur
{\quotedblbase}neue Kennzeichnung{\textquotedblleft}, da fr\"uher die
Flaschen blau gekennzeichnet wurden. 

Wichtige Begriffe

\begin{itemize}
\item F\"ullgr\"o{\ss}e 

\item F\"ullvolumen 

\item Durchflu{\ss} 

\end{itemize}
Die Druckflaschen haben unterschiedliche F\"ullgr\"o{\ss}en (Volumina),
g\"angige Flaschengr\"o{\ss}en sind 0,5~l, 1~l, 1,5~l, 2~l, 5~l, 10~l.
Der Sauerstoff wird unter Druck gelagert, durch diesen Trick kann die
Menge des zur Verf\"ugung stehenden Sauerstoffs um ein Vielfaches
gesteigert werden. Der Druck einer gef\"ullten Flasche betr\"agt
normalerweise ca. 200~bar. 

Jede Druckflasche verf\"ugt \"uber ein Hauptventil und ein
Standardgewinde. An diesem Gewinde ist enteder ein Druckschlauch oder
ein Druckminderer angeschlossen. Der  Druckminderer hat ein
Abgabeventil welches den Druck drosselt und eine Abgabe an den
Patienten erm\"oglicht. Je nach Bauart kann der Druckminderer \"uber
einen Regler, ein Flowmeter und ein Barometer verf\"ugen. Mit dem
Regler kann man die Durchlu{\ss}- bzw. Abgaberate ein\-stellen, das
Flowmeter (Durchflu{\ss}anzeige) zeigt die aktuelle Durchflu{\ss}rate
in Liter pro Minute (~l~/~min~)an. Das Barometer zeugt den Druck in der
Sauerstofflasche an wenn das Haupt\-ventil ge\"offnet ist. 

\begin{center}
 [Warning: Image ignored] % Unhandled or unsupported graphics:
%\includegraphics[width=7.189cm,height=4.94cm]{AASdev20090228-img9.png}

\end{center}
Das effektive F\"ullvolumen erh\"alt man wenn man die F\"ullgr\"o{\ss}e
(in~l) mit dem Flaschendruck (in bar) multipliziert.  

Wenn man wissen m\"ochte wie lange die vorhandene Sauerstoffmenge bei
einer bestimmten Abgabemenge reicht dividiert man einfach die
vorhandene Menge durch die Abgaberate. Achtung: Es ist eine
Ausreichende Reserve einzuplanen, und es ist zu Bedenken dass in der
Flasche ein Restdruck verbleiben muss! 

Mit Sauerstoff l\"asst sich gut rechnen ...

\[\begin{gathered}\boldsubformula{\text{Enthaltene
}}\boldsubformula{O}_{2}\boldsubformula{\text{{}-Menge }}=\text{
Volumen der Flasche }\times \text{
Druck}\\\boldsubformula{\text{Verf\"ugbare Zeit }}=\frac{\text{
Enthaltene Menge }}{\text{ Durchflu{\ss} }}\end{gathered}\]
Beispiel:

Flasche 10l mit 150bar: 1500l  O2

Patient ben\"otigt 5l O2 pro Minute: 

Die F\"ullung reicht f\"ur 300 Minuten, abz\"uglich eines
Sicherheitsvolumens.

Die Sauerstoffflaschen sind pfleglich zu behandeln

Sauerstoff ist zusammen mit Fett selbstentz\"undlich! Die Armaturen und
Gewinde sind fettfrei zu halten! Das Hantieren mit Feuer und offenem
Licht ist in der N\"ahe von Sauerstoffflaschen verboten! 

Der gro{\ss}e Druck birgt weitere Gefahren: Sollte ein Ventil abbrechen
kann die Flasche eine enorme Kraft entwickeln und zu einer Rakete
werden. Die Kraft einer mit 200~bar gef\"ullten Flasche reicht aus um
Mauern zu druchbrechen!  

Achtung! Vorsicht beim Hantieren -  Explosionsgefahr!

Flaschen stehen unter Druck, ein Abbrechen der Ventile ist katastrophal!

Sauerstoff ist zusammen mit Fett selbstentz\"undlich ${\rightarrow}$
Ventile nicht mit fetten/\"oligen H\"anden bedienen.

\paragraph{Verabreichungsarten}
\subparagraph{Sauerstoffbrille}
Nur bei vorhandener Spontanatmung!

zwei kurze, in die Nasen\"offnungen  eingef\"uhrte Kunst\-stoff\-stutzen

bei ca. 4 l O2 /min bis 40\% O2  in Atemgas  m\"oglich

\begin{center}
\tablehead{}\begin{supertabular}{|m{6.649cm}|m{6.651cm}|}
\hline
\centering Vorteile\par

 &
\centering Nachteile\par

\\\hline
gute Toleranz

fehlendes Engegef\"uhl

Sprechen\&Husten m\"oglich

 &
ungenaue Dosierung

Austrocknung  der Schleimh\"aute

\\\hline
\end{supertabular}
\end{center}
Bei behinderter Nasenatmung kann die  Sauerstoffbrille im Mundwinkel
plaziert werden. 

\begin{center}
\tablehead{}\begin{supertabular}{|m{3.225cm}|m{3.225cm}|m{3.225cm}|m{3.225cm}|}
\hline
\centering Was\par

 &
\centering Wann\par

 &
\centering Vorteile\par

 &
\centering Nachteile\par

\\\hline
\centering Sauerstoffbrille\par



\begin{center}
 [Warning: Image ignored] % Unhandled or unsupported graphics:
%\includegraphics[width=2.74cm,height=2.74cm]{AASdev20090228-img10.png}

\end{center}
 &
Spontanatmung

O2-Flow {\textless} 5l/min

 &
gute Toleranz

fehlendes Engegef\"uhl

Sprechen\&Husten m\"oglich

 &
ungenaue Dosierung

Austrocknung  der Schleimh\"aute

\\\hline
\centering Inhalationsmaske\par



\begin{center}
 [Warning: Image ignored] % Unhandled or unsupported graphics:
%\includegraphics[width=2.437cm,height=2.707cm]{AASdev20090228-img11.png}

\end{center}
 &
Spontanatmung

O2-Flow ab 5l/min

 &
h\"ohere O2{}-Konzentration

 &
h\"ohere O2{}-Konzentration

CO2{}-R\"uckatmung  wenn O2{}-Durchflu{\ss} {\textless} 5 l/min!)

\\\hline
  [Warning: Image ignored] % Unhandled or unsupported graphics:
%\includegraphics[width=2.652cm,height=2.476cm]{AASdev20090228-img12.png}
 

 &
 &
Noch h\"ohere O2{}-Konzentration

 &
Wie oben.

\\\hline
\centering Beatmungsbeutel mit  Reservoir\par

 &
Beatmung des Patienten (bei Reanimation oder im Ausnahmefall bei
Atemstillstand oder nicht ausreichender Eigenatmung)

Beatmung eines intubierten Patienten

 &
 &
\\\hline
\end{supertabular}
\end{center}
\paragraph{Beatmung mit Beatmungsbeutel}
JEDE BEATMUNG SOLLTE MIT RESERVOIRBEUTEL UND 10-15 l O2/min
DURCHGEF\"UHRT WERDEN!

Beatmungsfrequenz beim Erwachsenen: 12-15 x / min (Eigenfrequenz)

Achtung! Zuviel ist schlecht! Nicht zu schnell beatmen ${\rightarrow}$
Hyperventilationsgefahr!

\paragraph{Sauerstofftherapie}
Prinzipiell sollte sich die gew\"ahlte Therapie an der Diagnose und am
Schweregrad orientieren!

Bedenke: 

\begin{itemize}
\item Diagnose!

\item Schockbek\"ampfung

\item Isch\"amie (bei Akutem Koronoarsyndrom, Insult, ..)

\item Atemst\"orungen (Fremdk\"orper, Atemnot, {\dots})

\end{itemize}
\paragraph[Pulsoxymetrie]{Pulsoxymetrie%
%
}
\index{Pulsoxymetrie}\begin{center}
 [Warning: Image ignored] % Unhandled or unsupported graphics:
%\includegraphics[width=3.516cm,height=3.516cm]{AASdev20090228-img13.png}

\end{center}
Das Pulsoxymeter (kurz: Pulsoxy) mi{\ss}t dden Puls und die
Sauerstoffs\"attigung (\index{SpO2}SpO2)des Blutes (Verh\"altnis
O2-beladene/ nichtbeladene Erythrozyten). 

Normalwert 96-99\%, 

[F078?] Achtung: Ein {\quotedblbase}guter{\textquotedblleft}
\index{SpO2}SpO2{}-Wert bedeutet nicht {\quotedblbase}keinen Sauerstoff
geben{\textquotedblleft}. Die Entscheidung bez\"uglich Sauerstoffgabe
ber\"ucksichtigt die Diagnose und den Zustand des Patienten, der
\index{SpO2}SpO2{}-Wert hat wenig bis gar nichts zu sagen! 

Achtung: Das Pulsoxymeter verwechselt O2 und CO-Molek\"ule, bei einer
CO-Vergiftung zeigt das Pulsoxy falsch-hohe S\"attigungswerte an!

Weitere Fehlerquellen: kein Signal: Nagellack, Zentralisation

Das Pulsoxymeter misst die Herzfrequenz und die Sauerstoffs\"attigung
(=SpO2) des Blutes, aber es gilt:

[F078?] Ein {\quotedblbase}guter{\textquotedblleft} S\"attigungswert ist
keine Begr\"undung keinen Sauerstoff zu verabreichen!

\section{[RS XIII] Untersuchung und Anamnese}
Hier lernen Sie wie Sie die bereits gelernten therotischen Grundlagen in
der Praxis einsetzen, wie und in welcher Reihenfolge Sie im Einsatz
vorgehen, wie Sie Dringlichkeiten einsch\"atzen und Verdachtsdiagnosen
erarbeiten k\"onnen und wie Sie dabei dann auch noch den \"Uberblick
behalten k\"onnen.

\subsection{{\quotedblbase}Hilfe, ein Patient!{\textquotedblleft} - Was
mach{\textquotesingle} ich jetzt?}
Blabla

\"Ubersicht 

Ersteinsch\"atzung ${\rightarrow}$  Erstma{\ss}nahmen ${\rightarrow}$ 
Plan zurechtlegen wie es weiter gehen soll

blabla

mehr blabla

Die Anamnese erz\"ahlt uns die Geschichte des Patienten

\paragraph{S{\textquotesingle}KAMEL denkt an alles}
\subparagraph{Symptome und Schmerzen}
\subparagraph{Krankheiten}
\subparagraph{Allergien}
\subparagraph{Medikamente}
\subparagraph{Ereignisse}
\subparagraph{Das Letzte ...}
Untersuchungen k\"onnen uns weiterhelfen

\subparagraph{Sehen -- H\"oren -- F\"uhlen }
\subparagraph{Vitalwerte}
\subparagraph{Atmung}
\subparagraph{Pulsoxymetrie}
\subparagraph{Blutzuckermessung}
\subparagraph{K\"orpertemperatur}
\subparagraph{Neurocheck}
\subparagraph{Traumacheck}
Suche nach Verletzungen oder Verletzungsanzeichen

\begin{itemize}
\item Patient orientiert? Unfallhergang erinnerlich? Schmerzen?

\end{itemize}
\begin{itemize}
\item Abtasten mit festem Griff: von Kopf bis Fu{\ss} (Sch\"adeldach,
Hinterkopf, Stirn, Wangenknochen, Nasenbein, Kiefer, Hals, Schultern,
Arme, H\"ande, Brustkorb vorne und seitlich (Prellmarken?
Atemmechanik?), Bauch in 4 Quadranten (Prellmarken?
Abwehrspannung/harte Bauchdecke?), Schambein, Becken (stabil?), Beine
(sichtbare Verletzungen?), F\"u{\ss}e: Schuhe und Socken ausziehen)

\item {\quotedblbase}DMS{\textquotedblleft}: Finger und Zehen: Kontrolle
von Motorik (bewegen lassen), Sensibilit\"at ({\quotedblbase}Sp\"uren
Sie das?{\textquotedblleft}) und Durchblutung:
Nagelbettprobe/Rekapillarisierungszeit: leichter Druck auf Nagelende
bis Nagelbett wei{\ss} ist, loslassen. Zeit bis Nagelbett wieder
rosarot ist sollte nicht l\"anger als 1-2 Sekunden sein.
Seitenvergleich! Eine verl\"angerte Rekapzeit deutet auf eine
Durchblutungsst\"orung hin.

\item Kontrolle von Mund, Nase und Ohren auf Blutungen (mit einer
Lampe). Tupfertest: etwas Blut aus Ohren mit einem Tupfer auffangen.
Sollte Liquor im Blut enthalten sein bildet sich um den roten Blutfleck
ein heller, bernsteinfarbener Hof.

\end{itemize}
bei Sch\"adelverletzungen ist ein Neurocheck auf jeden Fall
erforderlich!

Tarumavortr\"age m\"ussen hierher verweisen

\subparagraph{Abtasten des Abomens }
Vorgehen in der Praxis

\paragraph{Szene\"uberblick}
\paragraph[Die ersten \ \ 30 Sekunden ]{Die ersten  30 Sekunden }
Ersteinsch\"atzung ${\rightarrow}$  Erstma{\ss}nahmen ${\rightarrow}$ 
Plan zurechtlegen wie es weiter gehen soll

Erster Schritt: Ersteinsch\"atzung

\begin{center}
\tablehead{}\begin{supertabular}{|m{3.158cm}m{10.142cm}|}
\hline
\centering Ersteindruck\par

 &
Aussehen, Gesichtsausdruck, Haltung, Sprache

Bl\"asse, Zyanose, Schwei{\ss}, Schohnhaltung, Atemnot, verwaschene
Sprache, Des\-orientiertheit, Kr\"ampfe 

\\\hline
\centering Bewusstsein\par

 &
Reaktion auf verbale Ansprache, Weckversuch, Schmerzreiz

klar, somnolent, sopor\"os, komat\"os/bewusstlos

\\\hline
\centering Atemweg\par

 &
Atemger\"ausch, Gestik

frei oder verlegt

\\\hline
\centering Atmung\par

 &
Frequenz, Atemmuster, Anstrengung beim Atmen, Atemger\"ausche,
Heimsauerstoff

\\\hline
\centering Kreislauf\par

 &
Bewusstsein, Hautfarbe, Radialispuls, Rekap-Zeit

\\\hline
\centering Traumacheck 1\par

 &
Schwere Verletzungen: Backen-, Oberschenkel-, Serienrippenfraktur,
massive Blutung

\\\hline
\centering Hauptbeschwerde\par

 &
Grund der Berufung, Schmerzen und andere f\"ur den Patienten dringliche
Symptome

\\\hline
\end{supertabular}
\end{center}
Nun stellen sich folgende Fragen:

\begin{itemize}
\item Was ist das Problem des Patienten?

\begin{center}
\begin{minipage}{3.863cm}
Alarmzeichen

Massive St\"orungen von 

Bewusstsein

Atemweg / Atmung 

Kreislauf 

typische Symptome 

Atemnot die sich nicht bessert 

Brodelndes Atemger\"ausch 

Thoraxschmerz  

Schocksymptome 

entgleiste Vitalwerte

schwere Verletzungen

{\dots} 

\end{minipage}
\end{center}
\item Ist der Patient akut vital bedroht?

\end{itemize}
\begin{itemize}
\item Anzeichen einer vitalen Bedrohung k\"onnen u.a. sein: St\"orungen
von Be\-wusstsein, Atemweg, Atmung oder Kreislauf sowie typische
Symptome wie sich nicht bessernde Atemnot, brodelndes
Atem\-ge\-r\"ausch, Thoraxschmerz, Schocksymptome und entgleiste
Vital\-werte.

\end{itemize}
W\"ahrend des weiteren Verlaufs (Anamnesegespr\"ach, etc.) kann man auf
weitere Befunde sto{\ss}en welche auf eine vitale Be\-drohung hinweisen
(zB. massiv erh\"ohter Blutdruck bei bei vor\-ge\-sch\"adigtem Herzen,
pl\"otzliche oder schleichende Zustands\-verschlechterung, etc.). 

Die Ersteinsch\"atzung muss st\"andig \"uberpr\"uft werden. Neue Befunde
und der Zustand des Patienten k\"onnen die Situation \"andern.

Zweiter Schritt: Erstma{\ss}nahmen

Welche Ma{\ss}nahmen sind f\"ur den Patienten am wichtigsten und
m\"ussen zuerst erfolgen? 

\begin{center}
\begin{minipage}{4.2cm}
Bei der Pr\"ufung {\quotedblbase}blind{\textquotedblleft} auf Nummer
sicher gehen und alle gleich f\"ur {\quotedblbase}vital
bedroht{\textquotedblleft} erkl\"aren gilt nicht! 

Viele Patienten bei der Pr\"ufung sind nicht vital bedroht und brauchen
keinen Notarzt. 

Erkl\"art man einen Patienten f\"ur vital bedroht muss man das
begr\"unden k\"onnen!

\end{minipage}
\end{center}
Bei vital bedrohten Patienten gibt es 5 Standardma{\ss}nahmen die sofort
zu erfolgen haben. Bei allen anderen Patienten muss individuell
\"uberlegt werden welche Ma{\ss}nahmen zuerst zu treffen sind.

\begin{center}
\tablehead{}\begin{supertabular}{|m{5.3770003cm}m{7.9230003cm}|}
\hline
\multicolumn{2}{|m{13.500001cm}|}{\centering Erstma{\ss}nahmen bei vital
bedrohten Patienten\par

}\\\hline
Situationsgerechte Lagerung

 &
\\\hline
Sauerstoffgabe

 &
je nach Indikation, Vorsicht bei COPD und Hyper\-ventilation

\\\hline
Notarzt{}-Nachforderung

 &
mit kurzer Begr\"undung

\\\hline
Reanimationsbereitschaft

 &
Platz schaffen, Beatmungsbeutel zu\-sammen\-bauen, ggfs. Defi und
Absaugung holen und in Griffweite stellen

\\\hline
Vitalwerte erheben und Monitoring

 &
RR, HF, Pulsoyxymetrie wenn vor\-handen

\\\hline
\end{supertabular}
\end{center}
Dritter Schritt: Durchatmen, nachdenken, planen

Nach den Erstma{\ss}nahmen muss man nun planen welche Ma{\ss}nahmen,
Untersuchungen und Fragen am wichtigsten sind und zuerst durchgef\"uhrt
werden m\"ussen und was eher in der Abfolge nach hinten ver\-schoben
wird. Die folgende Checkliste gibt einen \"Uberblick \"uber die
Ma{\ss}nahmen und das Anamnese\-ge\-spr\"ach, die Reihenfolge ist bei
jedem Patienten anders!

\paragraph{Checkliste}
\begin{flushleft}
\tablehead{}\begin{supertabular}{m{2.467cm}|m{1.0209999cm}m{6.926cm}m{2.483cm}|}
\hline
\multicolumn{1}{|m{2.467cm}|}{} &
 &
\multicolumn{2}{m{9.609cm}|}{m\"ogliche Ma{\ss}nahmen, die Reihenfolge
ist je nach Patient unterschiedlich

}\\\hline
\multicolumn{1}{|m{2.467cm}|}{\centering SELBSTSCHUTZ\par

\begin{center}
 [Warning: Image ignored] % Unhandled or unsupported graphics:
%\includegraphics[width=1.005cm,height=0.9cm]{AASdev20090228-img14.gif}

\end{center}
} &
 &
\multicolumn{2}{m{9.609cm}|}{\begin{itemize}
\item Schutzausr\"ustung (Handschuhe, Uniform, Helm)

\item Gefahrenbereich, Unfallstelle absichern, Hunde wegsperren lassen

\item ggfs. R\"uckzug, Spezialkr\"afte anfordern, {\dots} 

\item {\dots} 

\end{itemize}
}\\\hline
\multicolumn{1}{|m{2.467cm}|}{\centering \par

\begin{center}
 [Warning: Image ignored] % Unhandled or unsupported graphics:
%\includegraphics[width=1.12cm,height=0.63cm]{AASdev20090228-img15.gif}

\end{center}
\centering ERKENNEN\par

} &
 &
\multicolumn{2}{m{9.609cm}|}{\begin{itemize}
\item Situation/Milieu (zB Alter, soziales Umfeld, SMV, Alkohol,
Pensionistenheim, ...)

\item offensichtliche Symptome (zB Zyanose, Bewusstlosigkeit, starke
Blutung, ...)

\item Schockgefahr (Zyanose, kalter Schwei{\ss}, Radialispuls
tastbar/nicht tastbar, Puls beschleunigt)

\end{itemize}
}\\\hline
\multicolumn{1}{|m{2.467cm}|}{\centering \par

\begin{center}
 [Warning: Image ignored] % Unhandled or unsupported graphics:
%\includegraphics[width=1.12cm,height=1.12cm]{AASdev20090228-img16.png}

\end{center}
\centering ERSTE HILFE\par

} &
 &
\multicolumn{2}{m{9.609cm}|}{\begin{itemize}
\item Lagerung

\item Helmabnahme

\item Blutstillung

\item psychischer Beistand

\item W\"armeerhaltung

\end{itemize}
}\\\hline
\multicolumn{1}{|m{2.467cm}|}{\centering   [Warning: Image ignored]
% Unhandled or unsupported graphics:
%\includegraphics[width=1.18cm,height=1.18cm]{AASdev20090228-img17.png}
 \par

\centering ANAMNESE\par

} &
\centering S\par

 &
\multicolumn{2}{m{9.609cm}|}{Symptome \& Schmerzen 

(Zeit/Verlauf, St\"arke, Art/Qualit\"at, Lokalisation/Ausstrahlung,
Provokation)

}\\\hline
 &
\centering K\par

 &
\multicolumn{2}{m{9.609cm}|}{Krankheiten (Bestehen Grunderkrankungen, zB
Diabetes?)

}\\\hhline{~---}
 &
\centering A\par

 &
\multicolumn{2}{m{9.609cm}|}{Allergien/Allergische Reaktion (auch
Allergien auf Medikamente)

}\\\hhline{~---}
 &
\centering M\par

 &
\multicolumn{2}{m{9.609cm}|}{Medikamente (Nimmt Patient Medikamente?
Wogegen sind diese?)

}\\\hhline{~---}
 &
\centering E\par

 &
\multicolumn{2}{m{9.609cm}|}{Ereignisse vor Notfalleintritt (zB
Anstrengung vor Brustschmerz, Sturz, ...)

}\\\hhline{~---}
 &
\centering L\par

 &
\multicolumn{2}{m{9.609cm}|}{Letzte ... (zB letzte Mahlzeit, letzte
Regelblutung, letzter Spitalsaufenthalt)

}\\\hhline{~---}
\multicolumn{1}{|m{2.467cm}|}{\centering   [Warning: Image ignored]
% Unhandled or unsupported graphics:
%\includegraphics[width=1.679cm,height=1.82cm]{AASdev20090228-img18.jpg}
 \par

\centering BEFUNDE\par

} &
\centering RR\par

 &
\multicolumn{2}{m{9.609cm}|}{auskultatorische und palpatorische Messung

}\\\hline
 &
\centering HF\par

 &
\multicolumn{2}{m{9.609cm}|}{Herzfrequenz (Puls), Puls rhythmisch oder
arrhytmisch

}\\\hhline{~---}
 &
\centering AF\par

 &
\multicolumn{2}{m{9.609cm}|}{Atemfrequenz ausz\"ahlen

}\\\hhline{~---}
 &
\centering SpO2\par

 &
\multicolumn{2}{m{9.609cm}|}{Sauerstoffs\"attigung im Blut, Monitoring

}\\\hhline{~---}
 &
\centering BZ\par

 &
\multicolumn{2}{m{9.609cm}|}{Blutzucker (Wann war die letzte
Mahlzeit/Insulinspritze? Diabetes bekannt?)

}\\\hhline{~---}
 &
\centering Temp.\par

 &
\multicolumn{2}{m{9.609cm}|}{K\"orpertemperatur/Fieber

}\\\hhline{~---}
 &
\centering Neuro status\par

 &
\multicolumn{2}{m{9.609cm}|}{Bewusstsein (ev. GCS), Orientierung,
Pupillen, Halbseitenzeichen, Herdblick, Meningismus, retrograde Amnesie

}\\\hhline{~---}
 &
\centering Trauma check\par

 &
\multicolumn{2}{m{9.609cm}|}{abnorme Beweglichkeit, DMS an allen
Extremit\"aten (Nagelbettprobe), Wunden, paradoxe Atmung,
Abwehrspannung des Bauches, ...

}\\\hhline{~---}
\multicolumn{1}{|m{2.467cm}|}{\centering SANIT\"ATSHILFE\par

\begin{center}
 [Warning: Image ignored] % Unhandled or unsupported graphics:
%\includegraphics[width=1.349cm,height=1.319cm]{AASdev20090228-img19.png}

\end{center}
} &
 &
\begin{itemize}
\item Entscheidung Notarzt-Nachforderung

\item Fortf\"uhrung der Ersten Hilfe

\item Schienung

\item Sauerstoffgabe

\item Reanimationsbereitschaft

\item Verlaufskontrolle

\item ...

\end{itemize}
 &
\\\hline
\multicolumn{1}{|m{2.467cm}|}{\centering NOTARZT-ASSISTENZ\par

\begin{center}
 [Warning: Image ignored] % Unhandled or unsupported graphics:
%\includegraphics[width=1.349cm,height=1.319cm]{AASdev20090228-img20.png}

\end{center}
} &
 &
\multicolumn{2}{m{9.609cm}|}{\begin{itemize}
\item Infusion/peripherven\"osen Zugang vorbereiten

\item Intubation vorbereiten

\item Medikamente vorbereiten

\item EKG anlegen

\item ...

\end{itemize}
}\\\hline
\multicolumn{1}{|m{2.467cm}|}{\centering   [Warning: Image ignored]
% Unhandled or unsupported graphics:
%\includegraphics[width=1.88cm,height=1.88cm]{AASdev20090228-img21.png}
 \par

\centering KRANKENHAUS\par

} &
 &
\multicolumn{2}{m{9.609cm}|}{\begin{itemize}
\item Welches Spital/Bett buche ich ab? (Arbeitsunfall? Welche Abteilung
fahre ich an? Ist der Patient bereits in einem Spital in Behandlung?)

\item Brauche ich eine Spezialabteilung? (Schockraum, Stroke Unit, PTCA,
Verbrennungsstation)

\item Arzt\"ubergabe (Anamnese, Befunde, Ma{\ss}nahmen, Hinweise)

\end{itemize}
}\\\hline
\end{supertabular}
\end{flushleft}
\paragraph[Ein Beispiel: Herr Sepp hat Atemnot]{Ein Beispiel: Herr Sepp
hat Atemnot}
Berufung: Atemnot

Sie kommen in eine Wohnung. Die Umgebung erscheint sicher. Die Wohnung
sieht sauber und ordentlich aus. Auf der Couch liegt flach ein
50-j\"ahriger Mann mit weit aufgerissenen Augen und sichtlicher
Atemnot. Er wirkt bla{\ss} und schwei{\ss}ig und fasst ich mit der
rechten Hand an den Brustbereich. 

Sie begr\"u{\ss}en den Patienten, stellen sich vor und fragen welche
Beschwerden er hat. Der Patient ringt nach Luft und klagt \"uber
Schmerzen in der Brust. Diese haben vor einer Viertelstunde angefangen.


Welche Informationen haben Sie bis jetzt gewonnen?

Der Patient ist bei Bewusstsein. Er hat Atemnot und einen pl\"otzlich
eingetretenen Thoraxschmerz. Er ist bla{\ss} und schwei{\ss}ig,
vielleicht steckt ein Schockgeschehen dahinter. 

K\"onnen Sie schon etwas tun?

Ja.  Bei erhaltenem Bewusstsein, Atemnot und Thoraxschmerz wird der
Patient mit erh\"ohtem Oberk\"orper gelagert. Weiters sprechen die
Symptome f\"ur den Einsatz von Sauerstoff. Es muss noch schnell
gekl\"art werden ob eine Hyperventilationstetanie (nicht
wahrscheinlich) oder COPD-Erkrankung (erfragen) vorliegt. 

Ist der Patient vital bedroht?

Ja. Zwar wurden noch keien Vitalwerte erhoben, aber bereits jetzt gilt
der Patient als vital bedroht aufgrund der pl\"otzlich einsetzenden
Thoraxschmerzen und der Atemnot. 

\begin{itemize}
\item Der Oberk\"orper des Patienten wird hochgelagert 

\item Der Notarzt wird nachgefordert (Berufungsgrund: akuter
Thoraxschmerz mit Atemnot und schlechtem AZ).

\item Der Patient wird gefragt ob er an COPD leidet, dies wird verneint.
Eine Sauerstoffmaske mit einem Flow von 10l/min wird angelegt. 

\item Reanimationsbereitschaft wird hergestellt

\end{itemize}
Was sind die n\"achsten Schritte?

Wichtig sind nun die Vitalwerte und das fr\"uhe Beginnen des
Monitorings.Dann muss man genaueres \"uber das Beschwerdebild in
Erfahrung bringen. Der weitere Plan h\"angt davon ab.

Die Vitalwerte werden gemessen und ein Pulsoxy angeh\"angt (RR 100/60,
HF 110, SpO2 94\%). Der Patient gibt an dass er st\"arkste
stechende/brennende Schmerzen hinter dem Brustbein hat die ins Kiefer
ausstrahlen. Ausserdem f\"uhle es sich so an als w\"urde jemand auf dem
Patienten drauf stehen. Die Schmerzen sind nicht atem- oder
bewegungsabh\"angig  und haben vor einer Viertelstunde
{\quotedblbase}wie aus dem Nichts heraus{\textquotedblleft} begonnen.
Ausserdem bekomme er schlecht Luft, sonst h\"atte er keine Beschwerden.


An welche Differentialdiagnosen denken Sie?

Leitsymptom: Thoraxschmerz 

\begin{center}
\tablehead{}\begin{supertabular}{|m{3.225cm}|m{3.225cm}|m{3.225cm}|m{3.225cm}|}
\hline
\multicolumn{4}{|m{13.5cm}|}{\centering Tabelle: Differentialdiagnosen
Thoraxschmerz unseres Patienten\par

}\\\hline
 &
\centering daf\"ur spricht\par

 &
\centering dagegen spricht\par

 &
Bewertung

\\\hline
Herzinfarkt, Angina pectoris

 &
Schmerzen hinter dem Brustbein, Ausstrahlung, nicht atemabh\"angig,
Druck-/Engegef\"uhl, Atemnot

 &
 &
sehr wahrscheinlich

\\\hline
Lungenembolie

 &
Schmerzen hinter dem Brustbein, Atemnot

 &
untypische Ausstrahlung, 

 &
gut m\"oglich

\\\hline
Pneumothorax

 &
Schmerzen im Brustbereich, Atemnot

 &
pl\"otzliche st\"arkste Schmerzen

 &
eher unwahrscheinlich

\\\hline
Trauma

 &
Schmerzen

 &
kein Hinweis

 &
unwahrscheinlich

\\\hline
Erkrankung des St\"utz- und Bewegungsapparats

 &
Schmerzen 

 &
untypische Ausstrahlung, Atemnot, nicht bewegungsabh\"angig

 &
sehr unwahrscheinlich

\\\hline
\end{supertabular}
\end{center}
Die wahrscheinlichsten Diagnosen sind der Herzinfarkt und die
Lungenembolie.

Fallen Ihnen bei diesen Erkrankungen wichtige weitere Untersuchungen
ein?

Bei beiden Erkrankungen ist es wichtig beide Beine zu begutachten. Im
Falle einer Lungenembolie sucht man nach einer Beinvenenthrombose
(einseitig \"uberw\"armtes, bl\"aulich-r\"otliches geschwollenes Bein).
Bei einer Herzinsuffizienz bzw. Herzinfarkt beurteilt man den
R\"uckstau des Blutes, dh. ob die Beine geschwollen sind.

Ausserdem, hat der Patient so etwas gehabt? Oder sind einschl\"agige
Erkrankungen bekannt?  

Sie begutachten beide Beine des Patienten und stellen mittels
Daumendruckprobe fest dass beide Beine leicht geschwollen sind. Auf
Nachfrage gibt der Patient an dass er sonst keine geschwollenen Beine
hat. 

Was sind die n\"achsten Schritte?

Die wichtigsten Untersuchungen wurden gemacht, nun muss man sich um die
Anamnese k\"ummern. 

Symptome: siehe oben.

Krankheiten: Der Patient leidet an nicht-Insulinpflichtigem Diabetes
mellitus und Bluthochdruck. Sonst sind keine Krankheiten bekannt.
Spitalsbriefe gibt es nicht. 

Sie merken sich dass Sie nachher den Blutzucker messen wollen!

Allergien: keine bekannt.

Medikamente: Glucophage 850mg f\"ur den Zucker, Tritace f\"ur den
Blutdruck. Medikamente wurden heute wie verschrieben eingenommen. 

Ereignisse vor Notfalleintritt: Schmerzen begannen pl\"otzlich als der
Patient W\"asche waschen wollte. In den letzten Wochen immer wieder
Stiche in der HErzgegend, hat er aber nicht ernst genommen. 

Letzte(r):

Spitalsaufenthalt: Vor 10 Jahren in der Rudolfstiftung wegen Rotlauf.

Arztbesuch: Vor einer Woche wegen Neuaustellung eines Rezeptes.

Mahlzeit: Mittagessen, Gulasch

Regel: Der Patient ist m\"annlich!

pro:

Lungenembolie 

Wichtig sind nun die Vitalwerte und das fr\"uhe Beginnen des
Monitorings.Dann muss man genaueres \"uber das Beschwerdebild in
Erfahrung bringen. Der weitere Plan h\"angt davon ab.

\section{[RS IV] Anatomie, Physiologie}
Sebastian Gabriel, Michael Withofner

mit Beitr\"agen von Gerhard Gabriel

Changelog:

2009-02-25: Start Changelog.

\subsection{Das Skelett und der Bewegungsapparat}
Das Skelettsystem besteht aus

\begin{itemize}
\item R\"ohren knochen (z.B. Oberarmknochen)

\item platten  Knochen (z.B. Schulterblatt,  Sch\"adelknochen)

\end{itemize}
und dient  haupts\"achlich als

\begin{itemize}
\item St\"utzger\"ust  des K\"orpers

\item Ursprungs-  und Ansatz fl\"ache f\"ur die Muskulatur

\end{itemize}
und weiters als

\begin{itemize}
\item KalZiumspeicher (Ca)

\item Blutbildungsorgan (rotes Knochenmark!)

\end{itemize}
\paragraph{Aufbau eines Knochens}
\begin{center}
\tablehead{}\begin{supertabular}{|m{5.5150003cm}|m{7.7850003cm}|}
\hline
\index{Periost}Periost: \"uberzieht den Knochen und ist sehr gut mit
Nerven versorgt.  Verantwortlich f\"ur Dickenwachstum und
Frakturheilung (Kallusbildung)

\index{Kortikalis}Kortikalis: feste \"au{\ss}ere Schicht, \"au{\ss}ere
Form

\index{Spongiosa}Spongiosa: \index{Knochenb\"alkchen}Knochenb\"alkchen
(Trabekel) ${\rightarrow}$ Zug- und Druckfestigkeit

 &
[Warning: Draw object ignored]

\\\hline
\end{supertabular}
\end{center}
\paragraph[Das Knochenmark]{Das Knochenmark}
Rotes Knochenmark ist f\"ur die Blutbildung zust\"andig

Das rote Knochenmark ist das blutbildende Organ schlechthin. Rote
Blutk\"orperchen (Erythrozyzen), wei{\ss}e Blutk\"orperchen
(Leukozyten) sowie Blutpl\"attchen (Thrombozyten) werden hier gebildet.
Bei der Geburt findet man es in s\"amtlichen Knochen, im Laufe des
Lebens bildet es sich zur\"uck und verwandelt sich in
{\quotedblbase}Gelbes Fettmark{\textquotedblleft}. Rotes Knochenmark
findet man dann lediglich in den platten Knochen wie Becken, Brustbein
und Sch\"adel wo es f\"ur die Blutbildung verantwortlich ist und
normalerweise zeitlebens erhalten bleibt. 

Einige medizinische Behandlungsmethoden verursachen als Nebenwirkung
eine massive Sch\"adigung des Knochenmarks (z.B. Bestrahlung,
Chemotherapien bei Krebspatienten ({\quotedblbase}onkologsiche
Patienten{\textquotedblleft})). Diese Patienten leiden neben einer
Blutarmut auch an einem schwachen Immunsystem, da nicht gen\"ugend
Wei{\ss}e Blutk\"orperchen gebildet werden. Sie k\"onnen daher sehr
leicht schwerwiegende Infektionen bekommen. Diese Pateinten sind im
Krankentransport h\"aufig anzutreffen. 

Gelbes Knochenmark besteht aus Fett 

Im Laufe des Lebens bildet es sich das rote Knochenmark in den
R\"ohrenknochen zu fettreichem, gelben Knochenmark um. 

Das gelbe Fettmark kann bei schweren Knochenverletzungen von Bedeutung
sein, da verletzungs\-bedingt kleine Fettptr\"opchen in die Blutbahn
gelangen k\"onnen, die dann als Fettgerinnsel  zu ernsten
Komplikationen f\"uhren k\"onnen. 

\paragraph{Knochenwachstum}
\begin{center}
\tablehead{}\begin{supertabular}{|m{6.649cm}|m{6.651cm}|}
\hline
Das Skelett eines Neugeborenen besteht bei der Geburt zum gr\"o{\ss}ten
Teil aus Knorpel. Erst mit zunehmendem Lebensalter beginnt dieser von
Knochenkernen ausgehend, zu verkn\"ochern. Eine Wachstumszone
(\index{Epiphysenfuge}Epiphysenfuge) bleibt dabei allerdings erhalten.

Die Epiphysenfuge enth\"alt beim Kind /Jugendlichen bis etwa zum Ende
der Pubert\"at teilungsf\"ahiges Knorpelgewebe, daher ist ein
L\"angenwachstum m\"oglich. Beim Erwachsenen sind diese knorpeligen
Fugen hingegen vollst\"andig verkn\"ochert, soda{\ss} keinerlei
Wachstum mehr stattfinden kann.

 &
\centering   [Warning: Image ignored]
% Unhandled or unsupported graphics:
%\includegraphics[width=3.557cm,height=4.948cm]{AASdev20090228-img22.png}
 \par

\\\hline
\end{supertabular}
\end{center}
\paragraph[Gelenke - \ prinzipieller Aufbau]{Gelenke -  prinzipieller
Aufbau}
\begin{center}
\tablehead{}\begin{supertabular}{|m{6.649cm}|m{6.651cm}|}
\hline
Bewegliche Verbindung zwischen mindestens zwei Knochen

\begin{itemize}
\item Ein Gelenk ist von einer bindegewebigen Kapsel umgeben und durch
Sehnen und B\"ander stabil verbunden.

\item Die Gelenksfl\"achen sind zur Reibungs\-ver\-minderung mit 
Knorpel \"uberzogen und der Gelenksspalt mit einer d\"unnen Schicht 
Schmierfl\"ussigkeit (Synovialfl\"ussigkeit) gef\"ullt. 

\end{itemize}
 &
\centering [Warning: Draw object ignored]\par

\\\hline
\end{supertabular}
\end{center}
\subparagraph{Es gibt verschiedene Arten von Gelenken}
Klassische Beispiele: \index{Kugelgelenk}Kugelgelenk (H\"uftgelenk),
\index{Scharniergelenk}Scharniergelenk (Ellenbogengelenk),
\index{Sattelgelenk}Sattelgelenk (Daumengrundgelenk).  [Warning: Image
ignored] % Unhandled or unsupported graphics:
%\includegraphics[width=13.701cm,height=8.482cm]{AASdev20090228-img23.png}
 

Gliederung des Skeletts

\begin{itemize}
\item Kopf  (Caput)

\end{itemize}
\begin{itemize}
\item Gehirnsch\"adel

\item Gesichtssch\"adel

\end{itemize}
\begin{itemize}
\item Rumpf  

\end{itemize}
\begin{itemize}
\item Hals 

\item Brustkorb (Thorax )

\item Bauch (Abdomen )

\item Becken (

\end{itemize}
\begin{itemize}
\item obere Extremit\"aten 

\item untere Extremit\"aten

\end{itemize}
Der Sch\"adel 

\begin{center}
\tablehead{}\begin{supertabular}{|m{6.649cm}|m{6.651cm}|}
\hline
Kopf (Caput) / Sch\"adel (Cranium)

Der Sch\"adel gliedert sich in den Hirnsch\"adel und den
Gesichtssch\"adel

\captionof{figure}[Hirnsch\"adel vs. Gesichts\-sch\"adel]{Hirnsch\"adel
vs. Gesichts\-sch\"adel}
 &
\centering [Warning: Draw object ignored]\par

\\\hline
\end{supertabular}
\end{center}
\paragraph[Gesichtssch\"adel]{Gesichtssch\"adel}
\begin{center}
\tablehead{}\begin{supertabular}{|m{6.649cm}|m{6.651cm}|}
\hline
a ... Stirnbein (Os frontale)

b ... Nasenbein (Os nasale)

c ... Scheitelbein (Os parietale)

d ... Schl\"afenbein (Os temporale)

e ... Jochbein (Os zygomaticum)

f ... Oberkiefer (Maxilla)

g ... Unterkiefer (Mandibula)

 &
[Warning: Draw object ignored]

\\\hline
\end{supertabular}
\end{center}
\paragraph{Gehirnsch\"adel}
\begin{center}
\tablehead{}\begin{supertabular}{|m{6.649cm}|m{6.651cm}|}
\hline
\begin{itemize}
\item Sch\"adelbasis

\item Sch\"adeldach (Kalotte )

\item Die platten Knochen sind durch Sch\"adeln\"ahte  (Suturae)
verzahnt miteinander verkn\"ochert

\end{itemize}
a ... Stirnbein (Os frontale)

b ... Scheitelbein (Os parietale)

c ... Schl\"afenbein (Os temporale)

d .... Hinterhauptsbein (Os occipitale)

e ... gro{\ss}er Keilbeinfl\"ugel (Os sphenoidale)

 &
[Warning: Draw object ignored]

\\\hline
\end{supertabular}
\end{center}
\subparagraph[Fontanelle -- Sutur]{Fontanelle -- Sutur}
\index{Fontanelle}\index{Sutur}\begin{center}
\tablehead{}\begin{supertabular}{|m{8.6640005cm}|m{4.642cm}|}
\hline
Die Fontanellen sind knorpelige  Verbindungen  zwischen 
Sch\"adelknochen beim Neu\-geborenen (der Sch\"adel mu{\ss} sich f\"ur 
Geburt verformen k\"onnen!). Sie schlie{\ss}en  sich bis zum 18. LM,
verkn\"ochern  aber erst vollst\"andig mit zunehmendem  Lebensalter
(ca. 20 Lj. abgeschlossen). 

\begin{itemize}
\item gro{\ss}e Fontanelle (rot): zw. Stirnbein u. Scheitelbeinen

\item kleine Fontanelle (gr\"un): zw. Scheitelbeinen u.
Hinter\-haupts\-bein

\end{itemize}
 &
\centering [Warning: Draw object ignored]\par

\\\hline
\end{supertabular}
\end{center}
Die Sutur ist die endg\"ultig fest verkn\"ocherte Naht zwischen den
Sch\"adelknochen beim Erwachsenen ${\rightarrow}$ starrer  stabiler 
Sch\"adel, Schutz  f\"ur Gehirn.

\paragraph[Sch\"adelbasis]{Sch\"adelbasis}
\begin{center}
\tablehead{}\begin{supertabular}{|m{6.649cm}|m{6.651cm}|}
\hline
a ... Stirnbein (Os frontale)

b ... Siebbein (Os ethmoidale)

c ... Keilbein  (Os sphenoidale)

d ... Schl\"afenbein(e)  (Os temporale)

e ... Felsenbein  (Pars petrosa ossis temporalis)

f ... Hinterhauptsbein  (Os occipitale)

g ... gro{\ss}es Hinterhauptsloch (Foramen magnum)

 &
\centering [Warning: Draw object ignored]\par

\\\hline
\end{supertabular}
\end{center}
Gehirn liegt auf der Sch\"adelbasis auf (ausser wenn man auf dem Kopf
steht). 

Hirnnerven treten durch div. L\"ocher aus dem Sch\"adel aus.

Das gro{\ss}e Hinterhauptsloch (Foramen magnum) ist die
Durchtrittsstelle des R\"uckenmarks. 

\index{Hirnstammeinklemmung}Hirnstammeinklemmung

Das gro{\ss}e Hinterhauptsloch (Foramen magnum) ist die
Durchtrittsstelle des R\"uckenmarks. Bei einer Hirndrucksteigerung
(Hirn\"odem, Hirnblutung) ist es die einzige M\"oglichkeit zur
Ausdehnung des Gehirns! Dabei dr\"uckt sich das Gehirn allerdings
selber ab, es kommt zu einer t\"odlichen Hirnstammeinklemmung! 

${\rightarrow}$  Lebensgefahr!  

Die Wirbels\"aule (Columna vertebralis)

Bild von Wirbels\"aule gest\"ort!

[Warning: Draw object ignored]%
%
Die Wirbels\"aule besteht aus:

\begin{itemize}
\item 7 Halswirbeln  (Cervical-Wirbel)

\item 12 Brustwirbeln  (Thorakal-Wirbel)

\item 5 Lendenwirbeln  (Lumbal{}-Wirbel)

\item 5 verschmolzene Kreuzbein wirbel (Os sacrum)

\item 3-4 Stei{\ss}bein {\textquotedblleft}wirbel{\textquotedblleft} 
(Os coccygis) 

\end{itemize}
Doppelt S-f\"ormig gekr\"ummtes Achsenskelett

\begin{itemize}
\item erm\"oglicht den aufrechten Gang

\end{itemize}
besteht aus

\begin{itemize}
\item Wirbeln (vertebrae)

\item \index{Bandscheiben}Bandscheiben (Discus, Pl: Disci
intervertebrales): dienen der
{\quotedblbase}Sto{\ss}d\"ampfung{\textquotedblleft}

\item Im Wirbelkanal liegt das R\"uckenmark eingebettet.

\end{itemize}
\subparagraph{Die Wirbel: Aufbau und Arten}
Prinzipieller Aufbau:

\begin{itemize}
\item Wirbelk\"orper

\item Wirbelbogen  (umschlie{\ss}t das R\"uckenmark)

\item 2 Querforts\"atze

\item 1 Dornfortsatz

\end{itemize}
%
%
In den entsprechenden Abschnitten sehen die Wirbel unterschiedlich aus. 

Bild von Wirbelk\"orper hier einf\"ugen!

Der \index{Brustkorb}Brustkorb (\index{Thorax}Thorax)

Der  Brustkorb besteht vorne aus dem Brustbein (Sterrnum) und 12
Rippenpaaren, welche hinten mit der Wirbels\"aule in gelenkiger
Verbindung stehen. Die obersten 7 Rippen-Paare werden als 
{\quotedblbase}echte Rippen{\textquotedblleft} bezeichnet. Sie
erreichen mit ihren Spitzen das Brustbein und bilden praktisch
{\quotedblbase}Ringe{\textquotedblleft} um die Brusteingeweide. Die 5
unteren Rippenpaare bilden den bogenf\"ormigen unteren Rand des
Brustkorbes. Die Rippenpaar 8-10 sind dabei nur indirekt \"uber einen
Knorpel mit dem Brustbein verbunden ({\quotedblbase}unechte
Rippen{\textquotedblleft}), das 11. und das 12. Paar endet mit
knorpeligen Spitzen frei zwischen den
Bauchmuskeln({\quotedblbase}fliehende Rippen{\textquotedblleft})
bezeichnet werden.

Der Brustkorb enth\"alt und sch\"utzt die Brust\-organe (Lunge, Herz)
sowei die Oberbauchorgane (Leber, Magen). Zwischen den Rippen spannt
sich die Zwischenrippenmuskulatur (wichtig f\"ur Atemmechanik), in
diese sind segmentale Nerven- und Gef\"a{\ss}b\"undel eingelagert.

\begin{center}
\tablehead{}\begin{supertabular}{|m{6.649cm}|m{6.651cm}|}
\hline
  [Warning: Image ignored] % Unhandled or unsupported graphics:
%\includegraphics[width=4.388cm,height=5.019cm]{AASdev20090228-img24.png}
 

 &
Der Thorax besteht aus:

\begin{itemize}
\item Brustbein  (\index{Sternum}Sternum)

\item 12  Rippenpaaren 

\end{itemize}
\begin{itemize}
\item 7 {\quotedblbase}echte{\textquotedblleft}  

\item 3 {\quotedblbase}unechte{\textquotedblleft}

\item 2 {\quotedblbase}fliegenden{\textquotedblleft}

\end{itemize}
\begin{itemize}
\item Brustwirbels\"aule (12 Brust\-wirbel), sie Steht mit den Rippen in
gelenkiger Verbindung 

\end{itemize}
({\quotedblbase}echt{\textquotedblleft} bedeutet die Rippe hat eine
direkte Verbindung zum Sternum)

\\\hline
\end{supertabular}
\end{center}
Der \index{Schulterg\"urtel}Schulterg\"urtel

\begin{center}
\tablehead{}\begin{supertabular}{|m{6.649cm}|m{6.651cm}|}
\hline
... besteht aus (von vorne nach hinten) :

\begin{itemize}
\item \index{Brustbein}Brustbein  (\index{Sternum}Sternum)

\item \index{Schl\"usselbein}Schl\"usselbein  (Clavicula)

\item \index{Schulterblatt}Schulterblatt  (Scapula)

\end{itemize}
Der Schulterg\"urtel bildet die Gelenkspfanne f\"ur das Schultergelenk,
jeweils auf beiden Seiten.

Die Schulterbl\"atter sind nur durch Muskulatur mit dem Brustkorb
verbunden.

 &
\centering   [Warning: Image ignored]
% Unhandled or unsupported graphics:
%\includegraphics[width=4.379cm,height=4.453cm]{AASdev20090228-img25.png}
 \par

\\\hline
\end{supertabular}
\end{center}
Die obere Extremit\"at

\begin{center}
\tablehead{}\begin{supertabular}{|m{6.649cm}|m{6.651cm}|}
\hline
\begin{itemize}
\item Schultergelenk

\item Oberarmknochen  (Humerus)

\item Ellenbogengelenk

\item Unterarmknochen (2)

\end{itemize}
\begin{itemize}
\item Elle  (Ulna) 

\item Speiche  (Radius), daumenseitig!

\end{itemize}
\begin{itemize}
\item Handwurzelgelenk

\item (8) Handwurzelknochen

\item (5) Mittelhandknochen 

\item Finger: je 3  Endglieder

\item Daumen: nur 2  Endglieder

\end{itemize}
 &
[Warning: Draw object ignored]

\\\hline
\end{supertabular}
\end{center}
Achtung: Unterscheide: Hand -- Arm!

Der Beckeng\"urtel

\begin{center}
\tablehead{}\begin{supertabular}{|m{6.649cm}|m{6.651cm}|}
\hline
Besteht aus (von hinten nach vorne) :

\begin{itemize}
\item Kreuzbein (Os sacrum)

\item Beckenknochen, bestehend aus:

\end{itemize}
\begin{itemize}
\item \index{Darmbein}Darmbein (Os ilii)

\item \index{Sitzbein}Sitzbein (Os ischii)

\item \index{Schambein}Schambein (Os pubis)

\end{itemize}
Die \index{Symphyse}Symphyse stellt eine straffe Knorpelverbindung vorne
zwischen den beiden Beckenknochen dar (Dehnungs\-fuge f\"ur Geburt).

 &
\centering   [Warning: Image ignored]
% Unhandled or unsupported graphics:
%\includegraphics[width=5.527cm,height=3.555cm]{AASdev20090228-img26.png}
 \par

\\\hline
\end{supertabular}
\end{center}
Die untere Extremit\"at

\begin{center}
\tablehead{}\begin{supertabular}{|m{6.649cm}|m{6.651cm}|}
\hline
\begin{itemize}
\item Oberschenkel  (Femur)

\item Unterschenkelknochen (2)

\end{itemize}
\begin{itemize}
\item Schienbein  (Tibia) , dick

\item Wadenbein  (Fibula), d\"unn

\end{itemize}
\begin{itemize}
\item 8 Fu{\ss}wurzelknochen

\item 5 Mittelfu{\ss}knochen  (Ossa metatarsalia)

\item 5 Zehen:

\end{itemize}
\begin{itemize}
\item 4 Zehen: je 3  Endglieder

\item 1 Gro{\ss}zehe: nur 2 Endglieder

\end{itemize}
 &
\centering %
%
[Warning: Draw object ignored]\par

\\\hline
\end{supertabular}
\end{center}
Untere Extremit\"at Bild muss ausgetuascht werden.

Achtung: Unterscheide Fu{\ss} -- Bein!

Die Muskulatur

Prinzipiell gibt es 3  Arten von Muskelzellen:

\begin{itemize}
\item Quergestreifte (Skelett-)Muskulatur: willk\"urlich steuerbar,
trainierbar, er\-m\"ud\-bar

\item Glatte  Muskulatur: Nicht willentlich kontrollierbar,
unerm\"udlich, ausdauernd 

\item Herz muskel: Sonderstellung zwischen glatt und quergestreift,
Eigenschaften von beiden

\end{itemize}
Die quergestreifte Muskulatur heisst deswegen
{\quotedblbase}quergestreift{\textquotedblleft}, da man unter dem
Mikroskop tats\"achlich eine Streifung sehen kann, allerdings nicht mit
dem freien Auge) 

\begin{center}
\tablehead{}\begin{supertabular}{|m{6.649cm}|m{6.651cm}|}
\hline
Jeder Muskel hat eine vorgegebene  Zugrichtung. 

Muskeln k\"onnen sich nicht aktiv  entspannen, daher sind Gegenspieler 
(Agonist u. Antagonist) erforderlich (z.B. ein Muskel streckt den Arm,
ein anderer beugt ihn).

 &
%
%
[Warning: Draw object ignored]

\\\hline
\end{supertabular}
\end{center}
Ellbogen -- Muskelzug: Bild ersetzen (grau)

\subsection{Das Herz}
Das Herz ist ein etwa faustgro{\ss}er Hohlmuskel und hat die Funktion
einer Pumpe. Es liegt hinter dem Sternum\footnote{Brustbein} und
teilweise im linken Brustkorb. Unten liegt es auf dem Zwerchfell auf. 

\begin{center}
 [Warning: Image ignored] % Unhandled or unsupported graphics:
%\includegraphics[width=4.887cm,height=6.693cm]{AASdev20090228-img27.png}

\end{center}
Das Herz ist durch eine Scheidewand (Septum) in eine rechte und eine
linke Seite geteilt. Umgangssprachlich spricht man dabei jeweils vom
{\quotedblbase}rechten{\textquotedblleft} oder {\quotedblbase}linken
Herz{\textquotedblleft}. Jede H\"alfte wird durch Klappen in einen
Vorhof (Atrium) und eine Herzkammer (Ventrikel) geteilt. 

Herzklappen

Insgesamt 4 Herzklappen sorgen als Ventile daf\"ur dass das Blut nur in
eine Richtung str\"omt. Es gibt 2 Segelklappen und 2 Taschenklappen.
Die Segelklappen befinden sich zwischen den Vorh\"ofen und den
Herzkammern, die Taschenklappen befinden sich an den
Ausfluss\"offnungen der Kammern am \"Ubergang zur Lungenarterie (rechte
Seite) bzw. zur Hauptschlagader (linke Seite). 

Blutversorgung

Die Blutversorgung erfolgt \"uber die Herzkranzgef\"a{\ss}e
(Koronargef\"a{\ss}e) welche direkt aus der Hauptschlagader abgehen. 

Schichten

3 Schichten (au{\ss}en nach innen):

\begin{itemize}
\item Epikard (aussen)

\item Myokard (Muskel)

\item Endokard (innen)

\end{itemize}
\begin{minipage}{13.6cm}
 [Warning: Image ignored] % Unhandled or unsupported graphics:
%\includegraphics[width=13.6cm,height=6.168cm]{AASdev20090228-img28.jpg}
\captionof{figure}[Schema der Herz-Anatomie. ]{Schema der Herz-Anatomie.
}
Quelle: Wikipedia,  J\"org Rittmeister. G-FDL, CCAL

\end{minipage}

Herzfunktion

Der Herzmuskel kontrahiert und entspannt sich beim Erwachsenen 60 bis
100 mal in der Minute. Bei jeder Anspannung (Kontraktion) wird das Herz
ca. 70ml Blut in die Arterien aus. 

Das Herz pumpt in Ruhe ca. 4l Blut in der Minute. 

Blutdruck: systolisch/diastolisch  (Anspannung/Entspannung)

%
%
Reizleitungssystem

@ Johannes: Reizleitungssystem -- wie ausf\"uhrlich? Pr\"ufungsfrage?

\begin{center}
\tablehead{}\begin{supertabular}{|m{6.649cm}|m{6.651cm}|}
\hline
Eigenst\"andiges  Erregungs- und  Reizleitungssystem, streng
hierarchisch, aber jede  Station zur Erregungsbildung  f\"ahig 

bestehend aus:

\begin{itemize}
\item Sinusknoten  ({\quotedblbase}Schrittmacher{\textquotedblleft},
feuert mit einer Eigenfrequenz v. ca. 60/min. elektrische Impulse ab)

\item AV-(AtrioVentricular)-Knoten: Leitung vom Vorhof in die Kammern,
filtert zu hohe Frequenzen

\item His-B\"undel

\item Tawara-Schenkel (linker/rechter)

\item Purkinje-Fasern ({\quotedblbase}Endstrecke{\textquotedblleft})

\end{itemize}
 &
\centering   [Warning: Image ignored]
% Unhandled or unsupported graphics:
%\includegraphics[width=5.673cm,height=6.587cm]{AASdev20090228-img29.png}
 \par

\\\hline
\end{supertabular}
\end{center}
Der normale Schrittmacher ist der Sinusknoten, welcher die elektrischen
Impulse f\"ur die Vorh\"ofe und die Kammern erzeugt. Man spricht dann
vom Sinusrhythmus.

Herzrhythmus

m\"ogliche Rhythmen:

Sinusrhythmus: ({\textless}60/min ${\rightarrow}$ Sinusbradykardie,
{\textgreater}100/min ${\rightarrow}$ Sinustachy\-kardie):
Sinus\-knoten ist Schrittmacher, vor jedem QRS ein P, Abstand P-QRS
konstant, Abstand QRS {}-- QRS konstant, Abstand P -- P konstant

Arrhythmie: ({\textless}60/min ${\rightarrow}$ Brady arrhythmie,
{\textgreater}100/min ${\rightarrow}$ Tachyarrhythmie): ev. wechselnde
Schrittmacherzentren, nicht vor jedem QRS ein P, Abst\"ande QRS -- QRS
variabel

\subsection{Der Kreislauf }
Herz + Gef\"a{\ss}system + Blut = Kreislauf

Es wird der K\"orperkreislauf ({\quotedblbase}Gro{\ss}er
Kreislauf{\textquotedblleft},
{\quotedblbase}Hochdrucksystem{\textquotedblleft}) vom Lungenkreislauf
({\quotedblbase}Kleiner Kreislauf{\textquotedblleft},
{\quotedblbase}Niederdrucksystem{\textquotedblleft}) unterschieden. Das
Herz ist die f\"ur beide Systeme die gemeinsame Pumpe.

Kreislaufsteuerung

Das Stammhirn  steuert den Kreislauf  (RR + HF )*, Blutdruckrezeptoren
in den  Gef\"a{\ss}w\"anden (z.B.Carotis)

\begin{center}
\begin{tabular}{|m{5.328cm}m{7.977cm}}
\hline
\multicolumn{1}{|m{5.328cm}|}{\subparagraph{Kleiner (Lungen-)Kreislauf}
Die Hohlvenen m\"unden in den rechten Vorhof des Herzens und liefern
sauerstoffarmes Blut aus dem K\"orper an. Von hier wird das Blut in die
rechte Herzkammer gepumpt, um dann in die Lungenarterie zu gelangen. In
der Lunge ver\"asteln sich die Lungenarterien zu haard\"unnen
Gef\"a{\ss}en, den Kapillaren. Diese umspannen die Lungenbl\"aschen es
findet ein Gasaustausch statt: Sauerstoff wird aufgenommen und
Kohlendioxid abgegeben. Das jetzt mit Sauerstoff angereicherte Blut
wird \"uber die Lungenvenen in den linken Vorhof geleitet, der wiederum
das Blut in die linke Herzkammer pumpt. 

} &
\multicolumn{1}{m{7.977cm}|}{

\begin{center}
 [Warning: Image ignored] % Unhandled or unsupported graphics:
%\includegraphics[width=8.026cm,height=7.416cm]{AASdev20090228-img30.png}

\end{center}
\begin{center}
 [Warning: Image ignored] % Unhandled or unsupported graphics:
%\includegraphics[width=7.55cm,height=6.82cm]{AASdev20090228-img31.png}

\end{center}
Rechter Vorhof ${\rightarrow}$ rechte Herzkammer ${\rightarrow}$
Lungenarterie(n) ${\rightarrow}$ Lungenkappilaren (CO2-Abgabe,
O2-Aufnahme) ${\rightarrow}$ Lungenvenen ${\rightarrow}$ linker Vorhof 

${\rightarrow}$ linke Herzkammer ${\rightarrow}$ Aorta ${\rightarrow}$
abzweigende Arterien ${\rightarrow}$ Kapillaren (O2-Abgabe,
CO2-Aufnahme)  ${\rightarrow}$ Venen ${\rightarrow}$ Hohlvene
${\rightarrow}$ Rechter Vorhof

 ${\rightarrow}$ alles beginnt wieder von vorne

}\\\hline
\subparagraph{Gro{\ss}er Kreislauf}
Das Blut wird aus der linken Herzkammer ausgesto{\ss}en und gelangt
\"uber die die Hauptschlagader
({\quotedblbase}Aorta{\textquotedblleft}) in den K\"orperkreislauf. Aus
der Aorta entspringen Arterien
({\quotedblbase}Schlagadern{\textquotedblleft}), die sich immer weiter
verzweigen und an Durchmesser abnehmen. Sie versorgen den gesamten
K\"orper mit Blut. Der durch die Pumpleistung des Herzens erzeugte Puls
ist an oberfl\"achlich gelegenen Arterien tastbar (Radialis-Puls,
Carotis-Puls, ..). Daher r\"uhrt auch der Name
{\quotedblbase}Schlagader{\textquotedblleft}. 

Sie verzweigen sich, weiter an Durchmesser abnehmend, \"ahnlich dem
Ge\"ast eines Baumes, bishin zu Arteriolen und haard\"unnen
Gef\"a{\ss}en, den Haargef\"a{\ss}en
({\quotedblbase}Kapillaren{\textquotedblleft}). Hier findet wieder ein
Gasaustausch statt, diesmal allerdings spiegelbildlich zum
Lungenkreislauf: Sauerstoff wird an das Gewebe abgegeben, im Gegenzug
wird Kohlendioxid abtransportiert. 

Das nun sauerstoffarme Blut gelangt \"uber kleine Venolen in das
Venensystem und flie{\ss}t in den Venen herzw\"arts. Am Weg zum Herzen
nehmen die Venen an Umfang immer mehr zu, bis sie letztendlich in die
gro{\ss}en Venen ({\quotedblbase}obere und untere
Hohlvene{\textquotedblleft}) m\"unden. 

 &
\\\hhline{-~}
\end{tabular}
\end{center}
Klappen fehlen in der Beschreibung!

Wichtige Gef\"a{\ss}e

\subparagraph{Gef\"a{\ss}e}
Arterien: vom Herzen weg, dicke Muskelschicht

Venen: zum Herzen hin, Venenklappen

Kapillaren: d\"unne {\quotedblbase}Haargef\"a{\ss}e{\textquotedblleft}
in denen der Gasaustausch stattfindet, Verbindung zwischen Arterien und
Venen

Arterien sind durch eine dicke Muskelschicht gekennzeichnet. Venen haben
dagegen eine deutlich d\"unnere Wand als die muskul\"osen Arterien%
%


\begin{center}
\tablehead{}\begin{supertabular}{|m{2.728cm}|m{2.728cm}|m{2.728cm}|m{4.721cm}|}
\hline
\multicolumn{4}{|m{13.505cm}|}{\centering \"Ubersicht wichtiger
Blutgef\"a{\ss}e\par

}\\\hline
Hauptschlagader

 &
Aorta

 &
Arterie

 &
Geht von der linken Herzkammer ab, von dort zweigen alle anderen
K\"orperarterien ab

\\\hline
Halsschlagader

 &
Karotis

 &
Arterie

 &
1 Paar (2 Stk.), versorgen den Sch\"adel, Puls tastbar seitlich vom
Kehlkopf

\\\hline
Herzkranzgef\"a{\ss}e

 &
Koronararterien

 &
Arterie

 &
Versorgen das Herz mit Blut, gehen aus der Aorta ab

\\\hline
Lungenarterie

 &
Pulmonalarterie

 &
Arterie -- O2{}-arm!

 &
Geht von der rechten Herzkammer weg in die Lunge zu den Lungenkapillaren

\\\hline
Oberarmarterie

 &
Brachialarterie

A. brachialis

 &
Arterie

 &
Taststelle zum Puls-f\"uhlen beim Baby, Blutdruckmessung

\\\hline
Radialarterie

 &
A. radialis

 &
Arterie

 &
Verl\"auft am Unterarm speichenseitig (dort wo der Daumen ist). Puls
f\"uhlbar.

\\\hline
Obere Hohlvene

 &
Vena cava sup.

 &
Vene

 &
Obere Vene die in den rechten Vorhof m\"undet

\\\hline
Untere Hohlvene

 &
Vena cava inf.

 &
Vene

 &
Untere Vene die in den rechten Vorhof m\"undet. 

\\\hline
Lungenvene

 &
V. pulmonalis

 &
Vene -- O2{}-reich! 

 &
Kommt von den Lungenkapillaren und m\"undet in den linken Vorhof

\\\hline
Pfortader

 &
Vena portae

 &
Vene

 &
F\"uhrt O2-armes, aber n\"ahrstoffreiches Blut von Teilen des Darms zur
Leber (zur Entgiftung)

\\\hline
Haargef\"a{\ss}e

 &
Kapillare

 &
Kapillare

 &
D\"unne, dicht verzwweigte Gef\"a{\ss}e, Verbindung zwischen Arterien
und Venen. Hierfindet der Gas- und N\"ahrstoffaustausch statt. Es gibt
K\"orperkapillaren (K\"orperkreislauf) und Lungenkapillaren
(Lungenkreislauf)

\\\hline
\end{supertabular}
\end{center}
\paragraph{Blutdruck }
Abk\"urzung: RR. Gemessen immer in Herzh\"ohe!

systolisch:  Maximaldruck in der Arterie am H\"ohepunkt der
Kammerkontraktion

diastolisch: Druck in der Arterie w\"ahrend die Kammer erschlafft ist,
beeinflu{\ss}t durch  Gef\"a{\ss}wandspannung

Normbereich : zwischen 90/60 mmHg und 140/90 mmHg  

Hypotonie:  {\textless} 90/60 mmHg

Hypertonie: {\textgreater} 140/90mmHg  

Achtung: Eine Hypertonie ist zwar eine Erkrankung, aber nicht immer im
Einsatz von Bedeutung. Viele Patienten sind an ihren erh\"ohten
Blutdruck gewohnt und haben keine Beschwerden. Es darf nie nur nach dem
Wert gehandelt werden, es muss auch immer bedacht werden: 

\begin{itemize}
\item Zustand des Patienten.

\item Der sonst \"ubliche Blutdruck des Patienten (erfragen,
Blutdruckprotokoll).

\item Der Verlauf des Blutdruckes w\"ahrend der Behandlung
(Traumapatienten !!!).

\item Der Erregungszustand des Patienten.

\end{itemize}
Beispiele

Ein 65-j\"ahriger, m\"annlicher Patient klagt \"uber Schwindelgef\"uhl,
sie messen einen Blutdruck von 100/70. Im Blutdruckprotokoll finden sie
f\"ur die vergangenen vier Tage folgende Eintr\"age: 170/110, 180/115,
175/113, 183/112, ... .

Wie beurteilen Sie den Blutdruck? 

Der Blutdruck ist f\"ur den Patienten zu niedrig, der K\"orper ist auf
eine so pl\"otzliche Umstellung nicht vorbereitet gewesen. Im
Anamnesegespr\"ach erfahren sie dass er heute das erste Mal einen
{\quotedblbase}Nitro-Spray{\textquotedblleft} verwendet hat. Sie
vermerken dies am Einsatzprotokoll. Im Krankenhaus erfahren Sie dass
wahrscheinlich eine \"Uberdosierung dieses Medikaments die Ursache
dieses Blutdruckabfalls war. ${\rightarrow}$ Obwohl der Patient laut
Lehrbuch einen {\quotedblbase}sch\"onen{\textquotedblleft} Blutdruck
hatte, war dieser dennoch zu niedrig!

Ein 56-j\"ahriger Patient klagt \"uber starke kolikartige Schmerzen im
linken Unterbauch. Sie messen einen Blutdruck von 155/95. Der Patient
kennt seinen {\quotedblbase}gew\"ohnlichen{\textquotedblleft} Blutdruck
nicht.  

Wie beurteilen Sie den Blutdruck? 

Der Blutdruck ist laut Lehrbuch erh\"oht, aber dies ist erkl\"arbar die
starken Schmerzen und den damit verbundenen Erregungszustand. Die
Beschwerden stehen wahrscheinlich in keinem urs\"achlichen Zusammenhang
mit den Schmerzen. 

Eine 50-j\"ahrige Patientin hat Nasenbluten. Sie messen einen Blutdruck
von 220/120. Auf Nachfrage gibt sie an etwas Kopfsschmerzen zu haben,
und etwas schwindlig sei sie auch. 

Wie beurteilen Sie den Blutdruck? 

Der Blutdruck ist sehr stark erh\"oht, ausserdem zeigt die Patientin
Zeichen einer hypertensiven Krise (Nasenbluten, Kopfweh, Schwindel).
Hier ist der erh\"ohte Blutdruck Ursache f\"ur die Beschwerden und muss
mittels Medikamente unbedingt gesenkt werden. 

Das \index{Blut}Blut

ca. 7-8\% des K\"orpergewichts = 5-6l

besteht aus

\begin{itemize}
\item festen Bestandteilen (Zellen, Blutk\"orperchen)

\item \index{Plasma}Plasma

\end{itemize}
  [Warning: Image ignored] % Unhandled or unsupported graphics:
%\includegraphics[width=13.664cm,height=5.432cm]{AASdev20090228-img32.png}
 

Blut - Aufgaben

\begin{itemize}
\item Sauerstofftransport zu den K\"orperzellen

\item CO2-Abtransport

\item Transport von N\"ahrstoffen und Abfallprodukten

\item W\"armeverteilung

\item pH-Pufferungssystem (S\"aure/Basen-Gleichgewicht) (pH-Wert mu{\ss}
konstant zwischen 7,35 und 7,45 liegen!)

\end{itemize}
Zellul\"are Bestandteile

\begin{itemize}
\item Rote Blutk\"orperchen (Erythrozyten)

\item Wei{\ss}e Blutk\"orperchen (Leukozyten)

\item Blutpl\"attchen (Thrombozyten)

\end{itemize}
\begin{multicols}{2}
\index{Erythrozyten}Erythrozyten 

\begin{itemize}
\item rote Blutk\"orperchen

\item O2 - Transport

\item tragen Blutgruppenmerkmale (Blutgruppe)

\item seeeeeehr viele, die h\"aufigsten Blutk\"orperchen (4.000.000/mm3
)

\end{itemize}
\index{Leukozyten}Leukozyten 

\begin{itemize}
\item wei{\ss}e Blutk\"orperchen 

\item K\"orperabwehr, Entz\"undungszellen

\item (4.000 -- 8.000/mm3 )

\end{itemize}
\index{Thrombozyten}Thrombozyten 

\begin{itemize}
\item \index{Blutpl\"attchen}Blutpl\"attchen

\item wichtige Rolle beim Wundverschlu{\ss}

\item (150.000 -- 300.000/mm3 )

\item Blutgerinnung

\end{itemize}
Das \index{Plasma}Plasma

\begin{itemize}
\item 10\% Proteine, Enzyme, Hormone, Elektrolyte, Glucose, Fette

\item 90\% H2O

\item Aufgaben:

\end{itemize}
\begin{itemize}
\item N\"ahrstofftransport

\item W\"armeverteilung

\end{itemize}
\begin{itemize}
\item enth\"alt Proteine und Enzyme f\"ur  Blutgerinnung (Serum )

\item Pufferung

\end{itemize}
\end{multicols}
\subsection{Das Nervensystem}
Unterteilung

Man kann das Nervensystem nach seinem Aufbau und nach seiner Funktion
unterteilen.

\begin{center}
\begin{minipage}{4.001cm}
NEU 03.12.08

\end{minipage}
\end{center}
Vom Blickwinkel des Aufbaus wird das Nervensystem grob unterteilt in ein


\begin{itemize}
\item Zentrales Nervensystem (ZNS) und ein 

\item Peripheres Nervensystem.

\end{itemize}
Das ZNS ist im wesentlichen f\"ur Steuer-, Regel- und Denkfunktionen
zust\"andig sowie Teilweise f\"ur die Weiterleitung von Impulsen und
Informationen. Das periphere Nervensystem ist fast ausschliesslich nur
f\"ur die Weiterleitung von Informationen und Impulsen zu den
Endorganen (oder umgekehrt, von den Endorganen zum ZNS) zust\"andig). 

Betrachtet man die Funktion kann man das periphere Nervensystem in ein 

\begin{itemize}
\item Motorisches (willk\"urliche Bewegung),

\item Sensorisches und ein

\item Vegetatives (autonomes) Nervensystem

\end{itemize}
unterteilen. 

Das Zentralnervensystem

\begin{center}
\tablehead{}\begin{supertabular}{|m{6.649cm}|m{6.651cm}|}
\hline
Das Zentralnervensystem (ZNS) besteht aus:

\begin{itemize}
\item Gehirn

\end{itemize}
\begin{itemize}
\item Gro{\ss}hirn

\item Kleinhirn

\item Hirnstamm

\end{itemize}
\begin{itemize}
\item R\"uckenmark

\end{itemize}
Aufgaben:

\begin{itemize}
\item Aufnahme  und Verarbeitung  von \"au{\ss}eren und internen  Reizen
(Schmerz, F\"uhlen, Sehen, H\"oren, {\dots})

\item Bewusstsein und {\quotedblbase}Bewusstmachung{\textquotedblleft}

\item Abgabe  von zweckentsprechenden Impulsen  (Steuerung, Reaktion
etc.)

\end{itemize}
 &
\centering [Warning: Draw object ignored]\par

\\\hline
\end{supertabular}
\end{center}
%
%
ZNS - Aufgaben

\begin{center}
\begin{minipage}{4.001cm}
ZNS: ein bissl umgestellt:

Bild

ZNS Aufgaben

Anatomie

{}- Querschnitt Hirnsch\"adel

{}- Aufbau ZNS mit den versch Hirnen und Liquor

{}- Blut-Hirn-Schranke

\end{minipage}
\end{center}
\index{Gro{\ss}hirn}Gro{\ss}hirn

\begin{itemize}
\item bewu{\ss}tes Denken

\item oberste Instanz

\end{itemize}
\index{Zwischenhirn}Zwischenhirn

\begin{itemize}
\item Schaltinstanz und Filter R\"uckenmark - Gro{\ss}hirn

\end{itemize}
\index{Kleinhirn}Kleinhirn

\begin{itemize}
\item Bewegungskoordination ({\quotedblbase}Makros{\textquotedblleft})

\end{itemize}
\index{Hirnstamm}Hirnstamm und \index{verl\"angertes
R\"uckenmark}verl\"angertes R\"uckenmark

\begin{itemize}
\item Atem-, Kreislaufzentrum

\end{itemize}
\index{R\"uckenmark}R\"uckenmark

\begin{itemize}
\item Weiterleitung von Impulsen

\item Vom R\"uckenmark gehen die peripheren Nerven aus

\item (Muskeleigen-)Reflexe

\end{itemize}
\subparagraph{Anatomie}
\begin{center}
\tablehead{}\begin{supertabular}{|m{4.346cm}|m{8.988cm}|}
\hline
Querschnitt Gehirnsch\"adel

\begin{itemize}
\item Kopfschwarte (Galea)

\item Knochen

\item \index{Epiduralraum}Epiduralraum

\item \index{Dura mater}Dura mater (\index{Harte Hirnhaut}harte
Hirnhaut)

\item \index{Subduralraum}Subduralraum

\item \index{Arachnoidea}Arachnoidea (Spinngewebshaut)

\item \index{Subarachnoidalraum}Subarachnoidalraum

\item eingebettete Gef\"a{\ss}e

\item Pia mater (weiche Hirnhaut)

\end{itemize}
 &
  [Warning: Image ignored] % Unhandled or unsupported graphics:
%\includegraphics[width=9.035cm,height=6.869cm]{AASdev20090228-img33.png}
 

\\\hline
\end{supertabular}
\end{center}
\begin{itemize}
\item \end{itemize}
Aufbau ZNS

\begin{itemize}
\item Gro{\ss}hirn (Telencephalon)

\end{itemize}
\begin{itemize}
\item \index{Graue Substanz}Graue Substanz (Rinde)

\item \index{Wei{\ss}e Substanz}Wei{\ss}e Substanz (Mark)

\item Balken (Verbindung li/re)

\end{itemize}
\begin{itemize}
\item Zwischenhirn  (Diencephalon)

\end{itemize}
\begin{itemize}
\item Thalamus, Hypothalamus (Hormone, Temperaturregulation etc.)

\item Hypophyse (Hormone)

\end{itemize}
\begin{itemize}
\item Hirnstamm (Truncus encephali)

\item Mittelhirn (Mesencephalon)

\item Br\"ucke (Pons)

\item verl\"angertes Mark (Medulla oblongata)

\item Ventrikelsystem

\end{itemize}
\begin{itemize}
\item Liquorgef\"ullte Hohlr\"aume (Ventrikel)

\item haben nichts mit den Vnetrikel im Herz zu tun!

\end{itemize}
Liquor = Gehirnwasser, in Ventrikeln gebildet

Das Hirn und das R\"uckenmark schwimmt im Liquor, dadurch ist es ein
wenig gesch\"utzt. Tritt Liquor in die Umwelt aus (Nase, Ohr, ...)  ist
das ein Zeichen f\"ur eine schwerwiegende Verletzung. 

Tritt Liquor in die Umwelt aus ist das ein Zeichen f\"ur eine
schwerwiegende Verletzung. 

Blut-Hirn-Schranke 

Die Blutgef\"a{\ss}wand ist im Hirn besonders dicht. Sie filtert daher
besonders gut und nur f\"ur bestimmte Molek\"ule (O2,CO2, Wasser,...)
leicht passierbar. Sie bietet dem Hirn Schutz vor sch\"adlichen und
giftigen Substanzen. 

Die Blut-Hirn-Schranke bietet dem Hirn besonderen Schutz vor
sch\"adlichen Substanzen im Blut. 

PNS -- Peripheres Nervensystem

\begin{center}
\tablehead{}\begin{supertabular}{|m{6.649cm}|m{6.651cm}|}
\hline
\begin{itemize}
\item Kommunikation ZNS ${\leftarrow}$${\rightarrow}$ periphere  Organe

\item Spinalnerven (RM, segmental)

\item Hirnnerven (12 Nervi craniales)

\item elektrische Impulse, chemische Signale (Synapsen)

\end{itemize}
Schnell: 30-70 m/sec

Arten von Nerven

\begin{itemize}
\item motorisch

\item sensorisch

\item vegetativ

\end{itemize}
\begin{itemize}
\item sympathisch

\item parasympathisch

\end{itemize}
 &
\centering   [Warning: Image ignored]
% Unhandled or unsupported graphics:
%\includegraphics[width=5.28cm,height=4.405cm]{AASdev20090228-img34.png}
 \par

R\"uckenmark mit abgehenden peripheren Nerven. Beachte dass die Nerven
geordnet nach R\"uckenmarkssegmenten abgehen. 

\\\hline
\end{supertabular}
\end{center}
\paragraph[Vegetatives (autonomes) \ Nervensystem]{Vegetatives
(autonomes)  Nervensystem}
Sympathikus  und Parasymphatikus  sind  Gegenspieler

Sympathikus: {\quotedblbase}Fight or flight{\textquotedblleft} -
{\quotedblbase}Kampf und Flucht{\textquotedblleft}

Parasympathikus: {\quotedblbase}Rest and digest{\textquotedblleft} -
{\quotedblbase}Rast und Verdauung{\textquotedblleft}

z.B.: RR + HF

Sympathikus:  HF ${\uparrow}$, RR ${\uparrow}$

Parasymphatikus: HF ${\downarrow}$, RR ${\downarrow}$

  [Warning: Image ignored] % Unhandled or unsupported graphics:
%\includegraphics[width=6.864cm,height=4.481cm]{AASdev20090228-img35.png}
   [Warning: Image ignored] % Unhandled or unsupported graphics:
%\includegraphics[width=6.73cm,height=4.321cm]{AASdev20090228-img36.png}
 

  [Warning: Image ignored] % Unhandled or unsupported graphics:
%\includegraphics[width=13.29cm,height=16.484cm]{AASdev20090228-img37.png}
 

Nur zur Illustration: Das vegetative Nervensystem hat Einflu{\ss} auf
viele Organe.

Diese Bild NICHT LERNEN. Es dient nur der Verdeutlichung wie wichtig und
einflussreich das vegetative Nervensystem im K\"orper ist. 

\subsection{Die Haut}
gr\"o{\ss}tes Sinnesorgan

ca.1/6 des K\"orpergewichtes

ca. 1,6 m{\texttwosuperior} Oberfl\"ache

Grenzfl\"ache zur Au{\ss}enwelt

Aufgaben

\begin{itemize}
\item Schutz vor \"au{\ss}eren Einfl\"ussen (mechanisch, Erreger,
UV-Strahlung etc.)

\item Schutz vor Wasserverlust/Austrocknung

\item Schutz vor W\"armeverlust /  Temperaturregulation

\item Sekretionsorgan (Schwei{\ss}, Talg)

\item Sinnesorgan (Tastsinn, Temperatursinn)

\end{itemize}
Die Haut ist in Schichten aufgebaut%
%
  [Warning: Image ignored] % Unhandled or unsupported graphics:
%\includegraphics[width=12.952cm,height=9.712cm]{AASdev20090228-img38.gif}
 

Haut: Bild Querschnitt: muss vereinfacht werden, lat. Ausdr\"ucke weg
bzw. grau

\begin{itemize}
\item \index{Oberhaut}Oberhaut (Epidermis)

\end{itemize}
\begin{itemize}
\item Hornschicht, geschichtetes Plattenepithel

\end{itemize}
\begin{itemize}
\item \index{Lederhaut}Lederhaut (Corium, Dermis)

\end{itemize}
\begin{itemize}
\item Hautanhangsgebilde

\end{itemize}
\begin{itemize}
\item Schwei{\ss}dr\"usen

\item Talgdr\"usen

\item Haarfollikel, Nagelmatrix

\end{itemize}
\begin{itemize}
\item Tastk\"orperchen, freie Nervenendigungen (Schmerz)

\item Kollagenfasern (St\"utznetzwerk)

\item Arteriolen, Venolen (Versorgung der Haut, Temperaturregulation)

\end{itemize}
\begin{itemize}
\item \index{Unterhaut}Unterhaut (Subcutis)

\end{itemize}
\begin{itemize}
\item subkutanes Depotfett, W\"armeisolierung

\end{itemize}
Die Hautist wichtig f\"ur die Thermoregulation

Neben der Lunge dient vor allen Dingen  die Haut als wichtigstes Organ
zur Thermoregulation: Die Blutgef\"a{\ss}e  der  Haut sind stark
ineinander verwickelt und k\"onnen sich bei Bedarf stark erweitern. Die
stark durchblutete  Haut strahlt somit W\"arme ab. Weiters wird durch
die Verdunstungsk\"alte  des Schwei{\ss}es  ebenfalls W\"arme
abgegeben. 

\subsection{Atmung}
Die Atemwege sind gegliedert

\begin{center}
\tablehead{}\begin{supertabular}{|m{6.649cm}|m{6.651cm}|}
\hline
Obere Atemwege:

\begin{itemize}
\item Nasenrachenraum mit Nasennebenh\"ohlen 

\item Mundrachenraum

\item Rachen (gemeinsamer Weg von Nahrung und Atem!)

\item \index{Kehlkopf}Kehlkopf (Larynx)

\end{itemize}
Untere Atemwege:

\begin{itemize}
\item \index{Luftr\"ohre}Luftr\"ohre (\index{Trachea}Trachea)

\item \index{Hauptbronchus}Hauptbronchus (li./re.)

\item \index{Bronchioli}Bronchioli

\item \index{Lungenbl\"aschen}Lungenbl\"aschen
(\index{Alveolen}Alveolen)

\end{itemize}
 &
\centering   [Warning: Image ignored]
% Unhandled or unsupported graphics:
%\includegraphics[width=5.633cm,height=7.116cm]{AASdev20090228-img39.png}
 \par

\\\hline
\end{supertabular}
\end{center}
\begin{multicols}{2}
Obere Atemwege

Der Weg der Einatemluft beginnt in der Nase bzw. im Mund. Die Nase ist
mit Schleimhaut ausgekleidet, kn\"ocherne Vorspr\"unge sorgen f\"ur
eine Vorw\"armung der Luft. In der Schleimhaut be\-finden sich die
Sinneszellen f\"ur den Geruchssinn. 

In den Sch\"adelknochen befinden sich luft\-ge\-f\"ullte Hohlr\"aume,
die so\-genannten Neben\-h\"ohlen (Stirn-, Kiefer\-h\"ohlen,
Siebbeinzellen). Diese stehen in Verbindung mit dem Nasenrachenraum.
Entz\"undet sich die Schleimhaut spricht man von der
Neben\-h\"ohlen\-ent\-z\"undung, welche oft langwierig und schmerzhaft
sein kann. 

Die Luft gelangt nun in den Mund-Rachen\-raum und den Rachen. Hier
kreuzt der Weg der Nahrung mit dem Weg der Atemluft! Die Nahrung
gelangt vom Mundraum in die hinten gelegene Speiser\"ohre, die Luft
m\"ochte vom Nasen\-rachen\-raum in die vorne gelegene Luftr\"ohre. 

Um ein Verschlucken von Nahrung in die Lunge zu verhindern
verschlie{\ss}t der am Kehlkopf gelegene Kehldeckel w\"ahrend des
Schluckaktes die Luftr\"ohre. Dieser Reflex ist lebenswichtig! 

Im Kehlkopf  befinden sich die Stimmb\"ander, welche f\"ur die
Lautbildung beim Sprechen wichtig sind.

Untere Atemwege

In der Luftr\"ohre (Trachea) wird die vom Rachen und Kehlkopf kommende
Luft zur Lunge weitergeleitet. Die Luftr\"ohre teilt sich ein einen
rechten und einen linken Hauptbronchus, welche zu den jeweiligen
Lungenfl\"ugeln ziehen. Jeder Hauptbronchus verzweigt sich immer weiter
bis schlu{\ss}endlich sehr d\"unne Bronchioli in die Lungenbl\"aschen
(Alveolen) m\"unden. Hier findet der Gasaustausch statt. 

Die Abschnitte der Atemwege in denen kein Gasaustausch stattfindet und
der nur dem Transport dient bezeichnet man als
{\quotedblbase}\index{Totraum}Totraum{\textquotedblleft}.

\end{multicols}
weiter

\paragraph[Die Lunge ist f\"ur die Atmung zust\"andig]{Die Lunge ist
f\"ur die Atmung zust\"andig}
Die Lunge hat 2 Fl\"ugel, welche sich jeweils in 3 (rechts) bzw. 2
(links) Lappen gliedern: 

\begin{center}
\tablehead{}\begin{supertabular}{|m{6.649cm}|m{6.651cm}|}
\hline
\multicolumn{2}{|m{13.5cm}|}{\centering 1 Lunge\par

}\\\hline
\centering Rechter Fl\"ugel \par

 &
\centering Linker Fl\"ugel\par

\\\hline
\begin{itemize}
\item 3 Lappen

\end{itemize}
 &
\begin{itemize}
\item 2 Lappen

\end{itemize}
\\\hline
\begin{itemize}
\item Oberlappen

\item Mittellappen

\item Unterlappen

\end{itemize}
 &
\begin{itemize}
\item Oberlappen

\item Unterlappen

\end{itemize}
\\\hline
\end{supertabular}
\end{center}
Die Oberfl\"ache ist mit dem Brustfell, der lungenseitigen Pleura
\"uberzogen. In einem Lungefl\"ugel verzweigt sich der Hauptbronchus in
immer weitere kleinere Bronchien bis schlie{\ss}lich die Bronchioli in
die Lungenbl\"aschen, die Alveolen, enden. 

\subparagraph[Die Alveolen sind die Endst\"ucke des Atemwegs und sind
f\"ur den Gasaustausch zust\"andig]{Die Alveolen sind die Endst\"ucke
des Atemwegs und sind f\"ur den Gasaustausch zust\"andig}
In den Lungenbl\"aschen (Alveolen) hier findet der
\index{Gasaustausch}Gasaustausch zwischen der Einatemluft und dem Blut
statt. Sauerstoff gelangt dabei aus den Alveolen in das Blut und
Kohlendiovid wird im Gegenzug vom Blut an die Alveolen abgegeben. 

Dabei besteht zwischen Luft und Blut kein direkter Kontakt, die Gase
m\"ussen Hindernisse wie die Alveolenwand und die Blutgef\"a{\ss}wand
\"uberwinden. Vergr\"o{\ss}ert sich dieser Weg aufgrund einer
Erkrankung ist die Sauerstoffaufnahme vermindert! 

Sauerstoff (O2):   Alveolen ${\rightarrow}$ Blut

Kohlendioxid (CO2):  Blut ${\rightarrow}$ Alveolen

  [Warning: Image ignored] % Unhandled or unsupported graphics:
%\includegraphics[width=13.597cm,height=7.554cm]{AASdev20090228-img40.png}
 

\paragraph{Die Pleura}
\begin{center}
\tablehead{}\begin{supertabular}{|m{6.649cm}|m{6.651cm}|}
\hline
\begin{itemize}
\item Thoraxseitige Pleura%
%
 (Rippenfell, fest mit  Brustwand und Zwerchfell verwachsen, kleidet
damit die Brusth\"ohle aus)

\item Lungenseitige Pleura (Lungenfell, \"uberzieht die
Lungenoberfl\"ache)

\item Aufgrund eines Unterdrucks haften die beiden Pleuren aneinander,
sind aber gegeneinander verschieblich (wie 2 Glasplatten)

\item Geringe Menge Pleurafl\"ussigkeit ${\rightarrow}$ Adh\"asion der
beiden Bl\"atter!

\item Normalerweise keine Verwachsungen der beiden Bl\"atter

\end{itemize}
 &
\centering   [Warning: Image ignored]
% Unhandled or unsupported graphics:
%\includegraphics[width=5.422cm,height=5.249cm]{AASdev20090228-img41.gif}
 \par

\\\hline
\end{supertabular}
\end{center}
Welche Funktion hat aber nun die Pleura?

Da die eine Seite mit dem Brustkorb, die andere Seite mit der Lunge
verwachsen ist und der Raum zwischen den beiden Seiten luftdicht
abgechlossen ist haften die beiden Seiten aneinander, aber sind
gegeneinander verschieblich (Es gibt ja keine Verwachsungen der beiden
eiten und die Pleurafl\"ussigkeit fungiert sogar ein bisschen als
{\quotedblbase}Schmiermittel{\textquotedblleft}). 

Wenn sich nun der Brustkorb oder das Zwerchfell bei der Atmung bewegt
bewegt sich nat\"urlich das Rippenfell mit (es ist ja mit dem Brustkorb
und dem Zwerchfell verwachsen). Da der Pleurazwischenraum luftdicht ist
muss sich Zwangsl\"aufit aber auch das Lungenfell (an dem wiederum die
Lunge angewachsen ist) mitbewegen. Und wenn sich das Lungenfell
mitbewegt wird automatisch auch die Lunge wie eine Ziehharmonika
{\quotedblbase}Aufgezogen{\textquotedblleft} und Luft
{\quotedblbase}angesaugt{\textquotedblleft}. 

Frage: Was w\"urde passieren wenn der Zwischenraum zwischen den
Pleurabl\"attern aufgrund einer Verletung (Z.B. ein Loch nach einer
Stichwunde in die Brust) nicht mehr luftdicht ist? Antwort: s.
\footnote{Antwort: Die Lunge w\"urde in sich zusammenfallen, da sie ja
nicht mehr aufgespannt wird. Siehe
{\quotedblbase}Pneumothorax{\textquotedblleft}.} 

Atemmechanik%
%


  [Warning: Image ignored] % Unhandled or unsupported graphics:
%\includegraphics[width=13.648cm,height=7.52cm]{AASdev20090228-img42.png}
 

Der wichtigste Atemmuskel ist das Zwerchfell! 

Atemmechanik: Atemhilfsmuskulatur fehtl!

Luft

\begin{center}
\tablehead{}\begin{supertabular}{|m{4.367cm}|m{4.367cm}|m{4.367cm}|}
\hline
\centering Raumluft\par

 &
 &
\centering Ausatemluft\par

\\\hline
21 \% O2

 &
\centering Sauerstoff\par

 &
16 \% O2

\\\hline
79 \% N

 &
\centering Stickstoff\par

 &
79 \% N

\\\hline
0,03 \% CO2

 &
\centering Kohlendioxid\par

 &
4 \% CO2

\\\hline
1 \% Edelgase

 &
 &
1 \% Edelgase

\\\hline
\end{supertabular}
\end{center}
Atmungvolumina



\begin{center}
 [Warning: Image ignored] % Unhandled or unsupported graphics:
%\includegraphics[width=10.655cm,height=6.386cm]{AASdev20090228-img43.png}

\end{center}
+ 150ml in den oberen Luftwegen (Trachea, Rachenraum) ${\rightarrow}$
Dieser Teil nimmt nicht am Gasaustausch ${\rightarrow}$
\index{Totraum}Totraum 

Atmung

\begin{itemize}
\item \"au{\ss}ere  Atmung : Gasaustausch Alveolen /Lungenkapillaren

\item innere  Atmung : Gasaustausch Kapillaren /Gewebe

\end{itemize}
normale Atemfrequenzen sind altersabh\"angig:

\begin{center}
\begin{tabular}{|m{4.367cm}|m{4.367cm}|m{4.367cm}|}
\hline
 &
\centering AF\par

 &
\centering AZV\par

\\\hline
Erwachsener

 &
12 -- 16/min

 &
ca.  500 ml

\\\hline
Neugeborenes

 &
40/min 

 &
20-40 ml

\\\hline
Kleinkind

 &
25/min

 &
100-200 ml

\\\hline
Schulkind

 &
20/min

 &
200-400 ml

\\\hline
\end{tabular}
\end{center}
AMV = AZV x AF, beim Erwachsenen  ca. 6-8l

AF  ... Atemfrequenz

AMV  ... Atemminutenvolumen

AZV  ... Atemzugsvolumen ... Volumen bei einem normalen Atemzug

\subsection{Der Verdaungstrakt}
{\quotedblbase}GastroIntestinaltrakt{\textquotedblleft} (GIT)

\"Ubersicht

\begin{center}
\tablehead{}\begin{supertabular}{m{3.238cm}|m{3.674cm}|m{6.189cm}|}
\hline
\multicolumn{1}{|m{3.238cm}|}{Mundh\"ole

} &
 &
Zerkleinerung der Nahrung und Vermischung mit Speichel

\\\hline
 &
\multicolumn{1}{m{3.674cm}}{Zunge

} &
\\\hhline{~-~}
 &
\multicolumn{1}{m{3.674cm}}{Z\"ahne

} &
\\\hhline{~-~}
\multicolumn{1}{|m{3.238cm}|}{Rachen

} &
 &
Kreuzung des Speisebreis mit dem Luftweg

Kehldeckel verschliesst beim Schlucken die Luftr\"ohre

\\\hline
\multicolumn{1}{|m{3.238cm}|}{Speiser\"ohre (\"Osophagus)

} &
 &
\\\hline
\multicolumn{1}{|m{3.238cm}|}{Magen

} &
 &
Magens\"aure zerlegt Nahrung

Vermischung

Pf\"ortner gibt Portionsweise Speisebrei an den Zw\"olffingerdarm ab

\\\hline
\multicolumn{1}{|m{3.238cm}|}{D\"unndarm

} &
 &
\\\hline
 &
Zw\"olffingerdarm

 &
Gallenfl\"ussigkeit (${\leftarrow}$Leber) und Bauchspeichel
(${\leftarrow}$Pankreas) wird zugesetzt

\\\hhline{~--}
 &
Krummdarm

 &
Spaltung von N\"ahrstoffen (Eiwei{\ss}, Kohlehydrate, Fett) und
N\"ahrstoffaufnahme

\\\hhline{~--}
 &
\multicolumn{1}{m{3.674cm}}{Leerdarm

} &
\\\hhline{~-~}
\multicolumn{1}{|m{3.238cm}|}{Dickdarm (Colon)

} &
 &
\\\hline
 &
Blinddarm

 &
Sackgasse

\\\hhline{~--}
 &
Wurmfortsatz (Appendix) 

 &
Anh\"angsel der Sackgasse, kann sich entz\"unden.

\\\hhline{~--}
 &
Aufsteigender Teil

 &
Wasserentzug, Eindickung, Bakterienbesiedlung

\\\hhline{~--}
 &
\multicolumn{1}{m{3.674cm}}{Querender Teil

} &
\\\hhline{~-~}
 &
\multicolumn{1}{m{3.674cm}}{Absteigender Teil

} &
\\\hhline{~-~}
 &
\multicolumn{1}{m{3.674cm}}{S-f\"ormig gebogener Teil (Sigma)

} &
\\\hhline{~-~}
\multicolumn{1}{|m{3.238cm}|}{Mastdarm (Rektum)

} &
 &
{\quotedblbase}Speicherung{\textquotedblleft} des ausscheidungsbereiten
Stuhls

\\\hline
\multicolumn{1}{|m{3.238cm}|}{Anus

} &
 &
Absetzen des Stuhles

\\\hline
\end{supertabular}
\end{center}
\begin{itemize}
\item \end{itemize}
Der Stoffwechsel

\begin{itemize}
\item Die Verdauung dient dem Stoffwechsel

\item Proze{\ss}:

\end{itemize}
\begin{itemize}
\item Nahrungsaufnahme

\item Weiterverarbeitung im GIT

\item Transport der N\"ahrstoffe und Energietr\"ager  ins Blut

\item Zellaufbau und Abbau von Abfallstoffen

\end{itemize}
\begin{itemize}
\item \"Uber die Verdauung gewinnt der K\"orper  also einerseits die
n\"otigen {\quotedblbase}Bausteine{\textquotedblleft},  andererseits
auch die n\"otige Energie um  Zellen aufbauen und in Funktion erhalten 
zu k\"onnen.

\end{itemize}
Mundh\"ohle

Z\"ahne, zerkleinern Nahrung

Speicheldr\"usen: produzieren Speichel um Nahrung gleitf\"ahig zu machen
und Nahrung mit Amylase (Verdauungsenzym f\"ur Zucker und St\"arke) zu
versetzen

Zunge: schiebt Nahrung in Schlund weiter, verschlie{\ss}t beim Schlucken
Kehldeckel (sonst  Aspiration!)

Die Speiser\"ohre (\"Osophagus)

\begin{center}
\tablehead{}\begin{supertabular}{|m{6.649cm}|m{6.651cm}|}
\hline
 &
\\\hline
ca. 24cm langer, dreischichtiger Muskelschlauch (oberes Drittel
quergestreifte Muskulatur, untere 2/3 glatt) 

Verbindung Mund/Magen

{\quotedblbase}schiebt{\textquotedblleft} Nahrung durch
Muskelkontraktionen (Peristaltik)
{\quotedblbase}weiter{\textquotedblleft}

hier findet keine Verdauung statt

 &
\centering   [Warning: Image ignored]
% Unhandled or unsupported graphics:
%\includegraphics[width=2.341cm,height=5.85cm]{AASdev20090228-img44.png}
    [Warning: Image ignored] % Unhandled or unsupported graphics:
%\includegraphics[width=2.103cm,height=5.707cm]{AASdev20090228-img45.png}
 \par

\\\hline
\end{supertabular}
\end{center}
Der Magen (Gaster, Ventriculus)

\begin{itemize}
\item 2-3 l Magensaft

\item Schleim (sch\"utzt  Magenschleimhaut)

\item Salzs\"aure (Bakterienfalle, pH 1,0)

\item Pepsin  (Enzym f\"ur Eiwei{\ss}-Verdauung)

\item Wandmuskulatur mischt den Speisebrei mit der Magens\"aure

\item portionsweise Abgabe von Speisebrei durch den Pylorus
(Magenpf\"ortner) in den Zw\"olffingerdarm

\end{itemize}
\begin{center}
\tablehead{}\begin{supertabular}{|m{6.649cm}|m{6.651cm}|}
\hline
 &
\\\hline
\centering   [Warning: Image ignored]
% Unhandled or unsupported graphics:
%\includegraphics[width=4.116cm,height=6.196cm]{AASdev20090228-img46.png}
 \par

 &
\centering   [Warning: Image ignored]
% Unhandled or unsupported graphics:
%\includegraphics[width=4.746cm,height=4.44cm]{AASdev20090228-img47.png}
 \par

\\\hline
Lage des Magens

 &
Der Magen, der Pf\"ortner und das erste St\"uck Zw\"olffingerdarm im
Durchschnitt

\\\hline
\end{supertabular}
\end{center}
Zw\"olffingerdarm (Duodenum)

  [Warning: Image ignored] % Unhandled or unsupported graphics:
%\includegraphics[width=6.774cm,height=5.079cm]{AASdev20090228-img48.png}
   [Warning: Image ignored] % Unhandled or unsupported graphics:
%\includegraphics[width=5.278cm,height=4.663cm]{AASdev20090228-img49.png}
 

\begin{itemize}
\item Beginn des D\"unndarmes, ca 25 cm lang

\item Verdauungsenzyme aus Leber und Pankreas f\"ur Fett, Eiwei{\ss},
Zucker m\"unden hier ein

\end{itemize}
Die \index{Leber}Leber 

{\quotedblbase}Hepar{\textquotedblleft}

\begin{center}
\tablehead{}\begin{supertabular}{|m{6.649cm}|m{6.651cm}|}
\hline
\begin{itemize}
\item Die Leber liegt im rechten Oberbauch und hat einen linken und
rechten Lappen. 

\item \item Abbau von Giftstoffen

\item {\quotedblbase}Verdauungsdr\"use{\textquotedblleft}

\end{itemize}
\begin{itemize}
\item produziert \index{Gallenfl\"ussigkeit}Gallenfl\"ussigkeit
({\quotedblbase}Galle{\textquotedblleft}), welche in der Gallenblase
gespeichert wird: Fettverdauung

\item Von der Leber gehen die Galleng\"ange, an den Galleng\"angen
h\"angt die 

\end{itemize}
\begin{itemize}
\item \index{Gallenblase}Gallenblase: Nur Speicher f\"ur
Gallefl\"ussigkeit

\end{itemize}
\begin{itemize}
\item speichert Fett \& Zucker

\item erh\"alt \"uber \index{Pfortader}Pfortader (Vena portae)
n\"ahrstoffreiches Blut vom Darm zwecks Entgiftung

\end{itemize}
 &
\centering   [Warning: Image ignored]
% Unhandled or unsupported graphics:
%\includegraphics[width=4.719cm,height=3.539cm]{AASdev20090228-img50.png}
 \par

\\\hline
\end{supertabular}
\end{center}
  [Warning: Image ignored] % Unhandled or unsupported graphics:
%\includegraphics[width=5.469cm,height=4.037cm]{AASdev20090228-img51.png}
   [Warning: Image ignored] % Unhandled or unsupported graphics:
%\includegraphics[width=5.045cm,height=4.005cm]{AASdev20090228-img52.png}
   [Warning: Image ignored] % Unhandled or unsupported graphics:
%\includegraphics[width=2.373cm,height=4.068cm]{AASdev20090228-img53.png}
 

\index{Bauchspeicheldr\"use}Bauchspeicheldr\"use
(\index{Pankreas}Pankreas)

15 -- 22 cm lang, Kopf und Schwanz

Ausf\"uhrungsgang m\"undet mit Gallengang in den Zw\"olffingerdarm
(Duodenum)

Aufgaben: 2 wesentlcihe Funktionen

\begin{itemize}
\item produziert Enzyme f\"ur Verdauung ${\rightarrow}$ Zw\"offingerdarm

\item sch\"uttet Hormone ins Blut aus

\item \index{Insulin}Insulin ={\textgreater} Blutzucker ${\downarrow}$

\item Glukagon ${\rightarrow}$ Blutzucker ${\uparrow}$

\end{itemize}
Beim Diabetiker produziert die Bauchspeicheldr\"use nicht genug Insulin.




\begin{center}
 [Warning: Image ignored] % Unhandled or unsupported graphics:
%\includegraphics[width=8.804cm,height=5.081cm]{AASdev20090228-img54.jpg}

\end{center}
\begin{center}
 [Warning: Image ignored] % Unhandled or unsupported graphics:
%\includegraphics[width=5.788cm,height=4.525cm]{AASdev20090228-img55.png}

\end{center}
\index{Krummdarm}Krumm- und \index{Leerdarm}Leerdarm 

\begin{center}
 [Warning: Image ignored] % Unhandled or unsupported graphics:
%\includegraphics[width=3.049cm,height=3.091cm]{AASdev20090228-img56.png}

\end{center}
(Jejunum, Ileum)

\begin{itemize}
\item ca. 3 m lang

\item Oberfl\"achenvergr\"o{\ss}erung durch Darmzotten

\item N\"ahrstoffaufnahme

\item beweglich, am Darmgekr\"ose (Mesenterium) aufgeh\"angt

\end{itemize}
Dickdarm (Colon)

Der Dickdarm (Colon) ist ca. 1,5 m lang und wird in verschiedene
Abschnitte unterteilt: 

\begin{center}
 [Warning: Image ignored] % Unhandled or unsupported graphics:
%\includegraphics[width=3.751cm,height=4.328cm]{AASdev20090228-img57.png}
\captionof{figure}[Lage des Dickdarms und des Appendix]{Lage des
Dickdarms und des Appendix}

\end{center}
\begin{center}
\tablehead{}\begin{supertabular}{m{1.604cm}|m{4.8500004cm}|m{6.6460004cm}}
\hline
\multicolumn{1}{|m{1.604cm}|}{Dickdarm (Colon)

} &
teilt sich in:

 &
\\\hline
 &
Blinddarm

 &
Sackgasse

\\\hhline{~--}
 &
Wurmfortsatz (Appendix) 

 &
Anh\"angsel der Sackgasse, kann sich entz\"unden.

\\\hhline{~--}
 &
Aufsteigender Teil

 &
Wasserentzug, Eindickung, Bakterienbesiedlung

\\\hhline{~--}
 &
\multicolumn{1}{m{4.8500004cm}}{Querender Teil

} &
\\\hhline{~-~}
 &
\multicolumn{1}{m{4.8500004cm}}{Absteigender Teil

} &
\\\hhline{~-~}
 &
\multicolumn{1}{m{4.8500004cm}}{S-f\"ormig gebogener Teil (Sigma)

} &
\\\hhline{~-~}
\end{supertabular}
\end{center}
Der Leerdarm m\"undet in den Dickdarm. Diese M\"undung ist wie eine
T-Kreuzung gebaut: Nach oben geht es in den aufsteigenden Teil des
Dickdarmes, nach unten in den Blinddarm. Wie der Name des Blinddarmes
schon sagt endet er blind, er ist daher eine Sackgasse. Am Blinddarm
angeh\"angt findet man den Wurmfortsatz (Appendix) welcher sich
h\"aufig entz\"undet (Appendizitis). Der Wurmfortsatz erf\"ullt beim
Menschen keine nenneswerte Funktion, ausser der Arbeitsplatzsicherung
von An\"asthesisten und Chirurgen. 

\index{Bauchfell}Bauchfell

\begin{itemize}
\item Die gesamte Bauchh\"ohle ist mit  Bauchfell
(\index{Peritoneum}Peritoneum) ausgekleidet

\item d\"unner \"Uberzug

\item {\quotedblbase}Aufh\"angung{\textquotedblleft} der Bauchorgane 

\item zus\"atzlich gro{\ss}es + kleines Netz

\item Kann sich Entz\"unden (Bauchfellentz\"undung = Peritonitis): zB
bei Darmverletzung, Bauchfelldialyse 

\end{itemize}
\begin{center}
\tablehead{}\begin{supertabular}{|m{6.649cm}|m{6.651cm}|}
\hline
\centering   [Warning: Image ignored]
% Unhandled or unsupported graphics:
%\includegraphics[width=4.982cm,height=6.767cm]{AASdev20090228-img58.png}
 \par

 &
Der Bauch {\quotedblbase}aufgeklappt: Nach oben geklappt sieht man das
gele {\quotedblbase}Gro{\ss}e Netz{\textquotedblleft}, auf die Seiten
und nach unten die Bauchmuskulatur. 

Die hellen, fleischfarbenen Schlingen sind D\"unndarmschlingen, die
dicken braunen Schlingen sind der aufsteigende und der quere Teil des
Dickdarmes.

\\\hline
\end{supertabular}
\end{center}
\subsection{Sonstige Bauchorgane}
Die Milz 

(Lien, Splen)

\begin{center}
 [Warning: Image ignored] % Unhandled or unsupported graphics:
%\includegraphics[width=3.61cm,height=2.202cm]{AASdev20090228-img59.png}

\end{center}
\begin{itemize}
\item hat nichts  mit der Verdauung  zu tun 

\begin{center}
 [Warning: Image ignored] % Unhandled or unsupported graphics:
%\includegraphics[width=4.301cm,height=6.628cm]{AASdev20090228-img60.png}

\end{center}
\item 12 cm, oval, von Kapsel  umgeben

\item im linken Oberbauch unter den Rippen gelegen

\item Teil des Immunsystems

\item entfernt alte Erythrozyten, beim Embryo auch Blut\-bildung

\item sehr gut durchblutet, Blutspeicher (Gefahr:
\index{Milzri{\ss}}Milzri{\ss})

\end{itemize}
Achtung: sog. {\quotedblbase}\index{Zweizeitiger Milzriss}Zweizeitiger
Milzriss{\textquotedblleft}: Zuerst reisst die Milz ein, die Kapsel
bleib erhalten, es blutet zuerst in die Kapsel und der Druck steigt.
Nach einer gewissen Zeit reisst auch die Kapsel und es blutet nun
ungehindert aus der inneren Wunde. Der Patient wirkt zuerst
{\textquotedblleft}ganz normal{\textquotedblleft}, unerwartet und
pl\"otzlich verf\"allt er. 

Die Niere(n) (Ren)

\begin{itemize}
\item paarig angelegtes Organ

\begin{center}
 [Warning: Image ignored] % Unhandled or unsupported graphics:
%\includegraphics[width=3.873cm,height=4.654cm]{AASdev20090228-img61.png}

\end{center}
\item stark durchblutet, empfindlich auf RR- Abfall ${\rightarrow}$
Funktion ${\downarrow}$ 

\item reguliert

\end{itemize}
\begin{itemize}
\item Wasser- und Elektrolythaushalt

\item pH (S\"aure/Basen)

\end{itemize}
\begin{itemize}
\item Geht im Alter, besonders bei Diabetikern, kaputt ${\rightarrow}$
{\quotedblbase}Chronische Niereninsuffizienz{\textquotedblleft}
${\rightarrow}$ {\quotedblbase}Blutw\"asche{\textquotedblleft}
(\index{Dialyse}Dialyse)

\end{itemize}
Aufbau

\begin{itemize}
\item Rinde, Mark

\item Hilus

\item Gef\"a{\ss}e

\item Nierenbecken mit Anschluss an den Harntrakt

\end{itemize}
Niere: Anatomie und Funktion muss noch ausgearbeitet werden.

%
%
Mikroskopischer Aufbau

\begin{itemize}
\item Nephron (ca. 1 Mio/Niere)

\item Glomerulus: Hier wird das Blut gefiltert und Plasmawasser
abgepresst (\~{}150~l~/~Tag)

\item Schleifenapparat (Tubulus): Das abgepresste Plasma\-wasser wandert
den Schleifen\-apparat entlang, hier werden die meisten wichtigen
Stoffe und der Gro{\ss}\-teil des Wassers wiederum zur\"uckgeholt,
sodass am Ende nur mehr 1,5 -- 2~l Harn pro Tag in der Harnblasen
landen. 

\end{itemize}
Funktionsweise

\begin{itemize}
\item 1500 l Blut pro Tag wird
{\quotedblbase}bearbeitet{\textquotedblleft}

\item Glomerula: Filtration ${\rightarrow}$ 150  l
\index{Prim\"arharn}Prim\"arharn (enth\"alt Elektrolyte, etc.) wird
gewonnen und fliesst durch den Schleifenapparat

\item Schleifenapparat: {\quotedblbase}Recycling{\textquotedblleft},
Gutes kommt zur\"uck ins Blut, der Rest fliesst weiter Richtung
Harnblase ${\rightarrow}$ 1,5 -- 2  l (End-)Harn/Tag 

\end{itemize}
Die Nebenniere(n) (Adren)

paariges, jeweils am oberen Nierenpol gelegenes Organ

Hat nichts  mit dem Harntrakt oder der eigentlichen Nierenfunktion zu
tun.

Hormonproduktion

\begin{itemize}
\item Cortison

\item Sexualhormone (Androgene)

\item Adrenalin , Noradrenalin

\end{itemize}
Die Nebennieren sitzen wie {\quotedblbase}M\"utzen{\textquotedblleft}
auf den Nieren (s. Illustration \ref{seq:refIllustration4} unten). 

Harntrakt

\begin{center}
\tablehead{}\begin{supertabular}{|m{6.649cm}|m{6.651cm}|}
\hline
Der Harntrakt besteht aus 

\begin{itemize}
\item Nieren inkl. Nierenbecken

\item \index{Harnleiter}Harnleiter  (Ureter ), 30 cm lang, m\"undet  in
die

\item \index{Blase}Blase  (Vesica )

\item Harnr\"ohre  (Urethra, [2642?] 25cm, [2640?] 4cm)

\end{itemize}
Die Harnr\"ohre verl\"auft beim Mann von der Harnblase durch die
Prostata und den Penis und m\"undet an der Eichel in die Aussenwelt.
Bei einer Verg\"o{\ss}erung der Prostata kann der Harnabflu{\ss} bis
zur Unpassierbarkeit erschwert sein.  

Bei der Frau ist die Harnr\"ohre k\"urzer und m\"undet oberhalb des
Scheideneinganges. 

 &
\centering   [Warning: Image ignored]
% Unhandled or unsupported graphics:
%\includegraphics[width=5.85cm,height=7.43cm]{AASdev20090228-img62.jpg}
 \par

\captionof{figure}{Harntrakt mit Nieren (linke Niere im Querschnitt mit
Nierenbescken), Harnleiter, Blase und Harnr\"ohre. Auf den Nieren sieht
man wie eine Zipfelm\"utze die Nebennieren.}
\label{seq:refIllustration4}
\\\hline
\end{supertabular}
\end{center}
\subsection{Geschlechtsorgane}
M\"annliche Geschlechtsorgane

\begin{itemize}
\item innen:

\end{itemize}
\begin{itemize}
\item Hoden (Testes)

\end{itemize}
\begin{itemize}
\item Spermienproduktion

\item Geschlechtshormon-Produktion

\end{itemize}
\begin{itemize}
\item Nebenhoden: Spermienreifung und Lagerung

\item Samenstrang, besteht aus 

\end{itemize}
\begin{itemize}
\item Samenleiter

\item Nerven

\item Gef\"a{\ss}e

\item durchzieht Leistenkanal

\end{itemize}
\begin{itemize}
\item \index{Prostata}Prostata (Vorsteherdr\"use)

\end{itemize}
\begin{itemize}
\item produziert Samenfl\"ussigkeit (Beweglichkeit der Spermien)

\item Vergr\"o{\ss}ert sich im Alter und kann die Harnr\"ohre
abdr\"ucken, es kommt zum Harnstau ${\rightarrow}$ h\"aufger Notfall!

\end{itemize}
\begin{itemize}
\item au{\ss}en:

\end{itemize}
\begin{itemize}
\item Glied (Penis): Schwellk\"orper, Harnsamenr\"ohre, Begattungsorgan

\item Hodensack (Skrotum): umh\"ullt Hoden und Nebenhoden

\end{itemize}
\begin{itemize}
\item Hodenstiel kann sich verdrehen ${\rightarrow}$ Notfall!

\end{itemize}
  [Warning: Image ignored] % Unhandled or unsupported graphics:
%\includegraphics[width=13.699cm,height=8.52cm]{AASdev20090228-img63.png}
 

Weibliche Geschlechtsorgane

\begin{itemize}
\item innen:

\end{itemize}
\begin{itemize}
\item Eierst\"ocke (Ovar): Reifung der Eizellen

\item Eileiter (Tubae): Transport des reifen Eis in Richtung
Geb\"armutter

\item Geb\"armutter (\index{Uterus}Uterus): dehnbarer Hohlmuskel, mit
Schleimhaut ausgekleidet, 

\end{itemize}
\begin{itemize}
\item Schleimhaut wird regelm\"a{\ss}ig aufgebaut und wieder
abgesto{\ss}en (Menstruation, Regelblutung)

\item hier nistet sich die befruchtete Eizelle ein. 

\item Die M\"undung der Geb\"armutter in die Scheide ist der Muttermund.


\end{itemize}
\begin{itemize}
\item Scheide (Vagina)

\end{itemize}
\begin{itemize}
\item au{\ss}en: 

\end{itemize}
\begin{itemize}
\item Schamlippen: go{\ss}e/\"au{\ss}ere, kleine/innere

\item Kitzler (Klitoris): entspricht der m\"annlichen Eichel

\end{itemize}
  [Warning: Image ignored] % Unhandled or unsupported graphics:
%\includegraphics[width=13.399cm,height=8.485cm]{AASdev20090228-img64.png}
 

\paragraph{Der weibliche Zyklus}
\begin{itemize}
\item Der weibliche Zyklus dauert 28 Tage

\item gerechnet von Menstruation zu Menstruation  (Monatsblutung )

\item Uterusschleimhaut wird aufgebaut und wieder abgesto{\ss}en

\item hormongesteuert

\end{itemize}
\begin{itemize}
\item \"Ostrogen, Progesteron

\end{itemize}
\begin{itemize}
\item ca. am 14. Tag: Eisprung

\item Befruchtung nein: 

\end{itemize}
\begin{itemize}
\item Zyklus l\"auft normal weiter, Menstruation

\end{itemize}
\begin{itemize}
\item Befruchtung ja:

\end{itemize}
\begin{itemize}
\item Eizelle nistet sich im Uterus ein, Zyklus gestoppt

\end{itemize}
Was kann dabei schief gehen? Wenn sich die Eizelle schon im Ei\-leiter
einnistet kommt es zu einer
\index{Eileiter\-schwanger\-schaft}Eileiter\-schwanger\-schaft die
lebens\-bedrohlich ist (Eileiter rei{\ss}t und die Frau verblutet).

 [Warning: Image ignored] % Unhandled or unsupported graphics:
%\includegraphics[width=13.576cm,height=7.541cm]{AASdev20090228-img65.png}
\captionof{figure}[Der weibliche Zyklus. Oben: Follikelreifung. Unten:
Auf- und Abbau der Uterusschleimhaut. Am Tag 0 des Zyklus wird die
Uterusschleimheut abgestossen, es kommt zur Regelblutung, danach wird
die Schleimhaut wieder aufgebaut. Am Tag 14 erfolgt der Eispung. Kommt
es zu keiner Befruchtung wird die Schleimheut erneut am Tag 28
abgesto{\ss}en. Nach einer Befruchtung wandert die Eizelle weiter in
den Uterus und nistet sich dort ein. ]{Der weibliche Zyklus. Oben:
Follikelreifung. Unten: Auf- und Abbau der Uterusschleimhaut. Am Tag 0
des Zyklus wird die Uterusschleimheut abgestossen, es kommt zur
Regelblutung, danach wird die Schleimhaut wieder aufgebaut. Am Tag 14
erfolgt der Eispung. Kommt es zu keiner Befruchtung wird die
Schleimheut erneut am Tag 28 abgesto{\ss}en. Nach einer Befruchtung
wandert die Eizelle weiter in den Uterus und nistet sich dort ein. }


\section{[RS V] Vitalfunktionen}
Sebastian Gabriel [R]; Michael Withofner

Changelog:

2009-02-25: Start Changelog. Kapitel mit Korr Koch korr.

\"Uberblick

\begin{center}
\tablehead{}\begin{supertabular}{|m{6.649cm}|m{6.651cm}|}
\hline
Vitalfunktionen I. Ordnung

 &
Vitalfunktionen II. Ordnung

\\\hline
\begin{itemize}
\item Bewu{\ss}tsein

\item Atmung

\item Kreislauf

\end{itemize}
wichtigste, grundlegende Funktionen des K\"orpers

Ausfall einer dieser Funktionen 

bedeutet immer akute Lebensgefahr!!! 

mu{\ss} akutmedizinisch behandelt werden!!!

 &
\begin{itemize}
\item Wasser-/Elektrolythaushalt

\item S\"aure-/Basenhaushalt

\item Temperaturhaushalt

\item Hormonsystem

\item Stoffwechsel

\item Immunsystem

\end{itemize}
\\\hline
\end{supertabular}
\end{center}
\paragraph[NACA -- Score]{NACA -- Score}
\index{NACA}National Advisory Committee for Aeronautics

Diente urspr\"unglich der Bewertung von  Erkrankungen / Verletzungen und
der Entscheidung welche Soldaten zuerst auszufliegen w\"aren. Heute
dient der NACA-Score der groben statistischen Erfassung.  

1 bis 7:

\begin{center}
\begin{tabular}{m{1.201cm}m{12.1cm}}
\centering  NACA\par

 &
\\
\centering I\par

 &
 keine \"arztliche Therapie erforderlich

\\
\centering II\par

 &
 ambulante Abkl\"arung

\\
\centering III\par

 &
 station\"are Abkl\"arung

\\
\centering IV\par

 &
 akute Lebensgefahr nicht ausgeschlossen

\\
\centering V\par

 &
 akute Lebensgefahr

\\
\centering VI\par

 &
 Reanimation

\\
\centering VII\par

 &
 Tod

\\
\end{tabular}
\end{center}
\subsection{Bewusstsein}
\"Uberbegriff u.a. f\"ur Wachheit, Orientierung, Aufmerksamkeit,
Auffassungsgabe, Denkverlauf und
Merkf\"ahigkeit\label{ref:JRcite1Pschyrembel259}[1] . 

Die {\quotedblbase}wahre{\textquotedblleft} Definition von Bewusstsein
ist seit tausenden von Jahren ungekl\"art. Da f\"ur die
Sanit\"aterausbildung aber nur beschr\"ankt Zeit zur Verf\"ugung steht
ist die obige Definition f\"ur den Alltag erstmal ausreichend.

Das Bewusstsein ist besonders wichtig hinsichtlich des Schutzes des
Menschen vor Bedrohung. Ein bewusstseinsklarer Mensch kann sich gegen
Gefahren wehren, ein eingetr\"ubter oder bewusstloser Mensch kann dies
schlecht oder gar nicht. 

Bewu{\ss}seinsst\"orung

St\"orung des Bewusstseins, besonders der Wachheit und Orientierung. 

Orientierung

Es ist wichig zu \"Uberpr\"ufen ob der Patient orientiert ist. Man
unterscheidet dabei: 

\begin{enumerate}
\item Orientierung zur Person: Name, Alter. 

\item \"Ortliche Orientierung: Wo befindet sich der Patient gerade
(Stadt, Land, Adresse, Wohnung, Auto, etc.)?

\item Zeitliche Orientierung: Tageszeit, Wochentag, Monat, Jahreszeit,
Jahr? 

\end{enumerate}
Allgemein

Wie bereits oben erw\"ahnt ist das Bewusstsein wesentlich an der Abwehr
von Gefahren beteiligt. Je mehr das Bewusstsein gest\"ort oder
reduziert ist desto gef\"ahrlicher ist das f\"ur den Patienten. 

\begin{itemize}
\item Ab einem gewissen Punkt sind auch so wesentliche Funktionen wie
der Schluck-, W\"urge- bzw. Hustenreflex ausfallen und es kann zum
verschlucken von Mageninhalt oder Blut in die Lunge kommen
(Aspiration). Man spricht auch vom Ausfall der Schutzreflexe. 

\item Bei schweren Bewusstseinsst\"orungen kann es auch zum Erschlaffen
wichtiger Muskeln wie z.B. Der Zunge kommen. F\"allt diese zur\"uck
kann sie den Atemweg blockieren oder ganz verschlie{\ss}en. 

\end{itemize}
Daher sind schwere Bewusstseinsst\"orungen lebensgef\"ahrlich!

Ursachen

${\rightarrow}$ Sch\"adigung des zentralen Nervensystems

\begin{center}
\tablehead{}\begin{supertabular}{|m{6.649cm}|m{6.651cm}|}
\hline
\centering prim\"are  Ursachen\par

 &
\centering sekund\"are  Ursachen:\par

\\\hline
direkt das ZNS betreffend

 &
\"Uber Umwege das ZNS betreffend

\\\hline
Verletzungen (SHT)

Blutungen, Hirnschwellung

Schlaganfall

epileptischer Anfall

Entz\"undungen

Tumor

 &
St\"orungen der O2-Versorgung :

Atemst\"orungen  (Hypoxie)

Herz-Kreislaufst\"orungen

Elektrolytentgleisungen

Vergiftungen:

Stoffwechselst\"orungen (Zuckerhaushalt, Leberversagen)

Alkohol

Medikamente, Drogen

Reizgase,  L\"osungsmittel ...

\\\hline
\end{supertabular}
\end{center}
Grade und Symptome

Symptome Bewusstseinsst\"orungen: Tabelle geh\"ort erg\"anzt.%
%


\begin{center}
\tablehead{}\begin{supertabular}{|m{3.225cm}|m{3.225cm}|m{3.225cm}|m{3.225cm}|}
\hline
bewu{\ss}tseinsklar

 &
 &
 &
\\\hline
bewu{\ss}tseinsgetr\"ubt

 &
desorientiert

 &
 &
\\\hline
 &
delirant

 &
 &
\\\hline
 &
agitiert

 &
 &
\\\hline
 &
somnolent

 &
Schl\"afrig, aber leicht erweckbar. 

 &
\\\hline
 &
sopor\"os

 &
 &
Gefahr!

\\\hline
bewu{\ss}tlos

 &
Koma

 &
 &
Lebensgefahr! 

\\\hline
\end{supertabular}
\end{center}
Gefahren

\begin{itemize}
\item Zur\"uckfallen der Zunge ${\rightarrow}$ Erstickung!

\item Aspiration  (Husten- bzw. W\"urgereflex sowie  s\"amtliche andere
Schutzreflexe ausgefallen,  Muskulatur komplett erschlafft
={\textgreater}  Erbrochenes gelangt ungehindert in die  Lunge!)

\end{itemize}
\index{Glasgow Coma Scale}Glasgow Coma Scale (\index{GCS}GCS)

... dient zur objektiven Einsch\"atzung der Bewu{\ss}tseinslage, Skala
mit 3 bis max. 15 Punkten.

Bewu{\ss}tseinsklarer Patient: 15 Punkte. Toter Patient: 3 Punkte (nicht
0!)

Untersucht werden: 

\begin{flushleft}
\tablehead{}\begin{supertabular}{m{3.224cm}m{0.874cm}|m{1.1459999cm}|m{7.649cm}|}
\hline
\multicolumn{1}{|m{3.224cm}|}{} &
\centering Max.\par

 &
\centering Punkte\par

 &
\centering Reaktion\par

\\\hline
\multicolumn{1}{|m{3.224cm}|}{\begin{itemize}
\item Augen\"offnung

\end{itemize}
} &
4

 &
4

 &
spontan

\\\hline
 &
 &
3

 &
auf Ansprache

\\\hhline{~~--}
 &
 &
2

 &
Schmerzreiz

\\\hhline{~~--}
 &
 &
1

 &
nicht

\\\hhline{~~--}
\multicolumn{1}{|m{3.224cm}|}{\begin{itemize}
\item Verbale Antwort

\end{itemize}
} &
5

 &
5

 &
klar

\\\hline
 &
 &
4

 &
verwaschen

\\\hhline{~~--}
 &
 &
3

 &
stammelnd

\\\hhline{~~--}
 &
 &
2

 &
unartikuliert

\\\hhline{~~--}
 &
 &
1

 &
nicht

\\\hhline{~~--}
\multicolumn{1}{|m{3.224cm}|}{\begin{itemize}
\item Motorische Reaktion

\end{itemize}
} &
6

 &
6

 &
spontan

\\\hline
 &
 &
5

 &
gezielte Abwehrbewegungen

\\\hhline{~~--}
 &
 &
4

 &
ungezielte Abwehrbewegungen

\\\hhline{~~--}
 &
 &
3

 &
Beugung

\\\hhline{~~--}
 &
 &
2

 &
Strecken

\\\hhline{~~--}
 &
 &
1

 &
bewegt gar nicht

\\\hhline{~~--}
\end{supertabular}
\end{flushleft}
Ma{\ss}nahmen bei getr\"ubten Patienten

Ma{\ss}nahmen bei Bewusstseinsst\"orungen: Neufassung, bitte besonders
genau durchsehen!%
%


\begin{itemize}
\item BAK -- Schema  (Bewu{\ss}tsein/Atmung/Kreislauf)

\item Lagerung: situationsgerecht, im Zweifel stabile Seitenlage 

\item Je nach Verdachtsdiagnose: 

\end{itemize}
\begin{itemize}
\item O2: nach Ermessen, im Zweifel ca. 6-8 l

\item Reanimationsbereichtschaft

\item Notarzt 

\end{itemize}
\begin{itemize}
\item Engmaschiges Monitoring (Atmung, Pulsoxymetrie, RR)

\item W\"armeerhalt

\item Ursachenforschung -- Warum ist der Patient eingetr\"ubt?

\item Neurcheck

\end{itemize}
\begin{itemize}
\item Immer auch an BZ denken!

\end{itemize}
\begin{itemize}
\item Traumacheck

\end{itemize}
Ma{\ss}nahmen bei bewu{\ss}tlosen und sopor\"osen Patienten

\begin{itemize}
\item BAK -- Schema  (Bewu{\ss}tsein/Atmung/Kreislauf)

\item Lagerung: stabile Seitenlage

\item O2:  10 l  O2 , bei Zyanose\footnote{Zyanose: Bl\"auliche
Verf\"arbung der Haut infolge einer unzureichenden
Sauerstoffversorgung. Besonders auff\"allig an den Lippen, Gesicht, und
Extremit\"aten.}  15 l

\item Reanimationsbereichtschaft!

\item Notarzt! 

\item Engmaschiges Monitoring (Atmung, Pulsoxymetrie, RR)

\item W\"armeerhalt

\item Ursachenforschung -- Warum ist der Patient bewusstlos oder
sopor\"os?

\item Neurcheck

\end{itemize}
\begin{itemize}
\item Immer auch an BZ denken!

\end{itemize}
\begin{itemize}
\item Traumacheck

\end{itemize}
\subsection{Atmung}
Kurze Wiederholung

normale Atemfrequenz: (Erwachsener): 12-15x / min, AZV = 5-7 ml / kg KG

Atemminutenvolumen (AMV): AZV x AF

\index{Totraumventilation}Totraumventilation: Luftmenge in oberen  und
teilweise unteren Atemwegen, die  nicht am Gasaustausch teilnimmt.
Achtung \index{Schnappatmung}Schnappatmung!Dabei wird nur so wenig Luft
bewegt sodass aufgrund des Totraumes keine frische Luft in die
Lungenbl\"aschen (Alveolen) kommt. H\"aufig wird die Schnappatmung mit
flacher Atmung verwechselt, aber:

Schnappatmung = keine Atmung

St\"orungen der Atmung

Ateminsuffizienz = eingeschr\"ankte Atmung ={\textgreater} O2-Mangel
(Hypoxie)   

Atemstillstand = Aussetzen d. Atmung (Apnoe)

Ursachen:

\begin{itemize}
\item \end{itemize}
\begin{center}
\tablehead{}\begin{supertabular}{|m{2.71cm}|m{10.59cm}|}
\hline
\centering St\"orung des/der\par

 &
\\\hline
O2-Angebotes

 &
\begin{itemize}
\item G\"arkeller (CO2 ${\uparrow}$, O2 ${\downarrow}$)

\item {\quotedblbase}d\"unne Luft{\textquotedblleft} (z.B. Hochgebirge,
pO2 ${\downarrow}$)

\item Defekte Gastherme (CO ${\uparrow}$m O2${\downarrow}$)

\item Ertrinken

\end{itemize}
\\\hline
O2-Diffusion

 &
\begin{itemize}
\item Lungen\"odem

\item Lungenentz\"undung

\item Lungenembolie

\end{itemize}
\\\hline
Atemmechanik

 &
\begin{itemize}
\item Verlegung der 

\end{itemize}
\begin{itemize}
\item oberen Atemwege (Zunge, Erbrochenes, Laryngospasmus,
Glottis\"odem, Bolus)

\item unteren Atemwege (Bronchitis, Asthma, Lungen\"odem)

\end{itemize}
\begin{itemize}
\item Verminderung der Dehnbarkeit der Thoraxwand oder des Lungengewebes
(Rippenfraktur, Pneumothorax, Pleuraergu{\ss}, Emphysem)

\end{itemize}
\\\hline
Neuromuskul\"aren Atemregulation

 &
\begin{itemize}
\item SHT (Sch\"adel Hirn Trauma)

\item Schlaganfall (Insult)

\item Vergiftungen

\item Tumore

\item hoher Querschnitt (RM durchtrennt)

\item Muskelerkrankungen

\item Nervenerkankung

\item Stoffwechselst\"orungen

\end{itemize}
\\\hline
\end{supertabular}
\end{center}
Symptome

\begin{center}
\begin{tabular}{|m{3.225cm}|m{3.225cm}|m{3.225cm}|m{3.225cm}|}
\hline
\centering Gruppe\par

 &
\centering Befund\par

 &
\centering Beschreibung\par

 &
\centering Anmerkung\par

\\\hline
Dyspnoe

 &
Atemnot

 &
Leitsymptom

 &
\\\hline
Atemger\"ausch

 &
Brodeln

 &
Blubbern, klassisch f\"ur Lunden\"odem

 &
\\\hline
 &
Stridor

 &
 &
\\\hline
 &
Brummen, Giemen

 &
 &
\\\hline
 &
Rasselger\"ausche

 &
 &
\\\hline
Frequenz

 &
Beschleunigt

 &
Tachypnoe

 &
\\\hline
 &
Verlangsamt

 &
Bradypnoe

 &
\\\hline
Atemzugsvolumen

 &
Schnappatmung

 &
 &
\\\hline
 &
Flache Atmung

 &
 &
\\\hline
 &
Tiefe Atmung

 &
 &
\\\hline
Hautfarbe

 &
Blass

 &
Normal, An\"amie, Blutverlust

 &
\\\hline
 &
Rosig

 &
Normal, CO-Vergiftung

 &
\\\hline
 &
Bl\"aulich 

 &
Zyanose: Hypoxie!

 &
\\\hline
K\"orperlich

 &
Einziehungen an den Rippen

 &
Einsatz der Atemhilfsmuskulatur

 &
\\\hline
 &
Aufst\"utzen

 &
Einsatz der Atemhilfsmuskulatur

 &
\\\hline
 &
Nasenfl\"ugeln

 &
Atemnot, besonders bei Kleinkindern

 &
\\\hline
 &
Aufrechte Position

 &
 &
\\\hline
 &
{\quotedblbase}Paradoxe Atmung{\textquotedblleft}

 &
Brustkorb senkt sich bei Einatmung

 &
\\\hline
\end{tabular}
\end{center}
[F0EA?] Gefahr der Ateminsuffizienz

\begin{itemize}
\item Hypoxie

\end{itemize}
\begin{itemize}
\item in weiterer Folge Unterversorgung lebenswichtiger Organe wie Herz,
Hirn.

\end{itemize}
\begin{itemize}
\item Atmungsbedingte \"Ubers\"auerung (Respiratorische Azidose, s. H.IV
)

\end{itemize}
Ma{\ss}nahmen bei Atemstillstand

\begin{itemize}
\item Freimachen der Atemwege

\end{itemize}
\begin{itemize}
\item \"Offnen des Mundes (Esmarch-Handgriff)

\item Inspektion

\item fals n\"otig ausr\"aumen 

\item (falls n\"otig) Absaugen

\item \"Uberstrecken der HWS

\item Zungengrund wird angehoben -{\textgreater} Atemwege  sind frei!

\end{itemize}
\begin{itemize}
\item \"Uberpr\"ufung der Atemt\"atigkeit (bei  \"uberstreckter HWS)
durch

\end{itemize}
\begin{itemize}
\item SEHEN, H\"OREN, F\"UHLEN f\"ur mindestens 10  Sekunden!

\end{itemize}
\begin{itemize}
\item Sauerstoffgabe (wenn ausreichende Spontan-Atmung vorhanden)

\item Beatmung (keine ausreichende Spontanatmung bzw. bei Reanimation)

\end{itemize}
\begin{itemize}
\item Beatmungsbeutel mit Maske und Reservoir

\end{itemize}
Vgl. die entsprechenden Kapitel im Teil Erste Hilfe und
Sanit\"atstechnik bez\"uglich der genauen Durchf\"uhrung der
Ma{\ss}nahmen. 

Achtung: Bei einem mit Atemstillstand vorgefundenen Patienten ist
grunds\"atzlich von einer Reanimation auszugehen. 

\subsection[Herz- / Kreislauf]{Herz- / Kreislauf}
\label{bkm:RefHeading42514202}St\"orungen des Herz-/ Kreislaufsystems

\label{bkm:RefHeading42514602}Kreislaufstillstand

Kreislaufstillstand: Pumpversagen des Herzens aufgrund von

Asystolie (keine Herzt\"atigkeit)

Kammerflimmern (unkoordinierte Herzt\"atigkeit)

(Pulslose) Ventrikul\"are Tachykardie (HF zu schnell ${\rightarrow}$
diastolische F\"ullung nicht ausreichend ${\rightarrow}$ Auswurf
${\downarrow}$)

elektromechanische Entkopplung 

Ursachen

\begin{itemize}
\item Minderversorgung des Herzmuskels mit O2

\end{itemize}
\begin{itemize}
\item Herzinfarkt

\item Herzinsuffizienz

\end{itemize}
\begin{itemize}
\item St\"orungen der Atmung

\end{itemize}
\begin{itemize}
\item Asthma bronchiale

\item Lungenembolie

\item zentrale Atemregulationsst\"orungen

\end{itemize}
\begin{itemize}
\item sonstige St\"orungen

\end{itemize}
\begin{itemize}
\item Trauma

\item Vergiftungen

\item thermische Sch\"aden

\item Elektrolytverschiebungen,...

\end{itemize}
Pathophysiologie :

\begin{itemize}
\item Aussetzen des Herzschlages ={\textgreater} Patient ist sofort
pulslos!

\item Bewu{\ss}tlosigkeit (nach 15-20 sek.)

\item Atemstillstand/Schnappatmung (nach ca. 10 sek.)

\item Irreversible Hirnsch\"aden nach 3-5 min., bei K\"alte allerdings
l\"anger!

\end{itemize}
\paragraph{Tod}
Beachte:

[F0EA?] Nobody is dead until warm and dead!

Klinischer Tod vs. biologischer Tod

\begin{center}
\tablehead{}\begin{supertabular}{|m{6.649cm}|m{6.651cm}|}
\hline
\centering klinischer Tod\par

 &
\centering biologischer Tod\par

\\\hline
keine Atmung

kein Puls

Reanimation mu{\ss} innerhalb der ersten  Minuten erfolgen

(dann unter Umst\"anden) reversibel

 &
Organtod (irreversibel) des Gehirns

gleichzusetzen mit Tod des Individuums

\\\hline
\end{supertabular}
\end{center}
Wann ist es wirklich zu sp\"at?

Todesfeststellung: Grunds\"atzlich NUR durch den Arzt!

Todeszeichen: Nicht wie von Roman gefordert in Tabelle mit mit sicheren
und unsicheren Todeszeichen umgebaut, halte ich f\"ur un\"ubersichtlich
und bef\"urchte Verwechslungsgefahr mit Knochenbruchzeichen.[GAB]

Unter bestimmten Voraussetzungen allerdings auch durch den Sanit\"ater
(Unter\-lassung von  Reanimationsma{\ss}nahmen):

\begin{center}
 [Warning: Image ignored] % Unhandled or unsupported graphics:
%\includegraphics[width=4.19cm,height=3.028cm]{AASdev20090228-img66.png}

\end{center}
\begin{itemize}
\item %
%
Verwesung

\item Verletzungen, die nicht mit dem Leben  vereinbar sind

\end{itemize}
\begin{itemize}
\item Kopfamputation, {\dots}

\end{itemize}
\begin{itemize}
\item (Totenflecken (nicht 100\%ig feststellbar))

\item (Leichenstarre (nicht 100\%ig feststellbar))

\item Abbruch der Reanimationsma{\ss}nahmen aufgrund von Ersch\"opfung

\end{itemize}


\begin{center}
 [Warning: Image ignored] % Unhandled or unsupported graphics:
%\includegraphics[width=7.133cm,height=4.837cm]{AASdev20090228-img67.jpg}

\end{center}
Herzinsuffizienz, Herzversagen

Weitere, oft auch sehr bedrohliche, St\"orungen des Herzens werden im
Kapitel  Herz-Kreislaufst\"orungen 
(I.I{}-\pageref{bkm:RefHeading42342802}:
{\quotedblbase}Herz-Kreislaufst\"orungenunten{\quotedblbase})
besprochen.

Temperaturregulation [F0E1?]

W\"armeproduktion durch

\begin{itemize}
\item Zellaktivit\"at

\item Stoffwechselvorg\"ange

\item Bewegung (Muskelatmung)

\end{itemize}
W\"armeabgabe \"uber:

\begin{itemize}
\item Haut

\item Schwei{\ss}

\end{itemize}
\begin{itemize}
\item Wirkung durch Verdunstungsk\"alte

\item ca. 0,5 l / Tag

\item bis max. 1,5 l / Std. bei schwerer k\"orperlicher  Arbeit

\end{itemize}
physiologische Kerntemperatur: 37{\textdegree}C

Regulation durch den Hypothalamus

Stoffwechsel

 Definition 

Abl\"aufe von Stoffumsetzung und  Energiegewinnung

N\"ahrstoffe werden in W\"arme und Energie  umgewandelt

Arten

\begin{itemize}
\item Zuckerstoffwechsel (Kohlenhydrate)

\end{itemize}
\begin{itemize}
\item Energielieferant 

\end{itemize}
\begin{itemize}
\item Fettstoffwechsel (Fetts\"auren, ges\"attigt und  unges\"attigt)

\end{itemize}
\begin{itemize}
\item Doppelt so hoher Energiegehalt wie  Kohlenhydrate

\item L\"osungsmittel f\"ur Vitamine

\item Bausteine der Zellmembran

\item Bestandteil von Hormonen

\end{itemize}
\begin{itemize}
\item Eiwei{\ss}stoffwechsel (Proteine)

\end{itemize}
\begin{itemize}
\item Zellmembranbausteine

\item Enzyme!

\end{itemize}
\subsection[Wasser- und Elektro\-lyt\-haus\-halt]{%
%
Wasser- und Elektro\-lyt\-haus\-halt}
Das K\"orperwasser macht ca. 60\% des K\"orpergewichtes aus und teilt
sich auf:

\begin{itemize}
\item IntraZellularRaum (IZR), 40\% KG 

\item ExtraZellularRaum (EZR), 20\% KG 

\end{itemize}
\begin{itemize}
\item \index{Zwischengewebsraum}Zwischengewebsraum, \~{}16\%:
{\quotedblbase}freier{\textquotedblleft} Raum zwischen den Zellen im
Gewebe

\item \index{Gef\"a{\ss}raum}Gef\"a{\ss}raum \~{}4\%: Raum in den
Blutgef\"a{\ss}en

\end{itemize}
Trennung: 

IZR -- EZR:    Zellmembran (semipermeable Membran, siehe unten)

Zwischengewebsraum -- Gef\"a{\ss}raum:  Blutgef\"a{\ss}wand

\index{Atome}Atome, \index{Ionen}Ionen und
\index{Elektrolyte}Elektrolyte 

Atome sind die kleinsten Grundbestandteile von Marterie.

Ionen sind Atome die elektrisch geladen sind.

Geladene Ionen sind gut in Wasser l\"oslich. Wenn sie in Wasser gel\"ost
sind heissen sie {\quotedblbase}Elektrolyte{\textquotedblleft}. 

Im Wasser gel\"oste, elektrisch geladene Teilchen hei{\ss}en
Elektrolyte.

Die wichtigsten Ionen sind Natrium (Na+), Kalium (K+) (positiv geladen)
und Chlorid (Cl-)(negativ geladen). 

Funktion:

\begin{itemize}
\item Osmose, Diffusion und Verschiebung von Wasser zwischen den
einzelnen R\"aumen (siehe unten)

\end{itemize}
\begin{itemize}
\item ${\rightarrow}$ Volumenregulation des K\"orpers, siehe unten.  

\end{itemize}
\begin{itemize}
\item Elektrische Ladung: Wie bei einer Batterie ist eine Zelle geladen
(innen Minuspol, aussen Pluspol). Das ist besonders wichtig f\"ur

\end{itemize}
\begin{itemize}
\item Herz (Erregungsleitung!)

\item Nerven

\item Muskel

\end{itemize}
Elektrolyte sind geladene Teilchen, die  bei Aufspaltung von S\"auren,
Laugen  oder deren Salzen in w\"assriger L\"osung  entstehen.

Kationen: positiv geladene Teilchen (K+ , Na+ ,Mg++ , Ca++ ,...)

Anionen: negativ geladene Teilchen (Cl- )

Ionenverteilung:

IZR: viel K+ , wenig Na+

EZR: viel Na+ , wenig K+

%
%
\index{Osmose}Osmose und \index{Diffusion}Diffusion

GAB: Osmose und Diffusion: Grunds\"atzlich wichtig und relevant,
allenfalls bessere Erkl\"arung, aber bleibt in Normalschrift, nicht
heller. [GAB]

KOCH: Sollte alles grau werden

Geladene Teilchen k\"onnen die Zellmembran nicht selbstst\"andig
durchdringen, Wasser dagegen schon. Man nennt die Zellmembran daher
eine {\quotedblbase}\index{semipermeable Membran}semipermeable
Membran{\textquotedblleft}. 

Die Natur m\"ochte dass in einer Fl\"ussigkeit \"uberall die gleiche
Konzentration von geladenen Teilchen herrscht. Da die Zellmembran (und
teilweise auch die Blutgef\"a{\ss}wand) nur das Wasser durchl\"asst,
nicht jedoch die geladenen Teilchen verschiebt sich nur das Wasser von
einem Raum in den anderen. Man kann sagen dass die Elekktrolyte einen
Druck auf das Wasser aus\"uben {\quotedblbase}die Seiten zu
wechseln{\textquotedblleft}. Man nennt diesen Druck den
{\quotedblbase}\index{osmotischer Druck}osmotischen
Druck{\textquotedblleft}. Mann kann auch sagen {\quotedblbase}das
Wasser folgt dem Konzentrationsgef\"alle{\textquotedblleft}.

\begin{center}
 [Warning: Image ignored] % Unhandled or unsupported graphics:
%\includegraphics[width=6.131cm,height=3.646cm]{AASdev20090228-img68.gif}

\end{center}
Daneben gibt es den {\quotedblbase}onkotischen Druck{\textquotedblleft}
der durch Eiwei{\ss}bestandteile ausge\"ubt wird.

Wie kommen die Elektrolyte in die einzelnen R\"aume wenn doch nur Wasser
die Wand passieren kann?

Daf\"ur gibt es aktive Transporter
({\quotedblbase}Pumpen{\textquotedblleft}) die die Ionen hin und her
bef\"ordern. Die Steuerung richtige Verteilung der verschiedenen Ionen
ist recht kompliziert aber lebenswichtig. Eine
{\quotedblbase}falsche{\textquotedblleft} Verteilung der Elektrolyte
kann schlimme Folgen haben (zB. Herzrhythmusst\"orungen). 

Die richtige Verteilung der Elektrolyte im K\"orper ist lebenswichtig.
Elektrolyst\"orungen k\"onnen lebensgef\"ahrlich sein. 

H\"aufige Ursachen von \index{Elektrolyst\"orungen}Elektrolyst\"orungen

\begin{itemize}
\item Niereninsuffizienz (gest\"orte Ausscheidung

\end{itemize}
\begin{itemize}
\item Besonders bei Dialysepatienten

\end{itemize}
\begin{itemize}
\item Durchfall, Erbrechen 

\item Medikamente

\item Vergiftung

\end{itemize}
Wasser- und Elektrolythaushalt

Blutvolumen = Plasmawasser + zellul\"are Bestandteile (Blutk\"orperchen)

Beim Erwachsenen 7-8\% KG

Beim Neugeborenen / S\"augling 8-9\% KG  (70-80\%  K\"orperwasser!)

Verluste bis 20\% ohne Funktionsbeeintr\"achtigung tolerierbar,
allerdings Sch\"aden in der  Mikrozirkulation

${\rightarrow}$ auch ohne klinische Schockzeichen  ist bei
Volumensverlust Volumensersatz  indiziert!

\begin{itemize}
\item osmotischer Druck

\end{itemize}
\begin{itemize}
\item Volumenregulation des K\"orpers

\end{itemize}
\begin{itemize}
\item Konzentrationsgef\"alle ${\rightarrow}$ Spannungsunterschied an
den Zellmembranen ={\textgreater} Grundlage f\"ur Erregungsleitung /
Muskelkontraktion

\item Regelung durch:

\end{itemize}
\begin{itemize}
\item Hormone (Renin/Angiotensin/Aldosteron)

\item Ausscheidung \"uber Niere

\end{itemize}
S\"aure-/Basenhaushalt 

\index{S\"aure}S\"auren = Verbindungen, die  Wasserstoffionen (H+)
abgeben k\"onnen

\index{Lauge}Laugen = \index{Base}Basen = Verbindungen, die H+ 
aufnehmen k\"onnen

\index{pH-Wert}pH-Wert

Der pH-Wert gibt an wie sauer oder basisch eine L\"osung ist.

= Menge der freien H+  in einer L\"osung.

pH 7 ist neutral, pH {\textgreater} 7 basisch, pH {\textless} 7 sauer

pH normal ca. um 7,4.

pH {\textless} 7,35  ={\textgreater} Azidose

pH {\textgreater} 7,45  ={\textgreater} Alkalose

Der richtige pH-Wert ist lebenswichtig. Der K\"orper ist da sehr
empfindlich, der pH-Wert muss sich in engen grenzen bewegen. Schon
geringe Abweichungen k\"onnen gro{\ss}e Folgen haben. 

Leider gibt es im Rettungsdienst normalerweise keine M\"oglichkeit den
pH-Wert zu messen, jedoch haben wir gro{\ss}en Einflu{\ss} auf ihn
durch unsere Ma{\ss}nahmen, siehe unten. 

Bei einer \"Ubers\"auerung sprechen wir von einer
\index{Azidose}Azidose, wenn der Patient basisch wird von einer
\index{Alkalose}Alkalose.

\paragraph[\ Regulationsmechanismen]{ Regulationsmechanismen}
\begin{itemize}
\item Blut mit Puffersystem

\end{itemize}
\begin{itemize}
\item  %
%
H+ + HCO3- ${\leftarrow}$${\rightarrow}$ H2 O + CO2  (Formel nicht
lernen!)

\item S\"aure + Bikarbonat ${\leftarrow}$${\rightarrow}$ Wasser + CO2

\end{itemize}
\begin{itemize}
\item Atmet der Patient nicht (genug) staut sich das CO2 und bildet mit
Wasser H+ ${\rightarrow}$ Atmungsbedingte \"Ubers\"auerung

\item Atmet der Patient zu viel wird zu viel CO2 abgeatmet und wird von
H+ und dem anderen Stoff nachgebildet ${\rightarrow}$ H+ fehlt
${\rightarrow}$ Atmungsbedingte Alkalose

\item F\"allt zu viel H+ (S\"aure) im K\"orper an (durch den
Stoffwechsel, Stoffwechselbedingte \"Ubers\"auerung) wird es zu (Wasser
und) CO2 umgewandelt, das CO2 wird abgeatmet

\end{itemize}
\begin{itemize}
\item Lunge 

\end{itemize}
\begin{itemize}
\item Abatmung von CO2 ${\rightarrow}$ Einfluss auf den Puffer (s.o.)

\end{itemize}
\begin{itemize}
\item Niere

\end{itemize}
\begin{itemize}
\item Ausscheidung (oder Wiederaufnahme)  von H+

\end{itemize}
\subparagraph[Atmungsbedingte \"Ubers\"auerung - Respiratorische
Azidose]{Atmungsbedingte \"Ubers\"auerung - Respiratorische Azidose}
\label{bkm:RefHeading35233321}\index{\"Ubers\"auerung}\index{Respiratorische
Azidose}durch Hypoventilation / verminderten  Gasaustausch:

CO2  Abatmung vermindert ${\rightarrow}$ Hyperkapnie,  Ans\"auerung des
Blutes

Ursachen

Asthma

Lungen\"odem

Atemstillstand / insuffiziente Atmung

Therapie

Erh\"ohung des Atem-Minuten-Volumens

\subparagraph[Stoffwechselbedingte \"Ubes\"auerung -Metabolische
Azidose]{Stoffwechselbedingte \"Ubes\"auerung -Metabolische Azidose}
\index{\"Ubes\"auerung}\index{Metabolische Azidose}Zunahme von
organischen S\"auren im  K\"orper

Ursachen

\begin{itemize}
\item Nierenleistung 

\item Schock

\item Durchfall

\item diabetisches Koma

\end{itemize}
Gefahren

\begin{itemize}
\item Sch\"adigung des Zellstoffwechsels

\item Sch\"adigung des Herzmuskels

\item respiratorische und metabolische  Alkalose

\end{itemize}
\subparagraph[Atmungsbedingte (respiratorische)
Alkalose]{Atmungsbedingte (respiratorische) Alkalose}
\index{Alkalose}Hyperventilationssyndrom

Elektrolytverschiebung: Kribbeln in den Fingern, Pf\"otchenstellung,
Kr\"ampfe

\subparagraph{Stoffwechselbedingte (metabolische) Alkalose}
z.B. durch anhaltendes Erbrechen: Verlust von S\"aure

Zusammenfassung [F055?]

Wer zu wenig atmet wird sauer.

Wer zu viel atmet wird basisch und sp\"urt ein Kribbeln in den Fingern. 

Wer Probleme mit dem Stoffwechsel hat kann auch sauer oder basisch
werden.

\subsection[SCHOCK ]{SCHOCK }
[F0EF?][F068?][F0F0?][F0EB?][F0E9?][F0EA?][F0E1?][F078?][F073?][F072?][F069?][F055?][F04F?][F04E?]

Definition

{\quotedblbase}Lebensbedrohliche Kreislaufst\"orung aufgrund von
Unterversorgung des K\"orpers mit Blut und Sauerstoff infolge des
Versagens des Kreislaufsystems.{\textquotedblleft} 

Es k\"onnen einzelne oder mehrere Teile des Kreislaufsystems betroffen
sein  (Blut, Gef\"a{\ss}e, Herz).

Ob ein Schock oder eine Schockgefahr besteht wird anhand von
Schocksymtpomen und anhand der Verdachtsdiagnose beurteilt.

Schock bedeutet:

\begin{itemize}
\item akute Kreislaufinsuffizienz = Mi{\ss}verh\"altnis Herzauswurf - 
Sauerstoffbedarf des K\"orpers  

\item M\"ogliche Ursachen: 

\end{itemize}
\begin{itemize}
\item Herz pumpt wenig

\item Zu wenig Blut vorhanden 

\item Gef\"a{\ss}e sind zu weit gestellt und Blut versackt

\end{itemize}
Im Endeffekt kommt immer zu wenig Blut=Sauerstoff zu den Organen,
dadurch ist die Ge\-webedurchblutung nicht ausreichend  f\"ur die
Versorgung aller Organe mit Sauer\-stoff.

K\"orpereigene Mechanismen k\"onnen gegensteuern und versuchen den
Schock zu kompensieren (gegensteuern). Dies geht allerdings nur eine
gewisse Zeit und in einem gewissen Ausma{\ss} gut, daher ist eine
schnell  einsetzende Therapie wichtig! Bei l\"anger bestehendem Schock
kann dieser Zustand irreversibel werden. 

Kein Schock besteht bei kurzzeitigen  Bewu{\ss}tseinsst\"orungen wie
Kollaps/Synkope (kurzfristige Blutverteilungsst\"orung, kurzfristige
zerebrale Isch\"amie,  fl\"uchtige Symptome)

Schock -  Kompensationsmechanismen

Der K\"orper versucht selber dagegen etwas zu tun und gegenzusteuern
({\quotedblbase}kompensieren{\textquotedblleft}). Der wichtigste
Mechanismus ist die \index{Zentralisation}Zentralisation. Darunter
versteht man die Blutumverteilung zugunsten lebenswichtger Organe wie
Herz, Gehirn, Lunge und Leber. Dabei werden z.B. Haut, Muskel und Darm
sowie die Extremit\"aten (=Peripherie) kaum mehr  durchblutet.
Periphere Venen kollabieren und die Rekapillarisierungszeit
verl\"angert. Wenn die Kompensationsma{\ss}nahmen  versagen spricht man
vom {\quotedblbase}dekompensierter Schock{\textquotedblleft}. Es
besteht sp\"atestens dann akute Lebensgefahr![F0EA?] 

Beachte: 

\begin{center}
 [Warning: Image ignored] % Unhandled or unsupported graphics:
%\includegraphics[width=3.8cm,height=2.611cm]{AASdev20090228-img69.png}

\end{center}
Bei jedem Patienten ist zu erheben ob ein Schock oder eine realistische
Schockgefahr besteht oder nicht bzw. festzustellen um welche Art von
Schock es sich handelt (s.u.)!

Achtung!

Jeder Schockzustand ist grunds\"atzlich ein Fall f\"ur den NOTARZT!

Allgemeine Schocksymptomatik

\begin{center}
\begin{tabular}{|m{3.9750001cm}m{9.325cm}|}
\hline
 &
\centering Symptome\par

\\\hline
\centering \index{Schockgefahr}Schockgefahr\par

\centering [263A?]\par

 &
\begin{itemize}
\item Keine eindeutigen Anzeichen

\item Vorausschauend handeln!

\item Bedenke Verdachtsdiagnose.  Muss ich damit rechnen dass sich ein
Schock entwickelt?  Besteht Schockgefahr?

\end{itemize}
\\\hline
\centering Schock\par

\centering [E426?]\par

 &
\begin{itemize}
\item HF ${\uparrow}$

\item m\"a{\ss}iger RR-Abfall

\item RR noch {\textgreater} HF

\item neurologisch unauff\"allig, ev. Euphorisch oder agitiert

\item Haut bla{\ss}, evt. kaltschwei{\ss}ig

\item Venen sind kollabiert oder gestaut

\item Peripherie schlecht durchblutet

\end{itemize}
\\\hline
\centering Schwerer Schock\par

\centering Vollbild\par

\centering [E427?]\par

 &
\begin{itemize}
\item Tachykardie (HF ${\uparrow}$,
{\quotedblbase}fliehender{\textquotedblleft} Puls)

\item RR-Abfall (schwacher Puls)

\end{itemize}
\begin{itemize}
\item Herzfrequenz {\textgreater} systolischer RR

\end{itemize}
\begin{itemize}
\item Tachypnoe (schnelle, flache Atmung)

\item neurologisch auff\"allig, getr\"ubt

\item Haut bla{\ss}, zyanotisch (bl\"aulich), feucht

\item Venen kollabiert oder gestaut

\item nur mehr schwer behebbar! ${\rightarrow}$ Schock  m\"oglichst
fr\"uhzeitig erkennen \&  behandeln

\end{itemize}
\\\hline
\centering Endstadium\par

\centering [E429?]\par

 &
\begin{itemize}
\item Haut grau, klebrig

\item Versagen der Herzfunktion

\item RR kaum mehr me{\ss}bar

\item evt. Koma

\item evt. Schnappatmung

\item Patient ist dabei zu sterben.

\end{itemize}
\\\hline
\end{tabular}
\end{center}
Schockschere

Normalerweise ist die Herzfrequenz niedriger als der systolische
Blutdruck. Wenn der K\"orper versucht den Schock zu kompensieren kann
sich dieses Verh\"altnis umkehren. Man spricht dann von der
Schockschere, der Patient befindet sich dann bereits schon im Vollbild
eines Schocks! 

Achtung: Kinder haben andere Normalwerte f\"ur HF und RR!

[Warning: Draw object ignored]

Wenn sich das Verh\"altnis Herzfrequenz zu Systolischem Blutdruck
umkehrt befindet sich der Patient bereits im Schock. 

\label{bkm:RefHeading38104902}Allgemeine Ma{\ss}nahmen der
\index{Schockbek\"ampfung}Schockbek\"ampfung

Grundsatz: Der Schock muss bek\"ampft werden bevor man ihn bemerkt.

\begin{itemize}
\item Notarzt!

\item Situationsgerechte  Lagerung, je nach Diagnose und Schockart

\item Beengende Kleidung \"offnen

\item Sauerstoffgabe hochdosiert (ab 6l/min bei Schockgefahr, 10l/min
oder mehr bei bestehendem Schock, je nach Schwere%
%
)

\end{itemize}
Freigabe: Sauerstoffgabe hochdosiert (ab 6l/min bei Schockgefahr,
10l/min oder mehr bei bestehendem Schock, je nach Schwere

\begin{itemize}
\item Monitoring (Pulsoxy, RR-Manschette angelegt lassen)

\item Reanimationsbereitschaft 

\item W\"armeerhalt

\item Psychische Betreuung

\item Weitere Massnahmen je nach Schockart 

\end{itemize}
Schock - Arten

\begin{center}
\begin{tabular}{|m{4.8430004cm}|m{8.457cm}|}
\hline
\centering Mechanismus\par

 &
\centering Schockart\par

\\\hline
Herz pumpt zu wenig

 &
\begin{itemize}
\item Kardiogender Schock

\end{itemize}
\\\hline
Zu wenig Blut vorhanden 

 &
\begin{itemize}
\item Hypovol\"amischer Schock

\end{itemize}
\\\hline
Gef\"a{\ss}e inad\"aquat weitgestellt, Blut versackt

 &
\begin{itemize}
\item Anaphylaktischer Schock

\item Septischer Schock, toxischer Schock

\item Neurogener Schock

\end{itemize}
\\\hline
\end{tabular}
\end{center}
\paragraph{Hypovol\"amischer Schock}
Schock aufgrund eines Verlustes von Blutvolumen aus den Gef\"a{\ss}en.

Es kann dabei verloren gehen: 

\begin{itemize}
\item Blut

\item Plasma

\item Wasser und Elektrolyte

\end{itemize}
Ursachen

\begin{center}
\tablehead{}\begin{supertabular}{|m{6.649cm}|m{6.651cm}|}
\hline
Fl\"ussigkeitsverlust nach au{\ss}en

 &
Fl\"ussigkeitsverlust nach innen

\\\hline
\begin{itemize}
\item Blutung

\item gastrointestinal (Durchfall, Diarrhoe )

\item Verbrennungen (Hautbarriere versagt  gro{\ss}\-fl\"achig)

\item renal (Diabetes mellitus, Diuretika, Polyurie bei  Nierenversagen)

\end{itemize}
 &
\begin{itemize}
\item Weichteilblutung

\end{itemize}
\begin{itemize}
\item Nach Trauma, rupturierte Eileiter\-schwanger\-schaft, etc

\end{itemize}
\begin{itemize}
\item Wassereinlagerung ins Gewebe

\item Aszites ({\quotedblbase}Wasserbauch{\textquotedblleft} bei
Lebererkrankungen  etc.)

\item Ileus (Darmverschlu{\ss})

\item H\"amatothorax, H\"amatoperiton\"aum

\end{itemize}
\\\hline
\end{supertabular}
\end{center}
Bei akutem Fl\"ussigkeitsverlust kann ein Verlust von bis zu 20\% des
Blutvolumens  gut kompensiert werden. Auf gr\"o{\ss}ere Verluste  folgt
ein Absinken von Herzauswurf und RR!

Spezielle Ma{\ss}nahmen

\begin{itemize}
\item Allgemein: Allgemeine Ma{\ss}nahmen der Schockbek\"ampfung (siehe
oben) 

\item Blutstillung  (weitere Verluste verhindern!)

\item ggfs. Trauma-Vorgehen beachten

\item %
%
Wenn m\"oglich schonende Bergung

\item Lagerung: grunds\"atzlich %
%
Beine hoch  Wichtige Ausnahmen: {\textperiodcentered} kardiale Ursache,
{\textperiodcentered} Kopf- / Brust\-ver\-letzungen,
{\textperiodcentered} Hirndruck (SHT),  {\textperiodcentered}
Knochenbr\"uche der Beine oder des Beckens, {\textperiodcentered}
Bewusstlosigkeit ${\rightarrow}$ sonst F%
%
lachlagerung oder stabile Seitenlage

\item evt. Aviso (Ank\"undigung) im  Krankenhaus / Schockraum

\item Abtransport: rasch, schonend, nicht \"uberst\"urzt

\end{itemize}
Begriff {\quotedblbase}Schocklagerung{\textquotedblleft} ein f\"ur alle
mal in der Ausbildung streichen, nur Hinweis dass dieser Ausdruck
veraltet ist sollten sie ihn wo aufschnappen. Stattdessen
{\quotedblbase}Beine-hoch{\textquotedblleft}-Lagerung.

Wenn KI gegen Beine hoch: Flachlagerung bzw. stab. Seitenlage empfehlen?

\paragraph{Kardiogener Schock}
Schock infolge mangelnder Herzpumpleistung.

\begin{center}
\tablehead{}\begin{supertabular}{|m{6.649cm}|m{6.651cm}|}
\hline
Ursachen

\begin{itemize}
\item Herzversagen, Herzinsuffizienz

\end{itemize}
\begin{itemize}
\item Myokardinfarkt (h\"aufigste Ursache)

\item Herzrhythmusst\"orungen

\end{itemize}
\begin{itemize}
\item Lungenembolie

\end{itemize}
 &
charakterisiert durch

\begin{itemize}
\item RR-Abfall

\item eingeschr\"ankte Pumpleistung

\end{itemize}
\begin{itemize}
\item Blut staut sich zur\"uck

\end{itemize}
\\\hline
\end{supertabular}
\end{center}
Spezielle Ma{\ss}nahmen

\begin{itemize}
\item Allgemein: Allgemeine Ma{\ss}nahmen der Schockbek\"ampfung (siehe
oben) 

\item Lagerung: Oberk\"orper hoch! Evt. Beine h\"angen  lassen. 

\item medikament\"ose Therapie durch Notarzt  (Diuretika, Katecholamine,
Heparin)

\item Reanimationsbereitschaft  

\end{itemize}
\paragraph{Anaphylaktischer Schock}
Schock aufgrund einer \"uberschie{\ss}enden allergischen Reaktion.

{\quotedblbase}Allergischer Schock{\textquotedblleft}. Ausgel\"ost durch
\"uberschie{\ss}ende Reaktion der K\"orerabwehr (Immunreaktion) auf
bestimmte Stoffe (Antigene). 

\begin{center}
 [Warning: Image ignored] % Unhandled or unsupported graphics:
%\includegraphics[width=2.641cm,height=3.837cm]{AASdev20090228-img70.jpg}

\end{center}
Es werden s\"amtliche Gef\"a{\ss}e maximal erweitert, zus\"atzlich
steigt die Durchl\"assigkeit der Gef\"a{\ss}e
(Gef\"a{\ss}permeabilit\"at ${\uparrow}$) ${\rightarrow}$ 
Plasmaverlust in den Zwischengebwebsraum (Interstitium) 

RR f\"allt massiv ab!

M\"ogliche Antigene

\begin{itemize}
\item Pollen

\item Medikamente, Antibiotika

\item Kontrastmittel

\item tierische Gifte (Bienenstich etc.)

\item und vieles andere ...

\end{itemize}
Symptome%
%


\begin{center}
\tablehead{}\begin{supertabular}{|m{3.31cm}|m{9.99cm}|}
\hline
\centering Stadium\par

 &
\\\hline
Stadium I:

 &
Hautreaktion: evt. Ausschlag, R\"otung;
{\quotedblbase}Quaddeln{\textquotedblleft}, Flush, \"Odeme

Allgemeinsymptome: Juckreiz, Zittern, Schwindel, Kopfschmerz

\\\hline
Stadium II:

 &
RR ${\downarrow}$  , Tachykardie

\"Ubelkeit, Erbrechen, Durchfall 

\\\hline
Stadium III:

 &
Allgemeine Schocksymptome

Bronchospasmus

\\\hline
Stadium IV:

 &
Atem-/Kreislaufstillstand

Schwerer Schock

\\\hline
\end{supertabular}
\end{center}
Spezielle Ma{\ss}nahmen [F0EF?][F068?][F0F0?]

\begin{itemize}
\item Allgemein: Allgemeine Ma{\ss}nahmen der Schockbek\"ampfung (siehe
oben) 

\item Antigen  nach M\"oglichkeit entfernen!

\item Lagerung: situationsgerecht. Druckverlust im Vordergrund: Beine
hoch.%
%
 Atemprobleme im Vordergrund: stark erh\"ohter Oberk\"orper

\end{itemize}
\paragraph[Septischer Schock und Toxischer Schock]{Septischer Schock und
Toxischer Schock}
\label{bkm:RefHeading41284910}Siehe auch: Kap. [RS II] Hygiene 
{\guillemotright} Sepsis
[K.II{\guillemotright}\pageref{bkm:RefHeading41252510}, unten]

Schock infolge einer schweren, k\"orperweiten Infektion oder aufgrund
eines Gef\"a{\ss}-beeinflussenden Giftstoffes (Toxins).

Durch bakterielle Gef\"a{\ss}wandsch\"adigung, temperaturbedingt (hohes
Fieber {\textgreater} 38{\textdegree}C) sowie aufgrund von Giftstoffen
kommt es zu Fl\"ussigkeitsaustritt in den Zwischen\-gewebs\-raum sowie
zu Gef\"a{\ss}weitstellung.

Ursachen

\begin{center}
\tablehead{}\begin{supertabular}{|m{6.649cm}|m{6.651cm}|}
\hline
Diverse Erreger und Toxine

\begin{itemize}
\item Langanhaltender Entz\"undungsherd der pl\"otzlich Erreger
{\quotedblbase}ins Blut streut{\textquotedblleft} und
{\quotedblbase}explodiert{\textquotedblleft}

\item Besiedelte Fremdk\"orper (Diverse medizinische Sonden, Venflons,
etc.), siehe Sepsis
[K.II{\guillemotright}\pageref{bkm:RefHeading41252510}, unten]

\item Tampons

\item ...

\end{itemize}
 &
Symptome

\begin{itemize}
\item Schock

\item hohes Fieber oder kein Fieber

\item Haut hei{\ss}, trocken, rot oder bleich, grau,
{\quotedblbase}marmoriert{\textquotedblleft}

\end{itemize}
\\\hline
\end{supertabular}
\end{center}
Spezielle Ma{\ss}nahmen

\begin{itemize}
\item Allgemein: Allgemeine Ma{\ss}nahmen der Schockbek\"ampfung (siehe
oben) 

\item Lagerung: situationsgerecht, Beine hoch

\item Schutz vor Unterk\"uhlung/\"Uberw\"armung

\item Aviso im Krankenhaus / Intensivstation
({\quotedblbase}\"Uberwachungsbett{\textquotedblleft})

\item ben\"otigt Antibiotika (im KH)

\end{itemize}
\paragraph{Neurogener Schock}
Zusammenbruch der nerv\"osen Kreislaufregulation, inad\"aquate
arterielle und ven\"ose Gef\"a{\ss}erweiterung (Vasodilatation)

Ursachen 

\begin{itemize}
\item Neurologische Erkrankungen wie Querschnitt-Verletzung des
R\"uckenmarks

\item Die Gef\"a{\ss}e bekommen keine Nerven-Impulse und kontrahieren
sich nicht mehr.

\item ${\rightarrow}$ Blut versackt

\end{itemize}
Spezielle Ma{\ss}nahmen

\begin{itemize}
\item Allgemein: Allgemeine Ma{\ss}nahmen der Schockbek\"ampfung (siehe
oben) 

\item Lagerung: Flachlagerung (WS-Verletzung ist Kontraindikation gg.
Beine hoch!)

\end{itemize}
\section{[RS VI] Krankheiten und Notf\"alle}
Sebastian Gabriel, Michael Withofner

Changelog:

2009-02-25: Start Changelog. Bis inkl. bakt. Meningitis mit Korr. von
Koch verglichen und korr.

Eine willk\"urliche Einteilung

\begin{itemize}
\item St\"orungen des Bewu{\ss}tseins --  neurologische Notf\"alle

\item pulmonale Notf\"alle

\item Kreislaufst\"orungen

\item Gef\"a{\ss}verschl\"usse

\item Abdominelle Notf\"alle

\item Urologische Erkrankungen

\end{itemize}
Begriffe

Symptome: Krankheitserscheinungen

Komplikationen: m\"oglicherweise auftretende Schwierigkeiten

Therapie: Sinnvolle Ma{\ss}nahmen zur Verbesserung der  Symptome

Differentialdiagnose (DD): Erkrankung(en) mit \"ahnlichen Symptomen, 
Ver\-wechslungs\-gefahr

\subsection[Herz-Kreislaufst\"orungen]{Herz-Kreislaufst\"orungen}
\label{bkm:RefHeading42342802}Der Herz-Kreislaufstillstand wird unter 
Herz- / Kreislauf {\textgreater} Kreislaufstillstand 
(H.III{}-\pageref{bkm:RefHeading42514602}, oben) besprochen.

\begin{itemize}
\item Herz-/Kreislaufstillstand

\item Herzinsuffizienz

\end{itemize}
\begin{itemize}
\item Linksherz-Versagen

\item Rechtsherz-Versagen

\end{itemize}
\begin{itemize}
\item (Akutes) Koronarsyndrom

\end{itemize}
\begin{itemize}
\item Angina pectoris

\item Myokardinfarkt (MCI )

\end{itemize}
\begin{itemize}
\item Herzrhythmusst\"orungen

\item Hypertonie

\end{itemize}
\begin{itemize}
\item hypertensive Krise

\item hypertensiver Notfall

\end{itemize}
Herzinsuffizienz

Unzureichende Auswurfleistung des Herzens.

Leistungseinschr\"ankung des Herzmuskels f\"uhrt zu einer verminderten
Auswurfleistung, viele Ursachen sind m\"oglich. Manche Ursachen treten
akut auf, manche sind eher chronischer Natur.

Man unterteilt in \index{Rechtsherzinsuffizienz}Rechtsherzinsuffizienz,
\index{Linksherzinsuffizienz}Linksherzinsuffizienz  und
\index{Globalinsuffizienz}Globalinsuffizienz, je nachdem welcher
Herzteil betroffen ist (rechtes, linkes oder ganzes Herz). 

Beurteilung des Schweregrades

Es ist wichtig den Schweregrad einer Herzinsuffizienz zu beurteilen. Sie
muss nicht dramatisch sein, f\"ur viele Pateinten ist sie ein
chronischer Zusatnd mit dem sie gut leben k\"onnen. Problematisch wird
es jedoch wenn der Patient pl\"otzlich eine Symptomatik entwickelt und
die Kompenastionsmechanismen (k\"opereigene, Medikamente, {\dots})
versagen. Man spricht dann von der dekompensierten Herzinsuffizienz.
Man muss die Symptome generell anhand ihres Verlaufes beurteilen:
Bekannte, seit l\"angerer Zeit bestehende Symptome sind weniger wichtig
als pl\"otzlich auftretnde Symptome, diese sind meistens Alarmzeichen. 

Weitere klassiche Alamrzeichen sind Symptome wie Atemnot, pl\"otzlich
auftretende \"Odeme oder akute Schmerzen. Ein spezielles,
hochgef\"ahrliches, Alarmzeichen ist das brodelnde Atemger\"ausch beim
Lungen\"odem. 

Das kardiale Lungen\"odem wird im Kaptel
I.II{\guillemotright}\pageref{bkm:RefHeading36561325} (Lungen\"odem)
behandelt. 

\paragraph[Linksherzinsuffizienz]{Linksherzinsuffizienz}
\index{Linksherzinsuffizienz}Die Linksherzinsuffizienz entsteht aufgrund
mangelnder Auswurfleistung des linken Herzens. Durch den R\"uckstau von
Blut in den Lungenkreislauf kann es zu einem kardialen Lungen\"odem
kommen. 

Patienten mit einem kardialem Lungen\"odem sind lebensbedrohlich krank.
Das Herz vebraucht oft schon die allerletzten Reserven, schon die
Aufregung durch den Transport im Stiegenhaus zum Auto kann das Herz
endg\"ultig in den Herz-Kreislaufstillstand bringen. 

Es kommt immer wieder vor das solche Patienten im Stiegenhaus oder im
Auto reanimationspflichtig werden. 

Kein Transport von Patienten mit Lungen\"odem ohne Notarzt. Auch wenn
das Spital {\quotedblbase}um{\textquotesingle}s Eck{\textquotedblleft}
ist.

Symptome einer Linksherzinsuffizienz

\begin{itemize}
\item Dyspnoe !

\end{itemize}
\begin{itemize}
\item Orthopnoe: Patient bekommt nur mehr im Sitzen Luft, kann nicht
flach schlafen, etc.

\end{itemize}
\begin{itemize}
\item Tachypnoe

\item Tachykardie 

\item Hustenreiz !

\item Unruhe

\item Hypotonie !

\item Lungenstauung/Lungen\"odem (Blut staut sich zur\"uck in den
Lungenkreislauf)

\end{itemize}
Spezielle Ma{\ss}nahmen beim vital bedrohten Patienten

\begin{itemize}
\item Erstma{\ss}nahmen bei vital bedrohten Patienten: 

\end{itemize}
\begin{itemize}
\item Lagerung: Oberk\"orper hoch

\item O2 10~l, bei Bedarf bis zu 15~l

\item Notarzt

\item Reanimationsbereitschaft

\item Vitalwerte erheben \& Monitoring 

\end{itemize}
\begin{itemize}
\item Ursachenforschung

\end{itemize}
Spezielle Ma{\ss}nahmen bei nicht vital bedrohten Patienten

\begin{itemize}
\item Lagerung: Oberk\"orper hoch

\item O2 je nach Symptomen 6-8~l, bei Bedarf bis zu 15~l

\item Vitalwerte erheben \& Monitoring 

\end{itemize}
\paragraph{Rechtsherzinsuffizienz}
Die Rechtsherzinsuffizienz entsteht aufgrund mangelnder Auswurfleistung
des rechten Herzens. 

Symtome

\begin{itemize}
\item Bein-/Haut\"odeme

\item N\"achtliches Wasserlassen

\item gestaute Halsvenen 

\begin{center}
 [Warning: Image ignored] % Unhandled or unsupported graphics:
%\includegraphics[width=4.518cm,height=6.03cm]{AASdev20090228-img71.png}

\end{center}
\item Lebervergr\"o{\ss}erung, evt. mit Aszites
({\quotedblbase}Wasserbauch{\textquotedblleft})

\end{itemize}
Spezielle Ma{\ss}nahmen beim vital bedrohten Patienten

\begin{itemize}
\item Erstma{\ss}nahmen bei vital bedrohten Patienten: 

\end{itemize}
\begin{itemize}
\item Lagerung: Oberk\"orper hoch

\item O2 10~l, bei Bedarf bis zu 15~l

\item Notarzt

\item Reanimationsbereitschaft

\item Vitalwerte erheben \& Monitoring 

\end{itemize}
\begin{itemize}
\item Ursachenforschung

\end{itemize}
Spezielle Ma{\ss}nahmen bei nicht vital bedrohten Patienten

\begin{itemize}
\item Lagerung: Oberk\"orper hoch

\item O2 je nach Symptomen 6-8~l, bei Bedarf bis zu 15~l

\item Vitalwerte erheben \& Monitoring 

\end{itemize}
Koronare Herzkrankheit und Akutes Koronarsyndrom

\index{Koronare Herzkrankheit}Koronare Herzkrankheit (\index{KHK}KHK):
Erkrankung der Herzkranzgef\"a{\ss}e (Koronargef\"a{\ss}e) bei der es
zu einer Verengung der Gef\"a{\ss}e kommt.

\index{Akutes Koronarsyndrom}Akutes Koronarsyndrom: Sammelbegriff f\"ur
Herzinfarkt, stabile und instabile Angina pectoris. Bei diesen
Krankheitsbildern kommt es zu einer akuten Unterversorgung des
Herzmuskels mit Blut und Sauerstoff. Die Ursache sind in erster Linie
verengte oder verstopfte Herzkranzgef\"a{\ss}e . 

\paragraph[Angina pectoris]{Angina pectoris}
\index{Angina pectoris}Pl\"otzliches Engegef\"uhl in der Brust meist mit
Thoraxschmerz infolge einer Unter\-ver\-sorgung des Herzmuskels mit
sauerstoffreichen Blut. Der Anfall vergeht wieder von selbst ohne
bleibende Sch\"aden.

\begin{itemize}
\item Minderdurchblutung des Herzmuskels (Myokard) aufgrund einer oder
mehrerer Eineingung(en) von Herzkranzgef\"a{\ss}ren (Koronararterien)
oder

\item Erh\"ohter Sauerstoffbedarf des Herzens bei Anstregung

\end{itemize}
${\rightarrow}$ Mi{\ss}verh\"altnis Sauerstoffangebot - 
Sauerstoffbedarf des Herzens ${\rightarrow}$ Schmerz!

Der akute Angina-Pectoris-Anfall geh\"ort zu der Familie des Akuten
Koronarsyndroms. Dazu geh\"ort auch der Herzinfarkt, dem der gleiche
Mechanismus wie der der Angina Pectoris zu Grunde liegt, nur mit
schwer\-wiegenderen Auswirkungen. 

Bei den meisten Patienten ist die Angina Pectoris schon bekannt, bzw. es
besteht eine {\quotedblbase}\index{Koronare Herz-Krankheit}Koronare
Herz-Krankheit (\index{KHK}KHK){\textquotedblleft}. Die KHK ist eine
Erkrankung der Herzkranzgef\"a{\ss}e, diese werden durch Ablagerung von
Fetten und Kalk gesch\"adigt und eingeengt, dadurch kann weniger Blut
durchflie{\ss}en. So kommt es zu einer Mangelversorgung des Herzmuskels
und zu den Symptomen der Angina pectoris (oder zu einem Herzinfarkt). 

Symptome

\begin{center}
 [Warning: Image ignored] % Unhandled or unsupported graphics:
%\includegraphics[width=4.594cm,height=5.956cm]{AASdev20090228-img72.png}
\captionof{figure}[Patient mit einem akuten Koronarsyndrom, vermutlich
Angina Pectoris. Bei kaltem Wetter, beim Stiegen steigen mit schwerem
Gep\"ack: Ein pl\"otzliches Engegef\"uhl im Brustkorb und ein
stechender Schmerz im linken Thorax.]{Patient mit einem akuten
Koronarsyndrom, vermutlich Angina Pectoris. Bei kaltem Wetter, beim
Stiegen steigen mit schwerem Gep\"ack: Ein pl\"otzliches Engegef\"uhl
im Brustkorb und ein stechender Schmerz im linken Thorax.}

\end{center}
\begin{itemize}
\item Brustschmerz mit Ausstrahlung

\item Todesangst

\item Schw\"achegef\"uhl

\item Dyspnoe

\item prim\"ar nicht sicher von MCI unterscheidbar 

\item Symptome treten eher bei Belastung auf und vergehen oft bei
Entspannung. 

\begin{center}
 [Warning: Image ignored] % Unhandled or unsupported graphics:
%\includegraphics[width=6.173cm,height=4.067cm]{AASdev20090228-img73.png}
\captionof{figure}[Schmerzausdehnung beim isch\"amischen
Herzschmerz.]{Schmerzausdehnung beim isch\"amischen Herzschmerz.}

\end{center}
\end{itemize}
Spezielle Ma{\ss}nahmen

\begin{itemize}
\item Erstma{\ss}nahmen bei vital bedrohten Patienten: 

\end{itemize}
\begin{itemize}
\item Lagerung: Oberk\"orper hoch

\item O2 10~l, bei Bedarf bis zu 15~l

\item Notarzt

\item Reanimationsbereitschaft

\end{itemize}
\begin{itemize}
\item Vitalwerte erheben \& Monitoring 

\item striktes Bewegungsverbot!

\item %
%
Pat. darf eigenes Nitro nur  bei RR {\textgreater}100 syst. nehmen!

\end{itemize}
??? eigenes Nitro nehmen?

Vorschlag: Nur wenn sein Nitro leer war und er es deswegen nicht
genommen hat bei RRsyst{\textgreater}110

Wenn er es schon genommen hat und es nicht geholfen hat dann nicht

\paragraph{Herzinfarkt - Myokardinfarkt (MCI)}
Verschlu{\ss} eines oder mehrer Herzkranzgef\"a{\ss}e mit
anschlie{\ss}endem Absterben des betroffenen Herzmuskelgewebes. 

Symptome

Die Symptome entsprechen im Wesentlichen denen der Angina pectoris, doch
vergehen diese nicht mehr und es bleibt auch bei optimaler Behandlung
ein Schaden am Herzmuskel. Bei Diabetikern kann es aufgrund der
Nervensch\"adigungen auch zu einem schmerzlosen Infarkt kommen, man
spricht dabei auch vom {\quotedblbase}stummen
Infarkt{\textquotedblleft}.

\begin{center}
 [Warning: Image ignored] % Unhandled or unsupported graphics:
%\includegraphics[width=4.022cm,height=4.285cm]{AASdev20090228-img74.png}

\end{center}
\begin{itemize}
\item Brustschmerz mit Ausstrahlung 

\item Druckgef\"uhl / Beklemmung {\quotedblbase}als w\"urde jemand am
Brustkorb stehen{\textquotedblleft}

\item Todesangst

\item Dyspnoe

\item Kardiogener Schock, inbesonders: 

\end{itemize}
\begin{itemize}
\item Kaltschweissigkeit

\item Hypotonie

\end{itemize}
\begin{itemize}
\item %
%
Nicht ausreichende Besserung auf Nitro

\item Sonderform: stummer Infarkt bei Diabetes  mellitus

\end{itemize}
Spzielle Ma{\ss}nahmen

\begin{itemize}
\item Erstma{\ss}nahmen bei vital bedrohten Patienten: 

\end{itemize}
\begin{itemize}
\item Lagerung: Oberk\"orper hoch

\item O2 10~l, bei Bedarf bis zu 15~l

\item Notarzt

\item Reanimationsbereitschaft

\end{itemize}
\begin{itemize}
\item Vitalwerte erheben \& Monitoring 

\item Beengende Kleidung \"offnen

\item Striktes Bewegungsverbot!

\item \end{itemize}
Weitere Therapie

Herzinfarktpatienten werden {\quotedblbase}besonders{\textquotedblleft}
behandelt und normalerweise nicht auf eine beliebige Interne Station
gebracht. Meistens stehen dem Notarzt zur Auswahl:

\begin{itemize}
\item Lyse-Therapie mittels eines speziellen Medikaments, oder

\item Einweisung in ein \index{Herzkatheterlabor}Herzkatheterlabor
({\quotedblbase}PCI{\textquotedblleft},
{\quotedblbase}\index{PTCA}PTCA{\textquotedblleft} {}- Perkutane
Transluminale Coronar Angioplastie)

\end{itemize}
\begin{itemize}
\item In Wien gibt es einen Dienstplan f\"ur die Katheterlabore

\end{itemize}
\begin{itemize}
\item U.a. verf\"ugbar in: AKH, Donauspital, KH Hietzing

\end{itemize}
Kein Transport ohne Notarzt! Auch nicht wenn das Spital
{\quotedblbase}um{\textquotesingle}s Eck{\textquotedblleft} ist.

Der Patient braucht vermutlich nicht irgendein beliebiges Spital sondern
ein dienst\-bereites Herzkatheterlabor. Der Notarzt muss entscheiden ob
das notwendig und sinnvoll ist. 

Auch wenn der Patient ein Einweisungsformular (Spitalszettel) von einem
Arzt hat (\"Arztefunkdient, Hausarzt) oder ein Hausarzt meint
{\quotedblbase}es sei eh nicht dringend{\textquotedblleft} undbedingt
Notarzt beiziehen.

Komplikationen -- Womit muss ich rechnen?

\begin{itemize}
\item Herzrhythmusst\"orungen (alle Arten), bis hin zum
Atem-/Kresilaufstillstand und Reanimation

\item kardiogener Schock

\end{itemize}
Herzrhythmusst\"orungen 

\begin{center}
\tablehead{}\begin{supertabular}{|m{4.367cm}|m{4.367cm}|m{4.367cm}|}
\hline
 &
 &
\\\hline
Kammerflimmern

 &
Reanimation [F07E?]

 &
\\\hline
Pulslose Ventrikul\"are Tachykardie

 &
Reanimation [F07E?] 

 &
\\\hline
Asystolie (HF = 0)

 &
Reanimation  

 &
\\\hline
Bradykardie (HF {\textless} 60/min)

 &
 &
\\\hline
Tachykardie (HF {\textgreater}100/min)

 &
 &
\\\hline
Vorhofflimmern

 &
regellose Erregung im Vorhof

oft symptomlos

 &
Sehr viele Leute haben Vorhofflimmern und leben ganz gut damit.

Komplikationen:

Tachykardie-Anf\"alle

Thrombenbildung

deshalb Therapie mit Gerinnungshemmenden
({\quotedblbase}blutverd\"unnenden{\textquotedblleft}) Medikamenten:
Marcoumar ${\rightarrow}$ Patient ist Blutungsgef\"ahrdet

\\\hline
 &
 &
\\\hline
\end{supertabular}
\end{center}
\subparagraph[Kammerflimmern]{Kammerflimmern}
NICHT verwchseln mit Vorhofflimmern!

\begin{itemize}
\item Reanimationspflichtig.

\item Die elektrische Erregung im Herzen ist ungerichtet, als w\"urden
alle Autos eines Parkplatzes gleichzeitig in alle Richtungen losfahren
${\rightarrow}$

\item ${\rightarrow}$ Kein Auswurfleistung 

\end{itemize}
\subparagraph{Pulslose Ventrikul\"are Tachkardie}
\begin{itemize}
\item Reanimationspflichtig 

\item Vorstufe zum Kammerflimmern, Herz kontrahiert so schnell dass es
sich nicht mehr f\"ullen kann${\rightarrow}$ 

\item ${\rightarrow}$ keine Auswurfleistung 

\end{itemize}
\subparagraph{Asystolie}
\begin{itemize}
\item Reanimationspflichtig 

\item keine elektrische Erregung${\rightarrow}$ 

\item ${\rightarrow}$ keine Auswurfleistung 

\end{itemize}
\subparagraph{}
\subparagraph[Vorhofflimmern]{Vorhofflimmern}
\begin{itemize}
\item Regellose Erregung im Vorhof

\item oft symptomlos

\item Sehr viele Leute haben Vorhofflimmern und leben ganz gut damit.

\item Komplikationen:

\end{itemize}
\begin{itemize}
\item Tachykardie-Anf\"alle

\item Thrombenbildung

\end{itemize}
\begin{itemize}
\item deshalb Therapie mit Gerinnungshemmenden
({\quotedblbase}blutverd\"unnenden{\textquotedblleft}) Medikamenten:
Marcoumar ${\rightarrow}$ Patient ist Blutungsgef\"ahrdet

\end{itemize}
\paragraph{Zu schnell oder zu langsam: Tachy- und Bradykardie}
Denke bei der Beurteilung auch immer an die Umst\"ande und den
Erregungszustand des Patienten

\subparagraph{Bradykardie }
HF {\textless} 60/min (beim Erwachsenen)

Leistungssportler k\"onnen augrund ihres trainierten K\"orpers auch sehr
niedrige Ruhepulsfrequenzen haben. 

\subparagraph{Tachykardie }
HF {\textgreater}100/min (beim Erwachsenen)

Tachykarde Attacken sind ein h\"aufiges Notfallbild. Da oft chronische
Erkrankungen die Ursache sind  haben die Patienten diese Attacken immer
wieder und kennen die Symptome bereits (und oft auch die Behandlung). 

Symptome

\begin{itemize}
\item Herzklopfen

\item evt. Leichter Brustschmerz

\end{itemize}
Gefahr

Da das Herz pl\"otzlich viel mehr Arbeit leisten muss kann es sein dass
es an seine Grenzen st\"osst und versagt bzw. die Blutversorgung des
Herzens f\"ur diese Arbeit nicht mehr reicht und ein Herzinfarkt
entsteht. 

Ma{\ss}nahmen

\begin{itemize}
\item Erstma{\ss}nahmen bei vital bedrohten Patienten: 

\end{itemize}
\begin{itemize}
\item Lagerung: Oberk\"orper hoch

\item O2 10~l, bei Bedarf bis zu 15~l

\item Notarzt

\item Reanimationsbereitschaft

\end{itemize}
\begin{itemize}
\item Vitalwerte erheben \& Monitoring 

\item Beengende Kleidung \"offnen

\item Striktes Bewegungsverbot!

\item Psychische Betreuung! Oase der Ruhe schaffen. 

\end{itemize}
Hypotonie

RR {\textless} 100 mmHg syst.

\begin{multicols}{2}
Ursachen

Cardial

Volumenmangel, Exsikkose

Regulationsst\"orung (vasovagal)

Sportler, junge Frauen

Symptome

M\"udigkeit bis Kollaps

Kollaps

passagere Kreislaufregulationsst\"orung  mit RR-Abfall,
Blutverteilungsst\"orung  ohne Volumsverlust

RR-Abfall mit Bewu{\ss}tseinseintr\"ubung

{\quotedblbase}schwarzwerden vor Augen{\textquotedblleft}, Ohnmacht

Sturz = {\quotedblbase}Therapie{\textquotedblleft}

\end{multicols}
\paragraph{Kollaps, Synkope}
Symptome

\begin{itemize}
\item RR ${\downarrow}$

\item Schwindel (Vertigo )

\item ev. kurze Bewu{\ss}tlosigkeit

\item Patient(in) sackt zusammen

\item Tachykardie

\item Schwarzwerden vor Augen

\item kurze Bewu{\ss}tlosigkeit

\item Bradykardie

\item Ohnmacht

\item pl\"otzlicher RR-Abfal

\end{itemize}
Patient wacht auf dem Boden liegend  auf, wei{\ss} nicht, was geschehen
ist...

Ursachen:

\begin{itemize}
\item kardial

\item vaskul\"ar

\item zerebrovaskul\"ar

\item zerebral

\end{itemize}
ERKL\"AREN!

Therapie:

\begin{itemize}
\item sofern notwendig:

\item Beine hoch, 

\item evt. O2

\item Synkope

\item \item Hospitalisierung  zur Abkl\"arung  unbedingt erforderlich

\end{itemize}
Hypertonie

\index{Hypertonie}Hypertonie (\index{Bluthochdruck}Bluthochdruck): Oft
symptomlose chronische Erkrankung welche auf Dauer schwere Sch\"aden am
K\"orper anrichten kann und im Zuge von pl\"otzlichen Blutdruckkrisen
auch akut Probleme machen kann. 

\begin{itemize}
\item Chronische Hypertonie: Grenze derzeit bei {\textgreater}140/90 
(au{\ss}er bei Kindern)

\end{itemize}
Bei Schwangeren sind die Grenzwerte peinlich genau zu beachten! Denke an
die EPH-Gestose\footnote{Die EPH-Gestose wird in Kap.
I.VII{\guillemotright}\pageref{bkm:RefHeading39226112}: EPH-Gestose --
Pr\"aeklampsie - Eklampsie erkl\"art.}!

Ursachen

\begin{itemize}
\item Widerstandserh\"ohung (durch Engstellung der Gef\"a{\ss}e)

\item vermehrte Herzarbeit

\item Sympathikusaktivierung, Adrenalin+Cortison (Stre{\ss}hormone)
durch Dauerstre{\ss},  Bewegungsmangel, genetische Veranlagung

\end{itemize}
chronische Sch\"adigungen:

\begin{itemize}
\item Gef\"a{\ss}sch\"adigung

\end{itemize}
\begin{itemize}
\item Herzkranzgef\"a{\ss}e ${\rightarrow}$ Infarkt

\item Hirngef\"a{\ss}e ${\rightarrow}$ Schlaganfall

\item Nierengef\"a{\ss}e ${\rightarrow}$ Niereninsuffizienz
${\rightarrow}$ Dialyse ${\rightarrow}$ typischer KTW-Patient

\item Hauptschlagader ${\rightarrow}$ Aneurysma%
%


\end{itemize}
\begin{itemize}
\item Herzmuskel (Herz muss gegen gr\"o{\ss}eren Widerstand pumpen) 

\end{itemize}
\begin{itemize}
\item dilatative Kardiomyopathie

\end{itemize}
\paragraph[Hypertensive Krise / Hypertensiver \ Notfall]{Hypertensive
Krise / Hypertensiver  Notfall}
SEHR FRAGLICH !!!

\begin{center}
\tablehead{}\begin{supertabular}{|m{6.649cm}|m{6.651cm}|}
\hline
Hypertensive Krise

 &
Hypertensiver Notfall

\\\hline
RR {\textgreater} 230/130 mmHg

Patient f\"uhlt sich gut

 &
RR {\textgreater} 230/130 mmHg 

+ Symptome

\\\hline
Zufallsbefund?

Oder steht der Blutdruck doch in Zusammenhang mit der Berufung?

Patient hat \"uberhaupt keine Beschwerden? Warum hat er Sie dann
\"uberhaupt gerufen?

Wenn es wirklich nur ein Zufallsbefund ist, zum Beispiel wenn Sie
{\quotedblbase}aus Spa{\ss}{\textquotedblleft} jemand Blutdruck
gemessen haben ({\quotedblbase}Gesundheitstage{\textquotedblleft},
{\quotedblbase}Vorsorgeuntersuchung{\textquotedblleft}):

Ma{\ss}nahmen

\begin{itemize}
\item Lagerung: Oberk\"orper hoch

\item evt. O2 

\item Beengende Kleidung \"offnen

\item Monitoring

\item Hospitalisierung oder Notarzt zur Belassung

\end{itemize}
 &
Symptome:

\begin{itemize}
\item Kopfschmerz

\item Sehst\"orungen

\item Brustschmerz

\item evt. Nasenbluten

\item \"Ubelkeit, Erbrechen, Schwindel

\end{itemize}
Gefahr

\begin{itemize}
\item Herzversagen (Herz muss pl\"otzlich gegen gro{\ss}en Druck pumpen)

\end{itemize}
Ma{\ss}nahmen

\begin{itemize}
\item Lagerung: Oberk\"orper hoch

\item O2 

\item Beengende Kleidung \"offnen

\item Bewegungsverbot

\item Notarzt

\item Monitoring

\item Reanimationsbereitschaft

\end{itemize}
\\\hline
\end{supertabular}
\end{center}
Gef\"a{\ss}verschl\"usse - Extremit\"aten

Arteriell:  

Ven\"os:   \index{tVT}tVT (tiefe Venen Thrombose)

\paragraph{Arterieller Gef\"a{\ss}verschlu{\ss}}
Zuerst chronische Gef\"a{\ss}sch\"adigung: \index{PaVK}PaVK
(\index{periphere arterielle Verschlu{\ss}krankheit}periphere
arterielle  Verschlu{\ss}krankheit)

(Diabetes, Hypertonie, erh\"ohte Blutfette)

Symptome

\begin{flushleft}
\tablehead{}\begin{supertabular}{|m{1.8959999cm}|m{2.181cm}|m{9.016cm}|}
\hline
5 P  (englisch):

 &
 &
\\\hline
\centering 1.\par

 &
Pain

 &
schmerzend

\\\hline
\centering 2.\par

 &
Pale

 &
bla{\ss}

\\\hline
\centering 3.\par

 &
Pulsless

 &
pulslos

\\\hline
\centering 4.\par

 &
Paresthesia

 &
gef\"uhllos

\\\hline
\centering 5.\par

 &
Paralysis

 &
gel\"ahmt

\\\hline
\end{supertabular}
\end{flushleft}
Spezielle Ma{\ss}nahmen

\begin{itemize}
\item Lagerung: Extremit\"at h\"angen lassen, weich und warm lagern

\item Schonend transportieren

\end{itemize}
\paragraph{Tiefe Venenthrombose}
Blutgerinnsel (Thrombus ) verstopft Vene

Ursachen:

\begin{itemize}
\item Ven\"ose Insuffizienz (Krampfadern =  Venenaussackungen aufgrund
gest\"orten  Blutr\"uckflusses, Venenklappeninsuffizienz), R\"uckstau

\item Immobilisation

\item Bettl\"agrigkeit

\item Lange Resien ohne Bewegung: \index{Economy class Syndrom}Economy
class Syndrom 

\end{itemize}
Symptome:

\begin{itemize}
\item Schmerzen

\item Schwellung (einseitig, anderes Bein auch ansehen zum Vergleich!)

\item Dunkel-bl\"auliche (livide ) Verf\"arbung

\item \"Uberw\"armung

\end{itemize}
Gefahr

\begin{itemize}
\item Abschie{\ss}en des Thrombus ${\rightarrow}$ Lungenembolie (s.
I.II{\guillemotright}\pageref{bkm:RefHeading37612625})

\end{itemize}
Spezielle Ma{\ss}nahmen

\begin{itemize}
\item Lagerung:  liegend, Bein hoch!

\item Bewegungsverbot: Pat. nicht mehr gehen lassen  (Emboliegefahr!)

\end{itemize}
\subsection[Respiratorische / Pulmologische
\ Notf\"alle]{Respiratorische / Pulmologische  Notf\"alle}
St\"orungen der Atmung - \"Ubersicht

\begin{center}
 [Warning: Image ignored] % Unhandled or unsupported graphics:
%\includegraphics[width=3.754cm,height=4.22cm]{AASdev20090228-img75.png}

\end{center}
\begin{itemize}
\item mechanische Verlegung / Aspiration

\item Asthma bronchiale  / COPD

\item Lungenembolie

\item Lungen\"odem

\item Lungenentz\"undung (Pneumonie)

\end{itemize}
Mechanische Verlegung

Atemwege verlegt ${\rightarrow}$ Ventilation eingeschr\"ankt/aufgehoben

\index{Aspiration}Aspiration = Fremdk\"orper/-s\"afte (Magensaft,
Erbrochenes, Bolus) gelangt in die unteren Atemwege

besonders gef\"ahrlich f\"ur:

\begin{center}
 [Warning: Image ignored] % Unhandled or unsupported graphics:
%\includegraphics[width=6.871cm,height=4.395cm]{AASdev20090228-img76.png}

\end{center}
\begin{itemize}
\item Kinder

\item bewu{\ss}tseinseingeschr\"ankte Personen

\item Z.n. Bienenstich bei Allergiker (Kehldeckel\"odem)

\end{itemize}
Symptome:

\begin{itemize}
\item Dyspnoe, Husten

\item in/exspiratorisches Pfeifen (Stridor)

\item Apnoe, Asystolie

\item mechanische Verlegung

\end{itemize}
Therapie:

\begin{itemize}
\item Atemwege freimachen

\item Esmarch-Handgriff

\item Schlag zwischen Schulterbl\"atter

\item Heimlich-Griff

\item Erstma{\ss}nahmen bei vital bedrohten Patienten: 

\end{itemize}
\begin{itemize}
\item Lagerung: Oberk\"orper hoch

\item O2 10~l, bei Bedarf bis zu 15~l

\item Notarzt

\item Reanimationsbereitschaft

\end{itemize}
\begin{itemize}
\item Vitalwerte erheben \& Monitoring 

\item \end{itemize}
Asthma bronchiale

\begin{center}
\tablehead{}\begin{supertabular}{|m{6.078cm}|m{7.222cm}|}
\hline
Wiederholte Anf\"alle von Atemnot durch

\begin{center}
 [Warning: Image ignored] % Unhandled or unsupported graphics:
%\includegraphics[width=3.282cm,height=4.502cm]{AASdev20090228-img77.png}

\end{center}
\begin{itemize}
\item verengte Bronchien

\item vermehrte Schleimbildung

\end{itemize}
Symptome

\begin{itemize}
\item akute Atemnot mit trockenem Husten

\item exspiratorischer Stridor (Brummen, Pfeifen,  Giemen)

\item Einsatz der Atemhilfsmuskulatur

\item Verl\"angerte Ausatmungsphase

\end{itemize}
Komplikationen

\begin{itemize}
\item Zyanose

\item Status asthmaticus

\item Panik (${\rightarrow}$ Tachykardie)

\end{itemize}
 &
  [Warning: Image ignored] % Unhandled or unsupported graphics:
%\includegraphics[width=7.264cm,height=7.707cm]{AASdev20090228-img78.png}
 

\\\hline
\end{supertabular}
\end{center}
Spezielle Ma{\ss}nahmenx

\begin{itemize}
\item Erstma{\ss}nahmen bei vital bedrohten Patienten: Starke Atemnot
ohne Besserung

\begin{center}
 [Warning: Image ignored] % Unhandled or unsupported graphics:
%\includegraphics[width=6.378cm,height=4.784cm]{AASdev20090228-img79.jpg}

\end{center}
\end{itemize}
\begin{itemize}
\item Lagerung: Oberk\"orper hoch

\item O2 10~l, bei Bedarf bis zu 15~l

\item Notarzt

\item Reanimationsbereitschaft

\end{itemize}
\begin{itemize}
\item Vitalwerte erheben \& Monitoring 

\item Erstma{\ss}nahmen bei nicht vital bedrohten Patienten: Leichte
Atemnot, Besserung

\end{itemize}
\begin{itemize}
\item Lagerung: Oberk\"orper hoch

\item O2 ca. 6~l, bei Bedarf bis zu 15~l

\end{itemize}
\begin{itemize}
\item Vitalwerte erheben \& Monitoring 

\item Weiters immer:

\end{itemize}
\begin{itemize}
\item Beruhigen, voratmen (Lippenbremse)

\item Sprays, die Patient ev. bei sich hat d\"urfen nicht mehr genommen
werden.

\end{itemize}
Einnahme von verordneten Sprays analog zu Nitro?

O2-Gabe?

Achtung: In seltenen F\"allen kann es vorkommen dass bei hochdosierter
Sauerstoffgabe der Patient zu atmen aufh\"ort! 

COPD - Chronic obstructing pulmonary disease

Syn.: Chronisch-obstruktive Lungenerkrankung.

 [Warning: Image ignored] % Unhandled or unsupported graphics:
%\includegraphics[width=13.211cm,height=8.168cm]{AASdev20090228-img80.png}
\captionof{figure}[Eine COPD entsteht nicht pl\"otzlich: Ein
COPD{\textquotesingle}ler hat eine Karriere hinter sich.]{Eine COPD
entsteht nicht pl\"otzlich: Ein COPD{\textquotesingle}ler hat eine
Karriere hinter sich.}


Sammelbegriff f\"ur

\begin{itemize}
\item Lungenemphysem (\"Uberbl\"ahung der Lungenbl\"aschen)

\item chron. obstruktive Bronchitis

\end{itemize}
weiters:

\begin{itemize}
\item vermehrter Auswurf

\item Ausatmung (Exspiration ) erschwert

\item {\quotedblbase}Raucherhusten{\textquotedblleft}

\end{itemize}
Symptome

\begin{itemize}
\item Belastungsdyspnoe (bei schwererer COPD auch in Ruhe)

\item Husten mit Auswurf

\item Brummendes, giemendes Ausatemger\"ausch

\item Fa{\ss}thorax

\item Fa{\ss}thorax erkl\"aren

\item ben\"otigt evt. Heimsauerstoff

\end{itemize}
  [Warning: Image ignored] % Unhandled or unsupported graphics:
%\includegraphics[width=13.399cm,height=10.323cm]{AASdev20090228-img81.png}
 

Oben links: Bei der COPD sind die kleinen Luftwege verschleimt und
verengt.

Unten links: Die Lungenbl\"aschen (Alveolen) sind \"uberbl\"aht, weil
die Luft nur erschwert wieder entweichen kann. 

Therapie

\begin{itemize}
\item Wenn bei Eintreffen noch keine Besserung eingetreten ist Notarzt
erw\"agen!

\item Oberk\"orper hoch

\item voratmen, Lippenbremse

\item eigener %
%
Spray darf nicht mehr genommen werden

\end{itemize}
Spray bei COPD -- Regelung wie bei Nitro??

\begin{itemize}
\item Sauerstoff: VORSICHT:  O2 Vorsichtig dosieren, nur 2-3 l/ min. und
Atmung \"uberwachen. Wenn der Patient bereits Heimsauerstoff verordnet
hat 1-2 l/min h\"oher stellen und Atmung \"uberwachen! Heimsauerstoff
muss weiter gegeben werden!

\end{itemize}
Die CO2 - Konzentration im Blut ist bei COPD-Patienten chronisch hoch,
der K\"orper gew\"ohnt sich daran. Der K\"orper regelt dann den
Atemantrieb direkt \"uber den O2-Gehalt im Blut. 

Wird der O2-Gehalt pl\"otzlich k\"unstlich erh\"oht fehlt der
\index{COPD!Atemantrieb}Atemantrieb und der Patient kann einen
Atemstillstand erleiden. 

\label{bkm:RefHeading37612625}Lungenembolie

Einschwemmen eines Gerinnsels in die  Strombahn des Lungenkreislaufs

Herkunft des Embolus

\begin{center}
\tablehead{}\begin{supertabular}{|m{8.743cm}|m{4.563cm}|}
\hline
  [Warning: Image ignored] % Unhandled or unsupported graphics:
%\includegraphics[width=8.157cm,height=8.898cm]{AASdev20090228-img82.png}
 

 &
Tiefe (Bein-)Venenthrombose

Immobilisation (Gips, lange Reise)

Herz: Bei Rhythmusst\"orungen

Tumorpatienten

\\\hline
\end{supertabular}
\end{center}
Symptome:

\begin{itemize}
\item Atemnot

\item Akuter Thoraxschmerz

\end{itemize}
\begin{itemize}
\item eher atemabh\"angige Schmerzen (Atemstopp)

\end{itemize}
\begin{itemize}
\item Angst

\item Atemfrequenz und Herzschlag ${\uparrow}$

\end{itemize}
Eine Lungenembolie kann man leicht mit einem Herzinfarkt oder Angina
pectoris verwechseln. Dies hat aber bei der Erstversorgung kaum
Konsequenz, die Ma{\ss}nahmen sind die gleichen.

Typische Befunde

\begin{itemize}
\item Bein wie bei tiefer Beinvenenthrombose:  einseitg geschwollen,
\"uberw\"armt, rot/rot-bl\"aulich verf\"arbt, schmerzhaft

\begin{center}
\begin{minipage}{5.976cm}
${\leftarrow}$ wenn der Embolus aus dem Bein kommt

\end{minipage}
\end{center}
\item evt. gestaute Halsvenen

\begin{center}
\begin{minipage}{6.726cm}
${\leftarrow}$ Blut staut zur\"uck

\end{minipage}
\end{center}
\end{itemize}
Typische Anamnese

\begin{itemize}
\item Lange Reise, Flug, Auto, Bahn {\dots}

\item Bettl\"agrigkeit

\item Schwangerschaft

\item Raucher(in)

\item Verh\"utungspille

\end{itemize}
Ma{\ss}nahmen

\begin{itemize}
\item Erstma{\ss}nahmen bei vital bedrohten Patienten: Atemnot

\end{itemize}
\begin{itemize}
\item Lagerung: Oberk\"orper hoch

\item O2 10~l, bei Bedarf bis zu 15~l

\item Notarzt

\item Reanimationsbereitschaft

\end{itemize}
\begin{itemize}
\item Vitalwerte erheben \& Monitoring 

\item Striktes Bewegungsverbot

\item Beengende Kleidung \"offnen

\end{itemize}
\label{bkm:RefHeading36561325}Lungen\"odem

Fl\"ussigkeitsstau in den Alveolen und Interstitium

Ursachen:

\begin{itemize}
\item Durch R\"uckstau: akute Linksherzinsuffizienz, Linksherzversagen

\item z.B. bei MCI, zus\"aztliche Belastung bei vorgesch\"adigtem Herz
(Infekt)

\item Durch Entz\"undung: Inhalation toxischer Substanzen

\end{itemize}
Das gef\"ahrliche am Lungen\"odem sind dessen Ursachen!

Symptome:

\begin{itemize}
\item Atemnot

\item Zyanose

\item \index{Brodelnde Atemger\"ausche}Brodelnde Atemger\"ausche, oft
ohne Hilfsmittel h\"orbar, im Extremfall schon von der Wohnungst\"ur
aus.

\begin{center}
\begin{minipage}{4.2cm}
Brodelnde Atemger\"ausche: Als ob jemand mit einem Strohhalm in ein Glas
Wasser bl\"ast.

\end{minipage}
\end{center}
\end{itemize}
Spezielle Ma{\ss}nahmen

\begin{itemize}
\item Erstma{\ss}nahmen bei vital bedrohten Patienten

\end{itemize}
\begin{itemize}
\item Lagerung: Oberk\"orper hoch, wenn m\"oglich beine herabh\"angen
lassen

\item O2 10~l, bei Bedarf bis zu 15~l

\item Notarzt

\item Reanimationsbereitschaft

\end{itemize}
\begin{itemize}
\item Vitalwerte erheben \& Monitoring

\item Striktes Bewegungsverbot

\item Beengende Kleidung \"offnen

\item evt. Gabe von Nitroglycerin-Spray durch einen Notfallsanit\"ater

\end{itemize}
Patienten mit einem Lungen\"odem (v.a. bei kardialer Ursache) sind
lebensbedrohlich krank. Das Herz vebraucht oft schon die allerletzten
Reserven, schon die Aufregung durch den Transport im Stiegenhaus zum
Auto kann das Herz endg\"ultig in den Herz-Kreislaufstillstand treiben.


Es kommt h\"aufig vor das diese Patienten im Stiegenhaus oder im Auto
reanimationspflichtig werden. 

Kein Transport ohne Notarzt. Auch wenn das Spital
{\quotedblbase}um{\textquotesingle}s Eck{\textquotedblleft} ist.

Lungenentz\"undung (Pneumonie )

\begin{itemize}
\item durch Keime: Bakterien, Pilze (Pneumokokken,Chlamyden, Candida bei
 HIV/Immunsupression)

\item durch \index{Pneumonie!Aspiration}Aspiration

\end{itemize}
h\"oheres Lebensalter beg\"unstigt  Pneumonien

Symptome

Die Symptome sind oft nicht charakteristisch, eine Diagnose kann meist
erst im Krankenhaus gestellt werden. 

\begin{itemize}
\item Dyspnoe, Husten

\item Fieber ( mu{\ss} aber nicht sein)

\item Reduzierter Allgemeinzustand (AZ)

\end{itemize}
Ma{\ss}nahmen

\begin{itemize}
\item Symptomatisch

\end{itemize}
\begin{itemize}
\item W\"armeerhalt 

\end{itemize}
\begin{itemize}
\item Hospitalisierung, Abteilung: Innere Medizin

\end{itemize}
\subsection{Neurologische Notf\"alle}
St\"orungen des Nervensystems durch (\"Ubersicht):

\begin{itemize}
\item Schlaganfall

\item Vergiftungen

\item SHT (Sch\"adel-Hirn-Trauma)

\item metabolische Ursachen

\item Infektionen

\item Krampfanf\"alle

\item Erkrankeungen des St\"utzapparates (Wirbels\"aule, Bandscheiben)

\end{itemize}
Neurologische St\"orungen gehen oft mit einer St\"orung des Bewusstseins
oder Schmerzen einher.

Es fehlt: Lumbago, BWS/HWS

\index{Schlaganfall}Schlaganfall (\index{Insult}Insult)

\begin{enumerate}
\item Verschlu{\ss} eines Hirngef\"a{\ss}es mit anschlie{\ss}endem
Absterben der versorgten Hirnregion ({\quotedblbase}trockener
Insult{\textquotedblleft}), oder

\item Sponane Hirnblutung ({\quotedblbase}blutiger
Insult{\textquotedblleft})

\end{enumerate}
Syn.: \index{Apoplektischer Insult}\index{Apoplektischer
Insult}Apoplektischer Insult, \index{Apoplexie}Apoplexie

Arten:

\begin{center}
 [Warning: Image ignored] % Unhandled or unsupported graphics:
%\includegraphics[width=2.377cm,height=4.195cm]{AASdev20090228-img83.png}

\end{center}
\begin{itemize}
\item Gef\"a{\ss}verschlu{\ss} (isch\"amischer Insult)

\item Blutung (h\"amorrhagischer Insult)

\begin{center}
 [Warning: Image ignored] % Unhandled or unsupported graphics:
%\includegraphics[width=3.003cm,height=4.303cm]{AASdev20090228-img84.png}

\end{center}
\item Sonderform \index{TIA}TIA: \index{Transitorisch -- Isch\"amische
Attacke}Transitorisch -- Isch\"amische Attacke

\end{itemize}
\begin{itemize}
\item Symptome wie bei Insult, vergehen aber innerhalb von 24 Stunden.

\end{itemize}
  [Warning: Image ignored] % Unhandled or unsupported graphics:
%\includegraphics[width=9.526cm,height=5.428cm]{AASdev20090228-img85.png}
 

Symptome

Abh\"angig welche Hirnregion betroffen ist. Die Symptome k\"onnen sehr
unterschiedlich sein. 

Klassische Symptome sind:

\begin{itemize}
\item Halbseitenschw\"ache oder -l\"ahmung, h\"angener Mundwinkel

\item Sprachst\"orungen:

\end{itemize}
\begin{itemize}
\item verwaschene Sprache 

\item {\quotedblbase}wirre Sprache{\textquotedblleft}: Worte fehlen,
Silben werden vertauscht, ...

\end{itemize}
\begin{itemize}
\item Sehst\"orungen, pl\"otzliche Erblindung, Doppelbilder,
Gesichtsfeldausf\"alle

\item Schluckst\"orungen

\item pl\"otzlicher, rasender Kopfschmerz, Schwindel

\item in schweren F\"allen bis Bewu{\ss}tlosigkeit

\end{itemize}
Differentialdiagnose

\begin{itemize}
\item BZ ${\downarrow}$, ${\uparrow}$ (${\rightarrow}$ Bei
neurologischen Symptomen immer messen!)

\item Tumor

\end{itemize}
Ma{\ss}nahmen

\begin{itemize}
\item Erstma{\ss}nahmen bei vital bedrohten Patienten

\end{itemize}
\begin{itemize}
\item Lagerung: leicht erh\"ohter Oberk\"orper,

\item O2 10~l, bei Bedarf bis zu 15~l

\item Notarzt

\item Reanimationsbereitschaft

\end{itemize}
\begin{itemize}
\item Vitalwerte erheben \& Monitoring

\item Neurocheck inkl. BZ-Messung (DD!)

\item Stroke Unit!

\end{itemize}
Bakterielle \index{Meningitis}Meningitis

\begin{center}
\tablehead{}\begin{supertabular}{|m{6.649cm}|m{6.651cm}|}
\hline
= eitrige Entz\"undung der Hirnh\"aute

ansteckende Form:

Meningokokken-Meningitis ${\rightarrow}$ hoch ansteckend

Symptome:

\begin{itemize}
\item Kopfschmerzen, Lichtscheu

\item Erbrechen, \"Ubelkeit

\item Meningismus :

\end{itemize}
\begin{itemize}
\item Fieber

\item Nackensteifigkeit

\end{itemize}
\begin{itemize}
\item Ausschlag (Exanthem )

\end{itemize}
 &
\begin{center}
 [Warning: Image ignored] % Unhandled or unsupported graphics:
%\includegraphics[width=6.696cm,height=3.813cm]{AASdev20090228-img86.png}

\end{center}
\captionof{figure}[Endstadium einer bakteriellen Meningitis. Die Flecken
auf der Haut sind Blutungen in Folge einer nicht mehr funktionierenden
Blutgerinnung. Der Tod ist bei diesem Patienten innerhalb weniger
Stunden zu erwarten.]{Endstadium einer bakteriellen Meningitis. Die
Flecken auf der Haut sind Blutungen in Folge einer nicht mehr
funktionierenden Blutgerinnung. Der Tod ist bei diesem Patienten
innerhalb weniger Stunden zu erwarten.}
\\\hline
\end{supertabular}
\end{center}
FEHLT: Querverwis Meningimus

FEHLT: Ma{\ss}nahmen

Epilepsie%
%


Vom Gehirn ausgehende motorische Krampfanf\"alle aufgrund von akuten
oder  chronischen Erkrankungen des Gehirns.

Aufgrund von sich abnorm synchronisiert  entladenden Nervenzell-Gruppen
im  Gehirn

nach Art:

\begin{itemize}
\item tonisch

\item klonisch

\end{itemize}
nach Lokalisation:

\begin{itemize}
\item fokal - petit mal -- Anfall

\item generalisiert grand mal - Anfall 

\end{itemize}
Ursachen:

\begin{itemize}
\item Vergiftungen (Medikamente/Drogen/...)

\item Hirntumor

\item Narbengewebe nach Insult, Hirnverletzung

\item Hirnblutung, Isch\"amie

\item BZ 

\item angeboren, ohne organisch fa{\ss}bare  Ursache

\item Fieber bei Kleinkindern

\item Infektionen

\end{itemize}
\paragraph{Petit mal - Anfall}
fokale Anf\"alle:

\begin{itemize}
\item motorisch: Jackson-Anfall

\item {\quotedblbase}Jacksonian march{\textquotedblleft}

\item D\"ammerzustand (Absence )

\item Dauer: meist nur Sekunden

\item k\"onnen sich ev. zu grand mal-Anf\"allen  ausweiten

\end{itemize}
\paragraph{Grand mal - Anfall}
\begin{itemize}
\item Tonisch -- klonischer Anfall

\item zuerst kompletter Tonusverlust

\item Niederst\"urzen (Zungenbi{\ss})

\item Tonus ${\uparrow}$, unkontrollierte Zuckungen

\item Schaum vor dem Mund

\item Stuhl-/Urinabgang

\item nach Krampfphase Erholungsschlaf

\item ev. sogar tiefe Bewu{\ss}tlosigkeit

\end{itemize}
f\"ur den Anfall selbst besteht Amnesie! ${\rightarrow}$ Patient
aufkl\"aren was passiert ist, wo er sich befindet, etc.

Ma{\ss}nahmen

\begin{itemize}
\item Patient vor Verletzungen sch\"utzen

\item M\"obel etc. entfernen oder polstern

\item NICHT festhalten -- aber sichern, Eigenschutz beachten!

\begin{center}
\begin{minipage}{4.2cm}
Sichern: Patient darf nicht von der Trage oder aus dem Bett fallen, der
Kopf darf nicht am Boden oder Tischbein anschlagen, etc.

\end{minipage}
\end{center}
\item %
%
Notarzt wenn: 

\end{itemize}
\begin{itemize}
\item Krampf bei Eintreffen, 

\item {\textgreater} 2 Anf\"alle in 24 Stunden, 

\item erstmaliger Epi

\end{itemize}
\begin{itemize}
\item nach Anfall BAK-Schema, Atemwege freihalten

\item stabile Seitenlage  (Erholungsschlaf wie  Bewu{\ss}tlosigkeit)

\item Besondere Diagnostik: 

\end{itemize}
\begin{itemize}
\item Neurocheck inkl. BZ! 

\item Traumacheck 

\end{itemize}
\begin{itemize}
\item Reizabgeschirmter Transport

\end{itemize}
Epi {\textgreater} 3 min . bwz. erster  Epi ${\rightarrow}$ Notarzt

\paragraph{Status epilepticus }
\begin{itemize}
\item Anf\"alle {\textgreater} 15 Minuten oder

\item wiederholte Anf\"alle in kurzem Intervall

\end{itemize}
LEBENSGEFAHR!

Ma{\ss}nahmen

\begin{itemize}
\item NAW

\item reizabgeschirmter Transport

\item neurologische Notf\"alle -- Epilepsie

\end{itemize}
Indikationen f\"ur Spitalseinlieferung

\begin{itemize}
\item erstmaliges Auftreten eines Krampfanfalls

\item Status epilepticus

\item ABKL\"ARUNG  der Ursache in diesen  F\"allen wichtig!

\end{itemize}
Fieberkrampf

Krampfanfall bei Kleinkindern bei schnellem, hohen Fieberanstieg
({\textgreater} 39{\textdegree}C)

Sieht aus wie generalisierter grand mal-Anfall

DD

grand mal OHNE Fieber, Fieberkrampf  NUR bei Kindern {\textless} 4 Jahre

Tetanie 

(nicht verwechseln mit Tetanus)

neuromuskul\"are \"Ubererregbarkeit durch  St\"orungen im
Calcium-Spiegel

\paragraph{Hyperventilationstetanie }
\begin{itemize}
\item 90\% psychogen! (typisch: pubertierende  M\"adchen)

\item Lungenembolie

\item diabetische Stoffwechselentgleisung (Ketoazidose, kompensatorisch)

\end{itemize}
Symptome

\begin{itemize}
\item Tachypnoe mit Alkalose ={\textgreater} Verschiebung  beim freien
Ca++

\item Kribbeln, Unruhe

\item Kr\"ampfe (Gesichtsmuskulatur),  Pf\"otchenstellung

\item Panik

\end{itemize}
Therapie

\begin{itemize}
\item beruhigen!!!

\item Kommandoatmung

\item CO2 {}-R\"uckatmung  (Plastiksackerl,  Handschuh)

\item falls nicht erfolgreich, Notarzt f\"ur Sedierung  nachfordern

\item KEIN O2  verabreichen!!! Kontraindikation!!!

\end{itemize}
\subsection{Metabolische St\"orungen (Stoffwechselst\"orungen)}
Die relevanteste Stoffwechselst\"orung im Rettungsdienst ist der
\index{Diabetes mellitus}Diabetes mellitus. Dieser stellt eine
St\"orung des Kohlehydrat-(zucker-)stoffwechsels dar. Zucker ist ein
Kohlehydrat. 

Merke: Diabetes mellitus ist eine chronische Stoffwechselerkrankungen
mit  gest\"orter Blutzuckerregulation, daraus  resultierend:

\begin{itemize}
\item Hypoglyk\"amie ... Unterzucker

\item Hyperglyk\"amie ... {\quotedblbase}\"Uberzucker{\textquotedblleft}

\end{itemize}
Insulin 

Insulin ist ein Hormon und wird in den Beta-Zellen der
Bauchspeicheldr\"use (Pankreas) produziert. Es sorgt daf\"ur dass
Zucker (Glukose) aus dem Blut in die Zellen aufgenommen wird. 

Zwei St\"orungen sind denkbar:

\begin{enumerate}
\item Es ist zu wenig Insulin vorhanden.

\item Der K\"orper ist schon zu sehr an das Insulin gewohnt und reagiert
nicht mehr so wie er sollte, er braucht gr\"o{\ss}ere Mengen Insulin um
den gleichen Effekt wie ein Gesunder zu erreichen (Gew\"ohnungseffekt),
irgendwann kannn aber der K\"orper nicht mehr Insulin produzieren.

\end{enumerate}
Folglich unterscheidet man 2 Typen:

\begin{itemize}
\item Diabetes mellitus Typ I  (IDDM, insuline dependent) ... keine
Insulinproduktion des K\"orpers, muss  immer Insulin spritzen 

\item Diabetes mellitus Typ II (NIDDM, non insuline dependent) ...
Zuviel Zucker (schlechte Ern\"ahrung), folglich zu viel Insulin,
K\"orper gew\"ohnt sich, s.o. ${\rightarrow}$ braucht erst Sport,
Medikamente, im fort\-ge\-schrittenen Stadium oft Insulin

\end{itemize}
Sp\"atfolgen:

Ein hoher Blutzuckerspiegel \"uber Jahre sch\"adigt den gesammten
K\"orper. Es kommt zu:

\begin{itemize}
\item Gef\"a{\ss}sch\"aden: Verkalkung

\end{itemize}
\begin{itemize}
\item Gef\"a{\ss}verschl\"usse (kleine und gro{\ss}e Gef\"a{\ss}e)

\end{itemize}
\begin{itemize}
\item Beine: {\quotedblbase}Diabetischer Fu{\ss}{\textquotedblleft}, in
Folge der Mangeldurchblutung in der Peripherie heilen auch kleine
Wunden nicht mehr, in weiterer Folge muss immer weiter amputiert
werden.

\item Niere: Die sehr kleinen Gef\"a{\ss}e der Niere werde gesch\"adigt,
im weiteren Verlauf stellt die Niere ihre Arbeit ein ${\rightarrow}$
{\quotedblbase}Chronische Niereninsuffizienz{\textquotedblleft}. Diese
Erkrankung macht einen Gro{\ss}teil der Krankentransporte aus.

\item Herz: Die Herzkranzgef\"a{\ss}e (Koronar\-arterien) sind ebenso
betroffen. Diabetiker sind Risikopatienten f\"ur
Herz-Kreislauf-Erkrankungen wie zB Herz\-infarkt.

\item Hirn: Wie beim Herzen kommt es zu Gef\"a{\ss}sch\"aden, daher ist
das Risiko f\"ur Schlag\"anf\"alle erh\"oht.

\end{itemize}
\begin{itemize}
\item Nervensch\"aden (periphere Polyneuropathie ): Gef\"uhlslosigkeit
vor allem in den Extremit\"aten, aber auch am gesammten K\"orper: Ein
Herzinfarkt kann zum Beispiel schmerzlos ablaufen!

\item Augensch\"aden -- Erblindung

\end{itemize}
Krankheiten die sonst schmerzhaft sind k\"onnen beim Diabetiker aufgrund
der NErvensch\"aden ohne Schmerzen ablaufen. 

Achtung: Stummer Herzinfarkt beim Diabetiker!

 [Warning: Image ignored] % Unhandled or unsupported graphics:
%\includegraphics[width=11.387cm,height=16.225cm]{AASdev20090228-img87.png}
\captionof{figure}[Die h\"aufigsten Folgeerkrankungen des Diebetikers:
Nervensch\"adigungen, Herzinfarkt, Nierensch\"aden, Schlaganfall]{Die
h\"aufigsten Folgeerkrankungen des Diebetikers: Nervensch\"adigungen,
Herzinfarkt, Nierensch\"aden, Schlaganfall}


\subparagraph[Diabetes \ mellitus I]{Diabetes  mellitus I}
St\"orung des Immunsystems ${\rightarrow}$ Zerst\"orung der
Insulin-produzierenden ${\beta}$-Zellen des Pankreas ${\rightarrow}$
kein Insulin -- absoluter Mangel

Manifestationsalter: 0-25

5\% der Diabetiker, normalgewichtig ({\quotedblbase}jung,
sportlich{\textquotedblleft}), 2-6\% erblich

Therapie

Insulin (mu{\ss} gespritzt werden!)

Vorsicht: Hypoglyk\"amie-gef\"ahrdet (wenn zuviel Insulin gespritzt oder
zu wenig gegessen wurde) ${\rightarrow}$  Spritz/E{\ss}abstand
er\-fra\-gen

\subparagraph[Diabetes \ mellitus II]{Diabetes  mellitus II}
Relativer Insulinmangel, K\"orper ist \"ubers\"attig an Insulin und
reagiert kaum mehr daruf (Insulinresistenz), sp\"ater ${\beta}$-Zellen
ersch\"opft

\begin{itemize}
\item sp\"ate Manifestation: {\textgreater}35 Jahre

\item \"Ubergewicht

\item 95\% aller Diabetiker

\item 25-30\% erblich

\end{itemize}
Hypoglyk\"amie

Ursachen

\begin{center}
 [Warning: Image ignored] % Unhandled or unsupported graphics:
%\includegraphics[width=4.234cm,height=7.526cm]{AASdev20090228-img88.png}

\end{center}
\begin{itemize}
\item zuviel Insulin (spritzen ohne zu Essen)

\item fasten (spritzen ohne zu Essen)

\item Sport

\item Alkohol

\end{itemize}
BZ normalerweise 80-110 mg/dl

Symptome

\begin{itemize}
\item Unruhe, Zittern, Schwei{\ss}ausbruch

\item Hei{\ss}hunger, \"Ubelkeit

\item Tachykardie, Tachypnoe

\item Haut feucht/bla{\ss}

\item Verwirrtheit, Krampfanf\"alle, Bewu{\ss}tlosigkeit

\end{itemize}
Hypoglyk\"amie kann alle neurologischen  St\"orungen
aus\-l\"osen/imitieren!

Bei jeder  neurologischen Symptomatik  (v.a. Bewu{\ss}tseinsst\"orungen)
an  Hypoglyk\"amie denken

${\rightarrow}$ BZ-Messung bei jeder neurologischen St\"orung. 

Therapie

\begin{itemize}
\item Basistherapie bei Bewu{\ss}tlosigkeit

\item BZ messen! Wenn Patient ansprechbar,  Traubenzucker, danach
nachessen lassen.

\item evt. Notarzt

\end{itemize}
Hyperglyk\"amie

Diabetisches Koma 

(BZ {\textgreater} 400 mg/dl)

ketoazidotisches Koma

\begin{itemize}
\item BZ ${\uparrow}$

\item saure Ketonk\"orper ={\textgreater} Azidose (\"Ubers\"auerung)

\item kompensatorisch Ku{\ss}maul{\textquotesingle}sche Atmung (CO2 -
Abatmung)

\begin{center}
\begin{minipage}{4.2cm}
Erinnere Dich: Hyperventilationstetanie ${\rightarrow}$ Alkalose (pH
wird basisch). 

Bei der diab. Ketoazidose (sauer) versucht der K\"orper diesen
Mechanismus auszunutzen: 

{\quotedblbase}Er will sich vom sauren pH zum basischen pH
atmen.{\textquotedblleft}

${\rightarrow}$ Er hyperventiliert.

\end{minipage}
\end{center}
\item eher bei DM I

\end{itemize}
hyperosmolares Koma

\begin{itemize}
\item BZ ${\uparrow}$${\uparrow}$

\item zuviel Glucose ={\textgreater} rel. Fl\"ussigkeitsmangel

\item Niere gest\"ort

\item DM II

\end{itemize}
Ursachen

\begin{itemize}
\item Di\"atfehler (Kohlenstoff-reiche Ern\"ahrung)

\item Insulinzufuhr unzureichend

\item Stre{\ss}

\item Infektionen (Insulinbedarf ${\uparrow}$)

\end{itemize}
Symptome

\begin{itemize}
\item Durstgef\"uhl

\item vermehrter Harndrang, Harnmenge ${\uparrow}$

\item Bewu{\ss}tseinstr\"ubung bis Bewu{\ss}tlosigkeit

\item Ku{\ss}maul{\textquotesingle}sche Atmung (sehr tief und schnell)

\item Aceton-Geruch (Achtung, nicht mit Schnapsfahne verwechseln!)

\item DM I: ausgepr\"agte Symptome

\item DM II: schleichender Verlauf

\end{itemize}
Therapie

\begin{itemize}
\item Basistherapie Bewu{\ss}tlosigkeit

\item Notarzt

\item O2

\end{itemize}
\subsection{Abdominelle Notf\"alle}
Akutes Abdomen

Akutes Abdomen: Akute schmerzhafte Baucherkrankung (viele Ursachen 
m\"oglich!)

ferner: 

\begin{itemize}
\item harte Bauchdecke / Abwehrspannung  (D\'efense ) ${\rightarrow}$
Alarmzeichen

\item Schockgeschehen (septisch,  Volumenmangel, ...) ${\rightarrow}$
Notfall

\end{itemize}
Die Definition des Akuten Abdomens ist oft nicht ganz eindeutig:
Einerseits kann man da\-runter jede pl\"otzlich auftretende
Baucherkrankung verstehen\footnote{Longmore et.al.: Oxford Handbook of
Clinical Medicine, 7th edt., Oxford University Press, 2007},
\footnote{\cite{BLIM03} S.611}. \label{ref:JRcite1OHCM7}[2] Allerdings
ver\-stehen viele Leute darunter eher eine bedrohliche
Bauch\-erkrankung, bei der man akut (schnell) handeln muss
(Ab\-wehrspannung, harter Bauch).

Symptome

\begin{itemize}
\item Schonhaltung  (angezogene Beine, Patient  kr\"ummt sich)

\item diffuser Schmerz (ganzer Bauch), 

\item ev. lokalisierbarer Druckschmerz

\item harte  Bauchdecke (bretthart !!)

\item reflektorische Abwehrspannung  (D\'efense musculaire)

\item \"Ubelkeit, Erbrechen

\item Stuhl-, Windverhalten (ev. diagnostisch f.  Ileus) 

\item Fieber

\item Schock

\end{itemize}
Ursachen

\begin{itemize}
\item Entz\"undung, Drehung, Verlegung, Zerrei{\ss}ung

\item gyn\"akologisch

\item Verletzungen

\end{itemize}
Allgemeine Ma{\ss}nahmen (Ursachen-unabh\"angig)

\begin{itemize}
\item NA  bei Schocksymptomatik

\item O2

\item Lagerung: bauchdeckenentspannend  (Knierolle, Beine  angezogen),
nach Wunsch des Patienten

\item n\"uchtern lassen (in Hinblick auf m\"oglicherweise bevorstehende
Operation, nichts essen oder trinken lassen)

\item Hospitalisierung: ad Chirurgie, bei Kindern ${\rightarrow}$
Kinder, bei Frauen evt. Gyn\"akologie oder Chirurgie mit Gyn im Haus
(SMZ-O, Rudi, Wili, etc.)

\end{itemize}
M\"ogliche Ursachen f\"ur akutes Abdomen

\begin{itemize}
\item Bauchfellentz\"undung (Peritonitis )

\item Darmverschlu{\ss} (Ileus )

\item {\quotedblbase}Blinddarmentz\"undung{\textquotedblleft}
(Appendizitis )

\item Entz\"undung der Bauchspeicheldr\"use  (Pankreatitis )

\item Entz\"undung des Magens / Magengeschw\"ur  (Gastritis / Ulcus
ventriculi )

\item Gallenkolik

\item Bauchaortenaneurysma

\item Mesenterialinfarkt

\item abdominelle Erkrankungen

\item \"Osophagusvarizenblutung

\item Entz\"undung der Leber (Hepatitis )

\end{itemize}
\begin{itemize}
\item Hepatitis A

\item Hepatitis B / C

\end{itemize}
\begin{itemize}
\item Gelbsucht (Ikterus )

\item Gastroenteritis

\item Dehydratation, Exsikkose

\end{itemize}
Schmerzcharakter

Peritonitis

Peritonitis: Entz\"undung und Keimbesiedelung des  Bauchfells

Ursache

\begin{itemize}
\item Perforation eines Hohlorganes  (Magendurchbruch etc.)

\item Infektion \"uber Blut- / Lymphwege

\end{itemize}
Symptome

\begin{itemize}
\item Schmerzen (diffus)

\item {\quotedblbase}Totenstille{\textquotedblleft} (Peristaltik,
reflektorischer Ileus )

\item Abwehrspannung (D\'efense musculaire )

\item ev. Sepsiszeichen!

\end{itemize}
\index{Ileus}Ileus

Unterbrechung der normalen  Darmpassage

\index{Ileus}Ileus  (vollst\"andig)

\index{Subileus}Subileus  (unvollst\"andig)

Arten

\begin{itemize}
\item mechanisch  (Verschlingung, Abklemmung  etc.)

\item paralytisch  (toxisch, Medikamente)

\end{itemize}
Symptome

\begin{itemize}
\item Schmerzen

\item Erbrechen (ev. Koterbrechen, Miserere )

\item Stuhlverhalten

\item keine (paralytisch ) oder spritzende  (mechanisch )
Darmger\"ausche

\end{itemize}
Appendizitis

Entz\"undung des Wurmfortsatzes (Appendix vermiformis) des Blinddarms

Symptome

\begin{itemize}
\item akutes Abdomen

\item Schmerzen

\item Abwehrspannung (D\'efense musculaire )

\item \"Ubelkeit, Erbrechen

\item Schmerzen im rechten Unterbauch

\item Losla{\ss}schmerz li. Bauch

\item schmerzhafter McBurney-Punkt

\end{itemize}
Pankreatitis

Entz\"undung der Bauchspeicheldr\"use  (akut/chronisch)

Ursachen:

\begin{itemize}
\item Verlegung der Gallenwege

\item chron. Alkoholabusus

\item fettreiche Mahlzeit

\item Z.n. stumpfem Bauchtrauma

\end{itemize}
Symptome:

Schmerzen li. Oberbauch, g\"urtelf\"ormig  ausstrahlend

\"Ubelkeit, Erbrechen, Seifenst\"uhle, Bl\"ahungen

Gastritis / Ulcus ventriculi

Entz\"undung / Blutung (ev. mit  Durchbruch des Magens)

Ursachen

\begin{itemize}
\item Autoimmunerkrankungen

\item bakteriell (Helicobacter pylori )

\item chemisch

\item Gallens\"aurenreflux

\item Medikamente (NSAR)

\end{itemize}
Symptome

\begin{itemize}
\item Schmerzen (Oberbauch), \"Ubelkeit, Erbrechen

\item bei Blutung: Kaffesatz-Erbrechen,  Teerstuhl

\end{itemize}
Gallenkolik

Abflu{\ss}stau der Gallenwege

R\"uckstau

Gallenblase kontrahiert gegen Widerstand  (meist Stein)

Symptome

\begin{itemize}
\item krampfartige Schmerzen re. Oberbauch

\item Erbrechen nach fettreichen Mahlzeiten

\item ev. Gelbsucht (Ikterus )

\end{itemize}
Bauchaortenaneurysma

Zerrei{\ss}ung der Bauchaorta / Platzen  einer Aussackung
(Schwachstelle)

Symptome:

heftigste Schmerzen
({\quotedblbase}Vernichtungsschmerz{\textquotedblleft})

tastbare pulsierende Schwellung

herznah ${\rightarrow}$ hoher Blutverlust ${\rightarrow}$ schnell im 
Schock! 

Mesenterialinfarkt

Infarkt im Bereich des Darmgekr\"oses

Zelluntergang

Blutung

Austritt von Speisebrei in die freie Bauchh\"ohle  ={\textgreater}

zus\"atzlich Peritonitis

Symptome

pl\"otzliche starke Bauchschmerzen, lassen ev.  wieder nach

\"Osophagusvarizenblutung

{\quotedblbase}Krampfadern in der Speiser\"ohre{\textquotedblleft}

bei Lebererkrankungen 
({\quotedblbase}Umgehungskreislauf{\textquotedblleft})

schlechte Blutgerinnung

Lebensgefahr (hoher Blutverlust  m\"oglich)

Symptome

\begin{itemize}
\item Bluterbrechen im Schwall (H\"amatemesis)

\item Grunderkrankung:

\item Leberzirrhose

\item Ikterus

\end{itemize}
Therapie

\begin{itemize}
\item Notarzt: Sengstaken-Blakemore-Sonde

\item Spital: Ligatur (chirurgisch)

\end{itemize}
Hepatitis

Entz\"undung der Leber

Ursachen:

Gallen-Abflu{\ss}stau

Vergiftungen Lebensmittel(Pilze)/  Medikamente(Paracetamol)/  chron. 
Alkoholismus

Virusinfektionen (Hep A,  Hep B, Hep C )

Symptome:

grippale Symptome, Appetitlosigkeit

Schmerzen im re. Oberbauch

Gelbf\"arbung Augen/Haut (Ikterus )

\paragraph{Virushepatitis}
Hepatitis A

Infektion: f\"ako-oral

ungewaschene H\"ande

Meeresfr\"uchte

{\quotedblbase}Reise-Gelbsucht{\textquotedblleft}

akut, wird nie chronisch

Ikterus

\"Ubelkeit, Erbrechen

selten Komplikationen

Ikterus

Gelbsucht

zuerst an der Bindehaut (Auge) sichtbar

Ablagerung von Blutfarbstoff-Abbauprodukt  (Bilirubin )

pr\"ahepatisch : vermehrter Anfall von Bilirubin

intrahepatisch : gest\"orte Funktion der  Leberzellen

posthepatisch : Galleabflu{\ss} blockiert

Ausheilung

 oder

Chronifizierung...Leberzirrhose

Gastroenteritis

Brechdurchfall

Ursachen:

Vielf\"altig

Symptome:

\begin{itemize}
\item Erbrechen (Emesis )

\item Durchfall (Diarrhoe )

\item Wasser-/Elektrolytverlust

\end{itemize}
besonders gef\"ahrdet:

\begin{itemize}
\item Kinder

\item alte Menschen

\end{itemize}
Exsikkose-Gefahr!

Dehydratation, syn. Exsikkose

Ursachen:

\begin{itemize}
\item Fl\"ussigkeitszufuhr ${\downarrow}$

\item Fl\"ussigkeitsverlust ${\uparrow}$

\item forcierte Diurese  (Entw\"asserungsmedikamente)

\item starkes Schwitzen

\end{itemize}
Symptome:

\begin{itemize}
\item RR ${\downarrow}$

\item stehende Hautfalten

\item trockene, belegte Zunge

\item Elektrolytentgleisung

\item Verwirrtheit, Schwindel, Apathie

\end{itemize}
\subsection{Urologische Notf\"alle}
Nierenkolik

Ursache

\begin{itemize}
\item Nierenstein(e)

\end{itemize}
Symptome

\begin{itemize}
\item R\"ucken-/Flankenschmerzen

\item ev. Unterbauchschmerzen bei  Stein-/Sandabgang

\item {\quotedblbase}der Stein wandert, der Schmerz wandert, der Patient
wandert{\textquotedblleft}

\item Blut im Harn (H\"amaturie)

\end{itemize}
Therapie

\begin{itemize}
\item Lagerung nach Wunsch des Patienten

\item schonender Transport

\item unertr\"agliche Schmerzen ={\textgreater} Notarzt

\end{itemize}
Akutes Harnverhalten

Abflu{\ss}behinderung der Harnblase

Ursachen

\begin{itemize}
\item L\"ahmung (Querschnittl\"ahmung)

\item Schlie{\ss}muskel-Krampf

\item Verlegung der Harnr\"ohre

\item Medikamente

\end{itemize}
Gefahr

\begin{itemize}
\item Ruptur der Harnblase ${\rightarrow}$ Harn im kleinen Becken

\end{itemize}
Therapie

\begin{itemize}
\item Z\"ugiger Transport in eine Urologische Ambulanz

\item vorsichtig ablassen \"uber Harnkatheter oder Blasenpunktion
(Spital)

\end{itemize}
Niereninsuffizienz

Nierensch\"adigung ={\textgreater}  Funktionseinschr\"ankung

akut:

Schock

Entz\"undung

Steine

chronisch:

Chron. Entz\"undung

Zysten

Diabetes mellitus

Mi{\ss}bildungen

Tumore

Schmerzmittelabusus (NSAR)

Folge: Ur\"amie (Harnstoff/Harns\"aure im  Blut [F0E9?] )

Ur\"amie

Nierenversagen

Folgen:

renale An\"amie  (Erythropoietin  [F0EA?] )

Hypertonie

St\"orungen Wasser-/Elektrolythaushalt

Neuropathie

Osteopathie

Therapie:

Dialyse

Nierentransplantation (NTX )

Hodentorsion

Verdrehung eine Hodens im Hodensack  ={\textgreater}

Abschn\"urung von Samenstrang und  Blutgef\"a{\ss}en

Oft nach Eintritt in Pubert\"at, ev. {\textgreater} 20 J.

Symptome:

akuter heftiger Schmerz

Schwellung

Therapie:

hochlagern

umgehende urologisch-chirurgische  Intervention (ad URO!!!)

\subsection{Schwangerschaft und geburtshilfliche Notf\"alle}
Michael Withofner, Dr. med. univ.

Weiblicher Zyklus

Dauer: 28 Tage

Follikelreifung (Hormon: \"Ostrogen )

Eisprung (Ovulation )

Gelbk\"orperphase (Hormon: Progesteron )

2 M\"oglichkeiten:

\begin{enumerate}
\item Einnistung ={\textgreater} Schwangerschaft (Gestation )

\item keine Einnistung ={\textgreater} Regelblutung  (Menstruation )

\end{enumerate}
Schwangerschaft (Gestation)

Eizelle:

8-10 Std. nach Eisprung befruchtungsf\"ahig

ben\"otigt 4-6 Tage f\"ur Wanderung Eierstock- Eileiter-Uterus

befruchtet ={\textgreater} Einnistung

Placenta bildet sich: 

N\"ahrstoff- / O2 -Versorgung des Fetus

Hormon: ${\beta}$-HCG (in Blut und Urin der Mutter nachweisbar)

Schwangerschaft - Verlauf

Embryo : Frucht bis zur 12. SSW

Fetus : Frucht ab 12. SSW bis zur Geburt

erste Kindsbewegungen: 18.-20.SSW

Lungenreife: ca. ab 26. SSW 

Ab der SSW ~24 ist das Kind theoretisch \"uberlebensf\"ahig (und ggfs.
reanimationsplflichtig).

normale Schwangerschaftsdauer:

38-42 Wochen

9 (Kalender-)Monate

10 Lunarmonate

3 Schwangerschaftsdrittel (Trimenon )

Fr\"uhgeburt: bis zur vollendeten 37. SSW

Geburt 

Verlauf

3 Phasen:

\begin{center}
\begin{tabular}{|m{3.12cm}|m{8.7890005cm}|m{1.1919999cm}|}
\hline
\centering Phase\par

 &
\centering Beschreibung\par

 &
\centering Wehen\par

\\\hline
Er\"offnungsphase

 &
Er\"offnungswehen

Geb\"armutterhals (Cervix) verk\"urzt sich,  Muttermund \"offnet sich

Blasensprung

liegender Transport, sonst Beschleunigung der  Geburt und  Gefahr eines 
Nabelschnurvorfalls!! 

Kindskopf tritt in Geburtskanal ein

 &
\centering 10-15 min.\par

\\\hline
Austreibungsphase

 &
Pre{\ss}wehen

Pausen wichtig (Erholungsphase f\"ur Fetus)

Muttermund vollst\"andig er\"offnet

Pre{\ss}wehen: Kontraktion Bauch- plus  Uterusmuskulatur

Kopf wird sichtbar -- normalerweise  Hinterhauptslage (alle anderen
Lagen stellen  Komplikationen dar)

Dammschutz ${\rightarrow}$ sonst Dammri{\ss}, Schutz des Kindes vor
Stuhl der Mutter

Kopf geboren ${\rightarrow}$ K\"orper folgt schnell nach ${\rightarrow}$
 auffangen!

absaugen ${\rightarrow}$ abnabeln ${\rightarrow}$ warm einpacken
${\rightarrow}$ der Mutter auf die Brust legen

 &
\centering 2-3 min.\par

\\\hline
Nachgeburtsphase

 &
Placenta als Nachgeburt

nach 10-20 min. neuerliche Wehen

Geburt der Placenta  (Nachgeburt,  Mutterkuchen)

 &
\\\hline
\end{tabular}
\end{center}
\subparagraph{Neugeborenes - Vitalit\"atsbeurteilung}
APGAR-Score

\begin{enumerate}
\item A ussehen

\item P uls

\item G esichtsbewegung

\item A ktivit\"at

\item R espiration

\end{enumerate}
je 0-2 Pkt., jeweils nach 1min/5min/10min

Notf\"alle - Fr\"uhschwangerschaft

\paragraph{Eileiterschwangerschaft (Tubaria )}
Fehlgeburt (Abortus )

Notf\"alle - Tubaria

= extrauterine Schwangerschaft

Embryo hat sich in Eileiter (Tuba ) eingenistet

ab gewisser Gr\"o{\ss}e reicht Platz nicht mehr aus

Tube rei{\ss}t / platzt

Blutung, Schmerz, akutes Abdomen, Schock

${\rightarrow}$ Bei Frauen mit akutem Abdomen immer nach m\"oglicher
Schwangerschaft fragen!

\paragraph{Abortus}
Beendigung einer Schwangerschaft vor Beendigung der 28.SSW

Fr\"uhgeburt: 24.-37. SSW

Feten {\textgreater}24.SSW haben heute  \"Uberlebenschance!! (dank 
neonatologischer Intensivmedizin)

Symptome

\begin{itemize}
\item vaginale Blutung

\item Abgang von Klumpen

\item Wehenartige Schmerzen

\end{itemize}
Therapie

\begin{itemize}
\item sterile Vorlagen

\item Lagerung nach Fritsch

\item falls erforderlich Schocktherapie

\item je nach Zustand

\item NAW

\item direkt Gyn

\end{itemize}
Notf\"alle - Sp\"atschwangerschaft

\paragraph{Vorzeitige Plazental\"osung}
Aufgrund einer Einblutung zwischen Plazenta und Geb\"armutter kann es zu
einer vorzeitigen Abl\"osung der Plazenta mit dann ungehinderter
Blutung kommen.

\begin{itemize}
\item Lebensgefahr f\"ur Mutter und Kind! 

\end{itemize}
\begin{itemize}
\item Blutung der Mutter und Mangelversorgung des Kindes 

\end{itemize}
Die Diagnose ist durch en Sanit\"ater nicht sicher zu stellen. Im
Vordergrund steht daher die Behandlung der Symptome (Schock) und ein
z\"ugiger Transport. 

Symptome

\begin{itemize}
\item unspezifisch:

\end{itemize}
\begin{itemize}
\item Schock

\item akutes Abdomen

\end{itemize}
Therapie

\begin{itemize}
\item Allgemeine Ma{\ss}nahmen bei lebensbedrohlichen Zust\"anden

\item Schockbek\"ampfung

\item keine vaginale Untersuchung /  Scheidentamponade

\item Z\"ugiger Transport mit Voranmeldung

\end{itemize}
\paragraph{Vena-cava-Syndrom}
Es kann vorkommen dass das Kind die untere Hohlvene der Mutter
abdr\"uckt. Dadurch wird der Blutstrom zum Herzen vermindert und es
kann zu einem pl\"otzlichen Kollaps der Mutter und einer Unterversorung
des Kindes kommen. Die Therapie besteht in einer Linksseitenlagerung,
da die Hohlvene eher rechts der K\"orpermitte verl\"auft. 

Jede hochschwangere Pateintin sollte daher in Linksseitenlage und nicht
in R\"uckenlage transportiert werden. 

Ma{\ss}nahmen

\begin{itemize}
\item Linksseitenlage (Hohlvene verl\"auft eher rechts)

\end{itemize}
\paragraph[EPH-Gestose -- Pr\"aeklampsie - Eklampsie]{EPH-Gestose --
Pr\"aeklampsie - Eklampsie}
\label{bkm:RefHeading39226112}\index{EPH-Gestose}\begin{enumerate}
\item Edema (\"Odeme)

\item Proteinuria (Proteinurie = Eiwei{\ss} im Harn)

\item Hypertonus (RR  ${\uparrow}$)

\end{enumerate}
Synonym: \index{Pr\"aeklampsie}Pr\"aeklampsie

RR {\textgreater}140/90 bei einer Schwangeren ist zu  hoch!!!

Mutter-Kind-Pa{\ss} Untersuchungen zur  Fr\"uherkennung

\index{Eklampsie}Eklampsie = EPH + zerebrale  Krampfanf\"alle

Therapie:

\begin{itemize}
\item schonender Transport

\item \"Uberwachung, O2

\item Notarzt: Krampfdurchbrechung, RR-Senkung

\item (im Vollbild: vorzeitige Beendigung der  Schwangerschaft)

\end{itemize}
Peripartale Notf\"alle

vorzeitiger Fruchtwasserabgang

Nabelschnurvorfall

Placenta praevia

pathologische Geburtslagen

Uterusatonie

Uterusruptur

Neugeborenes atmet nicht (Asphyxie )

Plazenta l\"ost sich nicht

\paragraph[Vorzeitiger \ Fruchtwasserabgang]{Vorzeitiger 
Fruchtwasserabgang}
Fetus schwimmt w\"ahrend  Schwangerschaft in steriler Fl\"ussigkeit 
(Fruchtwasser)

vorzeitige Er\"offnung der Fruchtblase ={\textgreater}  Gefahren:

Nabelschnurvorfall

Infektion (Keimbesiedlung der Fruchth\"ohle)

Ma{\ss}nahmen: (auch f\"ur Transport ins KH)

\begin{itemize}
\item Linksseitenlagerung (sonst Vena-cava- Syndrom)

\item Becken erh\"oht lagern (Schwerkraft verz\"ogert  Geburt)

\item Nicht mehr aufstehen lassen! Liegender Transport!

\end{itemize}
Notf\"alle - Nabelschnurvorfall

Gefahren:

Abklemmung der Nabelschnur  (Sauerstoffversorgung  [F0EA?] )

Strangulation (Gehirnperfusion [F0EA?] )

Therapie:

Linksseitenlagerung, Beckenhochlagerung

Mutter soll NICHT pressen (Geburt w.m. erst  im KH!!!)

Wehenhemmung (hecheln, medikament\"os) 

Nabelschnurumschlingung des  kindlichen Halses w\"ahrend der Geburt 
sofort l\"osen!!!!

Placenta praevia

Plazenta liegt unterhalb des Fetus in  Geburtskanal

Geburtshindernis

massive Blutungen

Ma{\ss}nahmen:

Lagerung nach Fritsch

Schockbek\"ampfung

NAW

im KH: Notsectio

Notf\"alle -- pathologische  Geburtslagen

alle Lagen, bei denen nicht zuerst der  Kopf im Geburtskanal erscheint 

Stei{\ss}lage

Beckenendlage

Armvorfall

Gefahr

Abdr\"ucken der Nabelschnur

Geburtsstopp

Geburt nur im KH durchf\"uhrbar (Sectio- Bereitschaft!)

In jedem Fall Geburt  verhindern/hinausz\"ogern!!!

\paragraph{Notf\"alle - Uterusatonie}
normalerweise zieht sich in der  Nachgeburtsphase die Uterusmuskulatur 
zusammen

hier: Erschlaffung der Uterusmuskulatur  nach Plazental\"osung
={\textgreater} gro{\ss}fl\"achige  Blutung!!!!!

Ma{\ss}nahmen

\begin{itemize}
\item Allgemeine Ma{\ss}nahmen bei lebensbedrohlichen Zust\"anden

\item Schockbek\"ampfung

\item Lagerung: Fritsch, Beine hoch

\item Evt. Druck mit beiden F\"austen oberhalb  Schambein
${\rightarrow}$ Blutungsverminderung

\item Z\"ugiger Transport mit Voranmeldung

\end{itemize}
\paragraph{Uterusruptur}
pl\"otzlicher Schmerz w\"ahrend Wehen,  dann keine Wehen mehr
${\rightarrow}$ Geburtsstop

Lebensgefahr f\"ur Mutter (Blutung) und  Kind (O2 -Mangel!)

Therapie:

\begin{itemize}
\item Allgemeine Ma{\ss}nahmen bei lebensbedrohlichen Zust\"anden

\item Schockbek\"ampfung

\item Z\"ugiger Transport mit Voranmeldung

\end{itemize}
\paragraph{Asphyxie}
Neugeborenes atmet nicht oder nicht  ausreichend

siehe APGAR - Score

Therapie:

absaugen

Fu{\ss}sohlen und R\"ucken klopfen, reiben

O2

Reanimation, Notarzt

\paragraph{Plazenta l\"ost sich nicht}
Plazenta ist in die Uterusmuskulatur  eingewachsen (pathologisch) und
wird nicht nach einer halben Stunde nach Kindsgeburt geboren. 

Ma{\ss}nahmen

\begin{itemize}
\item Nabelschnur am Oberschenkel fixieren

\end{itemize}
NICHT ziehen!!!

\begin{itemize}
\item sterile Vorlage

\item ad Krankenhaus

\end{itemize}
\paragraph{Vaginale Blutungen}
verst\"arkte Regelblutung (Menorrhagie ): nie lebensbedrohlich

Verletzungen (Vergewaltigung, etc.)

in der Menopause: V.a. Malignom  (Tumorzerfallsblutung)

im Rahmen einer Schwangerschaft: Ursache  abkl\"aren

Therapie (ursachenunabh\"angig):

\begin{itemize}
\item Lagerung nach Fritsch

\item Gyn-Ambulanz

\end{itemize}
\subsection{P\"adiatrie - Kinderheilkunde}
Michael Withofner, Dr. med. univ.

Prinzipieller Umgang mit Kindern

\begin{itemize}
\item schonend

\item jegliche Aufregung vermeiden

\item Bezugsperson wichtig

\item Beruhigung

\item Kind w\"ahrend Transport am Besten eng bei  Vertrauensperson
belassen 

\end{itemize}
\paragraph[Anatomische und physiologische \ Besonderheiten]{Anatomische
und physiologische  Besonderheiten}
Kinder sind keine kleinen  Erwachsenen!

\begin{center}
\tablehead{}\begin{supertabular}{|m{5.235cm}|m{8.066cm}|}
\hline
 &
\\\hline
Atemfrequenz:

 &
Neugeborenes: 40x / min.

S\"augling: 20x / min.

Kleinkind: 15x / min.

Beatmung beim Kind immer in  Neutralposition (Kopf nicht \"uberstreckt)

Kehlkopf steht h\"oher, ist verkippt

\\\hline
Zeichen der Atemnot

 &
\begin{itemize}
\item Einziehungen zwischen Rippen

\item \index{Nasenfl\"ugeln}Nasenfl\"ugeln

\item beschleunigte Atmung

\item Atemger\"ausche

\end{itemize}
\\\hline
Pulstasten beim S\"augling:

 &
Oberarmarterie (Oberarm Innenseite)

Achsel

\\\hline
Herzfrequenz:

 &
Neugeborenes: 140 -- 160 / min.

S\"augling: 120 / min.

Kinder: 80-100 / min.

Bradykardie  ist fast immer  hypoxiebedingt ={\textgreater}
lebensgef\"ahrlich!!!

\\\hline
Wasser-, Elektrolyt- und W\"armehaushalt

 &
Kinder haben im Verh\"altnis zu ihrem  K\"orpergewicht eine
gr\"o{\ss}ere K\"orperoberfl\"ache  als Erwachsene!

h\"oherer Wasseranteil

Austrocknung und Unterk\"uhlung schneller  m\"oglich!!!

\\\hline
\end{supertabular}
\end{center}
Fremdk\"orperaspiration 

Verlegung der Atemwege durch  inhalierte Fremdk\"orper (Nahrungsmittel, 
Spielzeug)

Symptome:

\begin{itemize}
\item pl\"otzlich starker Husten

\item Atemnot

\item Stridor oder Keuchen

\item W\"urgen

\end{itemize}
Therapie:

\begin{itemize}
\item Schlag zwischen die Schulterbl\"atter tangential  (Kopf tief)

\item bei gr\"o{\ss}eren Kindern auch Heimlich-Maneuvre

\item keine Entfernungsversuche

\item Kind sitzen lassen, schonender Transport

\item Beatmung bei Atemstillstand

\end{itemize}
Laryngitis

Pseudokrupp / Kruppsyndrom

Obstruktion der oberen Atemwege

durch verschiedene respiratorische Viren

Glottis und Subglottis

Symptome:

\begin{itemize}
\item inspiratorischer  Stridor

\item trockener bellender Husten

\item Heiserkeit

\item Dyspnoe

\item kaum Fieber

\item typisches Alter: 3 Monate -- 6 Jahre

\end{itemize}
Therapie:

\begin{itemize}
\item kalte  Luft oder hei{\ss}er  Dampf

\item kaltes Wasser trinken lassen (DD Epiglottitis)

\item Notarzt: Cortison o.\"a.

\item immer Eltern und  Kind beruhigen

\end{itemize}
Epiglottitis

bakterielle Entz\"undung und daraus  folgende Schwellung des Kehldeckels

selten

ev. lebensbedrohlich

4.-5.Lj. Geh\"auft

schwere Erkrankung, mit richtiger  Behandlung aber meist gute Prognose

Symptome:

\begin{itemize}
\item hochfieberhafte Erkrankung

\item {\quotedblbase}hot potato voice{\textquotedblleft}

\item inspiratorischer Stridor

\item Mund offen, Speichelflu{\ss}, spricht kaum!!!

\item Angst!!!!!

\end{itemize}
Therapie:

\begin{itemize}
\item Notarzt

\item O2 , Luftbefeuchtung

\item Kinder sitzen lassen!

\end{itemize}
Laryngitis vs. Epiglottitis

Laryngitis:

wenig bedrohlich

Kind ist weinerlich

kann schlucken

schreit ev.

SIDS -- Sudden Infant Death Syndrome -- Pl\"otzlicher Kindstod

{\quotedblbase}\index{pl\"otzlicher Kindstod}pl\"otzlicher
Kindstod{\textquotedblleft}

\begin{itemize}
\item pl\"otzlich, meist im Schlaf

\item vorwiegend 1.-4.Lebensmonat

\item Saisonal (J\"an bis M\"arz)

\item famili\"ar und bei Fr\"uhchen geh\"auft

\end{itemize}
Ursachen

\begin{itemize}
\item letztendlich unbekannt

\item zentrale Atemregulationsst\"orung (unreifes  Atemzentrum?)

\end{itemize}
Warnzeichen:  lange Atempausen

Prophylaxe: Apnoemonitoring bei  Risikokindern

Vorgehensweise

Eltern sollen handeln k\"onnen: Gef\"uhl geben,  da{\ss} sie alles f\"ur
ihr Kind getan haben (sonst:  Schuldgef\"uhle)

im Zweifelsfall Pseudoma{\ss}nahmen (Transport  ins KH unter Reanimation
statt Belassung)

Zwilling unbedingt hospitalisieren!

\index{Ertrinkungsunfall}Ertrinkungsunfall

Tod in den ersten 24 h nach Unfall: Ertrinken

ansonsten: \index{Beinahe-Ertrinken}Beinahe-Ertrinken

Kleinkinder sind gef\"ahrdet (mangelnde  Aufsicht, k\"onnen nicht
schwimmen)!!!!!

Faktoren:

\begin{itemize}
\item Kleinkind (kann noch nicht schwimmen,  Panik!)

\item Gefahrenquellen:

\item Gartenteiche, Swimmingpools, schlecht  gesichert

\item Gew\"asser

\item unbeaufsichtigte Wannenb\"ader

\item mangelnde Aufsicht

\end{itemize}
Therapie:

\begin{itemize}
\item Vor Ort: CPR

\end{itemize}
Nobody is dead until warm and dead!

St\"orungen des Nervensystems

\paragraph{Krampfanf\"alle:}
\begin{itemize}
\item fokal

\item generalisiert

\item Gelegenheitskr\"ampfe:

\item Einfache Fieberkr\"ampfe

\item Komplizierte Fieberkr\"ampfe

\item Andere (Meningitis, Trauma, Synkope etc.)

\item Affektkr\"ampfe:

\item Pavor nocturnus

\end{itemize}
\paragraph{Fieberkr\"ampfe}
Einfach:

\begin{itemize}
\item erbliche Vorbelastung

\item schlaff oder generalisiert krampfend

\item selbstlimitierend

\item reifungsbedingt (Kleinkinder {\textless} 5 J.), kein 
Dauerschaden, kein  \"Ubergang in eine  Epilepsie

\end{itemize}
Kompliziert:

\begin{itemize}
\item dauern l\"anger

\item Bei Fieber \"ofters hintereinander

\item Bei prim\"ar neurologisch auff\"alligen Kindern  geh\"auft

\item Schlechtere Prognose

\item Neigung, in Epilepsie \"uberzugehen

\item IMMER medikament\"os unterbrechen  (Notarzt)

\item Kr\"ampfe {\textgreater} 20 min. ={\textgreater} Dauerschaden!!! 

\end{itemize}
Therapie:

\begin{itemize}
\item Wie {\quotedblbase}normaler{\textquotedblleft} Krampfanfall:

\item \item (Fiebersenkung unter  38,5{\textdegree}C (medikament\"os, 
mechanisch))

\item (Krampfdurchbrechung, falls notwendig  medikament\"os)

\end{itemize}
Kindesmi{\ss}handlung

Absichtliche psychische oder physische  Gewalt gegen Kinder, die zu 
Verletzungen, Tod und  Entwicklungsst\"orungen f\"uhren kann

Kommt in ALLEN sozialen Schichten vor

Hinweise:

Langes Intervall zwischen  Verletzungszeitpunkt und Verst\"andigung der 
Rettung

H\"aufiger Arzt-/KH-Wechsel

Anamnese und Verletzungsmuster  widersprechen einander!

Kindesmi{\ss}handlung

Symptome:

mehrere Verletzungen verschiedenen Alters  (R\"ontgen)

erkennbarer {\quotedblbase}verletzungsformender{\textquotedblleft} 
Gegenstand (Schuhabdruck etc.)

Sch\"uttelverletzungen (Subduralh\"amatom)

Zeichen der Vernachl\"assigung

Genitalregion (Verletzungen, Schmerzen ={\textgreater}  sexueller
Mi{\ss}brauch)

Verhaltensst\"orungen

Kindesmi{\ss}brauch

Ma{\ss}nahmen:

Daran denken!!! (Verbrennungen, SHT,  Knochenbr\"uche, ...)

Vorsicht mit voreiligen  Verdachts\"au{\ss}erungen!!!

ev. Vorwand f\"ur Hospitalisierung suchen

dem KH Verdacht vorsichtig mitteilen

\subsection{Psychiatrische Notf\"alle}
Michael Withofner, Dr. med. univ.

10\% aller psychiatrischen Erkrankungen  liegt eine organische Ursache
zugrunde.

Bei Erstmanifestation einer  psychiatrischen Erkrankung MUSS  daher eine
organische Ursache  ausgeschlossen werden!!!

Der Umgang mit psychiatrischen Patienten ist heikel und erfordert viel
Feingef\"uhl sowie die F\"ahigkeit, auf das Gegen\"uber einzugehen. 

Akute Intervention

Indikation:

Unmittelbare Lebensgefahr (Eigengef\"ahrdung)

Gesundheit und Leben anderer gef\"ahrdet (Fremdgef\"ahrdung)

erschwerend: unkooperativer Patient

2 M\"oglichkeiten:

\begin{enumerate}
\item Freiwillige Behandlung (Einwilligung des Patienten)

\item Unterbringung (Behandlung gegen oder ohne Willen des Patienten)

\end{enumerate}
Unterbringung (gem. Unterbringungsgesetz) 

fr\"uher: {\quotedblbase}Parere{\textquotedblleft}

Die \index{Unterbringung}Unterbringung wird durch das
\index{Unterbringungsgesetz}Unterbringungsgesetz geregelt. Auszug:

{\S} 3. In einer Anstalt darf nur untergebracht werden, wer

1. an einer psychischen Krankheit leidet und im Zusammenhang damit sein
Leben oder seine Gesundheit oder das Leben oder die Gesundheit anderer
ernstlich und erheblich gef\"ahrdet und

2. nicht in anderer Weise, insbesondere au{\ss}erhalb einer Anstalt,
ausreichend \"arztlich behandelt oder betreut werden kann.

Unterbringung ohne Verlangen

{\S} 8. Eine Person darf gegen oder ohne ihren Willen nur dann in eine
Anstalt gebracht werden, wenn ein im \"offentlichen Sanit\"atsdienst
stehender Arzt oder ein Polizeiarzt sie untersucht und bescheinigt,
da{\ss} die Voraussetzungen der Unterbringung vorliegen. In der
Bescheinigung sind im einzelnen die Gr\"unde anzuf\"uhren, aus denen
der Arzt die Voraussetzungen der Unterbringung f\"ur gegeben erachtet.

{\S} 9. (1) Die Organe des \"offentlichen Sicherheitsdienstes sind
berechtigt und verpflichtet, eine Person, bei der sie aus besonderen
Gr\"unden die Voraussetzungen der Unterbringung f\"ur gegeben erachten,
zur Untersuchung zum Arzt ({\S} 8) zu bringen oder diesen beizuziehen.
Bescheinigt der Arzt das Vorliegen der Voraussetzungen der
Unterbringung, so haben die Organe des \"offentlichen
Sicherheitsdienstes die betroffene Person in eine Anstalt zu bringen
oder dies zu veranlassen. Wird eine solche Bescheinigung nicht
ausgestellt, so darf die betroffene Person nicht l\"anger angehalten
werden.

(2) Bei Gefahr im Verzug k\"onnen die Organe des \"offentlichen
Sicherheitsdienstes die betroffene Person auch ohne Untersuchung und
Bescheinigung in eine Anstalt bringen.

(3) Der Arzt und die Organe des \"offentlichen Sicherheitsdienstes haben
unter m\"oglichster Schonung der betroffenen Person vorzugehen und die
notwendigen Vorkehrungen zur Abwehr von Gefahren zu treffen. Sie haben,
soweit das m\"oglich ist, mit psychiatrischen Einrichtungen
au{\ss}erhalb einer Anstalt zusammenzuarbeiten und erforderlichenfalls
den \"ortlichen Rettungsdienst beizuziehen.

Daraus ergibt sich bei fehlender Freiwilligkeit:

Voraussetzungen ({\S}3)

\begin{itemize}
\item Psychiatrische Erkrankung (urs\"achlich!),
(${\rightarrow}$~{\S}3~Abs.~1)

\item Eigen-  und/oder Fremdgef\"ahrdung, (${\rightarrow}$~{\S}3~Abs.~1)

\item Keine weniger einschr\"ankende M\"oglichkeit zur Abwendung der
Gef\"ahrdung, (${\rightarrow}$~{\S}3~Abs.~2)

\end{itemize}
Vorgehensweise ({\S}{\S}~8, 9)

Ein {\quotedblbase}im \"offentlichen Sanit\"atsdienst stehender Arzt
oder ein Polizeiarzt{\textquotedblleft} die Voraussetzungen zu
bescheinigen und zu Begr\"unden. In Wien ist das der
{\quotedblbase}\index{Amtsarzt}Amtsarzt{\textquotedblleft}, welcher
\"uber die Polizei verst\"andigt wird. 

Der Notarzt, auch wenn er von der Gemeinde Wien gestellt wird, hat nicht
die Kompetenz eine Unterbringung anzuordnen. (${\rightarrow}$~{\S}8)

F\"ur den Transport ist grunds\"atzlich die Polizei zust\"andig
(${\rightarrow}$~{\S}9~Abs.~1)! Die Polizei darf des Rettungsdienst
quasi als {\quotedblbase}Taxi{\textquotedblleft} beiziehen
(${\rightarrow}$~{\S}9~Abs.~3). 

Bei Gefahr im Verzug darf die Polizei auch ohne Bescheinigung eine
Person auf eine psychiatrische Abteilung zwecks Unterbringung
verbringen (${\rightarrow}$~Abs.~2). 

Im Zweifelsfall Polizei nachalarmieren...

Weiters: 

Eine Unterbringung kann nur in einer psychiatrischen Abteilung erfolgen!

Eine Unterbringung kann auch freiwillig erfolgen, da dann auch der
Transport freiwillig stattfindet sind dabei keine besonderen
Bestimmungen zu beachten.

Umgang mit Patienten

Eigenschutz!

ruhiger Kontakt

Situation:

Belassung? (ist Patient \"uberhaupt  reversf\"ahig?)

Unterbringung? (notwendig / rechtlich  zul\"assig?)

Polizei / Amtsarzt bzw. Notarzt? (stabil?  Eskalation?)

Leitsymptome

\begin{itemize}
\item Verwirrtheit, Bewu{\ss}tseinseintr\"ubung

\item Realit\"atsverlust (Wahn, Halluzinationen )

\item psychomotorische Unruhe, Aggressivit\"at

\item Depression , Verzweiflung, Suizidalit\"at

\item Angst

\item Antriebslosigkeit, Stupor

\item Verwirrtheit / Bewu{\ss}tseinstr\"ubung

\item einfache Verwirrung

\item Altersdemenz

\item delirante Verwirrung

\item Einfache Verwirrung

\end{itemize}
Kann bei allen {\quotedblbase}normalen{\textquotedblleft},
{\quotedblbase}k\"orperlichen{\textquotedblleft} Notf\"allen, die mit
Bewusstseins\-ein\-schr\"ankungen ein\-her\-gehen, auftreten!

Bsp.: Epilepsie , Commotio cerebri, Exsikkose etc.

Altersdemenz

Fortschreitende Einschr\"ankung der  geistigen Leistungsf\"ahigkeit des
Gehirns  im Alter

Betrifft

\begin{itemize}
\item Ged\"achtnis

\item Denkverm\"ogen

\item Sprache

\item Motorik

\item Pers\"onlichkeitsstruktur

\item delirante Verwirrung

\end{itemize}
Symptome

Desorientierung

zeitlich

\"ortlich

psychomotorische Unruhe

optische Halluzinationen

Personenverkennung

Wechselnde Symptomatik!!!!

delirante Verwirrung 

organische Symptome:

Tachykardie

Mydriasis (Pupillen weit)

trockene rote Haut

Kr\"ampfe

z.B. Entzugsdelir bei Alkohol, Drogen

Ma{\ss}nahmen

NAW, ad Interne \"Uberwachungsstation (RR  instabil!)

RR-Kontrollen

Realit\"atsverlust

M\"ogliche Ursachen

Endogene Psychosen (Schizophrenie,  schizoaffektive Psychose)

Organische Psychosen (Alkoholparanoia,  Drogenpsychose (z.B. Coke bugs
),  Encephalitis)

Realit\"atsverlust - Wahn

Wahn = falsche, allerdings in sich  logische Auffassung der Realit\"at,
an der  unkorrigierbar festgehalten wird  (paranoid = wahnhaft) 

Themen: Verfolgungswahn, Gr\"o{\ss}enwahn,  Verarmungswahn,...

Realit\"atsverlust - Halluzination

Halluzinationen: Trugwahrnehmungen,  f\"ur alle Sinne m\"oglich
(wei{\ss}e M\"ause,  Stimmen,
{\quotedblbase}Geisterber\"uhrungen{\textquotedblleft},...)

Psychomotorische Unruhe /  Aggressivit\"at

Organische Ursachen ausschlie{\ss}en!!!

Intoxikation durch Alkohol, Drogen,  Medikamente

Hypoglyk\"amie

St.p. SHT

Hirntumor

psychomotorische Unruhe /  Aggressivit\"at

Schwere Erregungszust\"ande:

Katatonie : Starre mit Sprechstereotypien, Grimassieren,
Befehlsautomatismen, Patient imitiert Untersucher

Manie: Distanzlosigkeit, aggressiv gef\"arbte Agitiertheit

psychoreaktiv: abnorme Pers\"onlichkeit

psychosoziale Krise: Reaktion auf Anla{\ss} / Stre{\ss}reaktion

psychomotorische Unruhe /  Aggressivit\"at

Hysterie (funktionelle Erregungszust\"ande):

anla{\ss}gebunden, entsprechende Grundpers\"onlichkeit

psychosomatischer Funktionsverlust (L\"ahmungen, Erblindung) ohne echte
organische Ursache

Entspricht psychischer Abwehrreaktion

hysterischer Krampfanfall:

keine Verletzungen, keine Apnoe, atypische Kr\"ampfe, kein
Stuhl-/Harnverlust

unbewu{\ss}t {\quotedblbase}gespielt{\textquotedblleft}

psychomotorische Unruhe /  Aggressivit\"at

Ma{\ss}nahmen:

ruhig bleiben! Achtung auf eigenen Tonfall!

NICHT provozieren!

falls Patient bewaffnet, ist Polizei zust\"andig

Selbstschutz besonders wichtig!!!

Manie

rastlose Unruhe - {\quotedblbase}energiegeladen{\textquotedblleft}

gereiztes Verhalten

\"Uberaktivit\"at

Distanzverlust

Verwirrungszust\"ande

Aggression

Selbst\"ubersch\"atzung

ev. als bipolare St\"orung abwechselnd mit  Depression

Depression

Traurigkeit

Schuldgef\"uhle

fehlender Antrieb

Schlafst\"orungen

Suizidgedanken

Depression / Verzweiflung / Suizidalit\"at

bei vielen psychiatrischen, aber auch  organischen Erkrankungen

Grundstimmung entweder:

antriebslos  oder

agitiert (h\"ohere Suizid gefahr!)

Ma{\ss}nahmen:

auf Gegenwart konzentrieren, soziale  Umgebung miteinbeziehen, 
Einf\"uhlungsverm\"ogen!!!

Depression / Verzweiflung / Suizidalit\"at

Selbstmordversuch ={\textgreater} station\"are  Aufnahme, falls
notwendig Parere!!!

Selbstmordphantasien: je konkreter,  desto gef\"ahrlicher!

Suizid = Selbstmord

Achtung, ev. sogar Mitnahmeselbstmord  m\"oglich!

Angst

Gef\"uhl der Bedrohung mit vegetativer  Begleitsymptomatik

oft ohne erkennbaren Grund

bei psychiatrischen und organischen  Krankheitsbildern

Angst

Panikattacke

\begin{itemize}
\item pl\"otzliches Angstgef\"uhl

\item Dauer: Sekunden bis Minuten

\item Thoraxbeschwerden (Infarktangst)

\item Hyperventilation

\item organisch mehrfach durchuntersuchte  Patienten, wirken gesund

\item Erwartungsangst

\item Vermeidungsverhalten

\item Antriebslosigkeit / Stupor

\end{itemize}
Stupor = Erschlaffung

Symptome:

\begin{itemize}
\item ausdrucksloser Patient

\item unbeweglich

\item starre Mimik

\item spricht nicht, i{\ss}t nicht ={\textgreater} ev. bis Verhungern!

\item Suchterkrankungen

\item prim\"ar psychisches Problem

\item sekund\"ar k\"orperliche und soziale Folgen

\item eigengesetzlicher Ablauf

\item zunehmender Kontrollverlust \"uber  eigenes Verhalten

\end{itemize}
Alkohol- und Medikamentenentzug

Symptome

\begin{itemize}
\item Schwitzen

\item Unruhe

\item Tremor  (Zittern der H\"ande)

\item Delirium  (Bewu{\ss}tseinsst\"orungen, Halluzinationen )

\item Krampfanf\"alle

\item RR-Entgleisung

\item {\quotedblbase}Cold turkey{\textquotedblleft} bei Heroinentzug

\end{itemize}
Anamnese

\begin{itemize}
\item Pl\"otzlich kein Zugang mehr zu Alkohol:

\item K\"urzlicher Einzug in Altersheim

\item k\"urzlich pflegebed\"urftig geworden

\item ...

\end{itemize}
Therapie

\begin{itemize}
\item Ruhige Gespr\"achsf\"uhrung, keine Kritik

\item Suchtanamnese

\end{itemize}
Notarztindikation

\begin{itemize}
\item Krampfanf\"alle

\item RR-Entgleisung

\item Transportverweigerung

\end{itemize}
\subsection{Systemische Hitzesch\"aden}
Der menschliche K\"orper h\"alt seine Kerntemperatur  im Normalfall
zwischen 36,5{\textdegree}  und 37,5{\textdegree} C  konstant.
\"Ubersteigt  die W\"armezufuhr die effektive  W\"armeabgabe,
resultiert ein Hitzeschaden : Hitzekollaps ,  Hitzeersch\"opfung ,
Hitzschlag  (bzw. bei  direkter Sonneneinstrahlung auf den Kopf ein 
Sonnen\-stich ).

Hitzekollaps

\begin{itemize}
\item W\"armestau  (zugef\"uhrte W\"arme  kurzfristig {\textgreater}
abgegebene)

\item Erweiterung der Blutgef\"a{\ss}e in der Haut  (maximal)

\item RR ${\downarrow}$ , Perfusion des Gehirns kurzzeitig
${\downarrow}$

\item kurze Ohnmacht, {\quotedblbase}schwarz vor
Augen{\textquotedblleft}

\end{itemize}
\begin{itemize}
\item Flachlagerung: ${\rightarrow}$ Gehirnperfusion wieder
${\uparrow}$, bewu{\ss}tseinsklar

\item nach Sekunden gebessert, {\quotedblbase}regelt sich von
selbst{\textquotedblleft}

\end{itemize}
\begin{itemize}
\item ${\rightarrow}$ Therapie: Patient sitzen lassen, Fl\"ussigkeit,
ev. BZ

\end{itemize}
Hitzeersch\"opfung

\begin{itemize}
\item W\"armeaufnahme {\textgreater}  W\"armeabgabe

\item durch starkes Schwitzen Wasser- und  Elektrolytverlust

\item starkes Durstgef\"uhl

\item RR ${\downarrow}$, Tachykardie

\item K\"orpertemperatur normal

\end{itemize}
Therapie

\begin{itemize}
\item K\"uhlung, Schatten,  Fl\"ussigkeitszufuhr, Vitalparameter, ev. 
BZ 

\end{itemize}
Hitzschlag

\begin{itemize}
\item Kopfschmerz, Schwindel, Erbrechen

\item Haut zun\"achst trocken, rot und hei{\ss},  sp\"ater grau-blau,
Kopf hochrot

\item stark  erh\"ohte  K\"orpertemperatur (bis  43{\textdegree}C
m\"oglich ={\textgreater} LEBENSGEFAHR !!!)

\item Bewusstseinsst\"orung, Bewusstlosigkeit

\item Schockzeichen

\item schnelle, flache Atmung

\end{itemize}
Ma{\ss}nahmen

\begin{itemize}
\item Flachlagerung an k\"uhlem Ort (Schatten), Beine und Kopf
hochlagern

\item beengende Kleidungsst\"ucke lockern

\item K\"uhlung von au{\ss}en (F\"acher, mit Eisw\"urfeln abreiben etc.)

\item K\"uhlung von innen: f\"ur ANSPRECHBARE Patienten viel k\"uhle
(alkoholfreie) Fl\"ussigkeit, ev. Elektrolyt-Getr\"ank

\item Schocklagerung u. -bek\"ampfung

\item Vitalparameter, O2, BZ

\end{itemize}
Sonnenstich

Der Sonnenstich ist eine Sonderform des  Hitzeschadens: Er entsteht
durch direkte  Sonneneinstrahlung auf den ungesch\"utzten  Kopf. 

\begin{itemize}
\item direkte  Sonneneinstrahlung auf den ungesch\"utzten 
Kopf${\rightarrow}$ Temperatur  im Gehirn  steigt  ${\rightarrow}$
Reizung der Hirnh\"aute  ${\rightarrow}$ Symptome wie bei 
Hirnhautentz\"undung (Meningitis) bzw.  neurologische Symptome durch 
Schwellung des ganzen Gehirnes

\item Kann kombiniert mit anderen Hitzesch\"aden auftreten

\item Gesichts- u. Kopfhaut hochrot u. hei{\ss}, restlicher K\"orper
k\"uhl

\item Abgeschlagenheit, heftiger Kopfschmerz, Schwindel

\item Unruhe,Verwirrtheitszust\"ande, Brechreiz

\item Nackensteifigkeit

\item Schwere F\"alle: Krampfanf\"alle, Bewusstlosigkeit

\end{itemize}
Ma{\ss}nahmen

\begin{itemize}
\item Flachlagerung an k\"uhlem Ort (Schatten), Kopf hochlagern

\item beengende Kleidungsst\"ucke lockern

\item K\"uhlung von au{\ss}en (F\"acher, mit Eisw\"urfeln abreiben,
kalte Umschl\"age auf die Stirn etc.)

\item K\"uhlung von innen: f\"ur ANSPRECHBARE Patienten viel k\"uhle
(alkoholfreie) Fl\"ussigkeit, ev. Elektrolyt-Getr\"ank

\item Schocklagerung u. -bek\"ampfung

\item Vitalparameter, O2, BZ

\end{itemize}
\subsection{Unterk\"uhlung}
\"Ahnlich wie beim Schock kommt es auch  bei der Unterk\"uhlung zur
Zentralisation  des Blutes, um die  K\"orperkerntemperatur aufrecht
erhalten  zu k\"onnen. Je niedriger diese  Kerntemperatur jedoch sinkt,
desto  schwerwiegender sind die Folgen!

Phasen

\begin{itemize}
\item Kampfphase:  36,5{\textdegree}-34{\textdegree}C, K\"altezittern,
Herzfrequenz steigt ({\textgreater}100 Schl\"age/Minute), schnelle
Atmung, Schmerzen

\item Ersch\"opfungsphase: 34{\textdegree}-30{\textdegree}C,
Bewusstseinstr\"ubung, K\"altestarre, Herzfrequenz sinkt wieder
({\textless}60/min), zu langsame Atmung, keine Schmerzen mehr

\item K\"altenarkose: 30{\textdegree}-27{\textdegree}C,
Bewusstlosigkeit, unregelm\"a{\ss}iger Herzschlag, Atempausen

\item {\textless}27{\textdegree}C ${\rightarrow}$ klinisch tot

\end{itemize}
Ma{\ss}nahmen

\begin{itemize}
\item nur soviel wie UNBEDINGT n\"otig ist  bewegen (BERGUNGSTOD)

\item Schutz vor weiterer Abk\"uhlung, KEINE AKTIVE ERW\"ARMUNG!!!

\item Woll- + Aludecken! M\"utze!

\item ev. Wiederbelebungsma{\ss}nahmen

\item Notarzt !

\end{itemize}
WICHTIG

Bei Unterk\"uhlung gilt:

{\quotedblbase}Nobody is dead until warm and dead.{\textquotedblleft}

Schwer unterk\"uhlte Patienten k\"onnen auch noch nach Stunden (!) ohne
Kreislauft\"atigkeit erfolgreich reanimiert werden!

\subsection{Vergiftungen}
Michael Withofner, Dr. med. univ.

Vergiftung - Definition

Vergiftung bedeutet eine (meist  dosisabh\"angige) organische 
Funktionsbeeintr\"achtigung / Sch\"adigung  des Organismus durch eine
(organische/ anorganische) chemische Substanz!

(z.B.: NaCl ={\textgreater} 20 dag t\"odlich,  Medikamente
={\textgreater} teilweise im Mikrogramm- Bereich t\"odlich!!!)

Vergiftungen - erkennen

(Au{\ss}en-)Anamnese besonders wichtig (Vergiftungsart/Ursache), oft
einziger Aufschlu{\ss}, da Patient ev. bewu{\ss}tseinsgetr\"ubt

6 W -Fragen:

\begin{enumerate}
\item Wer? (Alter, KG, Vorerkrankungen,... ={\textgreater}
Risikoabsch\"atzung!)

\item Was? (Art des Giftes ={\textgreater} ev. zu erwartende Symptome)

\item Wann? (Zeit seit Einnahme ={\textgreater} Magensp\"ulung ja/nein?)

\item Wie? (inhalativ/oral/perkutan/intraven\"os etc.)

\item Wieviel? (Dosisabsch\"atzung)

\item Warum? (akzidentiell/absichtlich)

\end{enumerate}
Ursachen:

akzidentiell: z.B.

(Arbeits-)Unfall: Reizgase, Chemikalien

Kinder: {\quotedblbase}Zuckerl{\textquotedblleft} = Medikamente

fahrl\"assig: falsch beschriftete Flaschen  (Putzmittel in
Getr\"ankeflasche etc.)

absichtlich:

Suizid: meist Medikamente (Benzodiazepine,  Barbiturate)

Drogen: \"Uberdosis (gewollt oder ungewollt), 
{\quotedblbase}Verschnitt{\textquotedblleft}

Mord: kriminelle Absicht, bei Verdacht Polizei  einschalten!!!

Aufnahme

oral : Verschlucken  der giftigen  Substanz 

inhalativ : Einatmen

rauchen/schnupfen: Kombination oral /  inhalativ

transkutan : fettl\"oslich e Substanz kann  Hautbarriere durchdringen
(z.B.:  organische L\"osungsmittel,  Insektenschutzmittel etc.)

intraven\"os/intramuskul\"ar/subkutan: spritzen

Vergiftungen -- unspezifische  Symptome

gastrointestinal:

Erbrechen

Durchfall (Diarrhoe )

Schmerzen im GIT

ZNS:

Bewu{\ss}tsein [F0EA?][F020?](Rausch)

ev. Koma

zerebrale Krampfanf\"alle

Vorgehensweise

5-Fingerregel:

\begin{enumerate}
\item Vitalparameter

\item Giftentfernung

\item Antidot-Therapie (Notarzt)

\item Gift-Asservierung

\item Transport

\end{enumerate}
Vergiftungen - Basisma{\ss}nahmen

allgemein:

Vitalparameter st\"andig \"uberwachen

O2  !!!

Asservierung (Gift bekannt -- ev. spezifische  Ma{\ss}nahmen!)

was liegt herum?

Mistk\"ubel durchsuchen

Erbrochenes (in Nierentasse mitnehmen!)

Vergiftungen -- Erstma{\ss}nahmen

prim\"are Giftentfernung / stoppen der Giftaufnahme

oral : 

Patient soll nichts essen / trinken

NICHT absichtlich Erbrechen herbeif\"uhren

im Spital: Aktivkohle (h\"aufig) bzw.  Magensp\"ulung (selten)

inhalativ :

Eigenschutz!!!

Entfernung aus Giftatmosph\"are

O2  / Frischluft

Vergiftungen -- Erstma{\ss}nahmen

prim\"are Giftentfernung / stoppen der Giftaufnahme

transkutan :

Eigenschutz!!!

Entfernung kontaminierter Kleidung

Sp\"ulung der betroffenen Hautpartien mit Wasser  (richtig abrinnen
lassen)

okul\"ar  (Auge):

Sp\"ulung mit Wasser / Ringerl\"osung (nach au{\ss}en abrinnen lassen!)

Kontaktlinsen entfernen

Vergiftungen -- spezifische  Ma{\ss}nahmen

bei bekanntem Gift ev. Antidot  (Gegengift)-Gabe m\"oglich (Notarzt / 
Spital)

Vergiftungszentrale:

406 43 43

bei bekanntem Giftstoff (6 W-Fragen) ev. spezifische Anweisungen
telefonisch

spezielle Vergiftungen

eventuell anhand von Symptomen  diagnostizierbar

{\quotedblbase}Toxidrome{\textquotedblleft} ={\textgreater} allerdings
eigene  Wissenschaft

dennoch gut, einige {\quotedblbase}typische{\textquotedblleft} 
Vergiftungen zu kennen

Medikamente / Drogen

Benzodiazepine  (Tranquilizer,  Schlafmittel):

Substanzen: Valium{\textregistered} , Rohypnol{\textregistered} ,
Dormicum{\textregistered} ,  Temesta{\textregistered}  ...

Symptome:

Bewu{\ss}tseinseintr\"ubung

RR [F0EA?]

Spontanatmung m\"a{\ss}ig [F0EA?]

Antidot: Anexate{\textregistered} 

Medikamente / Drogen

Schmerzmittel (NSAR, Nicht Steroidale AntiRheumatika)

Substanzen: 

Paracetamol: Mexalen{\textregistered}

Salicylate: Aspirin{\textregistered} , Brufen{\textregistered} ,
Voltaren{\textregistered}

Symptome / Folgen:

Paracetamol: 

Lebersch\"aden

Salicylate: 

Azidose

Hyperthermie

Kr\"ampfe

Therapie: Antidote im AKH / WSP lagernd,  H\"amofiltration!!!

Medikamente / Drogen

Opiate  (aus Milchsaft des Schlafmohns (Papaver somniferum))

Substanzen: Morphin (Fentanyl{\textregistered} etc.), Codein, Heroin
(illegal)

Unterliegen strengem Suchtgiftgesetz (Schmerzmittel)

Schwarzmarkt (Heroin), {\quotedblbase}Fixer{\textquotedblleft}, sehr
hohes Suchtpotential

Symptome:

Miose (Pupillen eng, stecknadelkopfgro{\ss}, nicht  lichtreaktiv)

Bewu{\ss}tsein [F0EA?]  bis Bewu{\ss}tlosigkeit

Atemdepression bis Atemstillstand (Brady- /Apnoe )

Therapie:

O2 , Beatmung, Vorgehen bei Bewu{\ss}tlosigkeit

Antidot: Narcanti{\textregistered}  (Wirkstoff: Naloxon )

Medikamente / Drogen

Achtung bei Kontakt mit drogenkranken  Patienten:

Immer  an Infektionserkrankung denken!!!

TBC

Hepatitis B + C

HIV

Immer Handschuhe tragen!!!

Blutkontakt unbedingt vermeiden

H\"andedesinfektion 

Medikamente / Drogen

\paragraph{Alkohol}
Substanzen: Methanol (Fuselalkohol,
{\quotedblbase}Vorlauf{\textquotedblleft}), Ethanol
({\quotedblbase}Alk{\textquotedblleft})

SEHR h\"aufige Vergiftung (gesellschaftlich akzeptiert),
{\quotedblbase}Kampfsaufen{\textquotedblleft}!

Symptome: 

Promilleabh\"angig

Euphorie

Sprach-/ Gangunsicherheit, Enthemmung

Bewu{\ss}tseinstr\"ubung,  Schmerzunempfindlichkeit, 
Temperaturregulationsst\"orung, Inkontinenz

Bewu{\ss}tlosigkeit, Atemdepression, Aspiration,  Unterk\"uhlung

Therapie:

Vitalparameter

stabile Seitenlage / Reanimation

BZ messen (Kinder!)

vorsichtige Beatmung (Alkoholisierte erbrechen  leichter!)

Alkohol -- chronische Sch\"aden

Lebersch\"aden

Zirrhose

Aszites

Krampfadern in der unteren Speiser\"ohre  (Blutungsgefahr)

Blutgerinnung [F0EA?]

Vitaminmangel ={\textgreater} 

intellektueller Abbau (Vit B)

An\"amie (Vit B12 , Fols\"aure)

Herzrhythmusst\"orungen

Pankreatitis (Entz\"undung der  Bauchspeicheldr\"use)

Alkohol - Entzug

\subparagraph{Entzugsdelir}
{\quotedblbase}Delirium tremens {\textquotedblleft}

Zittern (Tremor )

Unruhe, Angst, Psychose

Halluzinationen ({\quotedblbase}wei{\ss}e
M\"ause{\textquotedblleft},...)

Tachykardie

RR-Entgleisung

zerebrale Krampfanf\"alle

lebensbedrohlicher Zustand

{\quotedblbase}Akut-{\textquotedblleft}Therapie: Alkohol (Ethanol als
Schnaps  etc.)

im Spital: medikament\"os unterst\"utzter Entzug 

Fl\"ussigkeiten

\paragraph{Korrosiva (S\"auren/Laugen):}
Ver\"atzung

Lokale Sch\"adigung durch irreversible Zerst\"orung (Denaturierung) von
Eiwei{\ss}stoffen. 

betroffenen Organe

GIT

Haut

Augen

S\"auren/Laugen

\subparagraph{Symptome bei Ingestion:}
\"Atzspuren (Mund/Rachen), Abrinnspuren

W\"urgen

Erbrechen

Schmerzen

Vorsicht: Ruptur/Perforation/Blutungsgefahr!!!

Therapie:

\begin{itemize}
\item schluckweise Wasser trinken lassen

\item KEINE Neutralisationsversuche!!! (Milch  obsolet!!!)

\item NIEMALS absichtlich Erbrechen  herbeif\"uhren!!!!

\end{itemize}
S\"auren/Laugen

\subparagraph{Ver\"atzung der Haut}
Schorfbildung

glasige Verquellung

Nekrosen

Schmerzen

Therapie:

Eigenschutz! (Handschuhe)

kontaminierte Kleidung entfernen

betroffene Hautpartien mit Wasser absp\"ulen  (unverletzte Haut
sch\"utzen!)

locker, steril verbinden

S\"auren/Laugen

\subparagraph{Ver\"atzung der Augen}
Hornhauttr\"ubung, sp\"ater Narben

Visusverlust

Schmerzen

Therapie:

\begin{itemize}
\item Sofortige Sp\"ulung mit Wasser / Ringerl\"osung

\item danach beide Augen keimfrei abdecken!

\item Fl\"ussigkeiten

\end{itemize}
\paragraph[Schaumbildner (Wasch-/Putzmittel)]{Schaumbildner
(Wasch-/Putzmittel)}
Symptome:

verminderte Oberfl\"achenspannung:  Aspirationsgefahr ${\uparrow}$

Verlegung der Atemwege

\"Ubelkeit/Erbrechen

Therapie:

Patienten n\"uchtern lassen

NIEMALS absichtlich Erbrechen  herbeif\"uhren!!!!

Stickgase

Bei Gasunf\"allen hat der Selbstschutz Vorrang. Oft ist die Bergung nur
mit schwerem Atemschutz m\"oglich ${\rightarrow}$ FEUERWEHR!

Ein toter Helfer ist kein Helfer.

\paragraph{Kohlenmonoxid (CO)}
entsteht bei unvollst\"andiger Verbrennung (O2- Zufuhr ${\downarrow}$)

300x h\"ohere Affinit\"at zu H\"amoglobin (als O2), 

S\"attigung beim Pulsoxymeter falsch hoch!!!

Symptome:

Dyspnoe  ohne  Zyanose  (rosige Hautfarbe, kirschrote Lippen!!!)

Tachykardie ={\textgreater} Gefahr: MCI

Schwindel,Kopfschmerz, \"Ubelkeit, Erbrechen

ev. Kr\"ampfe bis Koma

Therapie:

Entfernung aus Gasatmosph\"are  (Eigenschutz!!!)

O2 hochdosiert , Hyperventilation, hyperbare  O2 -Therapie

\paragraph[Kohlendioxid (CO2) ]{Kohlendioxid (CO2) }
G\"argas, Verbrennungsgas, entsteht auch bei  Atmung

schwerer als Luft (sammelt sich in Mulden,  Weinkellern etc.)

Symptome:

Kopfschmerz, Schwindel, \"Ubelkeit, Kr\"ampfe, Ersticken

O2-Aufnahme ${\downarrow}$, stimuliert Atemzentrum (Hyperventilation)

Therapie

Bergung (Eigenschutz), Frischluft

10-15 l O2  

Reizgase

Soforttyp : (z.B. Tr\"anengas,...)

Konjunktivitis

starker Hustenreiz

Verz\"ogerungstyp: (z.B. Nitrosegase)

kaum Hustenreiz

Latenzphase (kaum Symptome)

3-24h sp\"ater: toxisches Lungen\"odem

bei Nichtbehandlung ev. irreversible Lungenfibrose

Reizgase

Vorkommen:

Chemische Industrie: Nitrosegase,  Schwefelwasserstoffe, Ammoniak

 Kunststoffindustrie: Chlorwasserstoffe,  Formaldehyd

D\"ungemittelherstellung: Ammoniak,  Nitrosegase, Schwefelwaserstoffe

Haushalt: Chlorgas (Kombination von best.  Haushaltsreinigern mit
Essig), Lacke,  Impr\"agniersprays, ...

Reizgase

Therapie:

Bergung (Eigenschutz!)

100\% O2 , ev. PEEP-Beatmung durch NFS /  Notarzt

wiederholt inhalatives Corticoid (Pulmicort{\textregistered} )  durch
NFS / Notarzt

Bei Reizgas / Stickgas immer auch nach  anderen bewu{\ss}tlosen Personen
fragen /  suchen!!! (unter Eigenschutz!!!)

\section{[RS VII] Traumatologische Notf\"alle}
Michael Withofner, Dr. med. univ.

\begin{itemize}
\item Sch\"adel-Hirn-Trauma (SHT)

\item Thorax trauma

\item Abdominal trauma / Bauchtrauma

\item Beckentrauma

\item Verletzungen der Wirbels\"aule

\item Polytrauma

\item Verletzungen der Extremit\"aten

\item Verbrennungen

\end{itemize}
Definition Trauma

Als Trauma  bezeichnet man eine {\textperiodcentered}~Sch\"adigung,
{\textperiodcentered}~Verletzung, oder {\textperiodcentered}~Wunde, die
durch eine Gewalteinwirkung physikalisch er oder chemisch er Art
verursacht wurde.

\subsection{Sch\"adel-Hirn-Trauma (SHT)}
Folge von Gewalteinwirkung auf den  Sch\"adel

Arten:

\begin{itemize}
\item gedecktes SHT: ohne  Er\"offnung der Dura  mater

\item offenes SHT: mit  Er\"offnung der Dura mater

\end{itemize}
begleitende Verletzungen:

\begin{itemize}
\item Verletzung der Kopfschwarte  (Ri{\ss}quetschwunde)

\item Frakturen von

\end{itemize}
\begin{itemize}
\item Sch\"adeldach

\item Sch\"adelbasis (Blutung aus Ohren, Liquoraustritt  aus der Nase)

\item Gesichtssch\"adel

\item Unterkiefer

\end{itemize}
Gedecktes SHT -- Klassifikation

\begin{itemize}
\item Gehirnersch\"utterung  {}-- Contusio cerebri

\item Gehirnprellung  {}-- Contusio cerebri

\item Gehirnquetschung

\end{itemize}
\paragraph[Gehirnersch\"utterung (Commotio
(cerebri))]{Gehirnersch\"utterung (Commotio (cerebri))}
\index{Gehirnersch\"utterung}\index{Commotio (cerebri)}\begin{itemize}
\item  Bewu{\ss}tlosigkeit /Bewu{\ss}tseinstr\"ubung {\textless}  1h

\item \"Ubelkeit, Erbrechen, Schwindel

\item evt. fl\"uchtige neurologische Ausf\"alle

\item retrograde  Amnesie  ({\quotedblbase}kann sich an den
Unfallhergang nicht erinnern{\textquotedblleft} )

\end{itemize}
\begin{itemize}
\item im CT/MR [${\rightarrow}$R{}-\pageref{bkm:LexComputertomographie}]
keine nachweisbaren Sch\"aden

\item keine  bleibenden Sch\"aden

\end{itemize}
gedecktes SHT - Klassifikation

\paragraph[Gehirnprellung (Contusio (cerebri))]{Gehirnprellung (Contusio
(cerebri))}
\index{Gehirnprellung}\index{Contusio (cerebri)}\begin{itemize}
\item Bewu{\ss}tlosigkeit  {\textgreater}  1h

\item Bewu{\ss}tseinstr\"ubung {\textgreater}24h

\item nachweisbare Gehirnsch\"adigung / bleibende  Ausf\"alle

\item ev. Gehirnnervenabrisse (z.B. Anosmie  bei  Riechnervabri{\ss})

\item  vegetative St\"orungen

\end{itemize}
\paragraph[Gehirnquetschung (SHT III{\textdegree})]{Gehirnquetschung
(SHT III{\textdegree})}
\index{Gehirnquetschung}\begin{itemize}
\item Bewu{\ss}tlosigkeit  {\textgreater}  24h

\begin{center}
 [Warning: Image ignored] % Unhandled or unsupported graphics:
%\includegraphics[width=7.09cm,height=2.2cm]{AASdev20090228-img89.png}

\end{center}
\item Atmungs-/ Kreislaufst\"orungen

\item Hirndruck symptomatik

\item evt. Einklemmung deAtemzentrums ${\rightarrow}$ zentraler
Atemstillsatnd

\item Behinderung der Hirndurchblutung

\item intrakranielle Blutungen

\end{itemize}
\begin{itemize}
\item epidural

\item subdural

\item subarachnoidal

\end{itemize}
\begin{itemize}
\item L\"ahmungen

\item zerebrale Kr\"ampfe

\end{itemize}
Intrakranielle Blutungen

\begin{itemize}
\item \index{Epidurale Blutung}Epidurale Blutung 

\end{itemize}
\begin{itemize}
\item Ri{\ss} einer Meningealarterie

\item Einblutung zwischen Kalotte und Dura mater

\item {\textquotedblleft}lichtes{\textquotedblright} (luzides )
Intervall: einige Stunden

\end{itemize}
\begin{itemize}
\item \index{Subdurale Blutung}Subdurale  Blutung 

\end{itemize}
\begin{itemize}
\item schleichender  (!) Verlauf: oft \"uber Tage

\item auch ohne  Trauma: Hypertoniker, Alkoholiker

\end{itemize}
\begin{itemize}
\item \index{Subarachnoidale Blutung}Subarachnoidale  Blutung
(\index{SAB}SAB)

\end{itemize}
\begin{itemize}
\item z.B. geplatztes Aneurysma, Br\"uckenvenenabri{\ss} bei
Decelerationstrauma

\end{itemize}
\begin{itemize}
\item kann spontan / ohne Trauma auftreten (Aneurysma)

\end{itemize}
\begin{itemize}
\item zuerst Herd symptomatik, schneller  Verfall, 
{\quotedblbase}W\"uhlblutung{\textquotedblleft}

\item Klassisch: pl\"otzlicher, schwerer,
{\quotedblbase}donnerschlagartiger{\textquotedblleft} Kopfschmerz

\end{itemize}
Offenes SHT

Hinweise:

\begin{itemize}
\item gro{\ss}e RQW mit tast-/sichtbaren  Knochensplittern

\item Hirnaustritt (schlechte Prognose)

\item Blutungen aus Nase / Ohren

\item Liquoraustritt, oft mit Blut vermischt 
({\quotedblbase}Tupfertest{\textquotedblleft})

\end{itemize}
SHT -- Allgemeine Diagnostik

Beurteilung von:

\begin{itemize}
\item A tmung / K reislauf  (BAK-Schema!!! )

\item neurologisch:

\end{itemize}
\begin{itemize}
\item B ewu{\ss}tseinslage: klar/getr\"ubt/bewu{\ss}tlos

\item Motorik/Reflexe

\item Pupillenreaktion 

\end{itemize}
\begin{itemize}
\item normal: rund + seitengleich + mittelweit + bds.  direkt und
indirekt lichtreaktiv

\item abnorm: entrundet, starr, weit oder eng nicht  lichtreaktiv

\end{itemize}
\begin{itemize}
\item Traumacheck (sonstige Verletzungen ={\textgreater}  Polytrauma?)

\item Glasgow-Coma-Scale (GCS)

\end{itemize}
Glasgow-Coma-Scale (GCS)

Bewertung von: {\textperiodcentered}~Augen\"offnung,
{\textperiodcentered}~beste verbale Reaktion,
{\textperiodcentered}~beste motorische Reaktion.

\begin{center}
 [Warning: Image ignored] % Unhandled or unsupported graphics:
%\includegraphics[width=13.667cm,height=20.134cm]{AASdev20090228-img90.png}

\end{center}
\paragraph{Bewu{\ss}tseinsst\"orung beim SHT}
\begin{itemize}
\item Sicherung der Vitalfunktionen

\item Atemwege freihalten

\item stabile Seitenlage bei Bewu{\ss}tlosigkeit

\item NOTARZT! ${\leftarrow}$${\rightarrow}$ Intubation und Beatmung ab
GCS {\textless}8 durch NA dringend  erforderlich

\item steriler Verband locker auf offene  Verletzungen

\item Knochen / Hirnteile NICHT  zur\"uckstopfen /  entfernen!

\item oft Kombination mit HWS-Verletzung  ={\textgreater}Stifneck    

\end{itemize}
\begin{itemize}
\item (Gefahr: hohe Querschnittl\"ahmung)

\end{itemize}
\begin{itemize}
\item Oberk\"orper/Kopf 30{\textdegree} erh\"oht lagern (nur wenn
Patient bewu{\ss}tseinsklar)

\item Sch\"adel/HWS/BWS bilden eine Linie

\end{itemize}
\begin{itemize}
\item Jedes SHT ist im Spital abzukl\"aren! (klinisch und  bildgebend)

\end{itemize}
Jedes SHT ist im Spital abzukl\"aren! Es besteht die Gefahr dass sich
der Patient gegenw\"artig in der {\quotedblbase}luciden
Phase{\textquotedblleft} befindet und sich sein Zustand erst in einigen
Stunden, dann aber schlagartig verschlechtert. Im Zweifelsfall wird der
Patient dort zur  Beobachtung (24h lang) aufgenommen.

\subsection{Thoraxtrauma}
\begin{itemize}
\item offenes

\end{itemize}
\begin{itemize}
\item offener Pneumothorax

\end{itemize}
\begin{itemize}
\item geschlossenes

\end{itemize}
\begin{itemize}
\item \index{Spontanpneumothorax}Spontanpneumothorax (Lungenri{\ss},
bull\"oses  Emphysem)

\item Lungenverletzung durch gebrochene Rippen 

\item Bronchusverletzungen

\end{itemize}
Thoraxtrauma -- allgemeine  Symptome

\begin{itemize}
\item Dyspnoe

\end{itemize}
\begin{itemize}
\item Tachypnoe

\item Zyanose

\end{itemize}
\begin{itemize}
\item atemabh\"angige  Schmerzen,  Inspirationsstop

\item ev. Hautemphysem (Luftansammlung in der Haut, typisches
{\quotedblbase}Knistern{\textquotedblleft})

\item Prellmarken (z.B. Lenks\"aule von PKW) 

\end{itemize}
\begin{itemize}
\item ={\textgreater} gr\"undlicher Traumacheck!

\end{itemize}
Thoraxtrauma -- spezifische  Symptome

\begin{itemize}
\item instabiler Thorax

\end{itemize}
\begin{itemize}
\item Sternumfraktur

\item Serienrippenfraktur (zus. paradoxe Atmung --  Brustkorb senkt beim
Einatmen)

\end{itemize}
\begin{itemize}
\item Dyspnoe + Prellmarken +  Bluthusten  (H\"amoptysis )

\end{itemize}
\begin{itemize}
\item Lungenprellung: Gasaustausch [F0EA?] , da  Lungenabschnitte
gesch\"adigt

\end{itemize}
\begin{itemize}
\item Brustschmerz, Schock 

\end{itemize}
\begin{itemize}
\item Rhythmusst\"orungen ${\rightarrow}$ Verletzung des  Herzens

\item schwerer Schock, schneller Verfall ${\rightarrow}$ Riss der  Aorta

\item Blut im Schwall aus dem Mund ${\rightarrow}$ 
Speiser\"ohrenverletzung

\end{itemize}
\index{Pneumothorax}Pneumothorax

\begin{itemize}
\item Lungen- (innerer) oder  Brustwandverletzung (\"au{\ss}erer)

\begin{center}
 [Warning: Image ignored] % Unhandled or unsupported graphics:
%\includegraphics[width=7.086cm,height=10.686cm]{AASdev20090228-img91.png}

\end{center}
\end{itemize}
\begin{itemize}
\item Luft gelangt in den Pleuraspalt

\item Unterdruck aufgehoben

\item Lunge kollabiert auf betroffener Seite

\end{itemize}
\begin{itemize}
\item {\quotedblbase}unkomplizierter{\textquotedblleft} Pneumothorax

\end{itemize}
\begin{itemize}
\item Luft kann ungehindert in Pleuraspalt und  wieder heraus

\item {\quotedblbase}nur{\textquotedblleft} Ausfall des betroffenen
Lungenfl\"ugels

\end{itemize}
\begin{itemize}
\item Ventil-, \index{Spannungspneumothorax}Spannungspneumothorax

\end{itemize}
\begin{itemize}
\item Luft kann zwar hinein, aber nicht wieder  heraus

\item Druck auf betroffener Seite steigt ${\rightarrow}$  Organe (Lunge,
Herz) werden auf die andere Seite gedr\"uckt

\end{itemize}
\begin{itemize}
\item gesunde Seite wird verdr\"angt + Abknickung  gro{\ss}er
Gef\"a{\ss}e!

\end{itemize}


\begin{center}
 [Warning: Image ignored] % Unhandled or unsupported graphics:
%\includegraphics[width=5.811cm,height=4.343cm]{AASdev20090228-img92.png}

\end{center}
Therapie

Pneumothorax, unkompliziert:

\begin{itemize}
\item Lagerung auf verletzte Seite (unverletzte kann  sich besser
entfalten)

\item O2

\item Notarzt, legt ev. vor Ort B\"ulaudrainage

\item (nicht versuchen, Wunde abzudichten ={\textgreater} ev. dadurch
Spannungspneumothorax)

\end{itemize}
Ventilpneumothorax:

\begin{itemize}
\item Entlastung durch Drainage!

\item offener Pneumothorax (Mittelfellpendeln)

\item steril abdecken

\end{itemize}
Stichverletzungen

NIEMALS  HERAUSZIEHEN!!!!!!!!

steril abdecken + fixieren

verletzte Strukturen??, dient ev. als 
{\quotedblbase}St\"opsel{\textquotedblleft}??

Fremdk\"orper soll nicht tiefer rutschen

\subsection[Bauchtrauma]{Bauchtrauma}
\index{Bauchtrauma}\begin{itemize}
\item offen (perforierend),

\end{itemize}
\begin{itemize}
\item  Er\"offnung der  Bauchh\"ohle

\item Stich-, Pf\"ahlungsverletzung

\item Schnittverletzung, {\quotedblbase}Platzbauch{\textquotedblleft}

\end{itemize}
\begin{itemize}
\item geschlossen (stumpf)

\end{itemize}
\begin{itemize}
\item Schlag, Tritt

\item Prellung durch Lenkstange bei PKW-Unfall

\end{itemize}
Offenes Bauchtrauma



\begin{center}
 [Warning: Image ignored] % Unhandled or unsupported graphics:
%\includegraphics[width=6.543cm,height=4.839cm]{AASdev20090228-img93.png}

\end{center}
\begin{itemize}
\item ev. Austritt von Darmschlingen

\end{itemize}
\begin{itemize}
\item NICHT  zur\"uckstopfen ${\rightarrow}$ sonst 
Darmri{\ss}/Darmverschlu{\ss}

\item steril abdecken, mit Ringerl\"osung feuchthalten

\item Bauchdecke entspannen , Knierolle ${\rightarrow}$ sonst 
Darmabklemmung

\end{itemize}
\begin{itemize}
\item Gefahr des Eindringens von Keimen  (Peritonitis/Sepsis!)

\end{itemize}
\begin{itemize}
\item steril  abdecken

\item Vitalparameter

\item O2

\item Notarzt

\end{itemize}
Geschlossenes Bauchtrauma

Gefahren:

\begin{itemize}
\item Gef\"a{\ss}verletzungen

\end{itemize}
\begin{itemize}
\item Aorta, Vena cava,...

\item massive innere Blutungen ${\rightarrow}$ Volumenmangelschock

\end{itemize}
\begin{itemize}
\item Zerrei{\ss}ung  parenchymat\"oser Organe (Milz, Leber)

\end{itemize}
\begin{itemize}
\item massive innere Blutungen ${\rightarrow}$ Volumenmangelschock

\end{itemize}
\begin{itemize}
\item Zerrei{\ss}ung  von Hohlorganen (Magen, Darm,Galleng\"ange )

\end{itemize}
\begin{itemize}
\item Austritt von Speisebrei, Magensaft, Stuhl, Gallensekret

\end{itemize}
\begin{itemize}
\item Bauchfellentz\"undung (Peritonitis) nach Std. bis  Tagen

\item Sepsis (hohe Lethalit\"at), septischer Schock

\end{itemize}
Symptome

\begin{itemize}
\item starke  Bauchschmerzen

\item brettharte  Bauchdecke  (Abwehrspannung, D\'efense )

\item Leitsymptom: Schockzeichen

\item beim Bewu{\ss}tlosen schwierig zu  diagnostizieren

\end{itemize}
\begin{itemize}
\item gr\"undlich nach \index{Prellmarken}Prellmarken suchen!!!

\end{itemize}
Immer nach Prellmarken suchen. Patienten entkleiden!

Bauchtrauma

\index{Milzruptur}Milzruptur

\begin{itemize}
\item prim\"ar: sofort, Kapsel durchgerissen

\item sekund\"ar: Sickerblutung, Kapsel anfangs noch intakt
${\rightarrow}$ Latenz bis Ri{\ss} der Kapsel

\end{itemize}
Cave! {\quotedblbase}\index{Zweizeitige Milzruptur}Zweizeitige
Milzruptur{\textquotedblleft}

\index{Aortenruptur}Aortenruptur

\begin{itemize}
\item bei VU selten, bei Sturz aus gro{\ss}er H\"ohe  h\"aufiger

\item Sehr schlechte Prognose

\end{itemize}
\index{Nierenruptur}Nierenruptur

\begin{itemize}
\item Leitsymptom: blutiger Harn, NICHT bei  Harnleiterabri{\ss}

\end{itemize}
Beckentrauma

\begin{itemize}
\item erhebliche Gewalteinwirkung  (Versch\"uttung, \"Uberrolltrauma)

\item massive  Blutungen (bis 5 l in wenigen  Minuten)

\item Gefahr des Ausblutens (ev. ohne sichtbare  Blutungsquelle!)

\item Mitverletzung von Organen im kleinen  Becken:

\end{itemize}
\begin{itemize}
\item Dickdarm ${\rightarrow}$ Blutung aus Anus

\item Harnsystem ${\rightarrow}$ Blut im Harn / Harn in der 
Bauchh\"ohle

\item (Ausnahme: Spontanruptur der gef\"ullten  Harnblase bei stumpfem
Bauchtrauma)

\item Hoden ${\rightarrow}$ H\"amatome am Skrotum und am  Damm

\end{itemize}
Therapie:

\begin{itemize}
\item Basisma{\ss}nahmen

\item Schockbek\"ampfung

\item Bei Beckenfraktur: \index{Beckengurt}Beckengurt, zur Not mit
G\"urtel%
%


\item O2  8-10 l mindestens

\item Volumenzufuhr i.V. (NA)

\item Schmerzbek\"ampfung

\end{itemize}
Prognose h\"angt vom Ausma{\ss} des  Blutverlustes ab.

\subsection{Wirbels\"aulenverletzungen}
\begin{itemize}
\item Typische Unfallmechanismen:

\end{itemize}
\begin{itemize}
\item Sturz aus gro{\ss}er H\"ohe

\item Motorradunfall (Genickbruch)

\item eingeklemmte Personen bei VU

\item Schlag auf Kopf (Stauchung der HWS)

\end{itemize}
\begin{itemize}
\item Gefahren:

\end{itemize}
\begin{itemize}
\item Kompression/Durchtrennung des  R\"uckenmarks

\item L\"ahmungen (Parese)

\item Sensibilit\"atsst\"orungen (Par\"asthesie)

\end{itemize}
\begin{itemize}
\item inkomplett: eines von beidem

\item komplett: beides

\end{itemize}
traumatische Einblutung ins R\"uckenmark

dadurch bedingtes \"Odem ={\textgreater} RM-Kompression

Wirbels\"aulenverletzung -- Schleudertrauma, Peitschenschlag

bei Auffahrunf\"allen

Zerrung des Halsmarks (bildgebend meist nicht nachweisbar)

Schmerzen im HWS-Bereich

Kopfschmerzen

Schwindel, Tinnitus

ev. neurologische Ausf\"alle (in schweren F\"allen)

Prophylaxe: Nackenst\"utze

Therapie: Immobilisation, Stifneck, (ev. O2 )

Wirbels\"aulenverletzung -  Querschnittssydrom

inkomplett: Motorik oder Sensorik  ausgefallen

komplett: beides ausgefallen

Symptome:

schlaffe L\"ahmung, Ausfall aller Reflexe  unterhalb des betroffenen
Segmentes

Blasen-, Darml\"ahmung ={\textgreater} Inkontinenz,  \"Uberlaufblase

Ausfall der Gef\"a{\ss}-und W\"armeregulation

Wirbels\"aulenverletzungen -- spinaler  Schock

pl\"otzliche spinale Durchtrennung mit  schlagartigem Wegfall aller
zentralen  erregenden Impulse...

Gef\"a{\ss}regulation versagt

Gef\"a{\ss}tonus [F0EA?]

RR [F0EA?]  

Wirbels\"aulenverletzungen -- {\quotedblbase}hoher{\textquotedblleft} 
Querschnitt

L\"asionen auf H\"ohe C1/C2 (1. bzw. 2.  HWK)

Abtrennung des Atemzentrums in der Medulla  oblongata

Anpralltrauma (Motorrad unter Lastwagen)

{\quotedblbase}K\"opfler{\textquotedblleft} in seichtes Gew\"asser

Erh\"angen

zentrale Ateml\"ahmung (keine  Spontanatmung!)

L\"asionen oberhalb C4 (4. HWK)

Segmente, die Zwerchfell und Kehlkopf  innervieren mitbetroffen

motorische Ateml\"ahmung ={\textgreater} auf Dauer  beatmungspflichtig

Wirbels\"aulentrauma -  Vorgehensweise

bei Verdacht  auf WS-Trauma  noch vor Ort  (bevor  der Patient bewegt 
wird) \"Uberpr\"ufung von Sensibilit\"at  und Motorik (und
Durchblutung) der Extremit\"aten !!!!!

(Merkabk\"urzung {\textquotedblleft}DMS {\textquotedblright}, Teil jedes
 Traumachecks)

Wirbels\"aulenverletzungen - Therapie

Strikte Immobilisierung (Vacuum-Matratze, Schaufeltrage, Stifneck)

Vacuummatratze immer gemeinsam mit  Stifneck

Bewegungen nur in der WS-Achse (keine  Beugung, Drehung,
Seitw\"artsneigung)

falls n\"otig, korrekte Helmabnahme (zu  zweit)

vorsichtiger L\"angszug (Stifneck anlegen)

Achtung beim Rautek-Griff: keine HWS- Stabilisierung!!!

Wirbels\"aulenverletzung -  Schocklagerung?

ja, spezielle Technik:

Patient auf Vacuummatratze + Trage

Trage schr\"ag: Kopf tief, Beine hoch

nein, wenn SHT oder kardiale  Komplikationen!!!

Anmerkungen zur Vacuummatratze

Immer mit Schaufeltrage

Immer mit Stifneck

Umbettungen durch mehrere Personen

\subsection{Polytrauma}
Gleichzeitig entstandene Verletzungen mehrerer K\"orperregionen
/Organsystemen, von  denen wenigstens eine allein oder die Kombination
mehrerer lebensgef\"ahrlich  ist.

\begin{itemize}
\item Bauchtrauma + Oberschenkel; SHT + WS + Thoraxtrauma etc.

\item gebrochener Finger + gebrochene Zehe ${\rightarrow}$ kein
Polytrauma

\end{itemize}
Problem:  Volumenmangelschock +  Ateminsuffizienz

Polytrauma - Vorgehensweise

\begin{itemize}
\item Schneller Traumacheck gleichzeitig  mit  Sichern der
Vitalparameter!

\item Bewu{\ss}tsein

\item Atmung

\item Kreislauf: Carotispuls, Pulse peripher  (orientierend!),
Kapillardurchblutung

\item gr\"undlicher Traumacheck erst sp\"ater

\item stabilisierende Ma{\ss}nahmen

\item Atemwege sichern (Intubation),  ven\"oser  Zugang (Notarzt!)

\item wenn einigerma{\ss}en stabil, gr\"undlicher  Traumacheck,
Verletzungsmuster (={\textgreater} zu  erwartende Komplikationen)

\end{itemize}
Polytrauma beim Kind

Schockindex unbrauchbar!!!

Kinder bleiben lange stabil, verfallen  dann aber schlagartig!!!

doppelte Vorsicht (lieber eine  Notarztalarmierung zu viel, als eine zu 
wenig!)

nach Unfallhergang : Patient gilt als polytraumatisiert  bis zum Beweis 
des Gegenteils! 

Polytrauma - Unfallort

Absichern  der Unfallstelle,  m\"oglichst auff\"allig machen (Blaulicht,
 Warnblinkanlage, Scheinwerfer)

bei mehreren Verletzten: Triage

Nachalarmierung : andere  Rettungsmittel, Feuerwehr, Polizei

Eingeklemmte im Fahrzeug stabilisieren,  nicht eigenm\"achtig bergen

Crash -Rettung nur  bei  Fremdgef\"ahrdung, HWS stabilisieren!

Polytrauma - Therapie

\begin{itemize}
\item Freimachen der Atemwege, 15l O2 , ggf.  Intubation

\item Kreislaufkontrolle, 2 gro{\ss}lumige Zug\"ange,  Volumen!

\item Blutstillung  (Druckverband, Abbindung wenn  indiziert)

\item Immobilisierung : Stifneck, Vacuummatratze,  Schaufeltrage

\item W\"armeschutz

\item BZ (Alkoholisierte)

\item Medikamente: 

\item Volumssubstitute (kristalloid und kolloidal)

\item Schmerzmittel

\end{itemize}
\subsection{Extremit\"atenverletzungen}
Gelenke betreffend

Verstauchung (Distorsion )

Verrenkung (Luxation )

Knochenbruch (Fractur )

\paragraph{Extremit\"atenverletzung - Verstauchung}
Zerrung / \"Uberdehnung des  Kapsel-/Bandapparates

Symptome

Schmerzen

Schwellung

Blutergu{\ss} (H\"amatom )

Therapie

ruhigstellen

hochlagern

k\"uhlen

ad Unfallabteilung (mu{\ss} nicht sofort sein)

\paragraph{Extremit\"atenverletzungen - Verrenkung}
partielle oder vollst\"andige Verrenkung  eines Gelenks

Symptome

Schmerzen

Schwellung

Blutergu{\ss}

abnorme Gelenkstellung (federnd fixiert), bewegungsunf\"ahig

Therapie

Stellung belassen

ruhigstellen

KEINE EINRENKVERSUCHE!!!

\paragraph[Extremit\"atenverletzungen -
\ Knochenbr\"uche]{Extremit\"atenverletzungen -  Knochenbr\"uche}
Arten:

\begin{itemize}
\item geschlossene Fraktur

\item offene Fraktur

\end{itemize}
jeweils

\begin{itemize}
\item komplett

\item inkomplett

\end{itemize}
Gefahren

\begin{itemize}
\item starke Blutungen

\item ev. bis Schock

\item begleitende Nervenverletzungen

\item Fettembolie (Fraktur langer R\"ohrenknochen)

\item Infektionsgefahr (Osteomyelitis bei offenen  Frakturen)

\end{itemize}
Immer gr\"undlichen Traumacheck  durchf\"uhren, um keine 
Begleitverletzungen zu \"ubersehen!!!

Extremit\"atenverletzungen -  Knochenbr\"uche

Frakturzeichen

\begin{center}
\begin{tabular}{|m{6.649cm}|m{6.651cm}|}
\hline
\multicolumn{2}{|m{13.5cm}|}{\centering
\index{Frakturzeichen}Frakturzeichen\par

}\\\hline
sichere

 &
unsichere

\\\hline
\begin{itemize}
\item abnorme Beweglichkeit
({\quotedblbase}Pseudogelenk{\textquotedblleft})

\item Fehlstellung

\item Krepitus (bitte NICHT absichtlich  ausprobieren!!)

\item Sichtbare freie Knochenteile (= offener Bruch)

\end{itemize}
 &
\begin{itemize}
\item Schmerz

\item Schwellung

\item H\"amatom

\item gest\"orte Funktion / Bewegungsunf\"ahigkeit

\end{itemize}
\\\hline
\end{tabular}
\end{center}
Therapie:

Traumacheck ${\rightarrow}$ je nach Bedarf

Blutstillung

Wundversorgung

Stifneck bei V.a. HWS-Verletzung

ev. NAW-Nachalarmierung (starke Schmerzen, gr\"o{\ss}erer Blutverlust,
Schock, relevanten Begleit\-verletzungen)

Ruhigstellung / Schienung (unter Zug)

\subsection{Verbrennungen}
Die Gr\"o{\ss}e des gesch\"adigten Areals wird in Prozent der
ver\-brannten K\"orperoberfl\"ache angegeben.

Faustregel - Neunerregel

Handfl\"ache (des Patienten) = 1~\% der K\"orperoberfl\"ache (KOF)

\begin{center}
 [Warning: Image ignored] % Unhandled or unsupported graphics:
%\includegraphics[width=6.591cm,height=7.03cm]{AASdev20090228-img94.png}

\end{center}
\begin{flushright}
\tablehead{}\begin{supertabular}{|m{2.134cm}|m{2.134cm}|m{2.134cm}|}
\hline
 &
\centering \%\par

 &
\\\hline
Kopf:

 &
\centering 9\par

 &
\\\hline
Arm, je:

 &
\centering 9\par

 &
\\\hline
Thorax vorne:

 &
\centering 18\par

 &
\\\hline
Thorax hinten:

 &
\centering 18\par

 &
\\\hline
Bein, je:

 &
\centering 18\par

 &
\\\hline
Genital: 

 &
\centering 1\par

 &
\\\hline
 &
 &
\\\hline
\end{supertabular}
\end{flushright}
Grade

Das Ausma{\ss} der Sch\"adigung durch Verbrennung wird in 3 Grade
eingeteilt:

\begin{flushleft}
\begin{tabular}{|m{1.836cm}|m{6.8980002cm}|m{4.3630004cm}|}
\hline
\centering Grad\par

 &
 &
\\\hline
\centering 1. Grad\par

 &
R\"otung (Juckreiz, Schmerz)

 &
\\\hline
\centering 2. Grad\par

 &
Blasenbildung, Schmerz

2a: oberfl\"achlich, Nerven intakt, 

2b: tiefer, Nerven zerst\"ort

 &
\\\hline
\centering 3. Grad\par

 &
trockene Hautfetzen (wei{\ss}-gelb-schwarz), Verkohlung, Schrumpf\-ung
des Gewebes durch Hitze\-ein\-wirkung

 &
\\\hline
\end{tabular}
\end{flushleft}
Ab 5\% verbrannter (mind. 2. Grad) KOF bei Kindern oder 10\% bei
Erwachsenen besteht Schockgefahr (Fl\"ussigkeitsverlust)

\begin{itemize}
\item immer  auch an Begleitverletzungen  denken

\item bestimmte Unfallmechanismen

\item Traumacheck !

\end{itemize}
Ma{\ss}nahmen

\begin{itemize}
\item Eigenschutz!

\item Kleiderbrand l\"oschen

\item Kleidung rasch entfernen (sofern nicht  eingeschmolzen!)

\item mind. 10 min. mit m\"a{\ss}ig kaltem Wasser  sp\"ulen -- NICHT
unterk\"uhlen!

\item keimfreier, nicht verklebender Verband, locker anlegen 

\item wenn vorhanden: Water-Jel%
%
\footnote{Streitfrage: Water-Jel: Umstritten} 

\item Schockbek\"ampfung

\item NAW wenn n\"otig

\end{itemize}
Was man besser nicht tun sollte...

\begin{itemize}
\item NEIN: Salbe, Puder, etc. auf frische  Brandwunde

\item NEIN: Blasen \"offnen

\item NEIN: festklebende Kleidung gewaltsam  entfernen

\item NEIN: aluminiumbeschichteten Verband  verwenden (Gefahr der
Hitzereflexion bei  nicht ausreichend gek\"uhlten  Verbrennungen!)

\end{itemize}
%
%
Inhalationstrauma

Leitsymptome

\begin{itemize}
\item angebrannte Gesichtshaare

\item geschwollene, ru{\ss}ige Lippen

\item Ru{\ss}partikel im Sputum

\item Heiserkeit

\item Stridor

\end{itemize}
${\rightarrow}$ Rauchgasvergiftung m\"oglich!

\subsection{Erfrierungen}
...sind Gewebsnekrosen, die durch lokale  Durchblutungsst\"orungen
aufgrund von  K\"alteeinwirkung (Eis, Wind, N\"asse,  K\"alte)
entstanden sind.

Besonders geff\"ahrdet sind k\"orperferne K\"orperteile wie Zehen,
Finger, Ohren oder Nase.

Einteilung erfolgt in 4 Stadien, diese haben am Notfallort keine
entscheidende Bedeutung f\"ur die zu setzenden Ma{\ss}nahmen. Den
Schweregrad einer Erfrierung kann man erst nach Tagen voll
einsch\"atzen. Deshalb ist es wichtig, die Warnsymptome ernst zu
nehmen!

Stadien

Symptome

zuerst 

Gef\"uhllosigkeit und Bl\"asse

starke Schmerzen und bl\"auliche Verf\"arbung

 einige Zeit sp\"ater

marmorierte Verf\"arbung der Haut mit  Blasenbildung

Taubheitsgef\"uhl, keine Empfindung bei  Ber\"uhrung

Bewegungsunf\"ahigkeit des betroffenen  K\"orperteils

starke Schmerzen

Ma{\ss}nahmen

\begin{itemize}
\item beengende Kleidungsst\"ucke lockern (Einschn\"urung erh\"oht
Risiko einer Erfrierung)

\item Keimfreier, gepolsterter Verband

\item Warme, gezuckerte Getr\"anke -- KEIN Alkohol

\item K\"orperstamm w\"armen, nie das erfrorene K\"orperteil abreiben!!

\end{itemize}
\subsection{Unf\"alle durch Strom- und Blitzschlag}
Niederspannung: bis 1000 Volt (Steckdose:  230 Volt)

Hochspannung: \"uber 1000 Volt, Stromst\"arke {\textgreater}  3 Ampere
(\"Uberlandleitungen) 

Effekte / Verletzungsmuster je nach  Spannung (Volt) / Stromst\"arke
(Ampere)

\begin{itemize}
\item Reizeffekte (Muskel-/Nervenerregung):

\end{itemize}
\begin{itemize}
\item Verkrampfung ({\quotedblbase}kann nicht
loslassen{\textquotedblleft})

\item Sensibilit\"atsst\"orungen

\item ZNS ${\rightarrow}$ Bewu{\ss}tseinsst\"orung

\item Herzrhythmusst\"orungen (Kammerflimmern, ev.  als Sp\"atfolge!)

\end{itemize}
\begin{itemize}
\item W\"armeproduktion:

\end{itemize}
\begin{itemize}
\item Strommarken (Ein-/Austrittsverbrennung)

\item Lichtbogen: oberfl\"achliche und tiefe  Verbrennungen

\end{itemize}
\begin{itemize}
\item Gleichstrom

\item Wechselstrom ${\rightarrow}$ eher  Herzrhythmusst\"orungen

\item Patient  kann Spannungsquelle ev. nicht  loslassen (Krampf der
Handmuskulatur)

\item Auch Defi stellt Strom-Gefahrenquelle  dar!!!

\item Blitzunfall

\end{itemize}
\begin{itemize}
\item mehrere Millionen (!) Volt

\item bis 100.000 Ampere

\end{itemize}
Symptome:

\begin{itemize}
\item Bewu{\ss}tlosigkeit

\item L\"ahmungen (vor\"ubergehend)

\item Blitzfiguren auf der Haut

\item Herzrhythmusst\"orungen m\"oglich

\end{itemize}
Cave: Das eigentliche Ausma{\ss} der Verletzung ist oft
{\quotedblbase}unsichtbar{\textquotedblleft}, das der Strom im
K\"orperinneren -- also {\quotedblbase}unterirdisch{\textquotedblleft}
- viel Schaden anrichtet. 

Ma{\ss}nahmen:

\begin{itemize}
\item Eigenschutz! 

\end{itemize}
\begin{itemize}
\item (Stromkreis  vorher  unterbrechen, FI bzw. bei Starkstrom durch 
Fachmann)

\end{itemize}
\begin{itemize}
\item Genaue  Patientenuntersuchung ${\rightarrow}$ Eintritts{}-  und
Austrittsmarke

\item auf ev. Herzrhythmusst\"orungen vorbereitet  sein ${\rightarrow}$
bis 24h Latenz

\item Vitalparameter ${\rightarrow}$ entsprechende  symptomatische
Versorgung

\item NAW ${\rightarrow}$ Hospitalisierung

\item Stromverbrennungen wie Verbrennungen  versorgen

\end{itemize}
\section[[RS II{]} Hygiene ]{[RS II] Hygiene }
\label{bkm:RefHeading32633609}Sebastian Gabriel [*]

\begin{center}
 [Warning: Image ignored] % Unhandled or unsupported graphics:
%\includegraphics[width=4.763cm,height=4.763cm]{AASdev20090228-img95.jpg}

\end{center}
Changelog:

2009-02-20: Fertig f\"ur Kurs 2009/02.

2009-02-25: Start Changelog.

\subsection{Grundlagen}
\begin{multicols}{2}
\index{Infektion}Infektion: Eindringen von krankmachenden Keimen in den
Organismus

\index{Pathogenit\"at}Pathogenit\"at: F\"ahigkeit eines Mikroorganismus,
eine Krankheit auszul\"osen

\index{Resistenz}Resistenz: Widerstandsf\"ahigkeit eines Organismus
gegen die  krankmachenden  Eigenschaften eines Erregers

\index{Inkubationszeit}Inkubationszeit: symptomlose Zeit zwischen
Infektion (Eintritt d. Erregers in den Organismus) und dem Auftreten
erster Symptome

\index{Prodromalstadium}Prodromalstadium  uncharakteristische Symptome
am Ende der Inkubationszeit (Fieber, Kopfschmerz, Erbrechen,
Gliederschmerzen)

\index{Epidemie}Epidemie: \"ortlich begrenztes, zeitlich vermehrtes
Auftreten einer  Infektionserkrankung (Grippe, Salmonellen)

\index{Pandemie}Pandemie: Epidemie die sich \"uber L\"ander u.
Kontinente ausbreitet (Pest im Mittelalter)

\index{Endemie}Endemie: Auf kleinere Gebiete beschr\"anktes, zeitlich
jedoch nicht begrenztes Auftreten einer Infektionserkrankung  (FSME,
Malaria, Lepra)

\index{Morbidit\"at}Morbidit\"at: Anteil der an XY Erkrankten einer
bestimmten Bev\"olkerungsgruppe (z.B. 5 pro 100.000 EW)

\index{Mortalit\"at}Mortalit\"at: Anteil der an XY Verstorbenen einer
bestimmten Bev\"olkerungsgruppe (z.B. 5 pro 100.000 EW)

\index{Letalit\"at}Letalit\"at:  Anteil der an XY Verstorbenen beozgen
auf die Gesamtzahl der Erkrankten ({\quotedblbase}1 von 5 Erkranten
stirbt{\textquotedblright})

\end{multicols}
\index{Entz\"undung}Entz\"undung

Reaktion des K\"orpers

\begin{enumerate}
\item R\"otung 

\item Schwellung 

\item \"Uberw\"armung 

\item Schmerz 

\item Funktionsverlust 

\end{enumerate}
Ausl\"oser:

\begin{itemize}
\item Fremdk\"orper

\item Krankheitserreger

\item Mechanische Beanspruchung

\item ...

\end{itemize}
Anh\"angsel
{\textquotedblleft}\index{{}-itis}{}-itis{\textquotedblright}:
Gastritis ${\rightarrow}$
{\quotedblbase}Magenentz\"undung{\textquotedblleft}, Hepatitis
${\rightarrow}$ Leberentz\"undung, Pankreatitis ${\rightarrow}$
Entz\"undung der Bauchspeicheldr\"use, Meningitis ${\rightarrow}$
Hirnhautentz\"undung, Rhinitis ${\rightarrow}$
{\quotedblbase}Schnupfen{\textquotedblleft}, Otitis ${\rightarrow}$
Ohrenentz\"undung, ...

\paragraph[Die wichtigsten \"Ubertragungswege]{Die wichtigsten
\"Ubertragungswege}
\begin{itemize}
\item f\"akal-oral (Schmierinfektion, zB. Badesee)

\item Aerogen (Tr\"opfcheninfektion)

\item Genital (Geschlechtsverkehr)

\item \"Uber/durch die Haut

\item Diaplazent\"ar (w\"ahrend Schwangerschaft)

\item W\"ahrend der Geburt

\item Eintrittspforten: alles was irgendwie offen ist

\end{itemize}
\begin{itemize}
\item Wunden, Magen-Darm-Trakt, Atmung, Auge, Penis/Scheide, Plazenta
...

\end{itemize}
\begin{itemize}
\item Fremdk\"orper sind Wohnh\"auser und Autobahnen f\"ur Keime

\end{itemize}
\begin{itemize}
\item Harnkatheter, ven\"ose Zug\"ange, Tubus, k\"unstlicher
Mageneingang, Dialysekatheter {\dots}

\end{itemize}
\paragraph[]{}
\paragraph[Kreuzinfektion]{Kreuzinfektion}
\begin{center}
\begin{minipage}{6.477cm}
${\leftarrow}$ DAVOR f\"urchten wir uns !!!

\end{minipage}
\end{center}
\index{Kreuzinfektion}Kreuzinfektion: Patient 1 ${\rightarrow}$
Sanit\"ater ${\rightarrow}$ Patient 2

Die wichtigsten Erreger

\begin{itemize}
\item Bakterien

\end{itemize}
\begin{itemize}
\item klein, im Mikroskop meist sichtbar;
{\textquotedblleft}Normale{\textquotedblright} Zellen

\end{itemize}
\begin{itemize}
\item Viren

\end{itemize}
\begin{itemize}
\item sicher klein, im Mikroskop i.d.R. nicht sichtbar, keine echten
Zellen, m\"ussen Zellen befallen

\end{itemize}
\begin{itemize}
\item Pilze

\item Tiere, Tierparasiten

\end{itemize}
\begin{itemize}
\item W\"urmer

\end{itemize}
\begin{itemize}
\item (Toxine)

\end{itemize}
\paragraph[Bakterien]{Bakterien}
\index{Bakterien}\begin{itemize}
\item {\quotedblbase}Echte{\textquotedblleft} Zellen, die leben und
aktiv etwas tun

\end{itemize}
\begin{itemize}
\item zB. Produktion von Substanzen (Toxine ...), Gasen, ...

\end{itemize}
\begin{itemize}
\item Viele Bakterien k\"onnen f\"ur den Menschen gut, egal oder auch
b\"ose sein

\item Haben die unterschiedlichsten Formen

\end{itemize}
\begin{itemize}
\item {\textquotedblleft}Kugeln{\textquotedblright}, St\"abchen, mit
Gei{\ss}eln ... 

\end{itemize}
Bakterien und Antibiotika

\begin{itemize}
\item Antibiotika t\"oten prinzipiell Bakterien, aber:

\end{itemize}
\begin{itemize}
\item Antibiotika haben ein bestimmtes Spektrum, dh. wirken nur auf
bestimmte Keime

\item Bakterien entwickeln Resistenzen gegen Antibiotika

\item Antibiotika helfen dann nicht mehr

\end{itemize}
Resistenzen sind ein enorm gro{\ss}es Problem in der heutigen Medizin!

... und damit auch im Rettungs- und Krankentransport.

{\textquotedblleft}\index{Spitalskeime}Spitalskeime{\textquotedblright}:
Keime mit {\textquotedblleft}angez\"uchteten{\textquotedblright}
Resistenzen

Unterscheide:
{\textquotedblleft}Wald-und-Wiesen{\textquotedblright}-Keime und
Spitalskeime 

${\rightarrow}$ Keim ist nicht gleich Keim, auch wenn sie den gleichen
Namen haben

Bakterien und Pathogenit\"at

\begin{flushleft}
\tablehead{}\begin{supertabular}{|m{7.605cm}|m{5.688cm}|}
\hline
\begin{itemize}
\item Die B\"osartigkeit von Bakterien h\"angt u.a. ab von:

\end{itemize}
\begin{itemize}
\item Welches Bakterium

\item Spezies, Stamm, Resistenzen

\item Wo ist die Infektion

\item Wie besiedelt ist die Umgebung

\item Wie ist der Immunstatus des Patienten

\item Kind, Chemotherapie, AIDS, ...

\item K\"orperliche Verfassung

\item Alter, Ern\"ahrung, Stress, ...

\end{itemize}
\begin{itemize}
\item Manche Keime machen immer krank

\item Manche Keime nur unter Umst\"anden

\item Opportunisten nur wenn der Wirt geschw\"acht ist, und das sind
viele Patienten im Rettungs- und Krankentransport!

\end{itemize}
 &
\centering   [Warning: Image ignored]
% Unhandled or unsupported graphics:
%\includegraphics[width=2.858cm,height=5.001cm]{AASdev20090228-img96.jpg}
 \par

\centering F\"ur manche Bakterien zahlt man freiwillig viel Geld, hier
am Beispiel des Lacotbacillus casei gezeigt. \par

\\\hline
\end{supertabular}
\end{flushleft}
\paragraph[Viren]{Viren}
\index{Viren}\begin{itemize}
\item Keine Zellen, sehr klein, im Lichtmikroskop nicht sichtbar

\item Brauchen Wirtszelle um sich zu vermehren

\item Schleusen eigenes genetisches Programm in die Wirtszelle ein

\item Nachweis:

\end{itemize}
\begin{itemize}
\item NICHT anz\"uchtbar auf N\"ahrb\"oden

\item Antigen- oder Antik\"orpernachweis

\item PCR (DNA/RNA-Nachweis)

\end{itemize}
\begin{itemize}
\item zB: viele Verk\"uhlungen, Grippe, Feuchtblattern, viele
Hepatitis-Formen ...

\end{itemize}
\subsection{Wichtige Krankheitsbilder}
\index{Grippe}Grippe (\index{Influenza}\index{Influenza}Influenza)

Die Grippe ist eine virale Erkrankung, welche je nach Patient auch einen
sehr schweren Verlauf nehmen kann und Komplikationen wie Lungen- oder 
Herzmuskelentz\"undungen oder weitere Infektionen nach sich ziehen
kann. Besonders gef\"ahrdet sind \"altere und immunsupprimierte
Personen. Gegen die Grippe wird eine Schutzimpfung angeboten. Diese
wirkt f\"ur ein Jahr und muss jede Saison widerholt werden. Die Impfung
bietet keinen hundertprozentigen Schutz, wird aber trotzdem f\"ur
Risikopatienten und Personen die einem erh\"oten Infektionsrisiko
ausgesetzt sind (darunter f\"allt auch Personal im Gesundheitsbereich)
empfohlen. 

Da es sich bei der Grippe um eine durch Viren verursachte Krankheit
handelt sind Antibiotika wirkungslos.

Die Grippe tritt in Wellen auf und kann t\"odlich enden

Die Grippe tritt besonders in der kalten Jahreszeit auf und tritt
regelm\"a{\ss}ig als \index{Epidemie}Epidemie auf. Es kann auch zum
Ausburch einer weltweiten \index{Pandemie}Pandemie kommen, z.B.
Spanische Grippe (1918, 25~Mio~Tote), Asiatische Grippe (1957,
1~Mio~Tote) und  Hongkong-Grippe (1968,
800.000~Tote)\footnote{http://de.wikipedia.org/wiki/Influenza}. 

In \"Osterreich erkranken pro Jahr ca. 380.000 Personen an der
Influenza, ca. 4.500 Personen m\"ussen station\"ar behandelt werden und
ca. 120 Personen sterben an den
Folgen\footnote{http://www.bmg.gv.at/cms/site/standard.html?channel=CH0742\&doc=CMS1075899626901}.

Symptome

Typisch f\"ur die Grippe ist ist das rasche, hohe anfiebern
(38{\textdegree}-40{\textdegree}C) mit Sch\"uttelfrost,  Kopf-, Hals-,
Muskel-, Gelenksschmerzen, Husten und stark reduziertem
Allgemeinszustand sowie langdauernde Abgeschlagenheit nach Ende der
Erkrankung. Bei schon vorher gschw\"achten Patienten kann der Verlauf
weniger spektakul\"ar ausshehen (kaum Fieber), das Risiko einer
Komplikationen ist allerdings nat\"urlich h\"oher. 

\index{grippaler Infekt}grippaler Infekt,
\index{Verk\"uhlung}Verk\"uhlung, \index{Erk\"altung}Erk\"altung

Die Verk\"uhlung ist eine Infektion der oberen Luftwege mit Viren,
allerdings mit anderen als bei der {\quotedblbase}echten
Grippe{\textquotedblleft} (Influenza). Die Symptome \"ahneln denen der
Grippe, allerdings sind sie bei weitem nicht so stark ausgepr\"agt. Im
Vordergrund stehen Halsschmerzen, Husten, rinnende oder verstopfte Nase
und eine eher leicht erh\"ohte Temperatur. 

\index{Tetanus}Tetanus (\index{Wundstarrkrampf}Wundstarrkrampf) [!!!]

Clostridium tetani

\begin{center}
 [Warning: Image ignored] % Unhandled or unsupported graphics:
%\includegraphics[width=5.984cm,height=3.539cm]{AASdev20090228-img97.gif}

\end{center}
\begin{itemize}
\item Der Erreger kommt \"uberall in der Umgebung vor

\item Produziert Toxin durch welches die Skelettmuskulatur tonisch
krampft

\item Inkubation: 4-21 Tage

\item Vorzeichen: M\"udigkeit, Inappetenz

\item Klinik: Tonische Kr\"ampfe der quergestreiften Muskulatur, oft
beginnend mit Kaumuskulatur

\item Problem: Schmerzen, Ateml\"ahmung

\item Letalit\"at ca. 50\%

\item Impfung!

\end{itemize}
Tetanus \index{Impfschutz}Impfschutz

\begin{itemize}
\item Auch bei kleinen Verletzungen nach dem Impfschutz fragen und
dokumentieren!

\item Patient soll unverz\"uglich im Impfpass nachschauen.

\item Wenn die letzte Impfung l\"anger als 5 Jahre zur\"uck liegt oder
der Patient noch nie geimpft wurde muss der Patient unverz\"uglich eine
Unfallabetilung aufsuchen. 

\item Am Revers best\"atigen lassen!

\end{itemize}
Eine mangelhafte Aufkl\"arung oder Dokumentation bez\"uglich Tetanus und
dessen Impfschutz ist grob fahrl\"assig und ein unverzeihbarer
Kunstfehler.

Tuberkulose 

Mykobacterium Tuberculosis, {\quotedblbase}Tbc{\textquotedblleft}

\begin{center}
 [Warning: Image ignored] % Unhandled or unsupported graphics:
%\includegraphics[width=4.615cm,height=6.153cm]{AASdev20090228-img98.jpg}

\end{center}
\begin{itemize}
\item Kann praktisch alle Organe befallen 

\item Der Verlauf wird vom Organbefall bestimmt.

\end{itemize}
\begin{itemize}
\item Lunge, Skelett, Darm, Haut, Genitalien, ...

\end{itemize}
\begin{itemize}
\item Ausgangspunkt ist fast immer die Lunge.

\item \"Ubertragung: aerogen, oral, Hautkontakt, plazent\"ar

\end{itemize}
\begin{itemize}
\item Abh\"angig von der Herdlokalisation

\end{itemize}
\begin{itemize}
\item Inkubationszeit: 4-12 Wochen

\item Klinik:

\end{itemize}
\begin{itemize}
\item Subfebrile Temperatur, Nachtschwei{\ss}, Husten, evt. Bluthusten,
Lymphknotenschwellung,  Pleuritis

\end{itemize}
\begin{itemize}
\item {\textquotedblleft}\index{Offene Tb}Offene Tb{\textquotedblright}:
Infektionsherd hat {\textquotedblleft}Luftkontakt{\textquotedblright}
${\rightarrow}$ Tr\"opfcheninfektion m\"oglich ${\rightarrow}$
Schutzmaske

\item Meldung an Leitstelle \& Betriebsleitung!

\end{itemize}
\index{HIV}HIV - \index{Humanes Immundefizienz Virus}Humanes
Immundefizienz Virus

\begin{itemize}
\item Bef\"allt vor allem die Zellen des Abwehrsystems 

\item Es vermehrt sich in ihnen, setzt sie au{\ss}er Funktion und
zerst\"ort sie schlie{\ss}lich.

\item Das k\"orpereigene Abwehrsystem kann HIV nicht aus dem K\"orper
entfernen.

\end{itemize}
\begin{itemize}
\item ${\rightarrow}$Immunschw\"ache

\end{itemize}
\begin{itemize}
\item HIV verursacht AIDS (Aquired Immune Deficiency Syndrom)

\item Begleiterkrankungen: 

\end{itemize}
\begin{itemize}
\item Infektionen mit normalerweise harmlosen Erregern

\end{itemize}
\begin{itemize}
\item Pilze, {\textquotedblleft}fakultative{\textquotedblright} Erreger,
{\textquotedblleft}Opportunisten{\textquotedblright}, seltene Arten der
Lungenentz\"undung

\end{itemize}
\begin{itemize}
\item Tumore

\end{itemize}
\begin{itemize}
\item \"Ubertragung

\end{itemize}
\begin{itemize}
\item K\"orperfl\"ussigkeiten

\item Blut, Blutprodukte, Sperma, Scheidensekret

\item Ungesch\"utzter Sex, Transfusionen

\item Nadelstichverletzungen

\item KEINE \"Ubertragung durch Sozialkontakte (H\"andesch\"utteln,
Haushalt, Geschirr, Toilettenanlagen)

\end{itemize}
\begin{itemize}
\item Erst nach 6-12 Wochen sind Antik\"orper nachweisbar

\end{itemize}
\begin{itemize}
\item HIV-Test nicht {\textquotedblleft}aktuell{\textquotedblright},
{\textquotedblleft}12 Wochen blind{\textquotedblright}

\end{itemize}
\begin{itemize}
\item Nach ca 5-10 Jahren bricht AIDS aus

\item Therapie:

\end{itemize}
\begin{itemize}
\item Keine Heilung m\"oglich, aber deutliche Lebensverl\"angerung

\end{itemize}
\begin{itemize}
\item Prophylaxe

\end{itemize}
\begin{itemize}
\item Postexpositionsprohylaxe bei Nadelstichverletzungen

\item Kondome

\end{itemize}
(Virale) \index{Hepatitis}Hepatitis - versch. Erreger (A-E)

Impfung nur f\"ur Hepatitis A und B verf\"ugbar
(Twinrix{\texttrademark})

Achtung: Non-Responder m\"oglich!

Virale Hepatitis - Klinik und Verlauf

\begin{itemize}
\item Grunds\"atzlich
{\textquotedblleft}Leberentz\"undung{\textquotedblright}

\item Vergr\"o{\ss}erte, schmerzhafte Leber 

\item Ikterus (Gelbsucht) 

\item Dunkler Urin (heller Sch\"uttelschaum),

\item evt. heller Stuhl (Acholie)

\item Unterscheide {\textquotedblleft}akut{\textquotedblright} und
{\textquotedblleft}chronisch{\textquotedblright}

\end{itemize}
\begin{itemize}
\item Chronisch: wenn Hepatitis nicht abheilt

\end{itemize}
\begin{itemize}
\item Ende: Leberversagen, Tod

\end{itemize}
Hepatitis-A - Hepatitis-A-Virus (HAV)

\begin{itemize}
\item Inkubationszeit 10 - 40 Tage, 

\item typische Reisekrankheit (Mittelmeerraum, S\"udamerika, Orient),
meist epidemisch, 

\item Dauer einige Wochen, 

\item kein \"Ubergang in chronischen Verlauf

\end{itemize}
Hepatitis-B - Hepatitis-B-Virus (HBV)

\begin{itemize}
\item Inkubationszeit 4 - 12 Wochen, 

\item meist sporadisches Auftreten, 

\item \"Ubertragung durch Speichel, Blut(produkte), Sperma und andere
K\"orperfl\"ussigkeiten, 

\item chronischer Verlauf in ca. 5 \%

\end{itemize}
Hepatitis C - Hepatitis-C-Virus (HCV)

\begin{itemize}
\item Inkubationszeit 6 - 12 Wochen, 

\item h\"aufig nach Bluttransfusion oder i.v. Drogenabusus, 

\item chronischer Verlauf (in 70 -- 80~\% der F\"alle)

\end{itemize}
F\"ur den Rettungsdienst sind besonders Hep B und C wichtig: Hohe
Infektiosit\"at, chronischer Verlauf!

Gegen Hep C gibt es keine Impfung!

\label{bkm:RefHeading41252510}Sepsis

Siehe auch: Septischer Schock und Toxischer Schock
[H.V{\guillemotright}\pageref{bkm:RefHeading41284910}, oben]

\begin{itemize}
\item Wenn Keime und deren Produkt ausser Kontrolle geraten erfolgt
Amoklauf des K\"orpers:

\item Ungehemmte Freisetzung von Stoffen des Entz\"undungs-, Gerinnungs-
und Immunsystems

\item Auch ohne Infektion m\"oglich:

\end{itemize}
\begin{itemize}
\item Systemic Inflammatory Response Syndrome (SIRS) 

\end{itemize}
\begin{itemize}
\item nach schwerem Trauma, Verbrennungen, Gewebshypoxie

\item Klinik abh\"angig vom Fortschritt und k\"orperlichen Abbau.

\end{itemize}
\begin{itemize}
\item Sehr stark reduzierter Allgemeinzustand

\item Fieber ({\textgreater}38{\textdegree} C) oder Hypothermie
({\textless}36{\textdegree} C)

\item Herzfrequenz {\textgreater}90/min

\item Atemfrequenz {\textgreater}20/min oder Hyperventilation

\item unkontrollierbare Gerinnung/Blutung

\item Hypotonie ${\rightarrow}$ Septischer Schock

\item Organversagen  

\end{itemize}
Multiresistente Keime

Vor der Entdeckung der Antibiotika waren bakterielle Infektionen sehr
gef\"urchtet, viele Menschen sind an -- nach heutiger Sicht simplen --
Infektionen gestorben. 

Heute kann man zwar solche Infektionen mit Antibiotika behandeln, jedoch
k\"onnen Bakterien Resistenzen gegen Antibiotika entwickeln. Die sonst
wirksamen Medikamente wirken dann nicht mehr oder nur noch sehr
abgeschw\"acht. Als behandelnder Arzt steht man dann hilflos vor dem
Patienten. 

Resistenzen machen Antibiotika wirkungslos. 

Es muss daher das Ziel aller sein die Ausbreitung von resistenten Keimen
zu verhindern. Dem Rettungs- und Krankentransport kommt dabei eine
Schl\"usselrolle zu, da es hier leicht zu Kreuzinfektionen kommen kann.


Resistenzen sind ein enorm gro{\ss}es Problem in der heutigen Medizin!

... und damit auch im Rettungs- und Krankentransport, da es hier leicht
zu Kreuzinfektionen kommen kann..

{\textquotedblleft}\index{Spitalskeime}Spitalskeime{\textquotedblright}:
Keime mit {\textquotedblleft}angez\"uchteten{\textquotedblright}
Resistenzen

Unterscheide:
{\textquotedblleft}Wald-und-Wiesen{\textquotedblright}-Keime und
Spitalskeime 

Der wohl bekannteste Problemkeim ist der MRSA (Multiresistenter
Staphylococcus aureus, siehe unten). Das liegt schlicht daran dass er
der erste war. In den letzten Jahren ist die Resistenzentwicklung aber
auch bei vielen anderen Keimen zu beobachten, Beispiele daf\"ur sind
die ESBL{\textquotesingle}s und VRE{\textquotesingle}s sowie der VRSA
(Extended Spectrum Beta-Laktamase-Bildner, Vancomycin-resistente
Enterokokken, Vancomycin-resistenter Staph. aureus). Das Problem der
multiresistenten Keime wird immer gr\"o{\ss}er.

Je wirkungsloser die Antibiotika werden desto wichtiger ist die Hygiene
und die Vermeidung von (Kreuz-)Infektionen. 

\begin{minipage}{13.67cm}
[Warning: Draw object ignored]

Zeichnung \stepcounter{Drawing}{\theDrawing}: Teufelskreis der
resistenten Keime. Oft das Bindeglied: Der Rettungs- und
Krankentransportdienst. Quelle: 

\end{minipage}

\paragraph[MRSA - Multi-Resistenter Staphylococcus aureus]{MRSA -
Multi-Resistenter Staphylococcus aureus}
Urspr\"unglich: Methicilin-Resistenter Staphylococcus Aureus

\begin{itemize}
\item Staphylococcus aureus ist ein
{\textquotedblleft}Allerweltskeim{\textquotedblright}, bis zu 50\% der
Bev\"olkerung sind besiedelt

\item Vor allem vordere Nasenh\"ole besiedelt, aber auch Haut
(Achselh\"ole, Analregion), Schleimh\"aute

\item Auch im Spital h\"aufiger Keim, aber dort oft
Antibiotika-Resistent: MRSA

\end{itemize}
\begin{itemize}
\item Multi-resistenter Staph. aureus

\item {\textquotedblleft}Spitalsinfektionen{\textquotedblright}

\item schwer bis gar nicht behandelbar

\item Wunden, die nicht abheilen wollen, Sepsis, Tod ...

\end{itemize}
\subparagraph{}
\subparagraph[MRSA-Transport - Allgemeine Ma{\ss}nahmen]{MRSA-Transport
- Allgemeine Ma{\ss}nahmen}
\label{bkm:RefHeading40262819}\begin{itemize}
\item Informationen \"uber MRSA-Tr\"agerschaft

\begin{center}
\begin{minipage}{4.941cm}
 gilt im Prinzip auch f\"ur die anderen
{\quotedblbase}Problemkeime{\textquotedblleft}

\end{minipage}
\end{center}
\item Art/Ort der MRSA-Besiedlung

\item Patientenvorbereitung und Transport

\item Wunden sind frisch verbunden und gut abgedeckt

\item Bei Besiedlung der Atemwege tr\"agt der Patient einen Mund-
/Nasenschutz

\item Vor dem Transport f\"uhrt der Patient eine hygienische
H\"andedesinfektion durch

\item Allgemeine Hygienema{\ss}nahmen

\end{itemize}
\begin{itemize}
\item Einmalhandschuhe und Schutzkittel.

\item Nach Ablegen sofortige hygienische H\"andedesinfektion  

\end{itemize}
\begin{itemize}
\item Nach Abschluss des Patiententransportes sind alle Materialien,
Ger\"ate, Instrumente und Fl\"achen, welche direkten Kontakt mit dem
Patienten hatten gem\"a{\ss} dem \"ublichen Hygieneplan zu
desinfizieren.

\end{itemize}
\begin{itemize}
\item Kugelschreiber!!!

\item Alle waagerechten Oberfl\"achen des Fahrzeuginnenraumes sind mit
einem Fl\"achendesinfektionsmittel in Form einer
Scheuerwischdesinfektion desinfizierend zu reinigen.

\item Abschlie{\ss}end f\"uhrt das Einsatzpersonal eine hygienische
H\"andedesinfektion durch.

\end{itemize}
\begin{itemize}
\item Nach Abschluss der vorangegangenen Hygiene- und
Desinfektionsma{\ss}nahmen ist das Fahrzeug wieder voll einsatzbereit.

\end{itemize}
F\"ur die anderen Problemkeime (VRSA, VRE, ESBL, {\dots}) gelten
sinngem\"a{\ss} die gleichen Verfahrensweisen. 

\subsection[Desinfektion]{Desinfektion}
\index{Desinfektion}\begin{center}
\tablehead{}\begin{supertabular}{|m{6.649cm}|m{6.651cm}|}
\hline
 &
\\\hline
DESINFEKTION

Ziel: Keimreduktion ${\rightarrow}$ nicht infekti\"os

Richtet sich gegen ausgew\"ahlte Keime

Aber: nicht keimfrei! nicht steril!

Chemische Substanzen

Einlegebad

 &
STERILISATION

Abt\"otung aller Mikroorganismen -
{\textquotedblleft}keimfrei{\textquotedblright}

Dampfsterilisation (\"Uberdruck)

Hei{\ss}luftsterilisation

\\\hline
\end{supertabular}
\end{center}
\index{Hygieneplan}Hygieneplan - Hygienebeauftragter

Der Hygieneplan regelt die konkreten Hygienema{\ss}nahmen und Was Wann
Wie Womit (Wo) zu geschehen hat, siehe allgemeiner Aushang.

Der derzeitige Hygienebeauftragte ist
\_\_\_\_\_\_\_\_\_\_\_\_\_\_\_\_\_\_\_\_\_\_\_\_\_\_\_\_\_\_\_\_\_\_.

Pers\"onliche Hygiene

\begin{itemize}
\item Schmuck, Uhren, Freundschaftsarmb\"ander sind perfekte
N\"ahrb\"oden ${\rightarrow}$ Vor dem Dienst ablegen!

\begin{center}
 [Warning: Image ignored] % Unhandled or unsupported graphics:
%\includegraphics[width=3.112cm,height=3.941cm]{AASdev20090228-img99.jpg}

\end{center}
\item Dienstkleidung regelm\"a{\ss}ig waschen!

\item Bei Kontamination wechseln!

\item Verwendung von Plastiksch\"urzen

\item Bei Dienstantritt

\end{itemize}
\begin{itemize}
\item H\"ande waschen

\item Nagelfalze reinigen

\item \index{Hygienische H\"andedesinfektion}Hygienische
\index{H\"andedesinfektion}H\"andedesinfektion

\end{itemize}
Wischdesinfektion

\begin{itemize}
\item Mit
\index{Fl\"achendesinfektionsmittel}Fl\"achendesinfektionsmittel 

\item Einwirkzeit lt. Packung beachten!

\item Fl\"ache muss w\"ahrend der Zeit feucht sein

\item in Autos nur Fl\"achendesinfektionsmittel ohne Alkohol!

\end{itemize}
\begin{itemize}
\item Alkohol zerst\"ort Acryl-Fl\"achen! ${\rightarrow}$
\index{Acryl-Des{\texttrademark}}Acryl-Des{\texttrademark}  

\end{itemize}
\begin{itemize}
\item Von au{\ss}en in Richtung Schmutzquelle wischen

\item Einmalt\"ucher wechseln

\end{itemize}
\begin{itemize}
\item nicht den Schmutz \"uberall verteilen

\end{itemize}
In den Fahrzeugen d\"urfen keine alkoholischen Desinfektionsmittel
verwendet werden da diese die Acryloberfl\"achen zerst\"oren k\"onnen!

Ger\"atedesinfektion

\begin{itemize}
\item Nach Kontamination:

\end{itemize}
\begin{itemize}
\item Grobe Reinigung vor Ort: Eintrocknen verhindern

\begin{center}
 [Warning: Image ignored] % Unhandled or unsupported graphics:
%\includegraphics[width=4.737cm,height=2.842cm]{AASdev20090228-img100.jpg}
\captionof{figure}[Desinfektionswannen, verschiedene
Gr\"o{\ss}en]{Desinfektionswannen, verschiedene Gr\"o{\ss}en}

\end{center}
\end{itemize}
\begin{itemize}
\item Wischdesinfektion

\item Ggfs. Schl\"auche durchsp\"ulen

\item Fl\"achendesinfektionsmittel in Nierenschale / Joghurtbecher,
saugen ...

\end{itemize}
\begin{itemize}
\item Meldung an Leitstelle: {\textquotedblleft}Nicht EB, m\"ussen
putzen{\textquotedblright}

\item Kontaminierte Teile gegen Ersatzmaterial tauschen

\end{itemize}
\begin{itemize}
\item Zerlegen

\item Nochmals sp\"ulen

\end{itemize}
\begin{itemize}
\item Wenn Desinfektionswanne leer ${\rightarrow}$ einlegen 

\end{itemize}
\begin{itemize}
\item Vollst\"andig zerlegt laut Anleitung

\item Luftblasenfrei!

\item Mit Niederhaltesieb unter die Oberfl\"ache dr\"ucken

\item Zeit auf Zettel notieren 

\end{itemize}
\begin{itemize}
\item {\textquotedblleft}was, wann, von wo, von wem
eingelegt{\textquotedblright}

\end{itemize}
\begin{itemize}
\item sonst nur bereit legen, in Plastiksack

\end{itemize}
\begin{itemize}
\item {\textquotedblleft}was, von wo, von wem, warum{\textquotedblright}

\end{itemize}
Hygienische H\"andedesinfektion

\begin{itemize}
\item Ziel: 

\end{itemize}
\begin{itemize}
\item Keimreduktion

\end{itemize}
\begin{itemize}
\item Vermeidung von Kreuzinfektionen

\item Patient1 ${\rightarrow}$ Sanit\"ater ${\rightarrow}$ Patient 2

\item {\textquotedblleft}Unterbrechung der
Infektionskette{\textquotedblright}

\end{itemize}
\begin{itemize}
\item Durchzuf\"uhren:

\end{itemize}
\begin{itemize}
\item Dienstantritt

\item nach jedem Patienten!!

\end{itemize}
\begin{itemize}
\item Womit?

\end{itemize}
\begin{itemize}
\item H\"andedesinfektionsl\"osung: Sterilium, Monopronto

\begin{center}
 [Warning: Image ignored] % Unhandled or unsupported graphics:
%\includegraphics[width=4.001cm,height=4.001cm]{AASdev20090228-img101.jpg}

\end{center}
\end{itemize}
Vorgehen 

\begin{enumerate}
\item Ringe entfernen 

\item Eine Portion alkoholisches H\"andedesinfektionsmittel (ca. 3ml =
1~Hohlhand) aus dem Spen\-der entnehmen. 

\item mittels Ellenbogen, vollkommene Benetzung beider H\"ande und
Handgelenke 

\item Handinnenfl\"achen gegeneinander reiben

\item geschlossene Finger 

\item Handgelenke einreiben

\item Mit rechter Handinnenfl\"ache linken Handr\"ucken reiben und
umgekehrt, dabei Finger ineinander verschr\"anken.

\item Mit ineinander verschr\"ankten Fingern Handinnenfl\"achen
gegeneinander reiben. 

\begin{center}
 [Warning: Image ignored] % Unhandled or unsupported graphics:
%\includegraphics[width=3.961cm,height=3.169cm]{AASdev20090228-img102.jpg}

\end{center}
\item H\"ande ineinander verhaken und Finger gegeneinander bewegen

\item Daumen mit gegen\"uberliegender Hand vollst\"andig umschlie{\ss}en
und rotierend reiben. 

\item Daumenkuppen nicht vergessen!

\item Fingerkuppen im Handteller kreisf\"ormig einreiben bis Alkohol
verdunstet ist

\end{enumerate}
 Einwirkzeit von mindestens 30 Sekunden beachten! 

[Warning: Draw object ignored]

\"Offnen von sterilem Material

Packung {\quotedblbase}aufsch\"alen{\textquotedblleft} 

  [Warning: Image ignored] % Unhandled or unsupported graphics:
%\includegraphics[width=6.656cm,height=4.993cm]{AASdev20090228-img103.jpg}
    [Warning: Image ignored] % Unhandled or unsupported graphics:
%\includegraphics[width=6.663cm,height=4.984cm]{AASdev20090228-img104.jpg}
 

Selbstschutz

\begin{itemize}
\item Immer Einmalhandschuhe beim Patientenkontakt tragen!

\item Handschuhe nach jedem Patienten wechseln

\item Nie mit schmutzigen Handschuhen in die Tasche greifen!

\item Ein Helfer hat Patientenkontakt, der andere reicht zu!

\item Achtung {\textquotedblleft}Kugelschreiber{\textquotedblright}

\item Einmalsch\"urzen bei Bedarf

\end{itemize}
\begin{itemize}
\item evt. Infektionsschutzanz\"uge, Maske

\end{itemize}
\begin{itemize}
\item Spitze Gegenst\"ande (Nadeln, Lanzetten, ... sofort in Gelbe Box!

\end{itemize}
\begin{itemize}
\item Ausnahme: ausdr\"uckliche Arztanweisung (Blutzuckerbestimmung aus
Nadel)

\item {\textquotedblleft}Wer sticht, der entsorgt!{\textquotedblright}

\end{itemize}
\paragraph{Kontamed, Sharp, Gelbe Box ...}
Spitze Gegenst\"ande (Nadeln, Lanzetten, ... sofort in Gelbe Box!

Ausnahme: ausdr\"uckliche Arztanweisung (Blutzuckerbestimmung aus Nadel)

Wer sticht, der entsorgt!



\begin{center}
 [Warning: Image ignored] % Unhandled or unsupported graphics:
%\includegraphics[width=3.377cm,height=4.552cm]{AASdev20090228-img105.jpg}
 [Warning: Image ignored] % Unhandled or unsupported graphics:
%\includegraphics[width=5.71cm,height=4.289cm]{AASdev20090228-img106.jpg}
\captionof{figure}[Kontamed-Boxen, verschiedene
Fabrikate.]{Kontamed-Boxen, verschiedene Fabrikate.}

\end{center}
NIE Nachstopfen.

NIE ausleeren.

Wenn die Box voll ist wird sie verschlossen und abgegeben.

Wundreinigung

\begin{itemize}
\item Bei frischen, nicht infizierten Wunden:

\end{itemize}
\begin{itemize}
\item Reinigung von innen nach au{\ss}en

\item Ziel: keine Keime in Wunde bringen

\end{itemize}
\begin{itemize}
\item Womit?

\end{itemize}
\begin{itemize}
\item Sterilen L\"osungen oder Octenisept

\item Sterile Tupfer

\end{itemize}
KEINE Papierhandt\"ucher oder unsterile Tupfer !!!

\begin{itemize}
\item Danach steriler Verband!

\end{itemize}
\begin{itemize}
\item Bei offenen Frakturen ist keimfreiheit be\-sonders wichtig !!!

\item Gefahr einer lebenslangen Knochen\-marks\-ent\-z\"undung
(Osteomyelitis)!

\item \index{OpSite}OpSite{}-Folie wenn vorhanden.

\end{itemize}
 [Warning: Image ignored] % Unhandled or unsupported graphics:
%\includegraphics[width=5.742cm,height=5.658cm]{AASdev20090228-img107.png}
\captionof{figure}[Reiningung einer sauberen Wunde: Von Innen nach
Aussen, der Tupfer wird anschliessend verworfen.]{Reiningung einer
sauberen Wunde: Von Innen nach Aussen, der Tupfer wird anschliessend
verworfen.}
 [Warning: Image ignored] % Unhandled or unsupported graphics:
%\includegraphics[width=4.219cm,height=4.219cm]{AASdev20090228-img108.jpg}
\captionof{figure}[Anbringen einer OpSite(TM)-Folie: Es ist wichtig die
Folie zu spannen damit sie sich nicht selbst verklebt. ]{Anbringen
einer OpSite(TM)-Folie: Es ist wichtig die Folie zu spannen damit sie
sich nicht selbst verklebt. }


\subsection{Besondere Verhaltensweisen}
%
%
\index{Nadelstichverletzungen}Nadelstichverletzungen

\begin{itemize}
\item Erstma{\ss}nahmen durchf\"uhren: 

\end{itemize}
\begin{itemize}
\item bei Stich-bzw.Schnittverletzungen: Wunde auspressen und 10 Minuten
desinfizieren

\item bei Spritzern mit Blut oder Sekreten in Auge oder Mund:
Schleimhaut mit Wasser u./o. Schleimhautdesinfektionsmittel sp\"ulen

\end{itemize}
\begin{itemize}
\item Leitstelle informieren

\item ${\rightarrow}$ Umgehend in ein Unfallspital der AUVA

\item Patient: Zielspital informieren:

\end{itemize}
\begin{itemize}
\item Blut sichern f\"ur HIV- und Hepatitis-Serologie

\end{itemize}
\begin{itemize}
\item  Unfallmeldung  (wichtig!), Blutabnahme 

\item ggf. HIV-Postexpositionsprohylaxe (PEP)

\end{itemize}
\begin{itemize}
\item kurzfristige Therapie (4 Wo) 

\item Einzelfall muss immer individuell beurteilt werden
(Nutzen/Risiko), starke Nebenwirkungen.

\item Muss innerhalb von 2 Std. begonnen werden ${\rightarrow}$ daher
SOFORT in ein Unfallspital der AUVA

\end{itemize}
\begin{itemize}
\item Meldung an Betriebsleitung

\item Unfallmeldung (wichtig!) an die AUVA

\end{itemize}
Unfallmeldung nicht vergessen!

Sofort in ein Unfallspital! 

\index{MRSA}MRSA und andere Resistente Keime

Die Ma{\ss}nahmen beim Transport von Patienten mit resistenten Keimen
werden unter {\quotedblbase}MRSA-Transport - Allgemeine
Ma{\ss}nahmen{\quotedblbase} auf Seite
K.II{\guillemotright}\pageref{bkm:RefHeading40262819} besprochen.

\section{[RS XI] Katastrophen, Gro{\ss}schadensereignisse,
Gefahrengutunfall}
Roman Koch [*]

Changelog:

2009-02-20: Fertig f\"ur Kurs 2009/02.

2009-02-25: Start Changelog.

\subsection{Katastrophen, Gro{\ss}schadensereignisse; Unfall}
\begin{minipage}{13.7cm}
  [Warning: Image ignored] % Unhandled or unsupported graphics:
%\includegraphics[width=13.68cm,height=6.929cm]{AASdev20090228-img109.jpg}
 

ICE-Ungl\"uck von Eschede am 3. Juni 1998

 [Warning: Image ignored] % Unhandled or unsupported graphics:
%\includegraphics[width=6.212cm,height=3.933cm]{AASdev20090228-img110.jpg}
Verstorbene:101

Schwerverletzte:88  Leicht- und Unverletzte:106

Foto Wikipedia, Nils Fretwurst; Lizenz: gemein frei
({\quotedblbase}public domain{\textquotedblleft})

\end{minipage}

Definition Unfall (Bagatelle, Notfall)

\begin{center}
\begin{tabular}{|m{3.31cm}|m{9.99cm}|}
\hline
Ereignis:

 &
pl\"otzlich, unerwartet

\\\hline
Dauer:

 &
Minuten bis 1 Std.

\\\hline
Infrastruktur:

 &
intakt

\\\hline
Verletzte:

 &
bis 14

\\\hline
Erstversorgung:

 &
innerhalb von Minuten durch Rettungsdienst nach Regeln der
Individual\-medizin

\\\hline
Abtransport:

 &
unbehindert Routinebetrieb mit \"ortlich vorhandenen Mitteln

\\\hline
Definitive Versorgung:

 &
in umliegende Krankenh\"auser nach Regeln der Individualmedizin

\\\hline
Material und Mannschaft:

 &
Normalbetrieb Rettungs- und Krankentransport

\\\hline
\end{tabular}
\end{center}
Definition \index{Gro{\ss}unfall}Gro{\ss}unfall
(\index{Gro{\ss}schadensereignis}Gro{\ss}schadensereignis)

{\quotedblbase}Ein Gro{\ss}unfall liegt vor, wenn anzunehmen ist, dass
das Ereignis mit den normalen \"ortlichen personellen und materiellen
Kr\"aften und Mitteln nicht bew\"altigt werden kann, aber keine
erkl\"arte Katastrophensituation vorliegt{\textquotedblleft}

\begin{center}
\begin{tabular}{|m{3.31cm}|m{9.99cm}|}
\hline
Ereignis:

 &
lokal begrenzt

\\\hline
Dauer:

 &
mehrere Stunden

\\\hline
Infrastruktur:

 &
intakt

\\\hline
Verletzte:

 &
ab 15

\\\hline
Erstversorgung:

 &
innerhalb von Minuten bis Stunden durch den Rettungsdienst nach den
Regeln der Katastrophenmedizin

\\\hline
Abtransport:

 &
relativ unbehindert nach der Erstversorgung, gen\"ugend vorhandene
Transportmittel

\\\hline
Definitive Versorgung:

 &
in umliegende Krankenh\"auser nach den Regeln der Individualmedizin

\\\hline
Material und Mannschaft:

 &
alle verf\"ugbaren Helfer, Normalbetrieb Rettungs- und Krankentransport
inkl. Sondereinsatzgruppe und zus\"atzlichem Material

\\\hline
\end{tabular}
\end{center}
Definition \index{Katastrophe}Katastrophe (au{\ss}ergew\"ohnliche Lage)

griech.:  {\quotedblbase}Umkehr der normalen
Verh\"altnisse{\textquotedblleft} 

Ein unvorhergesehenes Ereignis mit einem au{\ss}ergew\"ohnlichen
Schadensausma{\ss} das unmittelbar bevorsteht oder bereits eingetreten
ist, eine konkrete Gefahr f\"ur Menschen, Tiere, Umwelt, Kulturg\"uter
und Sachwerte sowie f\"ur die Infrastruktur zur Sicherstellung der
Versorgung mit lebensnotwendigen G\"utern und Dienstleistungen, welches
der koordinierten F\"uhrung durch die Beh\"orde bedarf. 

\begin{center}
\begin{tabular}{|m{3.31cm}|m{9.99cm}|}
\hline
Ereignis:

 &
(\"uber)-regionales Gebiet

\\\hline
Dauer:

 &
Tage, Wochen

\\\hline
Infrastruktur:

 &
gest\"ort, zerst\"ort

\\\hline
Verletzte:

 &
keine bis viele, je nach Art des Ereignisses

\\\hline
Erstversorgung:

 &
innerhalb von Stunden bis Tagen durch KHD\footnotemark{}{}-Einheiten
nach den Regeln der Katastrophenmedizin

\\\hline
Abtransport:

 &
stark behindert verz\"ogert, wenige/nicht vorhandene Transportmittel

\\\hline
Definitive Versorgung:

 &
auf Verbandspl\"atzen in Behelfsspit\"alern, in Lazaretten nach Regeln
der Katastrophenmedizin

\\\hline
Material und Mannschaft:

 &
\"Uberlebende und verf\"ugbare Helfer, einsatzbereite Fahrzeug,
Kat-Einsatz aller Hilfsorganisationen, \"uberregionale und
internationale Hilfe

\\\hline
\end{tabular}
\end{center}
\footnotetext{KHD: Katastrophenhilfsdienst}
Definition \index{Individualmedizin}Individualmedizin

Bedarfsgerechte und bestm\"ogliche Versorgung jedes Patienten 

Freie Arztwahl, welche auf das Vertrauensverh\"altnis zwischen Patient
und Arzt ausgerichtet ist

Definition \index{Katastrophenmedizin}Katastrophenmedizin

Eine Vielzahl von Hilfsbed\"urftigen stehen wenigen \"Arzten,
Sanit\"atern und sonstigen Helfern mit oft unzureichenden Mitteln
gegen\"uber. Ziel muss es sein, unter den gegebenen Umst\"anden so
viele Leben wie m\"oglich zu retten. Mit bewussten Verzicht auf zeit-
und personalaufwendige Maximalversorgung des Einzelnen zugunsten der
lebensrettenden Minimalversorgung Vieler.

Katastrophenmedizin bedeutet Massenversorgung von Verwundeten und
Kranken mit beschr\"ankten Mitteln und dem Zwang zur Einteilung nach
Dringlichkeit.

Bei Gro{\ss}unf\"allen und Katastrophen versucht man das beste
herauszuholen

\begin{center}
\tablehead{}\begin{supertabular}{|m{6.649cm}|m{6.651cm}|}
\hline
Ziele

\begin{itemize}
\item das Bestm\"ogliche

\item f\"ur die gr\"o{\ss}te Anzahl

\item zur rechten Zeit

\item am rechten Ort

\end{itemize}
 &
Die ist nur erreichbar wenn:

\begin{itemize}
\item der Bestgeeignete auf jeder Stufe f\"uhrt

\item alle Einsatzkr\"afte Ihre Aufgabe kennen 

\item das dazu ben\"otigte Material vorhanden, bereitgestellt und rasch
einsatzbereit ist

\item klare und eindeutige F\"uhrungsstrukturen vorhanden sind

\item die F\"uhrungspyramide von allen strikt eingehalten wird

\item die Organisation eingespielt/einge\"ubt ist

\item Eins\"atze geplant und vorbereitet sind

\end{itemize}
\\\hline
Daf\"ur werden Sonderausr\"ustungen zus\"atzlich erforderlich, da ein
erh\"ohter Bedarf besteht an:

\begin{itemize}
\item Material

\item Medikamente

\item Personal

\end{itemize}
 &
Wir m\"ussen rechnen mit

\begin{itemize}
\item einem gro{\ss}en Patientenanfall

\item besonderen Verletzungsmustern

\item einer l\"angeren Einsatzzeit vor Ort

\end{itemize}
\\\hline
\end{supertabular}
\end{center}
Es gibt verschiedene Katastrophenarten

\begin{center}
\tablehead{}\begin{supertabular}{|m{6.649cm}|m{6.651cm}|}
\hline
Naturkatastrophen

 &
Erdbeben,\"Uberschwemmungen, Erdrutsche/Lawinen etc.

\\\hline
Technische Katastrophen

 &
Flugzeugabsturz, Eisenbahnungl\"uck, Reaktor\-unfall, Schiffs\-untergang
etc.

\\\hline
Sekund\"arkatastrophen

 &
Hungersnot, Seuchen etc.

\\\hline
\end{supertabular}
\end{center}
Katastrophenhilfegesetze regeln die Zust\"andigkeiten

Die Katastrophenhilfsgesetze regeln die Aufgaben der L\"ander,
Organisationen und Zust\"andigkeiten in Ihrem Bereich.
Ausschlie{\ss}lich die Beh\"orde erkl\"art ein au{\ss}ergew\"ohnliches
Schadensereignis zur Katastrophe im Sinne des Gesetzes. Beh\"ordliche
Alarmpl\"ane dienen zur Festlegung der Ma{\ss}nahmen der
Katastrophenhilfe und zur Koordination der im Anlassfall t\"atig
werdenden Einrichtungen. Je nach Ausdehnung des Schadensgebietes gibt
es unterschiedliche Kompetenzen:

\begin{center}
\begin{tabular}{|m{2.375cm}m{0.45600003cm}m{6.1590004cm}|}
\hline
Gemeinde 

 &
\centering ${\rightarrow}$\par

 &
B\"urgermeister

\\
Bezirk 

 &
\centering ${\rightarrow}$\par

 &
Bezirkshauptmann

\\
Bundesland 

 &
\centering ${\rightarrow}$\par

 &
Landeshauptmann

\\
Bundesgebiet

 &
\centering ${\rightarrow}$\par

 &
Innenminister, Bundeskanzler

\\\hline
\end{tabular}
\end{center}
Die Katastrophenwarnung mittels Sirene alarmiert die Bev\"olkerung

Katastrophenschutz in Wien

\begin{center}
 [Warning: Image ignored] % Unhandled or unsupported graphics:
%\includegraphics[width=10.861cm,height=7.614cm]{AASdev20090228-img111.png}

\end{center}
\begin{center}
\tablehead{}\begin{supertabular}{|m{4.3190002cm}|m{8.981cm}|}
\hline
Wien ist f\"ur den Gro{\ss}\-schadens- und Katastrophenfall gut
ger\"ustet. Tausende Helfer sind dazu im Katastrophen\-schutz-Kreis
({\quotedblbase}K-Kreis{\textquotedblleft}) org\-anisiert.

Der K-Kreis ist ein Zu\-sammen\-schluss der beruflichen und freiwilligen
Hilfs- und Ein\-satz\-organisationen sowie wichtiger Beh\"orden.
Dreh\-scheibe und gemeinsame Anlaufstelle ist dabei die Stadt Wien,
Magistrats\-direktion - Katastrophen\-manage\-ment und
Sofort\-ma{\ss}\-nahmen und die Helfer Wiens.

 &
 [Warning: Image ignored] % Unhandled or unsupported graphics:
%\includegraphics[width=9.046cm,height=13.623cm]{AASdev20090228-img112.png}
\captionof{figure}[Der K-Kreis]{Der K-Kreis}


\\\hline
\end{supertabular}
\end{center}
Katastrophenhilfe

Hilfsma{\ss}nahmen die von uns als Sanit\"atsorganisation erwartet
werden:

\begin{itemize}
\item Bergung, Versorgung und Transport 

\item Pflegerische Betreuung 

\item Sorge f\"ur Unterkunft und Verpflegung 

\item Karitative Betreuung von Hilfsbed\"urftigen

\end{itemize}
\begin{center}
\tablehead{}\begin{supertabular}{|m{6.649cm}|m{6.651cm}|}
\hline
Sanit\"atseinsatz

 &
\begin{itemize}
\item Retten/Bergen

\item Triage

\item Lebens rettende Sofortma{\ss}nahmen

\item Erstversorgung

\item Sanit\"atshilfe

\item Abtransport

\end{itemize}
\\\hline
Pflegeeinsatz

 &
Wenn die pflegerische Versorgung durch die bestehenden Einrichtungen
nicht mehr sichergestellt ist.

\begin{itemize}
\item Erweiterung der vorhandenen Behandlungsm\"oglichkeiten

\item Errichtung von Hilfskrankenh\"ausern

\end{itemize}
\\\hline
Sozialeinsatz

 &
Wenn Menschen infolge einer Katastrophe ihre
Selbstversorgungsm\"oglichkeit verloren haben.

\begin{itemize}
\item Einrichtung/Betrieb von Notunterk\"unften/Lagern  

\item Beschaffung, Zubereitung und Ausgabe von Lebensmitteln, Kleidern,
Gebrauchsgegenst\"anden  

\item F\"ursorgliche Hilfe, Registrierung und Betreuung

\end{itemize}
\\\hline
\end{supertabular}
\end{center}
Grunds\"atze des Einsatzes

Grundkomponenten: F\"uhrung -- Zeit -- Raum 

Ziel ist es ohne \"uberst\"urztes Handeln die Ausdehnung der Krisenlage
auf benachbarte und andere R\"aume zu verhindern. Dabei wird
gem\"a{\ss} der gesetzlichen Bestimmungen von Kr\"aften der
Einsatzorganisationen, Beh\"orden und anderen Einrichtungen koordiniert
vorgegangen.

\paragraph{F\"uhrung}
Steuerndes Einwirken auf andere Menschen um ein bestimmtes Ziel zu
erreichen. 

Die F\"uhrung ist ein unentbehrliches Steuerungsinstrument zur
Erf\"ullung des Einsatzauftrages
({\quotedblbase}Ziel{\textquotedblleft}) und dem richtigen Einsatz von
Personal, Material und Transportmitteln 

Im Schadensraum gilt: {\quotedblbase}Ein einziger
Verantwortlicher{\textquotedblleft} mit Entscheidungsbefugnis und den
n\"otigen F\"uhrungsmittel ({\quotedblbase}Einheit der
F\"uhrung{\textquotedblleft}). Dieser wirkt mittels m\"oglichst
einfacher Ma{\ss}nahmen ({\quotedblbase}Einfachheit{\textquotedblleft})
auf den Einsatzverlauf ein
({\quotedblbase}Handlungsfreiheit{\textquotedblleft}). Mittel und
Mannschaften werden dabei gem\"a{\ss} ihrer Eigenart und
Leistungsf\"ahigkeit eingesetzt ({\quotedblbase}\"Okonomie der
Kr\"afte{\textquotedblleft}).Er muss rasch auf \"Anderungen der
Situation reagieren ({\quotedblbase}Beweglichkeit{\textquotedblleft})
und gleichzeitig vorausschauend handeln
({\quotedblbase}Reservenbildung{\textquotedblleft}).

{\quotedblbase}F\"uhren{\textquotedblleft} und
{\quotedblbase}Retten{\textquotedblleft} ist nicht gleichzeitig
m\"oglich!

Sanit\"ater/NA und Not\"arzte m\"ussen F\"uhrungsaufgaben \"ubernehmen

%
%
Der erste am Schadensort eintreffende Sanit\"ater muss:

\begin{enumerate}
\item Lagemeldung

\item den Standort der \"ortlichen Einsatzleitung festlegen

\item mit bereits anwesenden Einsatzkr\"aften die ersten Ma{\ss}nahmen
zur Rettung festlegen

\item wohin Patienten im Ein

\end{enumerate}
\paragraph[Zeit]{Zeit}
Jede Katastrophenereignis zeigt unabh\"angig von seiner Ursache
folgenden Zeitablauf

\begin{center}
\begin{tabular}{|m{4.367cm}|m{4.367cm}|m{4.367cm}|}
\hline
\centering 1. Phase\par

 &
\centering 2. Phase\par

 &
\centering 3. Phase \par

\\\hline
Isolation (Minuten)

 &
Retten (Stunden)

 &
Wiederherstellung (Tage)

\\\hline
Schaden begrenzen  {}- Lage erkennen {}- Melden

 &
Ordnung schaffen {}- Kommandoposten

 &
Lokale Einsatzleitung {}- Verbindung {}- Lage beurteilen

\\\hline
\"Uberleben

 &
Weiterleben

 &
Endbehandlung

\\\hline
Laien: {}- Nothilfe/Selbsthilfe {}- Retten/Erste Hilfe

 &
Geschulte Helfer: {}- LRSM {}- Transportf\"ahigkeit

 &
Spezialisten: {}- Spital nach Dringlichkeit behandeln

\\\hline
\multicolumn{3}{|m{13.501cm}|}{[Warning: Draw object ignored]

}\\\hline
\end{tabular}
\end{center}
\paragraph{Hilfe braucht Raum: Sanit\"atsdienstliche Organisation im
Schadensraum und in der Sanit\"atshilfsstelle (SanHist)}
Der Einsatz allgemein wird unterteilt in: 

\begin{itemize}
\item Schadensraum

\item Transportraum

\item Hospitalisationsraum

\end{itemize}
Der Hospitalisationsraum betrifft die Transportziele und
Aufnahmespit\"aler. 

Der Transportraum umfasst in diesem Zusammenhang alle Mittel und Wege um
die Patienten vom Schadens\-raum zun Hospitalisationsraum zu
verbringen. Dieser {\quotedblbase}allgemeine
Transportraum{\textquotedblleft} ist nicht zu ver\-wechseln mit dem
Transportraum der Sanit\"atshilfestelle (s.u.).  

Der Schadensraum ist der Ort des eigentlichen (schadens-)Geschehens. Es
gibt in ihm eine einheitliche Organisations\-form f\"ur den
Sanit\"atsdienst: Die Sanit\"atshilfsstelle (SanHiSt), welche aus den 3
R\"aumen Triageraum, Behandlungsraum und Transportraum besteht.
Definiert wird Sie in der ASB\"O Rahmen\-vor\-schrift f\"ur
Gro{\ss}unf\"alle und den Durchf\"uhrungsbestimmungen f\"ur
Gro{\ss}unf\"alle. 

\begin{minipage}{12.605cm}
 [Warning: Image ignored] % Unhandled or unsupported graphics:
%\includegraphics[width=12.605cm,height=14.03cm]{AASdev20090228-img113}
\captionof{figure}[Gliederung des Einsatzraumes]{Gliederung des
Einsatzraumes}
\begin{multicols}{2}
\"Aussere Absperrung: Gro{\ss}r\"aumige Verkehrsumleitung

Innere Absperrung: h\"alt Schaulustige fern, h\"alt Personen in Panik
zur\"uck

Sicherheitsring: Absperrung des Schadenplatzes, Sicherheit f\"ur
Einsatzkr\"afte

Pforte: Einbahnsystem f\"ur Einsatz-Kfz, keine Fremdfahrzeuge,
m\"oglichst nur eine Pforte 

Leiter Information (i) f\"ur Weiterleitung von Informationen an die
Au{\ss}enwelt (Presse, Angeh\"orige, ...)

Bergetriage (BT) 

Sanit\"atshilfsstelle (SanHist): mit Triageraum, Behandlungsraum und
Transportraum;

Einsatzleiter (EL) 

Leitender Notarzt (LNA)

Mobile Leitstelle (MLS) 

Material- und Meldestellet (M)

Sammelstelle Tote (t) und Unverletzte (U) 

\end{multicols}
\end{minipage}

\subparagraph{Triageraum}
Triage (frz.) = Sichtung, Auswahl

Triage ist das Begutachten aller Patienten mit geringen Mitteln und
Aufwand sowie das Einteilen in Triagegruppen mit dem Ziel m\"oglichst
viele Leben zu retten!

{\quotedblbase}Life before Limbs{\textquotedblleft} = Leben vor
Gliedma{\ss}en, Funktion vor Sch\"onheit!

Bergetriage

\begin{center}
\tablehead{}\begin{supertabular}{m{10.891cm}m{2.409cm}}
Nur wenn keine Gefahrenzone und gen\"ugend Personal/Zeit vorhanden ist!

\begin{itemize}
\item BT Team muss mind. aus einem NFS bestehen

\item Gefahren beachten

\item Erst Sichtung der Betroffenen und Festlegung der Dringlichkeit

\item Nicht behandeln!

\end{itemize}
 &
\raggedleft   [Warning: Image ignored]
% Unhandled or unsupported graphics:
%\includegraphics[width=2.396cm,height=2.149cm]{AASdev20090228-img114.png}
 \par

\\
\end{supertabular}
\end{center}
Triage (Sichtung)

\begin{center}
\tablehead{}\begin{supertabular}{m{10.891cm}m{2.409cm}}
\begin{itemize}
\item Bei Missverh\"altnis zwischen der Zahl der Betroffenen und den
Be\-handlungs\-m\"oglich\-keiten

\item Schnell, mit geringem Aufwand

\item Hoher Zeitdruck

\item Keine detaillierte Diagnostik m\"oglich

\item Triage ist dynamischer Vorgang

\item St\"andige Nachtriage ist erforderlich

\item Maximaler Zeitaufwand pro Patient gesamt:

\end{itemize}
\begin{itemize}
\item Gehender 1 Minute

\item Liegender 3 Minuten 

\end{itemize}
\begin{itemize}
\item Zentrale Registrierung aller Patienten mittels Patientenleitsystem
(PLS)

\item Zuordnung der Patienten zu Triagegruppen

\end{itemize}
\begin{itemize}
\item T1: Sofortbehandlung, lebensrettende Noteingriffe

\item T2: Dringende Behandlung, Schwerverletzte

\item T3: Minimalbehandlung

\item T4: Abwartende Behandlung

\end{itemize}
 &
\raggedleft   [Warning: Image ignored]
% Unhandled or unsupported graphics:
%\includegraphics[width=2.375cm,height=2.184cm]{AASdev20090228-img115}
 \par

\\
\end{supertabular}
\end{center}
\begin{center}
\begin{tabular}{|m{6.649cm}m{6.651cm}|}
\hline
\multicolumn{2}{|m{13.5cm}|}{\centering Zwei Triagepl\"atze, damit der
NA abwechselnd triagieren kann\par

}\\\hline
\multicolumn{1}{|m{6.649cm}|}{Vorbereitung durch RS/NFS

Entkleiden

Vitalparameter erheben (RR, HF, AF, SpO2)

Untersuchung durch den Notarzt

} &
\centering   [Warning: Image ignored]
% Unhandled or unsupported graphics:
%\includegraphics[width=5.132cm,height=3.783cm]{AASdev20090228-img116}
 \par

\\\hline
\end{tabular}
\end{center}
\subparagraph[Behandlungsraum]{Behandlungsraum}
\begin{flushleft}
\begin{tabular}{|m{3.328cm}|m{5.406cm}m{4.36cm}}
\hline
  [Warning: Image ignored] % Unhandled or unsupported graphics:
%\includegraphics[width=1.489cm,height=1.61cm]{AASdev20090228-img117.gif}
 

 &
\multicolumn{2}{m{9.966001cm}|}{Es werden drei Behandlungsstellen
eingerichtet

}\\\hline
  [Warning: Image ignored] % Unhandled or unsupported graphics:
%\includegraphics[width=1.629cm,height=1.489cm]{AASdev20090228-img118}
 

 &
\multicolumn{1}{m{5.406cm}|}{Sofortbehandlung vor Ort

T1 Lebensrettende Soforteingriffe

} &
\multicolumn{1}{m{4.36cm}|}{\centering   [Warning: Image ignored]
% Unhandled or unsupported graphics:
%\includegraphics[width=3.365cm,height=4.366cm]{AASdev20090228-img119.png}
 \par

}\\\hline
  [Warning: Image ignored] % Unhandled or unsupported graphics:
%\includegraphics[width=1.669cm,height=1.578cm]{AASdev20090228-img120.png}
 

 &
Dringende Behandlung

T2 Schwerverletzte

 &
\\\hhline{--~}
  [Warning: Image ignored] % Unhandled or unsupported graphics:
%\includegraphics[width=1.601cm,height=1.587cm]{AASdev20090228-img121.png}
   [Warning: Image ignored] % Unhandled or unsupported graphics:
%\includegraphics[width=1.587cm,height=1.641cm]{AASdev20090228-img122.png}
 

 &
Warten

T3 Warten Leichtverletzte

T4 Warten Schwerverletzte

 &
\\\hhline{--~}
\end{tabular}
\end{flushleft}
Sofortbehandlung vor Ort

Alle Patienten, welche ohne sofortige Behandlung sterben bzw. schwere
Sch\"aden erleiden

\begin{flushleft}
\tablehead{}\begin{supertabular}{m{5.642cm}m{4.379cm}m{3.076cm}}
Welche Patienten?

\begin{itemize}
\item Atemst\"orungen

\item Spannungspneumothorax

\item schwere \"au{\ss}ere Blutungen

\item schwerer traumatischer Schock

\end{itemize}
 &
Ma{\ss}nahmen

\begin{itemize}
\item \"Uberpr\"ufung der Triage

\item Volumenersatz

\item Intubation

\item Thoraxdrainage

\item Noteingriffe

\end{itemize}
 &
\centering   [Warning: Image ignored]
% Unhandled or unsupported graphics:
%\includegraphics[width=2.876cm,height=2.271cm]{AASdev20090228-img123}
 \par

\\
\end{supertabular}
\end{flushleft}
Dringende Behandlung

Behandlung und Herstellung der Transportf\"ahigkeit von schwer
verletzten/erkrankten Personen, eine weitere Behandlung im Krankenhaus
ist i.d.R. notwendig.

\begin{flushleft}
\tablehead{}\begin{supertabular}{m{5.642cm}m{4.379cm}m{3.076cm}}
Welche Patienten?

keine unmittelbare Lebensbedrohung

\begin{itemize}
\item SHT

\item Wirbelverletzungen mit L\"ahmungen

\item Bauchtrauma, innere Blutungen

\item schwere Thoraxtraumen

\item gro{\ss}e Gef\"a{\ss}verletzungen

\item offene Knochenbr\"uche und  Gelenks\-verletzungen

\item Augenverletzungen

\item gro{\ss}e Weichteilverletzungen

\item Verbrennungen 15 - 30\%

\item geschlossene Knochenbr\"uche und Gelenks\-ver\-letzungen

\end{itemize}
 &
Ma{\ss}nahmen

\begin{itemize}
\item Verb\"ande

\item Schockbehandlung

\item Fixierung/Schienung

\item Schmerzbek\"ampfung

\item pflegerische Ma{\ss}nahmen

\item LRSM wenn n\"otig

\end{itemize}
 &
\centering   [Warning: Image ignored]
% Unhandled or unsupported graphics:
%\includegraphics[width=2.775cm,height=2.44cm]{AASdev20090228-img124.png}
 \par

\\
\end{supertabular}
\end{flushleft}
Warten Leichtverletzte

Alle Patienten, welche nur einer Minimalbehandlung bed\"urfen. 

\begin{flushleft}
\tablehead{}\begin{supertabular}{m{5.642cm}m{4.379cm}m{3.076cm}}
Welche Patienten?

\begin{itemize}
\item Verbrennungen bis 15\%

\item kleine Weichteilwunden

\item einfache Knochenbr\"uche

\item Prellungen, Zerrungen

\item leichtes SHT

\end{itemize}
 &
Ma{\ss}nahmen

\"Uberpr\"ufung der Triage

pflegerische Ma{\ss}nahmen

Betreuung (KIT Team)

Entlassung ambulant Behandelter zur SaStU

 &
\centering   [Warning: Image ignored]
% Unhandled or unsupported graphics:
%\includegraphics[width=2.677cm,height=2.299cm]{AASdev20090228-img125.png}
 \par

\\
\end{supertabular}
\end{flushleft}
Warten Schwerverletzte

Alle Patienten, welche mit dem zur Verf\"ugung stehenden Mitteln
(Personal, Material) noch nicht versorgt werden k\"onnen . Abwartende
Behandlung. 

\begin{center}
\tablehead{}\begin{supertabular}{m{5.702cm}m{4.0550003cm}m{3.3439999cm}}
Welche Patienten?

\begin{itemize}
\item schweres Polytrauma

\item schwerstes SHT

\item offene K\"orperh\"ohlentraumen

\item Mehrh\"ohlenverletzungen

\item Verbrennungen \"uber 50\%

\end{itemize}
 &
Ma{\ss}nahmen

\begin{itemize}
\item \"Uberpr\"ufung der Triage

\item Behandlung  (Schmerzbek\"ampfung)

\item Betreuung (KIT Team)

\item \"Uberwachung

\end{itemize}
 &
\centering   [Warning: Image ignored]
% Unhandled or unsupported graphics:
%\includegraphics[width=2.659cm,height=2.421cm]{AASdev20090228-img126.png}
 \par

\\
\end{supertabular}
\end{center}
Transportpriorit\"at

Nach einer erfolgten Behandlung kann eine Transporpriorit\"at zugewiesen
werden. Man unterscheidet hierbei: 

\begin{itemize}
\item A: Hohe Transportpriorit\"at: Der rasche Transport zur
fr\"uhzeitigen Fachbehandlung nach \"arztlicher Notversorgung ist
angesagt. 

\item B: Niedrige Transportpriorit\"at: Der Transport zur Fachbehandlung
kann verz\"ogert erfolgen. 

\end{itemize}
\subparagraph{Transportraum}
\begin{center}
\tablehead{}\begin{supertabular}{|m{2.365cm}m{5.8640003cm}m{4.8710003cm}|}
\hline
Verladestelle

 &
Sollte in der N\"ahe der Behandlungsstelle II errichtet werden.

Es muss protokolliert werden, welcher Patient mit welchem
Transportmittel in welches Krankenhaus eingeliefert wird.

Einbahnregelung!

 &
\centering   [Warning: Image ignored]
% Unhandled or unsupported graphics:
%\includegraphics[width=3.373cm,height=2.484cm]{AASdev20090228-img127.png}
 \par

\\\hline
Wagenhalte\-platz

 &
\begin{itemize}
\item Gro{\ss} genug

\item Schr\"agparkordnung 

\item Sortiert nach NAW, RTW usw.

\item Schl\"ussel stecken lassen

\end{itemize}
\centering   [Warning: Image ignored]
% Unhandled or unsupported graphics:
%\includegraphics[width=3.55cm,height=2.498cm]{AASdev20090228-img128}
 \par

\centering Falscher Wagenhalteplatz\par

 &
\centering   [Warning: Image ignored]
% Unhandled or unsupported graphics:
%\includegraphics[width=3.375cm,height=2.438cm]{AASdev20090228-img129}
 \par

\\\hline
Hubschrauber\-landestelle

 &
\begin{itemize}
\item Die allgemeinen Vorschriften f\"ur Hubschrauberlandepl\"atze sind
zu beachten

\item Freie ebene Fl\"ache 

\item Auf Stromleitungen achten

\item Weit genug weg vom Behandlungsraum 

\end{itemize}
 &
\centering [Warning: Draw object ignored]\par

\\\hline
Transport\-management

 &
\begin{itemize}
\item Spiralf\"ormige Verteilung

\item Zeitliche Staffelung

\item Fachliche Zuordnung

\item Katastrophe nicht ins Spital verlagern!

\end{itemize}
 &
\centering   [Warning: Image ignored]
% Unhandled or unsupported graphics:
%\includegraphics[width=4.954cm,height=3.88cm]{AASdev20090228-img130}
 \par

\\\hline
\end{supertabular}
\end{center}
\subparagraph{Sammelstellen}
\begin{center}
\tablehead{}\begin{supertabular}{|m{4.367cm}m{4.367cm}m{4.367cm}|}
\hline
Sammelstelle Unverletzte

\centering   [Warning: Image ignored]
% Unhandled or unsupported graphics:
%\includegraphics[width=2.151cm,height=2.931cm]{AASdev20090228-img131}
 \par

 &
Ziel: Betreuung Unverletzter und Betroffener

\begin{itemize}
\item Betreuen

\item Beruhigen

\item Informieren

\item KIT Teams (ABW) integrieren

\end{itemize}
 &
\centering   [Warning: Image ignored]
% Unhandled or unsupported graphics:
%\includegraphics[width=2.01cm,height=2.154cm]{AASdev20090228-img132.png}
 \par

\\\hline
Sammelstelle Tote

 &
Sollte diskret eingerichtet werden.

Abgeschirmt von der Sammel\-stelle Unverletzte, dem Be\-handlungs\-raum
und der Au{\ss}en\-welt

 &
\centering   [Warning: Image ignored]
% Unhandled or unsupported graphics:
%\includegraphics[width=1.906cm,height=1.9cm]{AASdev20090228-img133.png}
 \par

\\\hline
\end{supertabular}
\end{center}
\subparagraph{Material- / Meldestelle}
\begin{itemize}
\item Dient zur Registrierung des Personals und Materials

\begin{center}
 [Warning: Image ignored] % Unhandled or unsupported graphics:
%\includegraphics[width=2.176cm,height=2.126cm]{AASdev20090228-img134}

\end{center}
\item Eintreffende Einheiten haben sich bei dieser zu melden

\item Steht im engen Kontakt zum Leiter SanHist und EL  

\end{itemize}
\subparagraph{Mobile Leitstelle (MLS)}
\begin{itemize}
\item Die MLS beherbergt i.d.R. die Einsatzleitung sowie wichtige
Administrations-, Telekommunkations- und Funkeinrichtungen. 

\begin{center}
 [Warning: Image ignored] % Unhandled or unsupported graphics:
%\includegraphics[width=2.902cm,height=2.915cm]{AASdev20090228-img135}

\end{center}
\item Bei Bedarf werden u.U. hier Handfunkger\"ate an die Bereichsleiter
ausgegeben. 

\item MLS ist die Br\"ucke zwischen verschiedenen Funkkan\"alen und
zust\"andig. 

\item Wenn die MLS die Einsatzleitung beinhaltet wird sie mit einem
rotem Dreh- oder Blinklicht markiert. 

\end{itemize}
\paragraph{F\"uhrungsstruktur in der SanHist}
\begin{center}
\tablehead{}\begin{supertabular}{m{2.4459999cm}m{4.0030003cm}m{2.299cm}m{4.1520004cm}}
\centering   [Warning: Image ignored]
% Unhandled or unsupported graphics:
%\includegraphics[width=1.654cm,height=1.624cm]{AASdev20090228-img136.png}
 \par

 &
Einsatzleiter Rettungsdienst

 &
\centering   [Warning: Image ignored]
% Unhandled or unsupported graphics:
%\includegraphics[width=1.62cm,height=1.62cm]{AASdev20090228-img137}
 \par

 &
Leiter Material-/Meldestelle

\\
\centering   [Warning: Image ignored]
% Unhandled or unsupported graphics:
%\includegraphics[width=1.62cm,height=1.62cm]{AASdev20090228-img138.png}
 \par

 &
Leiter SanHist

 &
\centering   [Warning: Image ignored]
% Unhandled or unsupported graphics:
%\includegraphics[width=1.489cm,height=1.61cm]{AASdev20090228-img139.gif}
 \par

 &
Leiter Behandlung

\\
\centering   [Warning: Image ignored]
% Unhandled or unsupported graphics:
%\includegraphics[width=1.62cm,height=1.62cm]{AASdev20090228-img140}
 \par

 &
LNA Leitender Notarzt

 &
\centering   [Warning: Image ignored]
% Unhandled or unsupported graphics:
%\includegraphics[width=1.62cm,height=1.62cm]{AASdev20090228-img141}
 \par

 &
Leiter Transport

\\
\centering   [Warning: Image ignored]
% Unhandled or unsupported graphics:
%\includegraphics[width=1.62cm,height=1.62cm]{AASdev20090228-img142}
 \par

 &
Leiter Triageraum/Triage

 &
\centering   [Warning: Image ignored]
% Unhandled or unsupported graphics:
%\includegraphics[width=1.62cm,height=1.62cm]{AASdev20090228-img143.png}
 \par

 &
Leiter Information

\\
\end{supertabular}
\end{center}
Patientenleitsystem (PLS)

Das \index{Patientenleitsystem}Patientenleitsystem (\index{PLS}PLS)
erm\"oglicht eine eindeutige Kennzeichnung von Peteinten und Zuordnung
von Behandlungsma{\ss}nahmen und Befunden. Der wichtigste Bestandteil
ist die \index{Patientenleittasche}Patientenleittasche
(\index{PLT}PLT). 

Sie besteht aus einer orangenen H\"ulle welche fortlaufend nummeriert
sind. Auf dieser Kunststoffh\"ulle k\"onnen wichtige Notizen wie
Triagegruppe, etc., vorgenommen werden. In der H\"ulle befinden sich
das Behandlungsprotokoll (blau), Idenentifikationsprotokoll (rosa),
Zusatzkarten, Kontaminationsaufkleber und ein Bogen mit gleichlautend
nummerierten Klebeetiketten. Am unteren Ende befinden sich zwei
Abschnitte welche die R\"uckverfolgbarkeit des Transports
sicherstellen. Die PLT ist mit einer Gummischnur versehen damit sie dem
Patienten umgeh\"angt werden kann. 

Sp\"atestens an der Triagestelle bekommt jeder Patient oder Betroffene
{\quotedblbase}seine{\textquotedblleft} PLT. Auch Unverletzte und
Verstrorbene werden auch mittels PLS registriert.

Zusatzkarten

Bei einer eventuellen Bergetriage bekommt jeder Patient eine
Patientenleittasche (PLT). Petienten, welche bevorrangt gerettet werden
sollen erhalten zus\"atzlich auf die PLT die gelbe DRINGEND Karte. Tote
erhalten eine PLT mit der schwarzen VERSTORBEN Karte.

Behandlungsprotokoll / Identifikationsprotokoll

Die Protokolle werden im Behandlungsraum ausgef\"ullt wenn Zeit daf\"ur
ist. Das Ausf\"ullen der Protokolle darf den Abtransport unter keinen
Umst\"anden verz\"ogern. Die Protokolle k\"onnen auch leer bleiben.

\index{Kontaminationsaufkleber}Kontaminationsaufkleber

Wenn der Patient kontaminiert ist, wird ein Dreieck auf die PLT geklebt.
Dies soll alle nachfolgenden Personen, welche mit dem Patienten in
Kontakt treten, warnen.

PLT-Nummer

Es erfolgt eine eindeutige Identifizierung der Patienten durch eine
einmalige Nummernkennzeichnung der Patientenleittasche. Mittels
Selbstklebenummern werden die Patienten in allen verschiedenen
Bereichen in Listen erfasst (Triage, Behandlung, Transport). Ausserdem
k\"onnen alle Effekte des Patienten damit gekennzeichnet werden.

Die Patientenleittasche verbleibt solange am Patienten, bis das
Aufnahmeverfahren in Krankenhaus abgeschlossen ist, d.h. bis er im
Spital administrativ erfasst ist. Anschlie{\ss}end wird die PLT der
Krankengeschichte des Patienten beigef\"ugt.

  [Warning: Image ignored] % Unhandled or unsupported graphics:
%\includegraphics[width=13.361cm,height=7.796cm]{AASdev20090228-img144}
 

\begin{center}
\begin{tabular}{m{6.649cm}m{6.651cm}}
\multicolumn{2}{m{13.5cm}}{\centering PLT Vorderseite\par

}\\
[Warning: Draw object ignored][Warning: Draw object ignored][Warning:
Draw object ignored]

\begin{center}
 [Warning: Image ignored] % Unhandled or unsupported graphics:
%\includegraphics[width=5.771cm,height=10.83cm]{AASdev20090228-img145}

\end{center}
 &
Triagegruppen:

Bei Triage nur Einteilung in I, II, III oder IV

Transportpriorit\"at:

A: Hohe Transportpriorit\"at: Der rasche Transport zur fr\"uhzeitigen
Fachbehandlung nach \"arztlicher Notversorgung ist angesagt. 

B: Niedrige Transportpriorit\"at: Der Transport zur Fachbehandlung kann
verz\"ogert erfolgen. 

wird sp\"ater entschieden

Zweiter Triageblock:

\begin{itemize}
\item F\"ur Umtriage oder

\item Spitalstriage

\end{itemize}
\\
\end{tabular}
\end{center}
\begin{center}
\tablehead{}\begin{supertabular}{m{6.649cm}m{6.651cm}}
\multicolumn{2}{m{13.5cm}}{\centering PLS Umtriage\par

}\\
[Warning: Draw object ignored][Warning: Draw object ignored][Warning:
Draw object ignored][Warning: Draw object ignored][Warning: Draw object
ignored][Warning: Draw object ignored]

\begin{center}
 [Warning: Image ignored] % Unhandled or unsupported graphics:
%\includegraphics[width=4.457cm,height=9.22cm]{AASdev20090228-img146}

\end{center}
 &
Umtriage innerhalb der ersten Viertelstunde

\begin{itemize}
\item Mittels Pfeil von einem zum nun richtigen Triagefeld

\item In der gleichen Zeile

\end{itemize}
\\
[Warning: Draw object ignored][Warning: Draw object ignored][Warning:
Draw object ignored][Warning: Draw object ignored][Warning: Draw object
ignored]

\begin{center}
 [Warning: Image ignored] % Unhandled or unsupported graphics:
%\includegraphics[width=4.457cm,height=9.063cm]{AASdev20090228-img147}

\end{center}
 &
Umtriage nach l\"angerer Zeit:

\begin{itemize}
\item Neues Kreuz beim richtigen Triagefeld

\item in der unteren Zeile

\end{itemize}
\\
\end{supertabular}
\end{center}
\begin{center}
\begin{tabular}{m{6.649cm}m{6.651cm}}
\multicolumn{2}{m{13.5cm}}{\centering PLT R\"uckseite\par

}\\
[Warning: Draw object ignored][Warning: Draw object ignored][Warning:
Draw object ignored][Warning: Draw object ignored][Warning: Draw object
ignored][Warning: Draw object ignored][Warning: Draw object
ignored][Warning: Draw object ignored][Warning: Draw object ignored]

\begin{center}
 [Warning: Image ignored] % Unhandled or unsupported graphics:
%\includegraphics[width=6.653cm,height=12.217cm]{AASdev20090228-img148}

\end{center}
 &
Der Triagearzt tr\"agt auf der R\"uckseite die zu setzenden
Ma{\ss}nahmen ein:

\begin{itemize}
\item Atmung (O2, Intubation, usw..)

\item Kreislauf (Blutstillung usw...)

\item Medikamente

\item Lagerung

\end{itemize}
Der Durchf\"uhrende (Arzt, NFS, RS) best\"atigt die Durchf\"uhrung der
Ma{\ss}nahmen im rechten Feld durch eine Paraphe

\\
 &
Abriss f\"ur das Zielspital:

Wird bei der Einlieferung mit der Ankunftszeit versehen und verbleibt in
der Krankengeschichte.

\\
 &
Abriss f\"ur die SanHist -- Leiter Transport:

Eintragen: Zielspital, Kfz Nummer, Verladezeit

Verbleibt beim Leiter Transport/Protokollf\"uhrer

\\
\end{tabular}
\end{center}
\subsection{Gefahrengutunfall}
\begin{center}
\tablehead{}\begin{supertabular}{m{6.649cm}m{6.651cm}}
  [Warning: Image ignored] % Unhandled or unsupported graphics:
%\includegraphics[width=6.233cm,height=4.351cm]{AASdev20090228-img149.jpg}
 

 &
\raggedleft   [Warning: Image ignored]
% Unhandled or unsupported graphics:
%\includegraphics[width=6.314cm,height=4.41cm]{AASdev20090228-img150.jpg}
 \par

\\
\end{supertabular}
\end{center}
\begin{center}
\tablehead{}\begin{supertabular}{|m{6.649cm}|m{6.651cm}|}
\hline
M\"ogliche Gefahren f\"ur den Sanit\"ater

\begin{itemize}
\item Brand bzw. Explosionsgefahr

\item Entweichen unter Druck stehender Gase

\item Giftigkeit, Ansteckungsgefahr

\item \"Atzwirkung

\item Radioaktivit\"at

\item Umweltgef\"ahrdung

\end{itemize}
 &
Gefahrenquellen k\"onnen entstehen bei

\begin{itemize}
\item Unf\"alle bei der Herstellung

\item Unf\"alle bei der Verwendung von gef\"ahrlichen Stoffen

\item Transportunf\"alle

\end{itemize}
\\\hline
\end{supertabular}
\end{center}
Andere Gefahrenquellen:

\begin{flushleft}
\tablehead{}\begin{supertabular}{|m{6.393cm}m{6.9030004cm}|}
\hline
Landwirtschaft

 &
CO2 in G\"arkellern, Siloanlagen, Brunnensch\"achte

Pflanzenschutzmittel

\\\hline
Kl\"argruben

 &
CO2 Methan, Schwefelwasserstoff

\\\hline
Obstlagerhallen

 &
wenig O2, viel CO2 oder N2 halten Obst frisch

\\\hline
Br\"ande in Wirtschaftsgeb\"auden

 &
extrem gef\"ahrliche Rauchgase

\\\hline
Tiefgaragen und Tunnels

 &
erh\"ohte CO- und CO2-Werte

\\\hline
Fl\"ussiggas

 &
Campingpl\"atze, feste und mobile K\"uchen

\\\hline
Pools, Schwimmb\"ader

 &
Chlor (Wasserdesinfektion)

\\\hline
Industrie

 &
K\"uhlmittel in gro{\ss}en K\"alteanlagen (Ammoniak)

\\\hline
Speziallabors, Einrichtungen mit und zur Verarbeitung von infekti\"osen
Abf\"allen oder Tierkadaver

 &
Infektionsgefahr

\\\hline
Nuklearmedizinische Institute, Industrie, Forschungseinrichtungen

 &
Radioaktivit\"at

\\\hline
\end{supertabular}
\end{flushleft}
Kennzeichnung von gef\"ahrlichen Stoffen und Gefahrenbereichen

\begin{center}
\begin{tabular}{|m{3.225cm}m{3.225cm}|m{3.225cm}m{3.225cm}|}
\hline
\multicolumn{2}{|m{6.65cm}|}{\centering Hinweis auf Gefahr\par

} &
\multicolumn{2}{m{6.65cm}|}{\centering Warnzeichen\par

}\\\hline
\centering   [Warning: Image ignored]
% Unhandled or unsupported graphics:
%\includegraphics[width=1.475cm,height=1.452cm]{AASdev20090228-img151.png}
 \par

 &
Explosionsgef\"ahrlich

 &
\centering   [Warning: Image ignored]
% Unhandled or unsupported graphics:
%\includegraphics[width=1.541cm,height=1.426cm]{AASdev20090228-img152.png}
 \par

 &
Allgemeine Gefahr

\\\hline
\centering   [Warning: Image ignored]
% Unhandled or unsupported graphics:
%\includegraphics[width=1.452cm,height=1.414cm]{AASdev20090228-img153.png}
 \par

 &
Giftig

 &
\centering   [Warning: Image ignored]
% Unhandled or unsupported graphics:
%\includegraphics[width=1.673cm,height=1.414cm]{AASdev20090228-img154.png}
 \par

 &
Radioaktive Stoffe, ionisierende Strahlen

\\\hline
\centering   [Warning: Image ignored]
% Unhandled or unsupported graphics:
%\includegraphics[width=1.497cm,height=1.611cm]{AASdev20090228-img155.png}
 \par

 &
Hochentz\"undlich

 &
\centering   [Warning: Image ignored]
% Unhandled or unsupported graphics:
%\includegraphics[width=1.696cm,height=1.412cm]{AASdev20090228-img156.png}
 \par

 &
Gasflaschen

\\\hline
\centering   [Warning: Image ignored]
% Unhandled or unsupported graphics:
%\includegraphics[width=1.475cm,height=1.435cm]{AASdev20090228-img157.png}
 \par

 &
Reizend

 &
\centering   [Warning: Image ignored]
% Unhandled or unsupported graphics:
%\includegraphics[width=1.673cm,height=1.414cm]{AASdev20090228-img158.png}
 \par

 &
Laserstrahlen

\\\hline
\centering   [Warning: Image ignored]
% Unhandled or unsupported graphics:
%\includegraphics[width=1.488cm,height=1.599cm]{AASdev20090228-img159.png}
 \par

 &
\"Atzend

 &
 &
\\\hline
\end{tabular}
\end{center}
Transport gef\"ahrlicher G\"uter

Kennzeichnungspflicht besteht erst dann, wenn die transportierte Menge
eines Stoffes die gesetzlich erlaubte Mindestmenge \"uberschreitet.
Diese Mindestmenge gilt f\"ur jeden einzelnen Stoff und nicht f\"ur die
Gesamtmenge.

Es k\"onnen also erhebliche Mengen gef\"ahrlicher G\"uter ohne
Kennzeichnung des Fahrzeuges transportiert werden. 

Gefahrzettel

Bei St\"uckguttransporten wird jedes Versandst\"uck gekennzeichnet. An
Fahrzeugen, die kein St\"uckgut transportieren m\"ussen Gefahrenzettel
an den L\"angsseiten und an der R\"uckseite angebracht werden.

\begin{flushleft}
\tablehead{}\begin{supertabular}{m{2.1599998cm}m{4.289cm}m{2.1599998cm}m{4.2840004cm}}
\centering Beispiele\par

 &
 &
 &
\\
\centering   [Warning: Image ignored]
% Unhandled or unsupported graphics:
%\includegraphics[width=1.659cm,height=1.885cm]{AASdev20090228-img160.png}
 \par

 &
Feuergef\"ahrlich (fl\"ussige Stoffe)

 &
\centering   [Warning: Image ignored]
% Unhandled or unsupported graphics:
%\includegraphics[width=1.856cm,height=2.096cm]{AASdev20090228-img161.png}
 \par

 &
Selbstentz\"undlich

\\
\centering   [Warning: Image ignored]
% Unhandled or unsupported graphics:
%\includegraphics[width=1.835cm,height=1.902cm]{AASdev20090228-img162.png}
 \par

 &
Feuergef\"ahrlich (feste Stoffe)

 &
\centering   [Warning: Image ignored]
% Unhandled or unsupported graphics:
%\includegraphics[width=1.909cm,height=2.162cm]{AASdev20090228-img163.png}
 \par

 &
Entz\"undbare Gase bei Ber\"uhrung mit Wasser

\\
\end{supertabular}
\end{flushleft}
Warntafel

\begin{center}
\begin{tabular}{|m{3.225cm}m{3.225cm}|m{3.225cm}m{3.225cm}|}
\hline
\multicolumn{2}{|m{6.65cm}|}{\centering AllgemeineKennzeichnung
(Sammeltransporte)\par

} &
\multicolumn{2}{m{6.65cm}|}{\centering SpezielleKennzeichnung\par

}\\\hline
\centering [Warning: Draw object ignored]\par

 &
 &
\centering [Warning: Draw object ignored]\par

 &
Gefahrnummer  (Kemler-Nummer)

\\\hline
 &
 &
 &
Stoffnummer  (UN-Nummer)

\\\hhline{~~~-}
\end{tabular}
\end{center}
Gefahrnummer

\begin{flushleft}
\tablehead{}\begin{supertabular}{|m{0.97999996cm}m{12.317cm}|}
\hline
\centering 2\par

 &
Entweichen von Gas durch Druck oder chemische Reaktion

\\\hline
\centering 3\par

 &
Entz\"undbarkeit von Fl\"ussigkeiten (D\"ampfen) und Gasen

\\\hline
\centering 4\par

 &
Entz\"undbarkeit fester Stoffe

\\\hline
\centering 5\par

 &
Oxidierende (brandf\"ordernde) Wirkung

\\\hline
\centering 6\par

 &
Giftig

\\\hline
\centering 7\par

 &
Radioaktivit\"at

\\\hline
\centering 8\par

 &
\"Atzwirkung

\\\hline
\centering 9\par

 &
Gefahr einer spontan heftigen Reaktion

\\\hline
\end{supertabular}
\end{flushleft}
Die Verdopplung einer Ziffer weist auf die Zunahme der entsprechenden 
Gefahr hin.

X  vor der Nummer zur Kennzeichnung der Gefahr ${\rightarrow}$ Stoff
reagiert in  gef\"ahrlicher Weise mit Wasser

Wenn die Gefahr eines Stoffes ausreichend von einer einzigen Ziffer 
angegeben werden kann, wird dieser Ziffer eine Null angef\"ugt.

Folgende Ziffernkombinationen haben eine besondere Bedeutung

\begin{flushleft}
\begin{tabular}{|m{0.97999996cm}m{12.317cm}|}
\hline
22

 &
tiefgek\"uhltes Gas

\\\hline
X323

 &
entz\"undbarer fl\"ussiger Stoff, der mit Wasser gef\"ahrlich reagiert,
wobei entz\"undbare Gase entweichen

\\\hline
X333

 &
selbstentz\"undliche Fl\"ussigkeit, die mit Wasser gef\"ahrlich reagiert

\\\hline
X423

 &
entz\"undbarer fester Stoff, der mit Wasser gef\"ahrlich reagiert, wobei
brennbare Gase entweichen

\\\hline
44

 &
entz\"undbarer fester Stoff, der sich bei erh\"ohter Temperatur in
Geschmolzenem Zustand befindet

\\\hline
539

 &
entz\"undbares organisches Peroxid

\\\hline
90

 &
verschiedene gef\"ahrlich Stoffe

\\\hline
\end{tabular}
\end{flushleft}
Einsatzrichtlinie

Gefahr erkennen

Absperrma{\ss}nahmen durchf\"uhren

Spezialkr\"afte anfordern

\begin{center}
\tablehead{}\begin{supertabular}{|m{2.222cm}|m{5.1460004cm}|m{5.7320004cm}|}
\hline
Gefahr erkennen

 &
Die Suche nach Warntafeln ist bei umgest\"urzten Fahrzeugen erschwert

 &
\centering   [Warning: Image ignored]
% Unhandled or unsupported graphics:
%\includegraphics[width=5.137cm,height=3.747cm]{AASdev20090228-img164.jpg}
 \par

\\\hline
Absperr\-ma{\ss}nahmen durch\-f\"uhren

 &
Die Absperrung sollte gro{\ss}r\"aumig  sein. Mit gen\"ugend Platz f\"ur
die technische Hilfeleistungen. Achtung  vor sich \"andernden
Situationen  (Windrichtung, Niederschlag)

 &
\centering   [Warning: Image ignored]
% Unhandled or unsupported graphics:
%\includegraphics[width=3.46cm,height=3.409cm]{AASdev20090228-img165.png}
 \par

\\\hline
Spezial\-kr\"afte anfordern

 &
 &
\centering   [Warning: Image ignored]
% Unhandled or unsupported graphics:
%\includegraphics[width=5.666cm,height=3.941cm]{AASdev20090228-img166.jpg}
 \par

\\\hline
\end{supertabular}
\end{center}
\section{Fallbeispiele und Arbeitsaufgaben}
Fallbeispiele und Arbeitsaufgaben f\"ur BLS

Fall 

\begin{flushleft}
\tablehead{}\begin{supertabular}{|m{2.0839999cm}|m{1.43cm}m{1.43cm}m{1.43cm}m{1.43cm}m{1.43cm}m{1.43cm}m{1.43cm}|}
\hline
\centering Berufung\par

 &
\multicolumn{1}{m{1.43cm}|}{\centering Typ\par

} &
\multicolumn{1}{m{1.43cm}|}{\centering RTW\par

} &
\multicolumn{1}{m{1.43cm}|}{\centering Code\par

} &
\multicolumn{1}{m{1.43cm}|}{\centering 26-A-1\par

} &
\multicolumn{3}{m{4.69cm}|}{}\\\hline
 &
\multicolumn{1}{m{1.43cm}|}{} &
\multicolumn{1}{m{1.43cm}|}{} &
\multicolumn{1}{m{1.43cm}|}{} &
\multicolumn{1}{m{1.43cm}|}{} &
\multicolumn{1}{m{1.43cm}|}{} &
\multicolumn{1}{m{1.43cm}|}{} &
\\\hline
\centering Umst\"ande\par

 &
\multicolumn{1}{m{1.43cm}|}{\centering Jahreszeit\par

} &
\multicolumn{1}{m{1.43cm}|}{\centering Tageszeit\par

} &
\multicolumn{1}{m{1.43cm}|}{\centering Wetter\par

} &
\multicolumn{1}{m{1.43cm}|}{} &
\multicolumn{1}{m{1.43cm}|}{\centering Ort\par

} &
\multicolumn{1}{m{1.43cm}|}{} &
\\\hline
 &
\multicolumn{1}{m{1.43cm}|}{\centering Sommer\par

} &
\multicolumn{1}{m{1.43cm}|}{\centering Vormittag\par

} &
\multicolumn{2}{m{3.06cm}|}{\centering schw\"ul, hei{\ss}\par

} &
\multicolumn{1}{m{1.43cm}|}{\centering Wohnung\par

} &
\multicolumn{1}{m{1.43cm}|}{} &
\\\hhline{~-------}
\centering Patient:\par

 &
\multicolumn{1}{m{1.43cm}|}{\centering Alter\par

} &
\multicolumn{1}{m{1.43cm}|}{\centering Geschl.\par

} &
\multicolumn{1}{m{1.43cm}|}{\centering AZ\par

} &
\multicolumn{1}{m{1.43cm}|}{\centering EZ\par

} &
\multicolumn{1}{m{1.43cm}|}{} &
\multicolumn{1}{m{1.43cm}|}{} &
\\\hline
 &
\multicolumn{1}{m{1.43cm}|}{\centering 65\par

} &
\multicolumn{1}{m{1.43cm}|}{\centering m\par

} &
\multicolumn{1}{m{1.43cm}|}{\centering red.\par

} &
\multicolumn{1}{m{1.43cm}|}{\centering gut\par

} &
\multicolumn{1}{m{1.43cm}|}{} &
\multicolumn{1}{m{1.43cm}|}{} &
\\\hhline{~-------}
Szene

 &
\multicolumn{7}{m{11.21cm}|}{Patient sitzt bla{\ss} in einem Lehnstuhl,
ansprechbar, wirkt nicht akut bedroht 

}\\\hline
Beschwerden

 &
\multicolumn{7}{m{11.21cm}|}{{\quotedblbase}Schwindel{\textquotedblleft},
erstmalig aufgetreten, keine Schmerzen

}\\\hline
Krankheiten

 &
\multicolumn{7}{m{11.21cm}|}{{\quotedblbase}Hypertonie,
Gef\"a{\ss}verkalkung, Angina pectoris{\textquotedblleft}

Blutdruckprotokoll: sie finden f\"ur die vergangenen vier Tage folgende
Eintr\"age: 170/110, 180/115, 175/113, 183/112, ...

}\\\hline
Allergien

 &
\multicolumn{7}{m{11.21cm}|}{keine bekannt

}\\\hline
Medikamente

 &
\multicolumn{7}{m{11.21cm}|}{Patient \"uberreicht
Liste:{\quotedblbase}Blutdruckmittel, aber die wirken nicht so richtig
sagt mein Arzt{\textquotedblleft}, Pat zeigt eine Sprayflasche
{\quotedblbase}Nitrolingual{\textquotedblleft} her den er gestern vom
Hausarzt {\quotedblbase}f\"ur sein Herz{\textquotedblleft} bekommen
hat.

}\\\hline
Ereignisse

 &
\multicolumn{7}{m{11.21cm}|}{{\quotedblbase}Ich habe das erste mal
diesen Nitro-Spray genommen und da ist mir damisch
word{\textquotesingle}n{\textquotedblleft}

}\\\hline
Letzte(r)

 &
\multicolumn{7}{m{11.21cm}|}{Krankenhausaufenthalt: vor 3 Jahren wg.
Blinddarmentz\"undung

Arztbesuch: gestern, wegen Herzschmerzen, Nitro-Spray verschrieben
bekommen

Mahlzeit: Fr\"uhst\"uck, 2 Marmeladesemmeln, Kaffee

}\\\hline
\multicolumn{8}{|m{13.494cm}|}{\centering Erstbefund\par

}\\\hline
Anblick

 &
\multicolumn{7}{m{11.21cm}|}{altersensprechend, gut gen\"ahrt, gepflegt,
jetzt schwach

}\\\hline
Haut

 &
\multicolumn{7}{m{11.21cm}|}{bla{\ss}, nicht schweissig

}\\\hline
Trauma grob

 &
\multicolumn{7}{m{11.21cm}|}{{}-{}-{}-

}\\\hline
 &
\multicolumn{1}{m{1.43cm}|}{\centering RR\par

} &
\multicolumn{1}{m{1.43cm}|}{\centering HF\par

} &
\multicolumn{1}{m{1.43cm}|}{\centering Puls\par

} &
\multicolumn{1}{m{1.43cm}|}{\centering SpO2\par

} &
\multicolumn{1}{m{1.43cm}|}{\centering AF\par

} &
\multicolumn{1}{m{1.43cm}|}{\centering BZ\par

} &
\centering Temp\par

\\\hline
 &
\multicolumn{1}{m{1.43cm}|}{\centering 100/70\par

} &
\multicolumn{1}{m{1.43cm}|}{\centering 120\par

} &
\multicolumn{1}{m{1.43cm}|}{\centering rhythm.\par

} &
\multicolumn{1}{m{1.43cm}|}{\centering 98\%\par

} &
\multicolumn{1}{m{1.43cm}|}{\centering 12\par

} &
\multicolumn{1}{m{1.43cm}|}{\centering 105\par

} &
\centering 37,1\par

\\\hline
Trauma Detail

 &
\multicolumn{7}{m{11.21cm}|}{{}-{}-{}-

}\\\hline
Neuro

 &
\multicolumn{7}{m{11.21cm}|}{grob neurologisch unauff\"allig.

}\\\hline
Ma{\ss}nahmen, Verlauf

 &
\multicolumn{1}{m{1.43cm}|}{} &
\multicolumn{5}{m{7.95cm}|}{} &
\\\hline
Transport

 &
\multicolumn{1}{m{1.43cm}|}{\centering Lagerung\par

} &
\multicolumn{5}{m{7.95cm}|}{} &
\\\hline
\"Ubergabe

 &
\multicolumn{1}{m{1.43cm}|}{\centering Abteilung:\par

} &
\multicolumn{5}{m{7.95cm}|}{} &
\\\hline
\end{supertabular}
\end{flushleft}
Wie beurteilen Sie den Blutdruck des Patienten? 

Welche Informationen fehlen Ihnen noch?

Wie lautet Ihre Verdachtsdiagnose? 

Welche Differentialdiagnosen ziehen Sie in Betracht?

\section{ILS: Notfallrettung Grundlagen}
Entspricht SAN 3-2-1-Einschulung

\subsection{Anamnese und Untersuchung}
\subsection{Megacodetraining}
\subsection{MPG}
\section{ILS: NFS-Aufbaumodul}
\section{ILS: Arzneimittel}
\section{ILS: Minimalinvasive T\"atigkeiten und Behandlung}
\section{Begriffserkl\"arungen}
\label{bkm:LexComputertomographie}Computertomographie. Abk.: CT. 

CT. ${\rightarrow}$Computertomographie.

Kernspintomographie. ${\rightarrow}$Magnetresonanztomographie

Magnetresonanztomographie. Syn.: Kernspintomographie. Abk.: MRT

MRT. ${\rightarrow}$Magnetresonanztomographie

NAW. ${\rightarrow}$Notarztwagen.

Notarztwagen. Abk.: NAW.

R\"ontgen. 

\section{Impressum, Hinweise}
\begin{multicols}{2}
ARBEITER-SAMARITER-BUND \"OSTERREICHS GRUPPE FLORIDSDORF-DONAUSTADT

WALLENBERGGASSE 2

A-1220 WIEN

TEL. +43 1 26 33 144

FAX  +43 1 26 33 144 75

WWW.SAMARITER.AT

OFFICE@SAMARITER.AT

ZVR: 317 994 238  UID: ATU 163 722 08  HG WIEN

Betriebsgesellschaft und Rechnungsanschrift:

ARBEITER-SAMARITER-BUND \"OSTERREICHS FLORIDSDORF-DONAUSTADT KRANKEN-,
RETTUNGSTRANSPORT UND SOZIALE DIENSTE GEM. GESMBH Kontakt und Adresse
w. o. FN 287700t, HG Wien UID: ATU 63216914

Gesch\"aftsf\"uhrer: Berhnard Lehner, Ing. Alexander Prischl

\"Arztliche Leitung: Dr. Johannes Steuer, MSc,  Stv: Dr. Helmut Paul

Leitung Ausbildungszentrum: Ing. Roman Koch

Aufsichtsapotheke: Mag.pharm. Ulrike Urban, Stadtapotheke
Klosterneuburg. 

Hinweise

Alle Angaben sind rein informativ zu ver\-stehen und begr\"unden KEINE
Dienst\-an\-wei\-sung oder Lehr\-meinung. 

Wie jede Wissenschaft ist die Medizin st\"andigen Ent\-wicklungen
unterworfen. Soweit fachliche Angaben ge\-macht werden, darf der Leser
zwar darauf ver\-trau\-en, dass die Autoren gro{\ss}e Sorg\-falt
da\-rauf verwandt haben, dass diese An\-ga\-ben dem Wissens\-stand bei
Ferti\-gstel\-lung des Werkes entspricht. F\"ur Angaben, spe\-ziell
\"uber Dosierungsanweisungen und Ap\-plikations\-formen, kann jedoch
keine Ge\-w\"ahr \"ubernommen werden. Jeder Leser ist angehalten, durch
sorg\-f\"altige Pr\"u\-fung der Literatur, der Lehrmeinung so\-wie der
Beipackzettel der ver\-wendeten Pr\"aparate und gegebenenfalls nach
Konsultation eines Spezialisten fest\-zustellen, ob die gegebene
Aussage oder Empfehlung f\"ur Do\-sierungen oder die Beachtung von
Kon\-tra\-indi\-ka\-tionen gegen\"uber der Angabe in dieser
Mit\-tei\-lung abweicht. Jede Dosierung oder Ap\-pli\-kation erfolgt
auf eigene Gefahr des Be\-nutzers. Das AMS appelliert an jeden Le\-ser
ihm etwa auffallende Un\-ge\-nau\-igkei\-ten dem AMS umgehend
mitzuteilen.

Gesch\"utzte Waren\-namen und Wa\-ren\-zei\-chen werden nicht besonders
kenntlich ge\-macht. Aus dem Fehlen eines solchen Hin\-weises kann also
nicht ge\-schlossen werden, dass es sich um einen freien Waren\-namen
handelt.

\end{multicols}
Sprechstunden

Mittwoch nachmittag. 

\printindex
\listoffigures
SO LONG..........

AND THANKS FOR ALL THE FISH

\section[Literaturverzeichnis]{Literaturverzeichnis}
[1] , 2002. Pschyrembel Klinisches W\"orterbuch. de Gruyter, .

[2] Longmore, Wilkinson, T. C. e., 2007. Oxford Handbook of Clinical
Medicine. Oxford University Press, .

\end{document}
